\documentclass[11pt]{amsart}


\date{version of \today}
\bibliographystyle{unsrt}
\usepackage{amsmath}
\usepackage{amssymb}
\usepackage{amsxtra}
\usepackage{abstract}
\usepackage{mathrsfs}
\usepackage{mathtools}
\usepackage[all]{xy}
\usepackage{amscd}
\usepackage{amsthm}
\usepackage[dvips]{graphicx}
\usepackage{ulem}
\usepackage{color}
\usepackage{appendix}
\usepackage{tikz}
\usepackage{tikz-cd}
\usepackage{url}
\newcommand{\comment}[1]{{\color[rgb]{0.15,0.8,0.15}#1}}
\newcommand{\nql}[1]{\textcolor{cyan}{#1}}
\newcommand{\bhx}[1]{\textcolor{purple}{#1}}
\newcommand{\pyc}[1]{\textcolor{red}{#1}}
\newcommand{\yml}[1]{\textcolor{blue}{#1}}

\newtheorem{Thm}{Theorem}[section]
\newtheorem{Lem}[Thm]{Lemma}
\newtheorem{Def}[Thm]{Definition}
\newtheorem{Cor}[Thm]{Corollary}
\newtheorem{Prop}[Thm]{Proposition}
\newtheorem{Ex1}[Thm]{Example}
\newtheorem{Rem1}[Thm]{Remark}
\newtheorem{Conj}[Thm]{Conjecture}
\newtheorem{Prob}[Thm]{Problem}

\newenvironment{Rem}{\begin{Rem1}\rm}{\end{Rem1}}
\newenvironment{Ex}{\begin{Ex1}\rm}{\end{Ex1}}

% environment of pyc %%%%%%%%%%%%%%%%%%%%%%%%%%%%%%%%%%%%%%%%%%%%%%%%%%%%%%%
\usepackage{bbm}
\usepackage{enumitem}

\tikzset{every picture/.style={line width=0.75pt}} %set default line width to 0.75pt
\newcommand{\bb}[1]{\mathbbm{#1}}

\newcommand{\sHom}{\underline{\mathrm{Hom}}}
\newcommand{\Hom}{\mathop{\mathrm{Hom}}}
\renewcommand{\top}{\mathop{\mathrm{top}}}
\newcommand{\rad}{\mathop{\mathrm{rad}}}
\newcommand{\soc}{\mathop{\mathrm{soc}}}
\renewcommand{\dim}{\mathop{\mathrm{dim}}}
\renewcommand{\span}{\mathop{\mathrm{span}}}
\renewcommand{\mod}{\mathop{\mathrm{mod}}}

\newenvironment{theorem}{\begin{Thm}}{\end{Thm}}
\newenvironment{lemma}{\begin{Lem}}{\end{Lem}}
\newenvironment{definition}{\begin{Def}}{\end{Def}}
\newenvironment{corollary}{\begin{Cor}}{\end{Cor}}
\newenvironment{proposition}{\begin{Prop}}{\end{Prop}}
\newenvironment{example}{\begin{Ex1}}{\end{Ex1}}
\newenvironment{remark}{\begin{Rem1}}{\end{Rem1}}
\newenvironment{conjecture}{\begin{Conj}}{\end{Conj}}
\newenvironment{problem}{\begin{Prob}}{\end{Prob}}

\usepackage[bookmarks=true, colorlinks=true, citecolor=blue, linkcolor=black]{hyperref}
% environment of pyc %%%%%%%%%%%%%%%%%%%%%%%%%%%%%%%%%%%%%%%%%%%%%%%%%%%%%%%

\newcommand{\qa}{kQ/I}

\newcommand{\m}{m}

\newcommand{\uar}{\uparrow}
\newcommand{\dar}{\downarrow}
\newcommand{\lra}{\longrightarrow}
\newcommand{\lla}{\longleftarrow}
\newcommand{\la}{\leftarrow}
\newcommand{\ra}{\rightarrow}
\newcommand{\sdp}{\times\kern-.2em\vrule height1.1ex depth-.05ex}
\newcommand{\epi}{\lra \kern-.8em\ra}
\newcommand{\C}{{\mathbb C}}
\newcommand{\Q}{{\mathbb Q}}
\newcommand{\N}{{\mathbb N}}
\newcommand{\R}{{\mathbb R}}
\newcommand{\F}{{\mathbb F}}
\newcommand{\ul}{\underline}
\newcommand{\ol}{\overline}
\newcommand{\Sn}{{\mathfrak S}}
\newcommand{\cal}[1]{\mathcal{#1}}
\newcommand{\Z}{{\mathbb Z}}
\newcommand{\Pres}{{\mathbb P}}
\newcommand{\dickbox}{{\vrule height5pt width5pt depth0pt}}
\newcommand{\Scomp}{comproj}
% \newcommand{\rad}{\mbox{rad}}
% \newcommand{\soc}{\mbox{soc}}
\newcommand{\stmod}{\underline{\mathrm{mod}}}
\newcommand{\al}{\alpha}
\newcommand{\be}{\beta}

\setlength{\textwidth}{16.3cm} \setlength{\textheight}{23.5cm}
\setlength{\topmargin}{0.2cm} \setlength{\oddsidemargin}{-3mm}
\setlength{\evensidemargin}{-3mm} \setlength{\abovedisplayskip}{3mm}
\setlength{\belowdisplayskip}{3mm}
\setlength{\abovedisplayshortskip}{0mm}
\setlength{\belowdisplayshortskip}{2mm} \normalbaselines
\raggedbottom

\begin{document}

	\title{Brauer graph algebras are closed under stable equivalence of Morita type}
	\author{Pengyun Chen$^a$, Nengqun Li$^b$, Yuming Liu$^a$, and Bohan Xing$^{a,*}$}
	\maketitle
	
	\renewcommand{\thefootnote}{\alph{footnote}}
	\setcounter{footnote}{-1} \footnote{\it{Mathematics Subject
			Classification(2020)}: 16G10, 16D50.}
	\renewcommand{\thefootnote}{\alph{footnote}}
	\setcounter{footnote}{-1} \footnote{\it{Keywords}: Brauer graph algebra, Deformed loop, Exceptional band module, Stable equivalence of Morita type, Symmetric stably biserial algebra.}
	\setcounter{footnote}{-1} \footnote{$^a$Pengyun Chen, Yuming Liu, and Bohan Xing, School of Mathematical Sciences, Laboratory of Mathematics and Complex Systems, Beijing Normal University,
		Beijing 100875,  P. R. China.}
	\setcounter{footnote}{-1} \footnote{$^b$Nengqun Li, School of Mathematics, Liaoning Normal University,
		Dalian 116029,  P. R. China.}
	\setcounter{footnote}{-1} \footnote{E-mail addresses: pychen@mail.bnu.edu.cn (P. Chen); linengqun@lnnu.edu.cn (N. Li); ymliu@bnu.edu.cn (Y. Liu); bhxing@mail.bnu.edu.cn (B. Xing).}
	\setcounter{footnote}{-1} \footnote{$^*$Corresponding author.}
	
\noindent{\bf Abstract:} We study stable equivalences of Brauer graph algebras. In particular, we prove that Brauer graph algebras are closed up to semisimple summands under stable equivalence of Morita type. As a consequence, we reprove the result of Antipov and Zvonareva that Brauer graph algebras are closed under derived equivalence. As a byproduct, we get a solution of the reconstruction problem posed by Rickard and Rouquier for algebras stably equivalent of Morita type to Brauer graph algebras.

\section{Introduction}

As one of the basic equivalences, the stable equivalence plays an important role in the representation theory of finite dimensional algebras and finite groups. One early result on stable equivalence arising from modular group representation theory says that any group algebra of finite representation type is stably equivalent to a Nakayama algebra. Moreover, it is well-known that the symmetric algebras stably equivalent to symmetric Nakayama algebras are exactly those given by Brauer trees (see for example in the textbook \cite{ARS}).

Brauer graph algebras are generalization of Brauer tree algebras and they coincide with symmetric special biserial algebras over an algebraically closed field. In fact, the Brauer graph algebras of finite representation type are precisely Brauer tree algebras. Each Brauer graph algebra is defined by a Brauer graph, which is a ribbon graph endowed with a multiplicity function, and the representation theory information of a Brauer graph algebra is encoded in its Brauer graph. In particular, each Brauer graph algebra can be defined by quiver with relations. For a nice survey on Brauer graph algebras, we refer to Schroll \cite{S}.

According to a result claimed by Pogorzaly in \cite{P}, self-injective special biserial algebras can be stably equivalent to stably biserial algebras only, and these two classes coincide. However, Ariki, Iijima and Park showed in \cite[Example A.7]{AIP} that these two classes do not coincide. By modifying some ideas from \cite{P}, Antipov and Zvonareva \cite{AZ1} gave a complete proof for the facts that a self-injective special biserial algebra can be stably equivalent to a stably biserial algebra only (here we assume that the self-injective algebras have no summand isomorphic to a Nakayama algebra of radical square zero) and that the Auslander-Reiten conjecture is true for algebras stably equivalent to self-injective special biserial algebras. Moreover, they showed that each symmetric stably biserial algebra is associated with a Brauer graph $(\Gamma,m)$ together with a set $\mathcal{L}$ of loops, called deformed loops; and in characteristic different from $2$, the classes of symmetric special biserial algebras, that is, Brauer graph algebras, and symmetric stably biserial algebras coincide.

Recently, Antipov and Zvonareva \cite{AZ2} proved, using some subtle derived invariants, that Brauer graph algebras are closed under derived equivalence. More recently, Opper and Zvonareva \cite{OZ} gave a complete classification of Brauer graph algebras up to derived equivalence. Nevertheless, the question of whether Brauer graph algebras are closed under stable equivalence remains largely unexplored. In general, this question may have a negative answer. For example, the Brauer graph algebra $k[x]/(x^2)$ is stably equivalent to the hereditary algebra $k(\cdot\rightarrow\cdot)$. However, if we exclude the algebras of radical square zero, then we even do not know whether a Brauer graph algebra can be stably equivalent to some non-symmetric algebra. Under some additional assumptions on the algebras related by a stable equivalence, we obtain the following result.

\begin{Thm}\textnormal{(see Theorem \ref{thm:st-BGA})}\label{thm:1.1}
Let $A$ be a non-local Brauer graph algebra, and let $B$ be a basic symmetric algebra with no semisimple summands. If $B$ and $A$ are stably equivalent, then $B$ is also a Brauer graph algebra.
\end{Thm}

In order to prove the above result, we shall use (and generalize) some invariants of the (triangulated) stable categories of Brauer graphs algebras which are proved by Antipov in \cite{A}. For example, for a Brauer graph algebra $A$, he shows that
$$
\mathrm{dim}_k \underline{\operatorname{Hom}}_A(X,Y) \leq 2,
$$
where $X$ is either a maximal chain module (that is, chain summand of the radical of an indecomposable non-uniserial projective module) or an exceptional band module (Definition \ref{def:exceptional-band}) and $Y$ is a string module; see the discussion following \cite[Proposition~2.2]{A} and \cite[Proposition~2.8]{A}. For certain homogeneous tubes arising from stably biserial algebras, we prove an analogous result. This will serve as one of key lemmas in the proof of Theorem \ref{thm:1.1}.

\begin{Lem}\textnormal{(see Lemma \ref{lem:equal-dimension})}
	Let $\varLambda$ be a symmetric stably biserial algebra, and let $X$ be a string module lying at the mouth of a homogeneous tube corresponding to a deformed loop in the Brauer graph of $\varLambda$. Then $X$ is a maximal chain module and $\Omega_{\varLambda}(X)$ is an exceptional band module. For any string $\varLambda$-module $Y$, we have
$$\mathrm{dim}_k \sHom_{\varLambda}(X,Y)=\mathrm{dim}_k \sHom_{\varLambda}(\Omega_{\varLambda}(X),Y)
\leq 2.$$
\end{Lem}

Theorem \ref{thm:1.1} also motivates us to investigate the above question for a special kind of stable equivalences, the stable equivalences of Morita type. Note that the stable equivalence of Morita type plays a fundamental role in the representation theory of self-injective algebras, establishing connections between their stable module categories while preserving essential homological properties (see for example in \cite[Chapter 5]{Z}). In this setting, we obtain the following result.

\begin{Thm}\textnormal{(see Theorem \ref{thm:st.M-BGA})}\label{thm:intro}
 Let $A$ be a Brauer graph algebra. Then for any basic algebra $B$ with no semisimple summands, if $B$ and $A$ are stably equivalent of Morita type, then $B$ is also a Brauer graph algebra.
\end{Thm}

Since any derived equivalence between selfinjective algebras induces a stable equivalence of Morita type, this gives an alternative proof of the result that Brauer graph algebras are closed under derived equivalence, see Corollary \ref{closed-under-derived-equivalence}.

As a byproduct, Theorem \ref{thm:intro} together with the result \cite[Remark 2.10]{GL} gives a solution of the reconstruction problem (see \cite[Introduction]{RR}) posed by Rickard and Rouquier for algebras stably equivalent of Morita type to Brauer graph algebras.

We note that there is another reason on studying stable equivalences of Brauer graph algebras. For Brauer tree algebras, it is well-known that two such algebras are stably equivalent if and only if they are derived equivalent. Furthermore, it has been shown in \cite{Asa2} (see also \cite{LL2}) that any stable equivalence between Brauer tree algebras can be lifted to a derived equivalence. However, for general Brauer graph algebras, this liftability property fails: \cite[Example~6.4]{LL} provides a stable (auto-)equivalence (even of Morita type) that cannot be lifted to a derived equivalence. Therefore, compared with derived equivalences, some aspects of stable equivalences for Brauer graph algebras are still mysterious.

In this paper, our results on stable equivalences are mainly obtained by analyzing the behavior of modules lying at the mouths of tubes in the Auslander-Reiten quiver. In a subsequent paper \cite{CLLX}, using invariants under stable equivalence of Morita type, such as the stable center and maximal tori, we will classify Brauer graph algebras up to stable equivalence of Morita type, and give a short proof of Theorem~\ref{thm:intro} following some idea from \cite{AZ2}.


\medskip
	\textbf{Outline.}\; In Section \ref{sec:prelim}, we review fundamental concepts and key results concerning stably biserial algebras, Brauer graph algebras, and stable equivalence of Morita type. In Section~\ref{app:lem}, we analyze the stable Hom-spaces from the string modules at the mouth of the homogeneous tube induced by deformed loops in the Brauer graph of a symmetric stably biserial algebra, and establish the key lemma. In Section~\ref{sec:stb-closed}, we prove our main results.
	
\section{Preliminaries}\label{sec:prelim}

	Throughout we assume that $k$ is an algebraically closed field and all algebras considered are finite dimensional $k$-algebras. Unless stated otherwise, all modules will be finitely generated left modules. Furthermore, we say that $A$ is a quiver algebra, if $A$ is isomorphic to $\qa$, where $Q$ is a finite quiver and $I$ is an admissible ideal in the path algebra $kQ$. We denote by $s(p)$ the source vertex of a path $p$ and by $t(p)$ its terminus vertex. We will write paths from right to left, for example, $p=\alpha_{n}\alpha_{n-1}\cdots\alpha_{1}$ is a path with starting arrow $\alpha_{1}$ and ending arrow $\alpha_{n}$. A path is called a cycle (equally, a closed path) if $s(p)=t(p)$.
	By abuse of notation we sometimes view an element in $kQ$ as an element in the quotient $\qa$ if no confusion can arise.

	In this paper, we often write an indecomposable $A$-module $M$ via its Loewy structure, which is represented by a diagram where the $i$-th row corresponds to the simple summands of the completely reducible module $\rad^{i-1}(A)M/\rad^{i}(A)M$ with $\rad(A)$ the Jacobson radical of $A$. Each number in the diagram denotes a distinct simple module in $A$. For further details, see for example in \cite[Page 174]{Ben}.

	Recall (for example from \cite{KZ,Z}) that for an algebra $A$,
	\begin{itemize}
		\item the {bounded derived category $\mathcal{D}^b(A)$} is the localization of homotopy category {$\mathcal{K}^b(A)$} by inverting quasi-isomorphisms, {which is a triangulated category with suspension functor the shift functor $[1]$};
		\item the stable category $A\text{-}\underline{\mathrm{mod}}$ is the quotient of the module category $ A\text{-}{\mathrm{mod}}$ modulo the ideal of homomorphisms that factor through projective modules. For $A$-modules $X$ and $Y$, we denote by $\Hom_A(X, Y)$ (resp. $\underline{\Hom}_A(X, Y)$) the homomorphism space in $A$-mod (resp. $A\text{-}\underline{\mathrm{mod}}$); moreover, for $f\in \Hom_A(X,Y)$, denote by $\underline{f}$ the image of $f$ in $\underline{\Hom}_A(X, Y)$.
	\end{itemize}
	We say that two algebras are derived equivalent (resp. stably equivalent) if their derived (resp. stable) categories are equivalent as triangulated categories (resp. as $k$-categories).

Regarding the precise definition and properties of the stable equivalence of Morita type, we will give them in Section \ref{subsec:stb-and-center}.

	\subsection{Stably biserial algebras}
		\
	
Recall that a $k$-algebra $A$ is called self-injective if $_AA$ is an injective $A$-module; and Frobenius if $_AA\cong DA$ as $A$-modules where $D$ is the usual $k$-duality of $A$; and weakly symmetric if the injective hull of each simple module of $A$ is isomorphic to the projective cover; and symmetric if $_AA_A\cong DA$ as $A$-$A$-bimodules. Indeed, a basic self-injective algebra is a Frobenius algebra, and symmetric algebras (respectively, Frobenius algebras) are naturally self-injective. 

Note that for a self-injective algebra $A$, the syzygy functor $\Omega_{A}$ defines a stable auto-equivalence (even of Morita type) over the stable category
$A\text{-}\underline{\mathrm{mod}}$ and $A\text{-}\underline{\mathrm{mod}}$ is a triangulated category (with suspension functor $\Omega_{A}^{-1}$) which is a natural quotient of $\mathcal{D}^b(A)$ (see for example in \cite{KZ,Z}).

	\begin{definition}\textnormal{(\cite[Definition 1]{AZ1})}
		Let $Q$ be a quiver and $I$ an admissible ideal of the path algebra $kQ$. A self-injective algebra $A$ is said to be stably biserial if $A$ is isomorphic to $kQ/I$, where $Q$ and $I$ satisfy the following conditions:
		\begin{enumerate}
			\item for each vertex $i\in Q$, the number of outgoing and incoming arrows are less than or equal to $2$;
			
			\item for each arrow $\alpha\in Q$, there is at most one arrow $\beta\in Q$ such that $\alpha\beta\notin \alpha\rad(A)\beta+\soc(A)$;
			
			\item for each arrow $\alpha\in Q$, there is at most one arrow $\gamma\in Q$ such that $\gamma\alpha\notin \gamma\rad(A)\alpha+\soc(A)$.
		\end{enumerate}
	\end{definition}

	For other description of stably biserial algebras we refer to \cite[Proposition A.5]{AIP}.

An algebra $A$ is called special biserial if it is isomorphic to an algebra of the form $\qa$ where $kQ$ is a path algebra and $I$ is an admissible ideal such that the following properties hold:
	
	\begin{enumerate}
		\item At every vertex $i$ in $Q$, there are at most two arrows starting at $i$ and there are at most two arrows ending at $i$.
		
		\item For every arrow $\alpha$ in $Q$, there exists at most one arrow $\beta$ such that $\beta\alpha\notin I$ and there exists at most one arrow $\gamma$ such that $\alpha\gamma\notin I$.
	\end{enumerate}

	A special biserial algebra $A$ is called a string algebra if the defining ideal $I$ is generated by paths. Note that if $A$ is a stably biserial algebra or a special biserial algebra, then the quotient algebra $\overline{A}$ by factoring the socles of indecomposable projective $A$-modules which are not uniserial is a string algebra, and in this way, the classification of indecomposable $A$-modules is reduced to the classification of indecomposable $\overline{A}$-modules, see Subsection 2.2.2 for further explanation.

Thus, according to the above definitions, self-injective special biserial algebras form a special subclass of stably biserial algebras. In fact, algebras that are stably equivalent to self-injective special biserial algebras have the following characterization.
	
	\begin{theorem}\textnormal{(\cite[Theorem 1]{AZ1})}\label{thm:sta-to-BGA=StB}
		Let $A$ be {an indecomposable} self-injective special biserial $k$-algebra which is not isomorphic to the Nakayama algebra with $\rad^2(A)=0$. If $B$ is a basic algebra stably equivalent to $A$, then $B$ is stably biserial.
	\end{theorem}

	Note that when a stably biserial algebra $A$ is symmetric, there always exists some nice presentation $kQ/I$ of $A$, see Theorem \ref{thm:sym-StBA} for the details.

The Auslander-Reiten conjecture for stable equivalences states that if two algebras are stably equivalent, then they have the same number of non-projective simple modules up to isomorphism. In particular, \cite[Theorem 0.1]{P} and \cite[Theorem 2]{AZ1} establish its validity for algebras stably equivalent to special biserial algebras.

	\begin{theorem}\label{thm:AR-conj}
		Let $A$ and $B$ be two $k$-algebras such that $ A\text{-}\underline{\mathrm{mod}}\cong B\text{-}\underline{\mathrm{mod}}$ and $A$ is special biserial. Then the numbers of isomorphism classes of non-projective simple modules over $A$ and $B$ are the same.
	\end{theorem}

	
	\subsection{Brauer graph algebras}
	\

	\subsubsection{Basic definitions}\label{subsec:def-BGA}
	\
	
	We recall the basic knowledge about Brauer graph algebras defined in \cite{DF}. For the convenience of the following construction, we adopt the notations from \cite{OZ} to define these algebras.
	
	\begin{Def}\textnormal{(\cite[Definition 1.1]{OZ})}\label{ribbon-graph}
		A ribbon graph is a tuple $\Gamma=(V,H,s,\iota,\rho)$, where
		\begin{enumerate}
			\item $V$ is a finite set whose elements are called vertices;
			
			\item $H$ is a finite set whose elements are called half-edges;
			
			\item $s: H\rightarrow V$ is a function;
			
			\item $\iota: H\rightarrow H$ is an involution without fixed points;
			
			\item $\rho: H\rightarrow H$ is a permutation whose cycles correspond to the sets $H_v:=s^{-1}(v)$, $v\in V$.
			
		\end{enumerate}
	\end{Def}
	
	Therefore, every ribbon graph defines a graph with vertex set $V$ whose edges are the orbits of $\iota$. An edge $\{h, \iota(h)\}$ is incident to the vertices $s(h)$ and $s(\iota(h))$. In fact, every ribbon graph gives rise to an oriented surface (see for example in \cite[Section 1.1]{OZ}).
	
		If $\Gamma$ is a ribbon graph, we write $V(\Gamma)$ and $H(\Gamma)$ for its sets of vertices and half-edges as well as  $H_v(\Gamma)$ for all half-edges $h\in H(\Gamma)$ such that $s(h)=v\in V(\Gamma)$. To simplify the notation, we often write $h^{\pm}:=\rho^{\pm 1}(h)$ for the successor and predecessor of a half-edge $h$ and $\bar{h}$ for its associated edge. The set of all edges is denoted by $E(\Gamma)$. Denote by $val(v)$ the valency of the vertex $v\in V$, it is defined to be the number of edges in $G$ incident to $v$, with the convention that a loop is counted twice.	
	Unless stated otherwise, we will assume that $\Gamma$ is connected, which means that its underlying graph is connected.
	
		\begin{Def}\textnormal{(\cite[Definition 1.5]{OZ})}\label{BG}
		A Brauer graph is a pair $(\Gamma,m)$ consisting of a ribbon graph $\Gamma$ and a function $m: V(\Gamma)\rightarrow \mathbb{Z}_+$.
	\end{Def}
	
	The function $m$ in Definition \ref{BG} is referred to as the multiplicity function and its values as multiplicities. Frequently, we omit $m$ from the notation and refer to $\Gamma$ as a Brauer graph. We denote by $\textbf{n}$ any constant multiplicity function with value $n$ at each vertex and say that a Brauer graph $(\Gamma,m)$ is multiplicity-free if $m=\textbf{1}$. In particular,  we call a given vertex $v$ is truncated if $\m(v)val(v)=1$.
	
	To any Brauer graph $(\Gamma,m)$ one can associate a quiver $Q=Q_\Gamma$ and an admissible ideal of relations $I=I_\Gamma$ in the path algebra $kQ$.
If $\Gamma$ is a single edge with two truncated vertices, then we define $Q$ to be a loop $\alpha$ and $I$ to be the ideal generated by $\alpha^2$. Otherwise, we define $Q$ and $I$ as follows.
	
	\begin{enumerate}
		\item The vertices of $Q$ correspond to the edges of $\Gamma$ and for every $h\in H$ with $s(h)$ not truncated, there is an arrow $\alpha_h:\bar{h}\rightarrow\bar{h^+}$. The assignment $\alpha_h\mapsto \alpha_{h^-}$ defines a permutation $\sigma=\sigma_\Gamma$ of the arrows of $Q$ whose orbits are in bijection with vertices which are not truncated. Hence every arrow $\alpha$ defines a closed path
	$$C_\alpha=\alpha\sigma(\alpha)\cdots\sigma^l(\alpha)$$
	where $l+1$ denotes the cardinality of the $\sigma$-orbit of $\alpha$. Every vertex of $Q$ is the starting point of at most two cycles of the form $C_\alpha$. If $\alpha=\alpha_h$, set $\m(C_\alpha):=\m(s(h))$.
	
	\item The ideal $I_\Gamma$ is generated by the following set of relations:
	\begin{enumerate}
		\item $$C_\alpha^{\m(C_\alpha)}-C_\beta^{\m(C_\beta)},$$
		where $\alpha,\beta\in Q_1$ and $t(\alpha)=t(\beta)$, that is $\alpha$ and $\beta$ end at the same edge of $\Gamma$.
		
		\item $$\alpha C_{\sigma(\alpha)}^{\m(C_\alpha)},$$
		where $\alpha$ is an arbitrary arrow in $Q$.
		
		\item $$\alpha\beta,$$
		where $\alpha,\beta\in Q_1$ are composable and $\sigma(\alpha)\neq \beta$;
	\end{enumerate}
	\end{enumerate}
	
	The resulting (finite dimensional) $k$-algebra $kQ_\Gamma/I_\Gamma$ will be denoted by $A_\Gamma$. Note that every nonzero path of $A_\Gamma$ is a subpath of a cycle $C_\alpha^{\m(C_\alpha)}$ for some arrow $\alpha$. We call a given $\alpha\in Q_1$ is an arrow around $v\in V$ if $s(t(\alpha))=v$ (equivalently, if $s(s(\alpha))=v$).
	
	\begin{Def}\textnormal{(\cite[Definition 1.6]{OZ})}
		A $k$-algebra $A$ is called a Brauer graph algebra if there exists a Brauer graph $(\Gamma,\m)$ such that $A$ is isomorphic to $A_\Gamma$ as $k$-algebras.
	\end{Def}
	
	Indeed, we have the following result.
	
	\begin{Thm}\textnormal{(\cite[Theorem 1.1]{Sch})}\label{ssbBGA}
		Let $A=\qa$ be a $k$-algebra. Then $A$ is a symmetric special biserial algebra if and only if $A$ is a Brauer graph algebra.
	\end{Thm}

Note that since we assume that any Brauer graph $(\Gamma,\m)$ is connected, the resulting Brauer graph algebra $A$ is always indecomposable.

	\subsubsection{Strings and bands}\label{subsec:string-band}
	\

Following \cite{BR,E,D}, we use string and band modules to describe the indecomposable modules over Brauer graph algebras. These correspond to the modules of the first and the second kind in \cite{A}, respectively.

Let $A = kQ/I$ be a Brauer graph algebra associated with a Brauer graph $\Gamma$. Recall that for each arrow $\alpha \in Q_1$, denote by $s(\alpha)$ and $t(\alpha)$ the source and target of $\alpha$, respectively. We denote by $\alpha^{-1}$ the formal inverse of $\alpha$, that is, the symbolic arrow satisfying
\[
s(\alpha^{-1}) = t(\alpha), \qquad t(\alpha^{-1}) = s(\alpha).
\]
Let $Q_1^{-1}$ denote the set of all formal inverses of arrows in $Q_1$.
We set $(\alpha^{-1})^{-1}=\alpha$. If $w=\alpha_n\cdots\alpha_1$ is a word in $Q_1\cup Q_1^{-1}$, then its inverse is the word $w^{-1}=\alpha_1^{-1}\cdots\alpha_n^{-1}$.

\begin{Def}\label{def:string}
A string is a word $w = \alpha_n \alpha_{n-1} \cdots \alpha_1$ with $\alpha_i \in Q_1 \cup Q_1^{-1}$ such that
\begin{enumerate}
    \item $\alpha_{i+1} \neq \alpha_i^{-1}$ for $1\leq i<n$;
    \item $s(\alpha_{i+1}) = t(\alpha_i)$ for $1\leq i<n$; and
    \item neither $w$ nor $w^{-1}$ contains, as a consecutive subword, a path occurring as a term of one of the defining relations of $I$ listed in Subsection~\ref{subsec:def-BGA}.
\end{enumerate}
The length of $w$ is the integer $|w| = n$. We say that $w$ is direct (resp. inverse) if each $\alpha_i \in Q_1$ (resp. $\alpha_i \in Q_1^{-1}$). The associated vertex sequence of $w$ is $(x_0,x_1,\ldots,x_n)$, where $x_0=s(\alpha_1)$ and $x_i=t(\alpha_i)$ for $1\leq i\leq n$. We set $s(w)=x_0$ and $t(w)=x_n$.
\end{Def}

Note that a direct string $\alpha_n \cdots \alpha_1$ can also be viewed as a path in the quiver algebra $kQ/I$. We shall use these two viewpoints interchangeably.

Next, we include strings of length zero, which correspond to stationary paths and serve as identity elements in the combinatorial setting.

\begin{Def}
A string of length zero is called a zero string. For each vertex $x \in Q_0$, we associate a unique zero string $\varepsilon_x$, corresponding to the idempotent at $x$. We set $|\varepsilon_x|=0$, $s(\varepsilon_x)=t(\varepsilon_x)=x$, and $\varepsilon_x^{-1}=\varepsilon_x$; its vertex sequence is $(x)$.
Note that a zero string is considered both direct and inverse.
\end{Def}

If $w=\alpha_n\cdots\alpha_1$ has positive length, a substring of $w$ is a consecutive subword $\alpha_j\cdots\alpha_i$ with $1\leq i\leq j\leq n$. We also regard $\varepsilon_x$ as a substring of a string $w$ whenever $x$ occurs in the vertex sequence of $w$. A string $w$ is called closed if $|w|>0$ and $s(w)=t(w)$.

Strings form the fundamental combinatorial objects used to describe indecomposable modules over special biserial algebras.
Among them, some strings exhibit a cyclic behavior, leading to the following notion.

\begin{Def}\label{def:band}
A band in $A$ is a closed string $b$ such that $b^m$ is a string for every integer $m > 0$, and $b$ is not a proper power of another string.
\end{Def}

\begin{Def}\label{def:string-module}
Let $w = \alpha_n \alpha_{n-1} \cdots \alpha_1$ be a string in $A$,
where each $\alpha_i \in Q_1 \cup Q_1^{-1}$. Let $(x_0,x_1,\ldots,x_n)$ be the associated vertex sequence of $w$. The string module $M=M(w)$ associated to $w$ is defined as follows.

For each vertex $x \in Q_0$, set
\[
M_x = \bigoplus_{i \, : \, x_i = x} k e_i,
\]
that is, $M_x$ is a direct sum of one copy of the field $k$ for each occurrence of $x$ in the sequence $x_0, x_1, \dots, x_n$.

For each $\alpha\in Q_1$,
define the linear map
$$M_{\alpha} : M_{s(\alpha)}=\bigoplus_{i \, : \, x_i = s(\alpha)} k e_i \to M_{t(\alpha)}=\bigoplus_{j \, : \, x_j = t(\alpha)} k e_j$$
by
\begin{equation*}
M_{\alpha}(e_i)=\begin{cases}
e_{i+1}, & \text{ if } \alpha_{i+1}=\alpha; \\
e_{i-1}, & \text{ if } \alpha_i=\alpha^{-1}; \\
0, & \text{ otherwise}.
\end{cases}
\end{equation*}

This defines a representation of $Q$ satisfying the relations in $I$, hence an $A$-module.

If $w = \varepsilon_x$ is the zero string at vertex $x$,
we define $M(w) := S_x$, the simple $A$-module at $x$.
\end{Def}

Note that each string module is indecomposable (see also Section 3). It follows from \cite{BR} that two string modules $M(w_1)$ and $M(w_2)$ are isomorphic if and only if $w_1 = w_2$ or $w_1 = w_2^{-1}$.

Now we recall that there is also the notion of a band module. Each distinct band determines an infinite family of band modules, which can be viewed as indecomposable modules parametrized by a nonzero scalar and a positive integer.

\begin{Def}\label{def:band-module}
Let $b = \beta_n \beta_{n-1} \cdots \beta_1$ be a band in $A$. Fix a positive integer $m \geq 1$ and a nonzero scalar $\lambda \in k^*$. Write $x_0 = s(\beta_1)$ and $x_i = t(\beta_i)$ for $1 \leq i < n$.  The band module associated with $(b, m, \lambda)$, denoted by $M(b, m, \lambda)$, is defined as follows.

For each vertex $x \in Q_0$, define
\[
M_x = \bigoplus_{i \, : \, x_i = x} k^m e_i,
\]
that is, $M_x$ is a direct sum of one copy of $k^m$ for each occurrence of $x$ in the sequence $x_0, x_1, \dots, x_{n-1}$.

For each $1\leq i\leq n$ let $\gamma_i\in Q_1$ denote the underlying
arrow of the symbol $\beta_i$ (so $\beta_i=\gamma_i$ or $\beta_i=\gamma_i^{-1}$). We now define the linear map $M_\gamma:M_{s(\gamma)}\to M_{t(\gamma)}$ for each $\gamma\in Q_1$:

\begin{enumerate}
	\item For $1\leq i<n$ define:
\begin{itemize}
  \item if $\beta_i=\gamma_i$, then
  \[
    M_{\gamma_i}(v e_{i-1}) = v e_i \quad\text{for all }v\in k^m,
  \]
  and $M_{\gamma_i}$ sends other summands to $0$;
  \item if $\beta_i=\gamma_i^{-1}$, then
  \[
    M_{\gamma_i}(v e_i) = v e_{i-1} \quad\text{for all }v\in k^m,
  \]
  and $M_{\gamma_i}$ sends other summands to $0$.
\end{itemize}
\item For $i=n$, we insert the Jordan block $J_m(\lambda)$:
\begin{itemize}
  \item If $\beta_n=\gamma_n$, define
  \[
    M_{\gamma_n}(v e_{n-1}) = J_m(\lambda)\,v e_0 \quad\text{for all }v\in k^m,
  \]
  and $M_{\gamma_n}$ sends other summands to $0$;
  \item If $\beta_n=\gamma_n^{-1}$, define
  \[
    M_{\gamma_n}(v e_0) = J_m(\lambda)^{-1}\,v e_{n-1} \quad\text{for all }v\in k^m
  \]
  (note that $J_m(\lambda)$ is invertible since $\lambda\neq 0$), and $M_{\gamma_n}$ sends other summands to $0$.
\end{itemize}
\end{enumerate}
All arrow maps $M_\gamma$ not prescribed above are defined to be the zero map.
One checks that these maps satisfy the relations in $I$, hence define an $A$-module,
which we denote by $M(b,m,\lambda)$.
\end{Def}

Note that each band $b$ gives rise to the family $\{M(b,m,\lambda)\mid m\ge1,\ \lambda\in k^*\}$ of indecomposable modules, and that two band modules $M(b,m,\lambda)$ and $M(b',m',\lambda')$ are isomorphic if and only if $m' = m$, $\lambda' =\lambda$, and $b'=b$ or $b'=b^{-1}$ up to taking rotations.

We note that the indecomposable projective $A$-modules which are not uniserial cannot be represented as string or band modules, because the corresponding word contains a path lying in $\soc(A)$, which cannot form a string: it violates condition~(3) of Definition~\ref{def:string}. Equivalently, the string and band $A$-modules are in fact defined over the string algebra $\overline{A}$ which is defined as the quotient algebra of $A$ by factoring the socles of indecomposable projective  $A$-modules which are not uniserial. The same applies to symmetric stably biserial algebras. Indeed, if $A$ is a symmetric stably biserial algebra, then the quotient algebra $\overline{A}$ defined by factoring the socles of indecomposable projective  $A$-modules which are not uniserial is also a string algebra, and therefore all the indecomposable $A$-modules except the indecomposable projective $A$-modules which are not uniserial can also be represented as string or band $A$-modules.

We also note that the notions of string and band depend on the choice of presentation $kQ/I$ of the algebra. For instance, in Example~\ref{ex-stably-biserial}, we exhibit different quiver presentations under which a band module may be changed to a string module.

	\subsubsection{Exceptional band modules}\label{subsec:excep-band}
	\

We now turn to a more detailed study of band modules over Brauer graph algebras. The following notion is taken from \cite[Section 2]{A}.

\begin{Def}\label{def:max-string module}
Let $A=kQ/I$ be a Brauer graph algebra. A path $p$ from $v$ to $w$ of $Q$ with positive length is called a band line if $p\notin \mathrm{soc}(A)$ and the simple modules $S_v$, $S_w$ of $A$ corresponding to the vertices $v,w$ are not $\Omega_{A}$-periodic. Moreover, $p$ is called a minimal (resp. maximal) band line if it is a minimal (resp. maximal) element in the partial ordered set consisting of all band lines.
\end{Def}

For a Brauer graph $\Gamma$, define $\Gamma^{(1)}$ to be the ribbon graph obtained from $\Gamma$ by deleting all leaves (which correspond to the edges connected with some truncated vertices), and define $\Gamma^{(n)}:=(\Gamma^{(n-1)})^{(1)}$ for every positive integer $n$ inductively. Then there exists some positive integer $r$ such that $\Gamma^{(r)}=\Gamma^{(r+1)}=\cdots$. Denote $\widetilde{\Gamma}=\Gamma^{(r)}$. Moreover, consider the quivers corresponding to these ribbon graphs, it is clear to see that the vertex set of $Q_{\widetilde{\Gamma}}$ is a subset of the vertex set of $Q_{\Gamma}$. For example, $\widetilde{\Gamma_1}=\Gamma_2$ in Example \ref{exa:2-ribbon graph}.

To construct the band lines in a Brauer graph algebra $A$, we use the following lemmas from \cite[Section 2]{A}, which provide criteria for the $\Omega_A$-periodicity of simple modules.


\begin{Lem}\label{lem:paths-in-band=band-line}
Let $X$ be a band module of $A$. Then each simple $A$-module which is a summand of top$(X)$ or soc$(X)$ is not $\Omega_{A}$-periodic.
\end{Lem}

\begin{Lem}
Suppose that $A$ is given by the Brauer graph $(\Gamma,m)$. Then there is a one-to-one correspondence between the isomorphism classes of simple $A$-modules which are not $\Omega_{A}$-periodic and the set of edges of $\widetilde{\Gamma}$.
\end{Lem}

Note that a minimal band line $p$ in $A$ determines an arrow $\alpha_p$ with $s(\alpha_p)=s(p)$ and $t(\alpha_p)=t(p)$ in $Q_{\widetilde{\Gamma}}$, and this correspondence is a bijection.

The following lemma provides us a method to construct a band from a string whose ending vertices corresponding to simple modules which are not $\Omega$-periodic. 

\begin{Lem}\label{lem-string-to-band} {\rm(\cite[Lemma 2.5]{A})}
Let $A$ be a Brauer graph algebra, and let
\[
w = p_l p_{l-1} \cdots p_1
\]
be a string in $A$, where each $p_i$ is either a direct string or an inverse string, appearing alternately, and both $s(w)$ and $t(w)$ correspond to simple $A$-modules that are not $\Omega_A$-periodic. Then there exists a band
\[
b = p_m p_{m-1} \cdots p_l \cdots p_1
\]
in $A$, where again each $p_i$ is either a direct string or an inverse string, appearing alternately.
\end{Lem}

Now we define the exceptional band module in \cite{A}.

\begin{Def}\label{def:exceptional-band}
A band $b$ of $A$ is called exceptional if it can be expressed up to taking rotations and inverses in the form $$b = p_{2k} p_{2k-1}^{-1} \cdots p_{2}p_{1}^{-1},$$ where each $p_1, p_3, \ldots, p_{2k-1}$ is a minimal band line and each $p_2, p_4, \ldots, p_{2k}$ is a maximal band line.

A band module $X$ of $A$ is called exceptional if it is defined by an exceptional band $b$, with the vector space at each vertex of $b$ being the one-dimensional space $k$, and all but one of the arrows act as identity maps, while one arrow acts as multiplication by some scalar $\lambda\in k^{*}$, that is, $X{\cong} M(b,1,\lambda)$.

If $b$ is a non-exceptional band, then we call a band module of the form $M(b,1,\lambda)$ a non-exceptional band module.
\end{Def}

The construction of the following string module is given in \cite[Proposition 2.3]{A}. We provide a concrete proof to show that the string module constructed in \cite[Proposition 2.3]{A} in fact does not lie in a tube of rank $1$ of the Auslander-Reiten quiver.

\begin{proposition}\label{prop:non-exceptional-tube}
Let $A = kQ/I$ be a Brauer graph algebra, and let $M$ be a non-exceptional band module. For any $n \in \mathbb{Z}_+$, there exists a string module $Y$ that does not lie in a tube of rank $1$ of the Auslander-Reiten quiver of $A$, such that
\[
\mathrm{dim}_k \underline{\Hom}_A(M, Y) > n.
\]
\end{proposition}

\begin{proof}
Suppose $M$ is given by the band
\[
p_{2m} p_{2m-1}^{-1} \cdots p_{2} p_{1}^{-1}
\]
of $A$. Then by Lemma \ref{lem:paths-in-band=band-line}, each $p_i$ is a band line. Since $M$ is not a exceptional band module, without loss of generality, at least one of the following possibilities occurs:

\begin{itemize}
\item[(1)] $p_1$ and $p_{2k}$ are not minimal band lines.
\item[(2)] $p_1$ and $p_{2k}$ are not maximal band lines.
\item[(3)] $p_1$ is not a minimal band line, and $p_{2k+1}$ is not a maximal band line (they may be the same line).
\end{itemize}

The third case, in contrast to the first and second, arises for band modules in which all even band lines are either simultaneously maximal or simultaneously minimal.
Note that this proposition is invariant under the $\Omega_A$-shift, and since the maximal band lines are mapped to minimal band lines under the $\Omega_A$-shift, we do not need to consider case (2).

Consider case (1). Denote by $q_1$ (resp. $q_2$) the minimal band line that contained in $p_1$ (resp. $p_{2k}$) and has a common starting arrow with $p_1$ (resp. $p_{2k}$). Then there exists an epimorphism $f : M \to Z$, where $Z$ is a string module which is given by $$w_1=q_2p_{2k-1}^{-1}\cdots p_{2}q_1^{-1}.$$

Since the simple modules corresponding to $s(w_1)$ and $t(w_1)$ are not $\Omega_A$-periodic, it follows from Lemma \ref{lem-string-to-band} that there exists a closed string $w = w_2w_1$ of $A$, where $w_2$ starts with a formal inverse in $Q_1^{-1}$ and ends with an arrow in $Q_1$.
For any natural number $n$, consider the string module $Y$ corresponding to the string
$$w_1w^n$$
It is easy to see that a monomorphism $h : Z \to Y$ and a morphism $hf : M \to Y$ correspond to any occurrence of the substring $Z$ in $Y$. Since $hf$ induces a nonzero map $\operatorname{top}(M) \to \operatorname{top}(Y)$ (the image of $p_i$ is nonzero), it is nonzero in the stable category.

We now show that such a module $Y$ does not lie in a tube of rank $1$ of the stable Auslander-Reiten quiver of $A$. It suffices to prove that $\tau(Y)\not\cong Y$, where $\tau$ is the Auslander-Reiten translation of $A$. First, note that $q_1$ and $q_2$ are not maximal direct strings, since there exist paths $p_1$ and $p_{2k}$ containing them that are not in the socle of $A$. Hence, $\tau(Y)$ is a string module obtained by attaching two cohooks to each side of $w_1 w^n$ (see for example in \cite{BR}), which is given by
\[
p''^{-1}\beta w_1 w^n \alpha^{-1}p',
\]
where $p'$ and $p''$ are maximal direct strings and $\alpha, \beta \in Q_1$. Therefore, $\tau(Y)\not\cong Y$.

Now consider case (3).

Let $k \neq 0$. Let $q_1$ denote the minimal band line sharing a common starting arrow with $p_1$, and let $q_2$ denote the maximal band line sharing a common ending arrow with $p_{2k+1}$. Then there exists a morphism from $M$ to an arbitrary string module that contains the substring
\[
w_1 = q_2^{-1}p_{2k}\cdots p_2 q_1^{-1}.
\]
As in the previous case, there exists a string module $Y$ of the form
\[
w_1w^n,
\]
where $w$ is a closed string containing $w_1$, such that
\[
\mathrm{dim}_k \underline{\Hom}_A(M, Y) > n.
\]
Note that $q_1$ is not a maximal direct string, since there exists a path $p_1$ containing it that are not in the socle of $A$. Thus, $\tau(Y)$ is a string module obtained by attaching two cohooks to each side of $w_1w^n$, which will not isomorphic to $Y$.

Let $k = 0$, and let $q_1 = p_1q \in \soc(A)$. Note that $q$ is also a band line, and since $p_1$ is neither a maximal nor a minimal band line, $q$ is also neither a maximal nor a minimal band line. It follows that we have the following decompositions
\[
p_1 = p' p'', \quad q = q' q'',
\]
where $p', p'', q', q''$ are band lines. So we get a new band line $q_2 = p'' q'$.
For any string module $Z$ containing $q_2$ as a substring, there exists a morphism $f : M \to Z$. We claim that $f$ is nonzero in the stable category.
Indeed, suppose to the contrary that $f$ is zero in the stable category. Then there exists a morphism $g : I(M) \to Z$ such that $f = g i$, where $I(M)$ denotes the injective hull of $M$ and $i : M \to I(M)$ is the canonical embedding. However, such a morphism $g$ cannot simultaneously map the top of $I(M)$ (corresponding to the simple module associated with $s(q_1)$) to a composition factor corresponding to a vertex in $q'$, and at the same time map the simple module corresponding to $s(p_1)$ to that corresponding to $s(p'')$. This yields a contradiction.

As in the previous case, there exists a string module $Y$ of the form
\[
q_2^{-1}w^n,
\]
where $w$ is a closed string containing $q_2^{-1}$, such that
\[
\mathrm{dim}_k \underline{\Hom}_A(M, Y) > n.
\]
Note that $q_2$ is not a maximal direct string, since there exists a path $p'q_2$ containing it that are not in the socle of $A$. Thus, $\tau(Y)$ is a string module obtained by attaching two cohooks to each side of $q_2^{-1}w^n$, which will not isomorphic to $Y$.
\end{proof}

\subsubsection{Faces and double-faces}
\

 We now recall the connection between exceptional band modules (which corresponds to exceptional modules of second kind  in \cite{A}) in a Brauer graph algebra $A$ and double-faces in its Brauer graph $\Gamma_A$.
We begin by recalling the definitions of faces and double-faces in a Brauer graph in \cite{OZ}, which correspond to Green walks in \cite{GreenJA} (or $G$-cycles in \cite{A}) and double-stepped Green walks in \cite{D} (or $G^2$-cycles in \cite{A}), respectively.

\begin{Def}
Let $\Gamma$ be a ribbon graph. A face of perimeter $m$ of $(\Gamma,m)$ is an equivalence class of primitive cyclic sequences $F=(h_1,h_2,\cdots,h_m)$, where $h_1,\cdots,h_m\in H(\Gamma)$ such that $h_{i+1}=\iota\rho(h_{i})$ for all $i$. We call a sequence $F$ primitive, if $F$ is not a power of another sequence of this form. Two sequences of such kind are considered equivalent if they agree after a cyclic permutation of their entries.
\end{Def}

\begin{Def}
Let $\Gamma$ be a ribbon graph. A double-face of perimeter $m$ of $(\Gamma,m)$ is an equivalence class of primitive cyclic sequences $F=(h_1,h_2,\cdots,h_m)$, where $h_1,\cdots,h_m\in H(\Gamma)$ such that $h_{i+1}=(\iota\rho)^{2}(h_{i})$ for all $i$. Two sequences of such kind are considered equivalent if they agree after a cyclic permutation of their entries.
\end{Def}

So, in particular, the faces give a partition on the set $H(\Gamma)$ of half-edges and the double-faces also give a partition on $H(\Gamma)$. Note that both notions are free of the multiplicity $m$.
We give the following example to help readers understand these concepts.

\begin{example}\label{exa:2-ribbon graph}
	Let $\Gamma_1$ and $\Gamma_2$ be the Brauer graphs with multiplicity one and the ribbon graphs are given as follows (the orientation around each vertex is clockwise), respectively.
	\begin{center}
\tikzset{every picture/.style={line width=0.75pt}} %set default line width to 0.75pt

\begin{tikzpicture}[x=0.75pt,y=0.75pt,yscale=-1,xscale=1]
%uncomment if require: \path (0,235); %set diagram left start at 0, and has height of 235

%Shape: Ellipse [id:dp5351628263682573]
\draw   (61.1,94.92) .. controls (61.1,76.46) and (75.68,61.5) .. (93.67,61.5) .. controls (111.67,61.5) and (126.25,76.46) .. (126.25,94.92) .. controls (126.25,113.37) and (111.67,128.33) .. (93.67,128.33) .. controls (75.68,128.33) and (61.1,113.37) .. (61.1,94.92) -- cycle ;
%Shape: Ellipse [id:dp22668671113333483]
\draw   (228.86,94.92) .. controls (228.86,76.46) and (243.45,61.5) .. (261.44,61.5) .. controls (279.43,61.5) and (294.02,76.46) .. (294.02,94.92) .. controls (294.02,113.37) and (279.43,128.33) .. (261.44,128.33) .. controls (243.45,128.33) and (228.86,113.37) .. (228.86,94.92) -- cycle ;
%Straight Lines [id:da3132565897170516]
\draw    (126.25,94.92) -- (171.86,94.92) ;
%Shape: Circle [id:dp7255337491938683]
\draw  [fill={rgb, 255:red, 0; green, 0; blue, 0 }  ,fill opacity=1 ] (169.36,94.92) .. controls (169.36,93.54) and (170.48,92.42) .. (171.86,92.42) .. controls (173.24,92.42) and (174.36,93.54) .. (174.36,94.92) .. controls (174.36,96.3) and (173.24,97.42) .. (171.86,97.42) .. controls (170.48,97.42) and (169.36,96.3) .. (169.36,94.92) -- cycle ;
%Shape: Circle [id:dp9224289625323931]
\draw  [fill={rgb, 255:red, 0; green, 0; blue, 0 }  ,fill opacity=1 ] (123.75,94.92) .. controls (123.75,93.54) and (124.87,92.42) .. (126.25,92.42) .. controls (127.63,92.42) and (128.75,93.54) .. (128.75,94.92) .. controls (128.75,96.3) and (127.63,97.42) .. (126.25,97.42) .. controls (124.87,97.42) and (123.75,96.3) .. (123.75,94.92) -- cycle ;
%Shape: Circle [id:dp8545479082638043]
\draw  [fill={rgb, 255:red, 0; green, 0; blue, 0 }  ,fill opacity=1 ] (58.6,94.92) .. controls (58.6,93.54) and (59.72,92.42) .. (61.1,92.42) .. controls (62.48,92.42) and (63.6,93.54) .. (63.6,94.92) .. controls (63.6,96.3) and (62.48,97.42) .. (61.1,97.42) .. controls (59.72,97.42) and (58.6,96.3) .. (58.6,94.92) -- cycle ;
%Shape: Circle [id:dp29041330982278835]
\draw  [fill={rgb, 255:red, 0; green, 0; blue, 0 }  ,fill opacity=1 ] (226.36,94.92) .. controls (226.36,93.54) and (227.48,92.42) .. (228.86,92.42) .. controls (230.24,92.42) and (231.36,93.54) .. (231.36,94.92) .. controls (231.36,96.3) and (230.24,97.42) .. (228.86,97.42) .. controls (227.48,97.42) and (226.36,96.3) .. (226.36,94.92) -- cycle ;
%Shape: Circle [id:dp4212398147160299]
\draw  [fill={rgb, 255:red, 0; green, 0; blue, 0 }  ,fill opacity=1 ] (291.52,94.92) .. controls (291.52,93.54) and (292.63,92.42) .. (294.02,92.42) .. controls (295.4,92.42) and (296.52,93.54) .. (296.52,94.92) .. controls (296.52,96.3) and (295.4,97.42) .. (294.02,97.42) .. controls (292.63,97.42) and (291.52,96.3) .. (291.52,94.92) -- cycle ;

% Text Node
\draw (48.34,64.2) node [anchor=north west][inner sep=0.75pt]   [align=left] {$h_1$};
% Text Node
\draw (130,78.23) node [anchor=north west][inner sep=0.75pt]   [align=left] {$h_3$};
% Text Node
\draw (122,106.3) node [anchor=north west][inner sep=0.75pt]   [align=left] {$h_2$};
% Text Node
\draw (215,65.53) node [anchor=north west][inner sep=0.75pt]   [align=left] {$h_1'$};
% Text Node
\draw (292.66,100.29) node [anchor=north west][inner sep=0.75pt]   [align=left] {$h_2'$};
\end{tikzpicture}
	\end{center}

The graph $\Gamma_1$ contains two faces: one of perimeter $2$, $\{h_1,h_2\}$, and another one of perimeter $4$, $\{\iota(h_1),\iota(h_3),h_3,\iota(h_2)\}$. From these, we obtain four double-faces: $\{h_1\}$, $\{h_2\}$, $\{\iota(h_1),h_3\}$, and $\{\iota(h_3),\iota(h_2)\}$.

In contrast, $\Gamma_2$ has two faces, both of perimeter 2: $\{h_1',h_2'\}$ and $\{\iota(h_1'),\iota(h_2')\}$. From these, we obtain four double-faces: $\{h_1'\}$, $\{h_2'\}$, $\{\iota(h_1')\}$, and $\{\iota(h_2')\}$.
\end{example}

As the example above illustrates, a face of even perimeter gives rise to two distinct double-faces. It is worth noting that a face with an odd perimeter, however, contributes to exactly one.

The following proposition shows that every double-face $F = (h_1, h_2, \ldots, h_k)$ in the graph $\widetilde{\Gamma}$ determines an exceptional band $b = p_{2k}^{-1} p_{2k-1} \cdots p_{2}^{-1} p_{1}$ of $A$ such that  $p_1, p_3, \ldots, p_{2k-1}$ are minimal band lines in $A$ corresponding to the arrows $\alpha_{h_1}, \alpha_{h_2}, \ldots, \alpha_{h_k}$ in the quiver $Q_{\widetilde{\Gamma}}$.

\begin{Prop}\textnormal{(\cite[Lemma 2.7]{A})}\label{prop:double-faces and exceptional bands}
Let $A$ be a Brauer graph algebra with associated Brauer graph $\Gamma$. Let $M$ be an exceptional band $A$-module defined by an exceptional band
\[
b = p_{2k}^{-1} p_{2k-1} \cdots p_{2}^{-1} p_{1},
\]
where $p_1, p_3, \ldots, p_{2k-1}$ are minimal band lines. Suppose that for each $1 \leq i \leq k$, the path $p_{2i-1}$ goes from vertex $v_i$ to vertex $w_i$ in the quiver $Q_\Gamma$.

Then the following statements hold.

\begin{enumerate}
    \item There exists a double-face $F = (h_1, h_2, \ldots, h_k)$ in $\widetilde{\Gamma}$ such that for each $i$, the arrow $\alpha_{h_i}$ in the quiver $Q_{\widetilde{\Gamma}}$ is also from $v_i$ to $w_i$, matching the endpoints of $p_{2i-1}$.

    \item The syzygy $\Omega_A(M)$ is also an exceptional band module. After applying a suitable cyclic permutation to the band defining $\Omega_A(M)$, $\Omega_A(M)$ can be defined by an exceptional band
    \[
	b' = q_{2k}^{-1} q_{2k-1} \cdots q_{2}^{-1} q_{1},
    \]
    where $q_1, q_3, \ldots, q_{2k-1}$ are minimal band lines from $v'_i$ to $w'_i$ in $Q_\Gamma$ and which correspond to a double-face $F' = (h_1', h_2', \ldots, h_k')$ in $\widetilde{\Gamma}$.
    \item Furthermore, either $F=F'$ is a face in $\widetilde{\Gamma}$ or $G:= (h_1, h_1', \ldots, h_{k},h_{k}')$ is a face in $\widetilde{\Gamma}$ such that for each $1 \leq i \leq k$:
    \begin{itemize}
        \item The arrow $\alpha_{h_{i}}$ in $Q_{\widetilde{\Gamma}}$ is from $v_i$ to $w_i$, matching the endpoints of the minimal band line $p_{2i-1}$ in $b$.
        \item The arrow $\alpha_{h_{i}'}$ in $Q_{\widetilde{\Gamma}}$ is from $v_i'$ to $w_i'$, matching the endpoints of the minimal band line $q_{2i-1}$ in $b'$.
    \end{itemize}
\end{enumerate}
\end{Prop}

We end this subsection by the following example.

\begin{example}\label{ex:two-exceptional-band} (Example \ref{exa:2-ribbon graph} revisited)
	Denote the vertices in $Q_{\Gamma_1}$ corresponding to $\bar{h_1}, \bar{h_2},\bar{h_3}$ by $v_1,v_2,v_3$ and the vertices in $Q_{\Gamma_2}$ corresponding to $\bar{h_1'}, \bar{h_2'}$ by $v_1,v_2$. Then the quivers associated to the ribbon graphs $\Gamma_1$ and $\Gamma_2$ are given as follows.
	
	$$% https://tikzcd.yichuanshen.de/#N4Igdg9gJgpgziAXAbVABwnAlgFyxMJZABgBoBGAXVJADcBDAGwFcYkQBFAfWAB1eA4vQC2w+l3IBfRCEml0mXPkIoArBWp0mrdtz6CRYrgCZps+SAzY8BIuQ00GLNjPLmF15UQDMDrc-Zjd0tFGxVkY1JiTScdGW9gqyVbFAA2P1iXEDc5D2TwgA4M7SygyU0YKABzeCJQADMAJwhhJEiQHAgkXxA4AAssepw2mgAjGDAoJDJ-OJB+JjQ+8WA+iUkQGkZ6ccYABVCvGSwwbFhgppakdQ6uxHTegaHuxxL2BcYllbXyAHINrY7GD7Q4pEAnM5sXIgS6te40TrXGj9QbDRAAWgemXevEWyx4-HwOHoAAo1sZfgBKAEgba7A6eMGMGDPaGwl63EaPVFIdE3bEyD5fAm8Imk8nUza0oEgxkqWks4Zs5pw9qIxAAFleAUFuM++P0YrJEklyqumoRdx6AvmeuFqy43g2ZrhPXV7RRz0Q9ml9NB8saWCqfWG2rmQoN5OdFnZ8M5iBunrRPTpwIZ+XYEKw5zDWQj3xM-wuKqR8YeSd5D1TsozMmZz1zOLxK0JEGJxr+psokiAA
\begin{tikzcd}
              &                                                                                           & 3 \arrow[rd, "\alpha_{h_3}"] &                                                                                            &  &               &                                                                                                           &  &                                                                                                           \\
Q_{\Gamma_1}: & 1 \arrow[rr, "\alpha_{h_1}" description, shift right=2] \arrow[ru, "\alpha_{\iota(h_1)}"] &                              & 2 \arrow[ll, "\alpha_{\iota(h_2)}", shift left=5] \arrow[ll, "\alpha_{h_2}"', shift right] &  & Q_{\Gamma_2}: & 1 \arrow[rr, "\alpha_{h_1'}" description, shift right=3] \arrow[rr, "\alpha_{\iota(h_1')}", shift left=6] &  & 2 \arrow[ll, "\alpha_{\iota(h_2')}", shift left=6] \arrow[ll, "\alpha_{h_2'}" description, shift right=3]
\end{tikzcd}$$
	
	Let $M$ be an exceptional band $A_{\Gamma_1}$-module defined by the exceptional band $(\alpha_{h_3}\alpha_{\iota(h_1)})^{-1}\alpha_{h_1}$, which can be represented by
	\begin{center}
\begin{tikzcd}
v_1 \arrow[rd, "\alpha_{\iota(h_1)}", no head] \arrow[dd, "\alpha_{h_1}"', no head] &                                         \\
                                                                                    & v_3 \arrow[ld, "\alpha_{h_3}", no head] \\
v_2                                                                                 &
\end{tikzcd}
	\end{center}
	Therefore, $\alpha_{h_1}$ is a minimal band line in $M$, which correspond to the double-face $\{h_1'\}$ in $\widetilde{\Gamma_1}=\Gamma_2$, and $\alpha_{h_3}\alpha_{\iota(h_1)}$ is a maximal band line of $A_{\Gamma_1}$. Note that $\alpha_{h_3}\alpha_{\iota(h_1)}$ is also a minimal band line corresponding to the double-face $\{\iota(h_1')\}$.

	At the same time, $\Omega_{A_{\Gamma_1}}(M)$ can be defined by the exceptional band $\alpha_{\iota(h_2)}^{-1}\alpha_{h_2}$, which can be represented by
	\begin{center}
\begin{tikzcd}
v_2 \arrow[dd, "\alpha_{h_2}"', no head, bend right] \arrow[dd, "\alpha_{\iota(h_2)}", no head, bend left] \\
                                                                                                           \\
v_1
\end{tikzcd}
	\end{center}
Therefore, $\alpha_{h_2}$ is a minimal band line in $\Omega_{A_{\Gamma_1}}(M)$, which correspond to the double-face $\{h_2'\}$ in $\widetilde{\Gamma_1}=\Gamma_2$. Moreover, the face $\{h_1',h_2'\}$ corresponds to the minimal band lines of $M$ and $\Omega_{A_{\Gamma_1}}(M)$.
\end{example}

\subsection{Symmetric stably biserial algebras}
\

Recall that we have introduced the stably biserial algebras in Section 2.1. We now concentrate on the symmetric stably biserial algebras.	
Note that every symmetric stably biserial algebra can be represented by a Brauer graph $(\Gamma, m)$ together with a special set $\mathcal{L}$ of loops in $Q_\Gamma$. This result was proved in \cite[Section 5]{AZ1}, and a clearer exposition can be found in \cite[Theorem 2.4]{AZ2}.

\begin{Thm}\label{thm:sym-StBA}
	Any symmetric stably biserial algebra $A$ has a presentation $kQ/I$, where $Q=Q_\Gamma$ is a quiver associated with a Brauer graph $(\Gamma,m)$, and the ideal of relations $I$ is generated by
	\begin{enumerate}
		\item $\alpha\beta=0$ for all $\alpha,\beta\in Q_1$, $\beta\neq \sigma(\alpha)$ and $\alpha\notin\mathcal{L}$;
		\item $C_\alpha^{m(C_\alpha)}=C_\beta^{m(C_\beta)}$ for all $\alpha,\beta\in Q_1$ with $s(\alpha)=s(\beta)$;
		\item $\alpha_i^2=t_{\alpha_i} C_{\alpha_i}^{m(C_{\alpha_i})}$ for each $\alpha_i\in\mathcal{L}$;
		\item $C_\alpha^{m(C_\alpha)}\beta=0$ for all $\alpha,\beta\in Q_1$,
	\end{enumerate}
	where $\mathcal{L}=\{\alpha_1,\cdots,\alpha_n\}\subseteq Q_1$ such that each $\alpha_i$ is a loop with $\sigma(\alpha_i)\neq \alpha_i$, and $t_{\alpha_i}\in k^*$.

	Moreover, when $\mathrm{char}(k)\neq 2$, any symmetric stably biserial algebra is isomorphic to an algebra $\qa$ as above with $\mathcal{L}=\varnothing$, in which case it is a Brauer graph algebra.
\end{Thm}

Compare to the defining relations of a Brauer graph algebra, the above relations $(3)$ are new. Observe that the fact that $Q_\Gamma$ contains deformed loops implies that $\Gamma$ contains loops. Note also that the associated string algebra of the above symmetric stably biserial algebra $A=kQ/I$ and the associated string algebra of the Brauer graph algebra defined by $(\Gamma,m)$ are the same (cf. Subsection 2.2.2).

In the remaining part of this paper, when we talk about a symmetric stably biserial algebra $A$ defined by $kQ/I$, we always assume that this presentation is given by the form in Theorem \ref{thm:sym-StBA}, that is, we assume that $Q=Q_{\Gamma}$ and $I$ is generated by the relations $(1)-(4)$.

Here we provide an example of a symmetric stably biserial algebra which is not isomorphic to a special biserial algebra.

\begin{Ex1}\label{ex-stably-biserial} {\rm(cf. \cite[Example A.7]{AIP})}
Let $k$ be a field of characteristic $2$ and let $A$ be a $k$-algebra given by the quiver
$$\xymatrix{
1\ar@/^/[r]^{\alpha}\ar@(dl,ul)^{\gamma} & 2\ar@/^/[l]^{\beta} \\
 }$$
with relations $$\gamma\beta\alpha=\beta\alpha\gamma=\gamma^2, \quad \gamma^3=\alpha\gamma^2=\beta\alpha\gamma\beta=\alpha\beta=0.$$
Then $A$ is a symmetric stably biserial algebra (with linear form $\varphi: A\to k$ by
$\varphi(\gamma\beta\alpha)=1$, $\varphi(\alpha\gamma\beta)=1$ and $\varphi(p)=0$ for each path $p\notin \mathrm{Soc}A$) but not isomorphic to any symmetric special biserial algebra. Note that $A$ is given by the multiplicity-free Brauer graph
$$\begin{tikzpicture}
\draw (0,0) circle (0.5);
\fill (0.5,0) circle (0.5ex);
\node at(-0.7,0){$1$};
\node at(1.0,0.3){$2$};
\draw (0.5,0)--(1.5,0);
\fill (1.5,0) circle (0.5ex);
\end{tikzpicture}$$
with $\mathcal{L}=\{\gamma\}$ and $t_{\gamma}=1$.
We have the following structure of the indecomposable projective $A$-modules:
\xymatrix@R=0.5pc@C=0.8pc
{&&1\ar@{-}[dl]\ar@{-}[dr]&& \\
 P_1= & 1 \ar@{-}[d]\ar@{--}[ddr]& & 2\ar@{-}[d] & \\
& 2 \ar@{-}[dr] &  & 1 \ar@{-}[dl] & \\
&& 1 & ,& }
\xymatrix@R=0.5pc@C=0.8pc
{&2\ar@{-}[d]& \\
 P_2= & 1\ar@{-}[d] & \\
& 1 \ar@{-}[d] & \\
& 2 & .}

Moreover, let $B$ be the Brauer graph algebra whose Brauer graph is the same as the one of $A$. Then $B$ is given by the quiver
$$\xymatrix{
1\ar@/^/[r]^{\alpha}\ar@(dl,ul)^{\gamma} & 2\ar@/^/[l]^{\beta} \\
 }$$
with relations $$\gamma\beta\alpha=\beta\alpha\gamma, \quad \gamma^2=\beta\alpha\gamma\beta=\alpha\beta=0.$$
We have the following structure of the indecomposable projective $B$-modules: \\
\xymatrix@R=0.5pc@C=0.8pc
{&&1\ar@{-}[dl]\ar@{-}[dr]&& \\
 Q_1= & 1 \ar@{-}[d]& & 2\ar@{-}[d] & \\
& 2 \ar@{-}[dr] &  & 1 \ar@{-}[dl] & \\
&& 1 & ,& }
\xymatrix@R=0.5pc@C=0.8pc
{&2\ar@{-}[d]& \\
 Q_2= & 1\ar@{-}[d] & \\
& 1 \ar@{-}[d] & \\
& 2 & .}

Suppose that $A$ is isomorphic to some Brauer graph algebra $C$. Then the associated Brauer graphs of $A$ and $C$ are the same (see \cite[Lemma 3.1]{AZ2}). It follows that the Brauer graph algebras $B$ and $C$ are isomorphic. So, in order to prove that $A$ is not isomorphic to any Brauer graph algebra, it is enough to show that $A$ is not isomorphic to $B$. In fact, $A$ and $B$ are not stably equivalent. The reason is as follows. Each indecomposable $A$-module (resp. $B$-module) at the mouth of a tube of rank $1$ of the Auslander-Reiten quiver of $A$ (resp. $B$) is of the form
\[\begin{tikzcd}[row sep=7, column sep=0]
            & {1} & \\
            {X_t =} & {2} & t \\
            & {1} &
            \arrow[shift left=2, bend left=15, dashed, no head, from=1-2, to=3-2]
            \arrow[no head, from=1-2, to=2-2]
            \arrow[no head, from=2-2, to=3-2]
        \end{tikzcd}\]
with $t\in k$, where the dashed line means that \(\gamma-t\beta\alpha\) acts as zero on the elements of $\top(X_t)$. A calculation shows that $\Omega_{A}(X_t)\cong X_{1+t}$ and $\Omega_{B}(X_t)\cong X_{t}$ for each $t\in k$.
If $F:B\text{-}\mathrm{\underline{mod}}\rightarrow A\text{-}\mathrm{\underline{mod}}$ is a stable equivalence, then $F(X_0)$ is isomorphic to some $X_t$ with $t\in k$. We have
$$X_{t+1}\cong\Omega_{A}(X_t)\cong\Omega_{A}F(X_0)\cong F\Omega_{B}(X_0)\cong F(X_0)\cong X_t,$$
a contradiction.

However, when $\mathrm{char}(k)\neq 2$, the change of basis given by
\[
\alpha'=\alpha,\quad \beta'=\beta,\quad \gamma'=\gamma-\frac{\beta\alpha}{2}
\]
yields an isomorphism $A\cong B$. Under this isomorphism, certain band modules over $A$ are sent to string modules over $B$.
For instance, the band module $X=A(\gamma-\frac{\beta\alpha}{2})$ over $A$ is mapped to the string module $Y=B\gamma$ over $B$, with the following Loewy structures:
\[
\begin{tikzcd}[row sep=7, column sep=0]
            & {1} & \\
            {X =} & {2} & \frac{1}{2}, \\
            & {1} &
            \arrow[shift left=2, bend left=15, dashed, no head, from=1-2, to=3-2]
            \arrow[no head, from=1-2, to=2-2]
            \arrow[no head, from=2-2, to=3-2]
        \end{tikzcd}
\qquad
\begin{tikzcd}[row sep=7, column sep=0]
            & {1}  \\
            {Y =} & {2} \\
            & {1}
            \arrow[no head, from=1-2, to=2-2]
            \arrow[no head, from=2-2, to=3-2]
        \end{tikzcd}.
\]
\end{Ex1}

Note that although every symmetric stably biserial algebra admits a presentation of the form $kQ/I$ described in Theorem~\ref{thm:sym-StBA}, an algebra of this form need not be symmetric, as follows from \cite[Remark~3.2]{AZ2}. Moreover, even in characteristic $2$, there exist stably biserial algebras of the form $kQ/I$ given in Theorem~\ref{thm:sym-StBA} that are isomorphic to special biserial algebras, as mentioned in the final paragraph of the proof of \cite[Lemma~3.1]{AZ2}. The following lemma makes these observations explicit.

\begin{lemma}\label{lem:stb-local-iso}
	Let $0\neq\lambda_1\in k$, and consider the stably biserial algebra
$$B=k\langle x,y\rangle/\langle xy-yx,\,x^2-\lambda_1yx,\,y^2-\lambda_2yx,\,x^2y,\,y^2x\rangle.$$
The set $\mathcal{L}$ of deformed loops is given by
$$\mathcal{L}=\begin{cases}
	\{x\}, & \text{if }\lambda_2=0,\\
	\{x,y\}, & \text{if }\lambda_2\neq0.
\end{cases}$$
Then $B$ is symmetric if and only if $\lambda_1\lambda_2\neq1$. Moreover, if $B$ is symmetric, then $B$ is isomorphic to the special biserial algebra
$$A=k\langle x,y\rangle/\langle x^2-y^2,\,xy,\,yx\rangle.$$
\end{lemma}

\begin{proof}
Note that $B$ has $k$-basis $\{1, x, y, yx\}$, and
$$J(B)=kx\oplus ky\oplus kyx,\qquad J(B)^2=kyx.$$ Suppose first that $\lambda_1\lambda_2=1$. Then $z:=x-\lambda_1y$ satisfies
$$zx=x^2-\lambda_1yx=0,\qquad zy=xy-\lambda_1y^2=(1-\lambda_1\lambda_2)yx=0.$$
Thus $z,yx\in\operatorname{soc}(B)$. Since $z\notin J(B)^2$ whereas $yx\in J(B)^2$, they are linearly independent. Consequently, $\mathrm{dim}_k\operatorname{soc}(B)\ge2$. Therefore, if $B$ is symmetric, then $\lambda_1\lambda_2\neq1.$
Conversely, suppose that $\lambda_1\lambda_2\neq1$. Define a linear form $\varphi:B\to k$ by
$$\varphi(yx)=1,\qquad\varphi(1)=\varphi(x)=\varphi(y)=0.$$
Then the associative symmetric bilinear form $\langle u,v\rangle=\varphi(uv)$ is nondegenerate, so $B$ is symmetric.

Finally, assume that $B$ is symmetric. Then $\lambda_1\lambda_2\neq1$. Since $k$ is algebraically closed, there exists $0\neq c\in k$ satisfying
$$c^2(\lambda_1\lambda_2-1)=1.$$
Let
$$X=x,\qquad Y=c(x-\lambda_1y)$$
be two elements in $B$.
One easily verifies that
$$X^2=Y^2=\lambda_1yx,\qquad XY=YX=0.$$
Hence the assignment
$$x\mapsto X,\qquad y\mapsto Y$$
defines a homomorphism $A\longrightarrow B$. Moreover, the images of $X$ and $Y$ are linearly independent in $J(B)/J(B)^2$, so they generate $B$. Therefore the homomorphism is surjective. Since both $A$ and $B$ are $4$-dimensional, it is an isomorphism. Thus $A\cong B.$
\end{proof}

By comparing the tubes in the stable Auslander-Reiten quiver of a symmetric stably biserial algebra with the faces of its corresponding Brauer graph, Antipov and Zvonareva established in \cite[Section 4.2]{AZ2} the following result.

\begin{Prop}\label{prop:stbA-tube-and-face}
	Let $A$ be a representation-infinite symmetric stably biserial algebra with Brauer graph $\Gamma$. Then, for any face of $\Gamma$ with perimeter $p > 2$, the structure of the stable Auslander-Reiten quiver of $A$ reflects this combinatorial data. If $p$ is odd, the quiver contains a tube of rank $p$ that is invariant under the syzygy functor $\Omega_A$. If $p$ is even, it contains a pair of tubes, each of rank $p/2$, which are exchanged by $\Omega_A$.
\end{Prop}

Specifically, for a face $F=(h_1,h_2,\cdots,h_{p})$, let us denote by $\mathcal{M}_F$ the set of modules in the mouth of the $p$-tube
or two $p/2$-tubes, corresponding to $F$. They can be constructed as follows: take two consecutive
half-edges $h_i$ and $h_{i+1}=\iota\rho(h_{i})$ from the face $F$, they correspond to two vertices $\bar{h_i}$ and $\bar{h_{i+1}}$ in the quiver of
$A$ and an arrow $\alpha_{h_i}: \bar{h_i}\rightarrow \bar{h_{i+1}}$ (resp. a path $C_{\alpha_{\rho^{-1}(h_{i+1})}}: \bar{h_{i+1}}\rightarrow \bar{h_{i+1}}$) if $s(h_i)$ is not truncated (resp. if $s(h_i)$ is truncated); the modules in $\mathcal{M}_F$ are of the form $M_i:=Ae_{\bar{h_i}}/A\alpha_{h_i}$ (resp. $M_i:=Ae_{\bar{h_i}}/AC_{\alpha_{\rho^{-1}(h_{i+1})}}$). Thus, $\mathcal{M}_F=\{M_1,\cdots,M_p\}$. Note that the modules in the union of all $\mathcal{M}_F$ are exactly the simple modules whose projective covers are uniserial and the maximal uniserial quotients of indecomposable non-uniserial projective modules.

It should be noted that if one considers only Brauer graph algebras (that is, symmetric special biserial algebras), then for a face in $\Gamma$ with perimeter $1$ or $2$, Proposition \ref{prop:stbA-tube-and-face} also holds. More precisely, if $p=1$, then such a face corresponds to a tube of rank $1$ in the stable Auslander-Reiten quiver, and if $p=2$, then such a face corresponds to a pair of two tubes of rank $1$ in the stable Auslander-Reiten quiver which are exchanged by the syzygy functor. For instance, the $A_{\Gamma_1}$ in Example \ref{ex:two-exceptional-band} illustrates this correspondence for a face of perimeter $2$, which yields a pair of tubes that are not stable under the syzygy functor.

However, this correspondence fails for the broader class of symmetric stably biserial algebras. As a counterexample, consider a deformed loop in a symmetric stably biserial algebra $A$ (which gives rise to a face of perimeter $1$), as shown below:
$$\begin{tikzpicture}
\draw (-0.5,0) circle (0.5);
\fill (0,0) circle (0.5ex);
\node at(-1.2,0) {$1$};
\node at(-0.3,0) {$m$};
\draw[-] (0,0) -- (1,1);
\draw[-] (0,0) -- (1,-1);
\draw (-0.45,-0.15) rectangle (-0.15,0.15);
\node at(0.5,0.75) {$2$};
\node at(0.5,-0.75) {$n$};
\draw[dotted] (0.866,-0.5) arc (-30:30:1);
\end{tikzpicture}$$
Then there exist two tubes of rank 1 which are not stable under the syzygy functor $\Omega_A$; indeed, $\Omega_A$ sends one to the other. The indecomposable modules at the mouths of these homogeneous tubes are as follows:

$$\begin{tikzpicture}
	\begin{scope}[xshift=-60pt]
		\node at(0,0) {$1$};
\node at(0,-0.3) {$2$};
\draw[dotted] (0,-0.5) -- (0,-0.9);
\node at(0,-1.1) {$n$};
\node at(0,-1.4) {$1$};
\node at(0,-1.7) {$1$};
\draw[dotted] (0,-1.9) -- (0,-2.3);
\node at(0,-2.5) {$1$};
\node at(0,-2.8) {$2$};
\draw[dotted] (0,-3) -- (0,-3.4);
\node at(0,-3.6) {$n$};
\node at(0,-3.9) {$1$};
	\end{scope}
		\begin{scope}[xshift=60pt]
		\node at(0,0) {$1$};
\node at(0,-0.3) {$2$};
\draw[dotted] (0,-0.5) -- (0,-0.9);
\node at(0,-1.1) {$n$};
\node at(0,-1.4) {$1$};
\node at(0,-1.7) {$1$};
\draw[dotted] (0,-1.9) -- (0,-2.3);
\node at(0,-2.5) {$1$};
\node at(0,-2.8) {$2$};
\draw[dotted] (0,-3) -- (0,-3.4);
\node at(0,-3.6) {$n$};
\node at(0,-3.9) {$1$};
\node at(0.7,-2) {$t$\;.};
\draw[dashed] (0.15,-3.9) arc (-15:15:7.5);
	\end{scope}
	\node at(0,-2) {\text{and}};
\end{tikzpicture}$$
Note that the former one is a string module and the latter one is an exceptional band module. We will study the stable Hom-spaces from these modules in Section \ref{app:lem}.


\subsection{Stable equivalences of Morita type and centers}\label{subsec:stb-and-center}
\

According to \cite[Proposition~5.2]{BHS}, if an algebra is stably equivalent to a symmetric algebra whose radical square is nonzero, then it must be weakly symmetric. However, we even do not know a concrete example where a symmetric algebra is stably equivalent to some non-symmetric algebra over an algebraically closed field, although there do exist such examples over a non-algebraically closed field (see \cite{OTY}).

Therefore, in order to simplify our discussion, we consider those stable equivalences which preserve the property of being symmetric. This leads us to consider stable equivalences of Morita type. We recall their basic definition below and refer the reader to~\cite[Chapter 5]{Z} for further details.


\begin{Def}
	Let $A$ and $B$ be two $k$-algebras.
A stable equivalence of Morita type between $A$ and $B$
consists of bimodules
\[
{}_A M_B \quad \text{and} \quad {}_B N_A
\]
such that:
\begin{enumerate}
  \item $M$ is projective as a left $A$-module and as a right $B$-module;
  \item $N$ is projective as a left $B$-module and as a right $A$-module;
  \item There exist projective bimodules $P$ (an $A$-$A$-bimodule) and $Q$ (a $B$-$B$-bimodule) such that
  \[
  M \otimes_B N \cong A \oplus P \quad \text{as } A\text{-}A\text{-bimodules,}
  \]
  and
  \[
  N \otimes_A M \cong B \oplus Q \quad \text{as } B\text{-}B\text{-bimodules.}
  \]
\end{enumerate}

\noindent
In this situation, the tensor functor{
\[
M \otimes_B - \; : \; B\text{-}\mathrm{mod}\longrightarrow A\text{-}\mathrm{mod}
\]
induces a stable equivalence $B\text{-}\underline{\mathrm{mod}}\longrightarrow A\text{-}\underline{\mathrm{mod}}$,
with quasi-inverse induced by
\[
N \otimes_A - \; : \; A\text{-}\mathrm{mod} \longrightarrow B\text{-}\mathrm{mod}.
\]
}\end{Def}

Clearly, the above functor $M \otimes_B - \; : \; B\text{-}\mathrm{mod}\longrightarrow A\text{-}\mathrm{mod}$ is exact and commutes with the syzygy functors $\Omega_A$ and $\Omega_B$. In particular, if both $A$ and $B$ are self-injective algebras, then $M \otimes_B -$  induces a triangulated equivalence $B\text{-}\underline{\mathrm{mod}}\longrightarrow A\text{-}\underline{\mathrm{mod}}$. Moreover, a stable equivalence of Morita type preserves the property of being symmetric and indecomposability of algebras (see \cite[Proposition 2.1 and Corollary 2.4]{L}).

Now we recall some basic notions concerning the centers of $k$-algebras. Let $A$ be a $k$-algebra.
The center of $A$ is
  \[
  Z(A) = \{\, z \in A \mid za = az \text{ for all } a \in A \,\}.
  \]
  It is well-known that $Z(A) \cong \operatorname{End}_{A^e}(A)$.

Recall that for a $k$-algebra $A$, the Reynolds ideal $R(A)$ is defined by $R(A)=Z(A)\cap\soc(A)$.  For a symmetric stably biserial algebra $A$, we have $\soc(A)\subseteq Z(A)$ and $R(A)=\mathrm{soc}(A)$.

\begin{proposition}\textnormal{(\cite[Theorem 1.7]{ZZ})}\label{prop:ZZ}
	Let $A$ and $B$ be two symmetric indecomposable $k$-algebras which are stably equivalent of Morita type. If $k$ is of positive characteristic, then we have an isomorphism of algebras $Z(A)/R(A)\cong Z(B)/R(B)$.
\end{proposition}

\subsubsection{Centers of symmetric stably biserial algebras}
\

For symmetric stably biserial algebras, whose quiver presentations are given by Proposition \ref{thm:sym-StBA}, we now recall the results on their centers from \cite{AZ2}.

Assume that $A$ is a symmetric stably biserial algebra defined by a Brauer graph $(\Gamma, m)$ together with a set $\mathcal{L}$ of deformed loops.
Let $\{C_1, C_2, \ldots, C_r\}$ be the set of cycles of the form $C_\alpha$ in $Q_\Gamma$ (see Subsection~\ref{subsec:def-BGA} for the definition of $C_\alpha$), considered up to cyclic permutation, where each $C_i$ corresponds to a vertex $v_i$ of $\Gamma$. For each $1 \le i \le r$, consider a cyclic sequence $(\alpha_{i,1}, \alpha_{i,2}, \ldots, \alpha_{i,l_i})$ of arrows in the cycle $C_i$, where $\sigma(\alpha_{i,j}) = \alpha_{i,j+1}$ and $l_i$ denotes the length of $C_i$. Let $r' \leq r$ be an integer such that $m(C_i) > 1$ for $i = 1, \ldots, r'$ and $m(C_i) = 1$ for $i = r' + 1, \ldots, r$.


For each loop $\gamma = \alpha_{i,j}$ such that $\sigma(\gamma) \ne \gamma$, set
\[
q_{\gamma} = q_{\alpha_{i,j}}
= \alpha_{i,j-1} \cdots \alpha_{i,1} \alpha_{i,l_i} \cdots \alpha_{i,j+1}
(\alpha_{i,j} \alpha_{i,j-1} \cdots \alpha_{i,1} \alpha_{i,l_i} \cdots \alpha_{i,j+1})^{m(C_i)-1}.
\]
For each vertex $v$ of $Q$, let $\alpha_{i,j}$ be an arrow starting at $v$, and set
\[
s_v = (\alpha_{i,j-1} \cdots \alpha_{i,1} \alpha_{i,l_i} \cdots \alpha_{i,j})^{m(C_i)}.
\]
Note that $s_v$ does not depend on $\alpha_{i,j}$ according to the defining relation $(a)$ in $I_\Gamma$, and $s_v$ is a socle element of $A$ corresponding to the vertex $v$. Note also that the following result only depends on the associated Brauer graph $(\Gamma, m)$ of $A$ and is free of the characteristic of the field $k$.

\begin{Prop}\textnormal{(\cite[Proposition 4.1]{AZ2})}\label{prop:center-of-stBA}
Let $A=kQ/I$ be a symmetric stably biserial algebra defined by a Brauer graph $(\Gamma, m)$. As a vector space over $k$, the center $Z(A)$ is generated by $1$ together with the following elements:
\begin{itemize}
  \item[(a)] Elements
  \[
  p_{i,t} =
  (\alpha_{i,l_i}\cdots\alpha_{i,2}\alpha_{i,1})^t
  + (\alpha_{i,1}\cdots\alpha_{i,3}\alpha_{i,2})^t
  + \cdots
  + (\alpha_{i,l_i-1}\cdots\alpha_{i,1}\alpha_{i,l_i})^t,
  \]
  for $i = 1, 2, \ldots, r'$ and $t = 1, \ldots, m(C_i) - 1$.

  \item[(b)] Elements $q_{\gamma}$ for each loop $\gamma$ such that $\sigma(\gamma) \ne \gamma$.

  \item[(c)] Elements $s_v$ for each vertex $v \in Q_0$.
\end{itemize}
\end{Prop}



\section{Hom-spaces from the modules at the tube induced by a deformed loop}\label{app:lem}
	In this section we prove a dimension formula (see Lemma \ref{lem:equal-dimension}) that will be used in the proof of Theorem \ref{thm:st.M-BGA}. {The inequality part of this dimension formula can be seen as a generalization of some similar dimension formulas for Brauer graph algebras in \cite{A}.} To prove our result, we need a description of some nice basis for the space of homomorphisms between two string modules introduced by Crawley-Boevey in \cite{Cra89}.

	Let $\varLambda = kQ/I$ be a symmetric stably biserial algebra and let $\overline{\varLambda} = kQ/I'$ be the associated string algebra. Then $I'$ is generated by a set $\mathcal{R}$ of paths of length at least two in $kQ$. Let $X$ and $Y$ be two string $\varLambda$-modules associated with strings $w_X$ and $w_Y$, respectively. Let $S$ be a quiver of type $A_n$. We denote by $V_S$ the $kS$-module which is one-dimensional at each vertex and in which the linear maps corresponding to each arrow are the identity. As in \cite{Cra89}, we can depict a string $w_X$ as a quiver morphism $F: S\ra Q$ with $S$, which satisfies
\begin{itemize}
  \item[(1)] there are no subquivers of the form $\cdot\stackrel{a}{\longleftarrow}\cdot\stackrel{b}{\longrightarrow}\cdot$ or $\cdot\stackrel{a}{\longrightarrow}\cdot\stackrel{b}{\longleftarrow}\cdot$ in $S$ with $F(a)=F(b)$,
  \item[(2)] there is no path in $S$ whose image under $F$ is in $\mathcal{R}$,
\end{itemize}
the associated string module $X$ is $F_\lambda(V_S)$, where $F_\lambda$ is the push-down functor $F_\lambda: kS\mbox{-mod}\longrightarrow \overline{\varLambda}\mbox{-mod}$ defined by
$$F_\lambda(V)(x)=\bigoplus_{y\in F^{-1}(x)} V(y),$$
for $x$ either a vertex or an arrow in $Q$. {Note that the quiver morphism $F: S\ra Q$ factors through the topological universal cover $\pi: \widetilde{Q}\longrightarrow Q$ and therefore the functor $kF: kS\ra kQ/I'$ factors through the covering functor $k\pi: k\widetilde{Q}/\widetilde{I'}\longrightarrow kQ/I'$. It follows that the push-down functor $F_\lambda: kS\mbox{-mod}\longrightarrow \overline{\varLambda}\mbox{-mod}$ equals the composition of the push-down functors $kS\mbox{-mod}\longrightarrow k\widetilde{Q}/\widetilde{I'}\mbox{-mod}$ (which is fully faithful) and $k\widetilde{Q}/\widetilde{I'}\mbox{-mod}\longrightarrow\overline{\varLambda}\mbox{-mod}$. From this it is easy to see that the string module $F_\lambda(V_S)$ is indecomposable.}

Let $w_X$ be a string defined by a quiver morphism $F: S\ra Q$ as above. The quiver $S$ decorated by the vertices and the arrows in $Q$ through $F$ is called the diagram of the string module $F_\lambda(V_S)$ (cf. Example \ref{string module}). The vertices of $S$ correspond to a $k$-basis of $F_\lambda(V_S)$, which is called the diagrammatic basis of this string module (after \cite{A} and \cite{AZ1}). Note that diagrams provide an equivalent description of string modules: they are in bijection with equivalence classes of strings under the relation identifying a string with its inverse.

From now on, by abuse of notation, we will use the same symbol $X$ to denote a string module, its diagram, and a fixed representative string whenever no confusion can arise.

	{Now let $X=F_\lambda(V_S)$ and $Y=E_\lambda(V_T)$ be two string modules.} According to \cite{Cra89} (see also \cite{Sch99}), the space Hom$_{\varLambda}(X,Y)$ has a $k$-basis described as follows. Each basis element is given by a triple $(S',T',\alpha)$, where $S'$ is a connected subquiver of $S$ which is closed under taking predecessors of arrows, $T'$ is a connected subsquiver of $T$ which is closed under taking successors of arrows, and $\alpha$ is a quiver isomorphism $S'\longrightarrow T'$ such that $E\circ\alpha=F$. Following \cite{A}, we will call the elements of this basis as diagrammatic morphisms in Hom$_{\varLambda}(X,Y)$. To describe the homomorphism corresponding to the triple $(S',T',\alpha)$ we must specify a map
$$f(x): \bigoplus_{y\in F^{-1}(x)}k\longrightarrow\bigoplus_{z\in E^{-1}(x)}k$$
for each vertex $x$ in $Q$. This map should have component $f(x)_{yz}=1$ if $y\in S'$ and $z=\alpha(y)$, and otherwise $f(x)_{yz}=0$.

{Note that if $X$ is a uniserial module (this is the case we will use), then we have the following simple way to depict the diagrammatic morphisms in Hom$_{\varLambda}(X,Y)$: they are just determined by those uniserial submodules $M$ of $Y$ such that $M$ is isomorphic to a quotient module of $X$.}

It should be noted that Crawley-Boevey stated his result in \cite{Cra89} in a more general setting, that is, for two tree modules over any zero-relation algebra. As we mentioned before, string modules over a stably biserial algebra can be seen as string modules over a related string algebra (which is clearly a zero-relation algebra) and therefore our situation satisfies the condition in \cite{Cra89}.

We refer the reader to \cite{BC}, where the diagrammatic morphisms and the corresponding Hom spaces are interpreted via the surface model of {the associated string algebra $\overline{\varLambda}$.} However,  this geometric approach is not sufficient, since we need to work with stable Hom spaces, and the surface model loses the information of the non-uniserial projective {$\varLambda$}-modules.


We provide the following example to help the reader understand {the diagrammatic basis and the diagrammatic morphisms.}

\begin{Ex1}\label{string module}
Let $\varLambda$ be the symmetric stably biserial algebra in Example \ref{ex-stably-biserial}. {Using the Loewy diagrams, let}
\[\begin{tikzcd}[row sep=-4, column sep=-4]
            & 1 \\
            {X=} & 2 \\
            & 1
        \end{tikzcd}\]
and
\[\begin{tikzcd}[row sep=-4, column sep=-4]
            & 1 &&&& \\
            {Y=} && 2 &&1& \\
            &&& 1 && 2
        \end{tikzcd}\]
be two string $\varLambda$-modules defined by the strings $\beta\alpha$, $\alpha\gamma^{-1}\beta\alpha$, respectively. We depict them by the following diagrams respectively:
\[\begin{tikzpicture}[baseline=(current bounding box.center)]
  \node (x1) at (0,1) {$x_1$};
  \node (y1) at (0.5,0) {$y_1$};
  \node (x2) at (1,-1) {$x_2$};
  \node (Xeq) at (-1,0) {$X=$};

  \draw[->] (x1) -- (y1);
  \draw[->] (y1) -- (x2);
\end{tikzpicture},\]

\[\begin{tikzpicture}[baseline=(current bounding box.center)]
  \node (x3) at (0,1) {$x_3$};
  \node (y2) at (0.5,0) {$y_2$};
  \node (x4) at (1.5,0) {$x_4$};
  \node (x5) at (1,-1) {$x_5$};
  \node (y3) at (2,-1) {$y_3$};
  \node (Yeq) at (-1,0) {$Y=$};

  \draw[->] (x3) -- (y2);
  \draw[->] (y2) -- (x5);
  \draw[->] (x4) -- (x5);
  \draw[->] (x4) -- (y3);
\end{tikzpicture},\]
where the diagrammatic basis elements $x_i$ and $y_i$ correspond to the simple $\varLambda$-modules $S_1$ and $S_2$, respectively. Following \cite{AZ1}, we call a source vertex in the above {diagrams} a peak and a sink vertex a deep. For instance, $x_3$ and $x_4$ are peaks and $x_5$ and $y_3$ are deeps in the {diagram} of $Y$.

There are exactly two diagrammatic morphisms $f_1,f_2$ from $X$ to $Y$, where $f_1$ (which sends $x_1$ to $x_3$) is given by the quiver isomorphism $\alpha_1\colon (x_1\ra y_1\ra x_2)\ra (x_3\ra y_2\ra x_5)$ and $f_2$ (which sends $x_1$ to $x_5$) is given by $\alpha_2\colon(x_1)\ra (x_5)$, and $\{f_1,f_2\}$ forms a basis of $\Hom_{\varLambda}(X,Y)$. Alternatively, since $X$ is uniserial, we can read $f_1,f_2$ directly from the Loewy structures of $X$ and $Y$, they correspond to the uniserial submodules \begin{tikzcd}[row sep=-4, column sep=-4]
           1 & & \\
           & 2 &\\
           & & 1
        \end{tikzcd} and
 $\ 1$ of $Y$ respectively. Note that by passing to the stable category, $\underline{f_2}=0$ and $\underline{\operatorname{Hom}}_{\varLambda}(X,Y)$ is spanned by $\underline{f_1}$.
\end{Ex1}

The following lemma will be used in the proof of Theorem \ref{thm:st.M-BGA}. It should be noted that although we only need this lemma in characteristic $2$ for Theorem \ref{thm:st.M-BGA}, this lemma is free of characteristic of the base  field $k$.

    \begin{lemma}\label{lem:equal-dimension}
        Let $\varLambda$ be a symmetric stably biserial algebra {with a presentation $kQ/I$ as in Theorem \ref{thm:sym-StBA}}. Consider the subgraph
        \[\begin{tikzpicture}
            \draw (-0.5,0) circle (0.5);
            \fill (0,0) circle (0.5ex);
            \node at(-1.2,0) {$v$};
            \node at(-0.3,0) {$m$};
            \draw[-] (0,0) -- (1,1);
            \draw[-] (0,0) -- (1,-1);
            \draw (-0.45,-0.15) rectangle (-0.15,0.15);
            \node at(0.5,0.75) {$x_1$};
            \node at(0.5,-0.75) {$x_k$};
            \draw[dotted] (0.866,-0.5) arc (-30:30:1);
        \end{tikzpicture}\]
        of the Brauer graph of $\varLambda$. It consists of a deformed loop $v$ together with all half-edges $x_1,\cdots,x_k$ incident with the same vertex of multiplicity $m$, and {without loss of generality we can assume that} the orientation around this vertex is clockwise.
        Let $X$ be the string $\varLambda$-module and $\Omega(X)$ be the exceptional band $\varLambda$-module defined as follows:
		\[\begin{tikzcd}[row sep=-4, column sep=0]
            & v &&&& v \\
            & {x_1} &&&& {x_1} \\
            & \vdots &&&& \vdots \\
            & {x_k} &&&& {x_k} \\
            & v &&&& v \\
            & v &&&& v \\
            {X=} & \vdots & {,} & \quad & {\Omega(X)=} & \vdots \\
            & v &&&& v \\
            & v &&&& v \\
            & {x_1} &&&& {x_1} \\
            & \vdots &&&& \vdots \\
            & {x_k} &&&& {x_k} \\
            & v &&&& v
            \arrow["t", shift left=3, bend left=15, dashed, no head, from=1-6, to=13-6]
        \end{tikzcd},\]
        where the direct sequence $v\ra x_1\ra\cdots\ra x_k\ra v$ with $k+2$ vertices occurs $m$ times in $X$ and $\Omega(X)$, and the dashed line means that \(\alpha-t \alpha^{-1}C_\alpha^m\) acts as zero on the elements of $\top(\Omega(X))$, with $t\in k^*$, $\alpha\colon v\ra v$ being the deformed loop. Let $Y$ be a string $\varLambda$-module. Then
        \[\dim\nolimits_k \sHom_{\varLambda}(X,Y)=\dim\nolimits_k \sHom_{\varLambda}(\Omega_{\varLambda}(X),Y) \leq 2.\]

    \end{lemma}

	{\it Outline of the proof for the equality part:}
		
		For simplicity, we just write $\Hom$ and $\sHom$ for \(\Hom_{\varLambda}\) and \(\sHom_{\varLambda}\) respectively. We first prove that $\dim_k \sHom (X,Y) \leq \dim_k \sHom(\Omega(X),Y)$. Consider all diagrammatic morphisms from $X$ to $Y$. As we mentioned before, they forms a basis $\mathscr{B}$ of $\Hom(X,Y)$. For a fixed morphism $f$ in this basis, we define {an appropriate} subdiagram $N=N_f$ of $Y$ which contains the image of the starting vertex of $X$, and we classify the local configuration of $Y$ containing $N$ as a subdiagram by discussing all the possibilities of $N$ case by case. Moreover, we define an equivalence relation $\sim$ on $\mathscr{B}$: two morphisms $f,f'\in\mathscr{B}$ are equivalent if their corresponding subdiagrams $N_f,N_{f'}$ of $Y$ are the same. For each equivalence class $\mathscr{C}$ in $\mathscr{B}$, we define a finite set of morphisms $\mathscr{C}'$ in $\Hom(\Omega(X),Y)$.
		Let {$\underline{\mathscr{C}}$} (resp. {$\underline{\mathscr{C}'}$}) denote a basis of $\span\{\underline{f_i}\mid f_i\in {\mathscr{C}}\}$ (resp. $\span\{\underline{g_j}\mid g_j\in {\mathscr{C'}}\}$) in the Hom-spaces of the corresponding stable categories respectively.
		We prove the following statements:

			(1) {$\bigcup\limits_{\mathscr{C}\in\mathscr{B}/\sim}\underline{\mathscr{C}}$} forms a basis of \(\underline{\operatorname{Hom}}(X,Y)\).

			(2) {$\bigcup\limits_{\mathscr{C}\in\mathscr{B}/\sim}\underline{\mathscr{C}'}$} is linearly independent in \(\underline{\operatorname{Hom}}(\Omega(X),Y)\).

			(3) The two sets in (1) and (2) have the same cardinality. \\
		From (1)-(3) it follows that
		\[\begin{aligned}
			\dim\nolimits_k \sHom_{\varLambda}(X,Y)&={\sum_{\mathscr{C}\in\mathscr{B}/\sim}|\underline{\mathscr{C}}|}\\
			&={\sum_{\mathscr{C}\in\mathscr{B}/\sim}|\underline{\mathscr{C}'}|}\leq\dim\nolimits_k \sHom_{\varLambda}(\Omega_{\varLambda}(X),Y).
		\end{aligned}\]

		We then prove the reverse inequality by showing that $\Hom(\Omega(X), Y)$ is spanned by $\bigcup\limits_{\mathscr{C}\in\mathscr{B}/\sim}\mathscr{C}'$, which implies the fact that the defined $\bigcup\limits_{\mathscr{C}\in\mathscr{B}/\sim}\underline{\mathscr{C}'}$ forms a basis of \(\underline{\operatorname{Hom}}(\Omega(X),Y)\) and hence
        \[\dim\nolimits_k \sHom_{\varLambda}(X,Y)=\dim\nolimits_k \sHom_{\varLambda}(\Omega_{\varLambda}(X),Y).\]

		\medskip


	Now we give the detailed proof of Lemma \ref{lem:equal-dimension}.

\begin{proof}
		We first prove that $\dim\nolimits_k \sHom_{\varLambda}(X,Y) \leq 2$ and $\dim\nolimits_k \sHom_{\varLambda}(X,Y) \leq \dim\nolimits_k \sHom_{\varLambda}(\Omega_{\varLambda}(X),Y)$.

        Let $\eta\colon v\ra x_1\ra\cdots\ra x_k\ra v'\ra\cdots\ra v'\ra x_1\ra\cdots\ra x_k\ra v$ be the diagram of the string module $X$, where we denote those vertices $v$ which are not endpoints of $\eta$ by $v'$.
        Consider an infinite zigzag diagram $Z$ given by $\cdots\la\eta\la\eta\la\cdots\la\eta\la\cdots$. Then the vertices labelled by $v$ in the copies of $\eta$ are precisely all the peaks and deeps (see these notions in Example \ref{string module}) of $Z$.
		Take the basis of $\Hom(X,Y)$ formed by all diagrammatic morphisms, and denote by $f$ an element of this basis. Then $f$ sends the starting vertex of $X$ to some $v$ in $Y$, and this $v$ must lie in a unique common subdiagram {$N=N_f$} of $Z$ and $Y$ satisfying the following properties:
		\begin{enumerate}[label=(\alph*)]
			\item In the vertex notation of $Z$, the two endpoints of $N$ are both $v$, that is, either of the two endpoints is a peak or a deep of $Z$.
			\item In the vertex notation of $Z$, $f$ sends the starting vertex of $X$ to some $v$ in $N$.
			\item $N$ is maximal as a subdiagram of $Y$ with respect to the above two properties.
		\end{enumerate}
		
Note that we can construct $N_f$ as follows. We fix the vertex corresponding to \(\top(f(X))\) in $Y$, and this vertex `grows' to a diagram satisfying $(a)$ and $(b)$ along the common subdiagram of the diagrams $Z$ and $Y$ with the greatest extent possible. The idea behind this definition is that each $N_f$ controls some morphisms $g$ in $\Hom(\Omega(X),Y)$ such that \(\top(g(\Omega(X)))\) is generated by a linear combination of some diagrammatic basis elements contained in $N_f$, and all these morphisms together induces a basis of \(\underline{\operatorname{Hom}}(\Omega(X),Y)\).

[{\it Interlude.}] In Example \ref{string module}, {by the new notation as in Lemma \ref{lem:equal-dimension}, the diagram $\eta$ is given by $v_1\ra x_1\ra v_2$ where $v_1=v_2=v=1$, the diagram $Y$ is given by the following diagram:}
\[\begin{tikzpicture}[baseline=(current bounding box.center)]
  \node (x3) at (0,1) {$v_1$};
  \node (y2) at (0.5,0) {$x_1$};
  \node (x4) at (1.5,0) {$v_2$};
  \node (x5) at (1,-1) {$v_3$};
  \node (y3) at (2,-1) {$x_1'$};
  \node (Yeq) at (-1,0) {$Y=$};

  \draw[->] (x3) -- (y2);
  \draw[->] (y2) -- (x5);
  \draw[->] (x4) -- (x5);
  \draw[->] (x4) -- (y3);
\end{tikzpicture},\] $N_{f_1}$ and $N_{f_2}$ are the same (which we denote by $N$) and $N$ is given by the subdiagram
\[\begin{tikzpicture}
  \node (x3) at (0,1) {$v_1$};
  \node (y2) at (0.5,0) {$x_1$};
  \node (x4) at (1.5,0) {$v_2$};
  \node (x5) at (1,-1) {$v_3$};

  \draw[->] (x3) -- (y2);
  \draw[->] (y2) -- (x5);
  \draw[->] (x4) -- (x5);
\end{tikzpicture}\]
of $Y$. We write
 \[\begin{tikzcd}[row sep=4, column sep=0]
            & {x'_1} & \\
            {\Omega(X)=} & {y'_1} & t \\
            & {x'_2} & ,
            \arrow[shift left=2, bend left=15, dashed, no head, from=1-2, to=3-2]
        \end{tikzcd}\]
which means that $y_1' = \alpha x_1', x_2' = \beta y_1'$ and $tx_2' = \gamma x_1'$. {Observe that $\Hom(\Omega(X),Y)$ has a basis $\{g_1,g_2\}$, where $g_1(x'_1)=v_1 +v_2$ and $g_2(x'_1)=v_3$.} Clearly both $g_1$ and $g_2$ are `controlled' by $N$. By passing to the stable category, $\underline{g_2}=0$ and $\underline{\operatorname{Hom}}(\Omega(X),Y)$ is spanned by $\underline{g_1}$.

As we mentioned before, each diagrammatic morphism $f$ determines a subdiagram $N_f$ of $Y$. We classify the possible local configurations of $Y$ obtained by extending $N_f$ along $Y$ to the first peak or deep on both sides.

		In the following cases, we use $x,x'\in \{v,v',x_1,\cdots,x_k\}$ to denote arbitrary vertices. Note that the vertices denoted by $x$ or $x'$ may be the endpoints of $Y$. For convenience, we fix the arrows within $\eta$ to be directed from left to right, and the arrows connecting each $\eta$ to be directed from right to left, and we may take the (horizontal) flip of $Y$ for this purpose.

		We now divide the discussion into major cases and secondary cases. The major cases are classified according to the number of occurrences of $v$ in $N$, while the secondary cases describe the corresponding local configurations of $Y$.

        {\bf Case 1.} $N$ is a single vertex $v$. After possibly taking the flip of $Y$, the local configuration of $Y$ containing $N$ as a subdiagram can be classified into the following cases:
\\
        Case 1.1.
            \begin{tikzcd}
                v
            \end{tikzcd}.
\\
        Case 1.2.
% https://q.uiver.app/#q=WzAsMyxbMSwwLCJ4Il0sWzIsMywidiJdLFswLDNdLFswLDEsIlxcY2RvdHMiLDFdLFswLDIsIlxcY2RvdHMiLDEseyJzdHlsZSI6eyJib2R5Ijp7Im5hbWUiOiJub25lIn0sImhlYWQiOnsibmFtZSI6Im5vbmUifX19XSxbMCw0LCIiLDEseyJzaG9ydGVuIjp7InRhcmdldCI6MjB9LCJzdHlsZSI6eyJib2R5Ijp7Im5hbWUiOiJub25lIn0sImhlYWQiOnsibmFtZSI6Im5vbmUifX19XSxbMCw1LCIiLDEseyJsZXZlbCI6MX1dXQ==
\begin{tikzcd}[row sep=small, column sep=tiny,  every label/.append style = {font=\normalsize}]
	& x \\
	\\
	\\
	{} && v
	\arrow[""{name=0, anchor=center, inner sep=0}, "\cdots"{description}, sloped,,  ,  draw=none, from=1-2, to=4-1]
	\arrow["\cdots"{description}, sloped,,  ,  from=1-2, to=4-3]
	\arrow[""{name=1, anchor=center, inner sep=0}, draw=none, from=1-2, to=0]
	\arrow[from=1-2, to=1]
\end{tikzcd}, where
% https://q.uiver.app/#q=WzAsMixbMCwwLCJ4Il0sWzEsMywidiJdLFswLDEsIlxcY2RvdHMiLDFdXQ==
\begin{tikzcd}[row sep=small, column sep=tiny,  every label/.append style = {font=\normalsize}]
	x \\
	\\
	\\
	& v
	\arrow["\cdots"{description}, sloped,,  ,  from=1-1, to=4-2]
\end{tikzcd} is a {proper subdiagram of $\eta$ ending at $v$}.
\\
        Case 1.3.
% https://q.uiver.app/#q=WzAsMyxbMCwwLCJ2Il0sWzEsMywieCJdLFsyLDBdLFswLDEsIlxcY2RvdHMiLDFdLFsxLDIsIlxcY2RvdHMiLDEseyJzdHlsZSI6eyJib2R5Ijp7Im5hbWUiOiJub25lIn0sImhlYWQiOnsibmFtZSI6Im5vbmUifX19XSxbMSw0LCIiLDEseyJzaG9ydGVuIjp7InRhcmdldCI6MjB9LCJzdHlsZSI6eyJib2R5Ijp7Im5hbWUiOiJub25lIn0sImhlYWQiOnsibmFtZSI6Im5vbmUifX19XSxbNSwxLCIiLDEseyJsZXZlbCI6MX1dXQ==
\begin{tikzcd}[row sep=small, column sep=tiny,  every label/.append style = {font=\normalsize}]
	v && {} \\
	\\
	\\
	& x
	\arrow["\cdots"{description}, sloped,,  ,  from=1-1, to=4-2]
	\arrow[""{name=0, anchor=center, inner sep=0}, "\cdots"{description}, sloped,,  ,  draw=none, from=4-2, to=1-3]
	\arrow[""{name=1, anchor=center, inner sep=0}, draw=none, from=4-2, to=0]
	\arrow[from=1, to=4-2]
\end{tikzcd}, where
% https://q.uiver.app/#q=WzAsMixbMCwwLCJ2Il0sWzEsMywieCJdLFswLDEsIlxcY2RvdHMiLDFdXQ==
\begin{tikzcd}[row sep=small, column sep=tiny,  every label/.append style = {font=\normalsize}]
	v \\
	\\
	\\
	& x
	\arrow["\cdots"{description}, sloped,,  ,  from=1-1, to=4-2]
\end{tikzcd}
        is a {proper subdiagram of $\eta$ starting at $v$}.

		Note that in Case~1.2, if $x \ra \cdots \ra v$ equals $\eta$, then by the definition of $N$, $N$ contains at least two vertices that are peaks or deeps of $Z$.
		In this situation, the configuration is not included in Case 1, but is instead covered by cases with index $\geq 2$. Case 1.3 is analogous, {by considering $v \ra \cdots \ra x$ whether equals $\eta$ or not.}

        For convenience, in the following cases, we use $x \ra\cdots\ra v$, $x' \ra\cdots\ra v$ {to denote each proper subdiagram of $\eta$ ending at $v$, and use} $v \ra\cdots\ra x$, $v \ra\cdots\ra x'$ to denote each {proper subdiagram of $\eta$ starting at $v$,} whenever it appears in $Y$ as a subdiagram (these subdiagrams share only one endpoint with $N$). Note that the vertex corresponding to $\top f(X)$ may be the starting or the ending vertex of some $\eta$, or it may be a middle vertex of some $\eta$, and the latter case is included in Case 2.8.

        {\bf Case 2.} $N$ contains two vertices that are peaks or deeps of $Z$,
			that is, $N=$
			% https://q.uiver.app/#q=WzAsMixbMCwwLCJ2Il0sWzEsMywidiJdLFswLDEsIlxcY2RvdHMiLDFdXQ==
			\begin{tikzcd}[row sep=small, column sep=tiny,  every label/.append style = {font=\normalsize}]
				v \\
				\\
				\\
				& v
				\arrow["\cdots"{description}, sloped,,  ,  from=1-1, to=4-2]
			\end{tikzcd} or
            \begin{tikzcd}[row sep=small, column sep=tiny,  every label/.append style = {font=\normalsize}]
                & {} \\
                \\
                \\
                v\\
                {}
                \arrow[""{name=0, anchor=center, inner sep=0}, "v"{description}, draw=none, from=1-2, to=4-1]
                \arrow[""{name=1, anchor=center, inner sep=0}, draw=none, from=4-1, to=0]
                \arrow[from=1, to=4-1]
            \end{tikzcd}. After possibly taking the flip of $Y$, the local configuration of $Y$ containing $N$ as a subdiagram can be classified into the following cases: \\
        Case 2.1.
			\begin{tikzcd}[row sep=small, column sep=tiny,  every label/.append style = {font=\normalsize}]
				v \\
				\\
				\\
				& v
				\arrow["\cdots"{description}, sloped,,  ,  from=1-1, to=4-2]
			\end{tikzcd}. \quad
        Case 2.2.
			\begin{tikzcd}[row sep=small, column sep=tiny,  every label/.append style = {font=\normalsize}]
                & {} \\
                \\
                \\
                v\\
                {}
                \arrow[""{name=0, anchor=center, inner sep=0}, "v"{description}, draw=none, from=1-2, to=4-1]
                \arrow[""{name=1, anchor=center, inner sep=0}, draw=none, from=4-1, to=0]
                \arrow[from=1, to=4-1]
            \end{tikzcd}. \quad
        Case 2.3.
% https://q.uiver.app/#q=WzAsNCxbMiwzLCJ2Il0sWzEsMCwieCJdLFswLDNdLFszLDBdLFsxLDAsIlxcY2RvdHMiLDFdLFsxLDIsIlxcY2RvdHMiLDEseyJzdHlsZSI6eyJib2R5Ijp7Im5hbWUiOiJub25lIn0sImhlYWQiOnsibmFtZSI6Im5vbmUifX19XSxbMywwLCJ2IiwxLHsic3R5bGUiOnsiYm9keSI6eyJuYW1lIjoibm9uZSJ9LCJoZWFkIjp7Im5hbWUiOiJub25lIn19fV0sWzEsNSwiIiwxLHsic2hvcnRlbiI6eyJ0YXJnZXQiOjIwfSwic3R5bGUiOnsiYm9keSI6eyJuYW1lIjoibm9uZSJ9LCJoZWFkIjp7Im5hbWUiOiJub25lIn19fV0sWzAsNiwiIiwxLHsic2hvcnRlbiI6eyJ0YXJnZXQiOjIwfSwic3R5bGUiOnsiYm9keSI6eyJuYW1lIjoibm9uZSJ9LCJoZWFkIjp7Im5hbWUiOiJub25lIn19fV0sWzEsNywiIiwxLHsibGV2ZWwiOjF9XSxbOCwwLCIiLDEseyJsZXZlbCI6MX1dXQ==
\begin{tikzcd}[row sep=small, column sep=tiny,  every label/.append style = {font=\normalsize}]
	& x && {} \\
	\\
	\\
	{} && v
	\arrow[""{name=0, anchor=center, inner sep=0}, "\cdots"{description}, sloped,,  ,  draw=none, from=1-2, to=4-1]
	\arrow["\cdots"{description}, sloped,,  ,  from=1-2, to=4-3]
	\arrow[""{name=1, anchor=center, inner sep=0}, "v"{description}, draw=none, from=1-4, to=4-3]
	\arrow[""{name=2, anchor=center, inner sep=0}, draw=none, from=1-2, to=0]
	\arrow[""{name=3, anchor=center, inner sep=0}, draw=none, from=4-3, to=1]
	\arrow[from=1-2, to=2]
	\arrow[from=3, to=4-3]
\end{tikzcd}. \quad
        Case 2.4.
% https://q.uiver.app/#q=WzAsNSxbMiwxLCJ4Il0sWzEsNCwidiJdLFswLDddLFsyLDBdLFszLDRdLFswLDEsIlxcY2RvdHMiLDFdLFsxLDIsInYiLDEseyJzdHlsZSI6eyJib2R5Ijp7Im5hbWUiOiJub25lIn0sImhlYWQiOnsibmFtZSI6Im5vbmUifX19XSxbMCw0LCJcXGNkb3RzIiwxLHsic3R5bGUiOnsiYm9keSI6eyJuYW1lIjoibm9uZSJ9LCJoZWFkIjp7Im5hbWUiOiJub25lIn19fV0sWzEsNiwiIiwxLHsic2hvcnRlbiI6eyJ0YXJnZXQiOjIwfSwic3R5bGUiOnsiYm9keSI6eyJuYW1lIjoibm9uZSJ9LCJoZWFkIjp7Im5hbWUiOiJub25lIn19fV0sWzAsNywiIiwxLHsic2hvcnRlbiI6eyJ0YXJnZXQiOjIwfSwic3R5bGUiOnsiYm9keSI6eyJuYW1lIjoibm9uZSJ9LCJoZWFkIjp7Im5hbWUiOiJub25lIn19fV0sWzEsOCwiIiwxLHsibGV2ZWwiOjF9XSxbMCw5LCIiLDEseyJsZXZlbCI6MX1dXQ==
\begin{tikzcd}[row sep=small, column sep=tiny,  every label/.append style = {font=\normalsize}]
	&& {} \\
	&& x \\
	\\
	\\
	& v && {} \\
	\\
	\\
	{}
	\arrow["\cdots"{description}, sloped,,  ,  from=2-3, to=5-2]
	\arrow[""{name=0, anchor=center, inner sep=0}, "\cdots"{description}, sloped,,  ,  draw=none, from=2-3, to=5-4]
	\arrow[""{name=1, anchor=center, inner sep=0}, "v"{description}, draw=none, from=5-2, to=8-1]
	\arrow[""{name=2, anchor=center, inner sep=0}, draw=none, from=2-3, to=0]
	\arrow[""{name=3, anchor=center, inner sep=0}, draw=none, from=5-2, to=1]
	\arrow[from=2-3, to=2]
	\arrow[from=5-2, to=3]
\end{tikzcd}. \\
        Case 2.5.
% https://q.uiver.app/#q=WzAsNCxbMiwzLCJ2Il0sWzEsNiwieCJdLFszLDBdLFswLDNdLFswLDEsIlxcY2RvdHMiLDFdLFsyLDAsInYiLDEseyJzdHlsZSI6eyJib2R5Ijp7Im5hbWUiOiJub25lIn0sImhlYWQiOnsibmFtZSI6Im5vbmUifX19XSxbMywxLCJcXGNkb3RzIiwxLHsic3R5bGUiOnsiYm9keSI6eyJuYW1lIjoibm9uZSJ9LCJoZWFkIjp7Im5hbWUiOiJub25lIn19fV0sWzAsNSwiIiwxLHsic2hvcnRlbiI6eyJ0YXJnZXQiOjIwfSwic3R5bGUiOnsiYm9keSI6eyJuYW1lIjoibm9uZSJ9LCJoZWFkIjp7Im5hbWUiOiJub25lIn19fV0sWzEsNiwiIiwxLHsic2hvcnRlbiI6eyJ0YXJnZXQiOjIwfSwic3R5bGUiOnsiYm9keSI6eyJuYW1lIjoibm9uZSJ9LCJoZWFkIjp7Im5hbWUiOiJub25lIn19fV0sWzcsMCwiIiwxLHsibGV2ZWwiOjF9XSxbOCwxLCIiLDEseyJsZXZlbCI6MX1dXQ==
\begin{tikzcd}[row sep=small, column sep=tiny,  every label/.append style = {font=\normalsize}]
	&&& {} \\
	\\
	\\
	{} && v \\
	\\
	\\
	& x
	\arrow[""{name=0, anchor=center, inner sep=0}, "v"{description}, draw=none, from=1-4, to=4-3]
	\arrow[""{name=1, anchor=center, inner sep=0}, "\cdots"{description}, sloped,,  ,  draw=none, from=4-1, to=7-2]
	\arrow["\cdots"{description}, sloped,,  ,  from=4-3, to=7-2]
	\arrow[""{name=2, anchor=center, inner sep=0}, draw=none, from=4-3, to=0]
	\arrow[""{name=3, anchor=center, inner sep=0}, draw=none, from=7-2, to=1]
	\arrow[from=2, to=4-3]
	\arrow[from=3, to=7-2]
\end{tikzcd}. \quad
        Case 2.6.
% https://q.uiver.app/#q=WzAsNCxbMSwwLCJ2Il0sWzIsMywieCJdLFswLDNdLFszLDBdLFswLDEsIlxcY2RvdHMiLDFdLFswLDIsInYiLDEseyJzdHlsZSI6eyJib2R5Ijp7Im5hbWUiOiJub25lIn0sImhlYWQiOnsibmFtZSI6Im5vbmUifX19XSxbMSwzLCJcXGNkb3RzIiwxLHsic3R5bGUiOnsiYm9keSI6eyJuYW1lIjoibm9uZSJ9LCJoZWFkIjp7Im5hbWUiOiJub25lIn19fV0sWzAsNSwiIiwxLHsic2hvcnRlbiI6eyJ0YXJnZXQiOjIwfSwic3R5bGUiOnsiYm9keSI6eyJuYW1lIjoibm9uZSJ9LCJoZWFkIjp7Im5hbWUiOiJub25lIn19fV0sWzEsNiwiIiwxLHsic2hvcnRlbiI6eyJ0YXJnZXQiOjIwfSwic3R5bGUiOnsiYm9keSI6eyJuYW1lIjoibm9uZSJ9LCJoZWFkIjp7Im5hbWUiOiJub25lIn19fV0sWzAsNywiIiwxLHsibGV2ZWwiOjF9XSxbOCwxLCIiLDEseyJsZXZlbCI6MX1dXQ==
\begin{tikzcd}[row sep=small, column sep=tiny,  every label/.append style = {font=\normalsize}]
	& v && {} \\
	\\
	\\
	{} && x
	\arrow[""{name=0, anchor=center, inner sep=0}, "v"{description}, draw=none, from=1-2, to=4-1]
	\arrow["\cdots"{description}, sloped,,  ,  from=1-2, to=4-3]
	\arrow[""{name=1, anchor=center, inner sep=0}, "\cdots"{description}, sloped,,  ,  draw=none, from=4-3, to=1-4]
	\arrow[""{name=2, anchor=center, inner sep=0}, draw=none, from=1-2, to=0]
	\arrow[""{name=3, anchor=center, inner sep=0}, draw=none, from=4-3, to=1]
	\arrow[from=1-2, to=2]
	\arrow[from=3, to=4-3]
\end{tikzcd}. \quad
        Case 2.7.
% https://q.uiver.app/#q=WzAsOCxbNCwwLCJ4JyJdLFszLDMsInYiXSxbMiw2XSxbMSwwXSxbMiwzLCJPIixbMCwwLDEwMCwxXV0sWzEsMywiTyIsWzAsMCwxMDAsMV1dLFs1LDNdLFswLDZdLFswLDEsIlxcY2RvdHMiLDFdLFsxLDIsInYiLDEseyJzdHlsZSI6eyJib2R5Ijp7Im5hbWUiOiJub25lIn0sImhlYWQiOnsibmFtZSI6Im5vbmUifX19XSxbMyw0LCJ4IiwxLHsic3R5bGUiOnsiYm9keSI6eyJuYW1lIjoibm9uZSJ9LCJoZWFkIjp7Im5hbWUiOiJub25lIn19fV0sWzAsNiwiXFxjZG90cyIsMSx7InN0eWxlIjp7ImJvZHkiOnsibmFtZSI6Im5vbmUifSwiaGVhZCI6eyJuYW1lIjoibm9uZSJ9fX1dLFs1LDcsIiIsMSx7InN0eWxlIjp7ImJvZHkiOnsibmFtZSI6Im5vbmUifSwiaGVhZCI6eyJuYW1lIjoibm9uZSJ9fX1dLFsxLDksIiIsMSx7InNob3J0ZW4iOnsidGFyZ2V0IjoyMH0sInN0eWxlIjp7ImJvZHkiOnsibmFtZSI6Im5vbmUifSwiaGVhZCI6eyJuYW1lIjoibm9uZSJ9fX1dLFsxMCw5LCJcXGNkb3RzIiwxLHsic2hvcnRlbiI6eyJzb3VyY2UiOjIwLCJ0YXJnZXQiOjIwfSwic3R5bGUiOnsiYm9keSI6eyJuYW1lIjoibm9uZSJ9LCJoZWFkIjp7Im5hbWUiOiJub25lIn19fV0sWzQsOSwiIiwxLHsibGV2ZWwiOjEsInN0eWxlIjp7ImJvZHkiOnsibmFtZSI6Im5vbmUifSwiaGVhZCI6eyJuYW1lIjoibm9uZSJ9fX1dLFs0LDEwLCIiLDEseyJzaG9ydGVuIjp7InRhcmdldCI6MjB9LCJzdHlsZSI6eyJib2R5Ijp7Im5hbWUiOiJub25lIn0sImhlYWQiOnsibmFtZSI6Im5vbmUifX19XSxbMTAsNSwiIiwxLHsic2hvcnRlbiI6eyJzb3VyY2UiOjIwfSwic3R5bGUiOnsiYm9keSI6eyJuYW1lIjoibm9uZSJ9LCJoZWFkIjp7Im5hbWUiOiJub25lIn19fV0sWzAsMTEsIiIsMSx7InNob3J0ZW4iOnsidGFyZ2V0IjoyMH0sInN0eWxlIjp7ImJvZHkiOnsibmFtZSI6Im5vbmUifSwiaGVhZCI6eyJuYW1lIjoibm9uZSJ9fX1dLFsxMCwxMiwiXFxjZG90cyIsMSx7InNob3J0ZW4iOnsic291cmNlIjoyMCwidGFyZ2V0IjoyMH0sInN0eWxlIjp7ImJvZHkiOnsibmFtZSI6Im5vbmUifSwiaGVhZCI6eyJuYW1lIjoibm9uZSJ9fX1dLFsxLDEzLCIiLDEseyJsZXZlbCI6MX1dLFs0LDE1LCIiLDEseyJsZXZlbCI6MX1dLFsxNiw0LCIiLDEseyJsZXZlbCI6MSwic3R5bGUiOnsiaGVhZCI6eyJuYW1lIjoibm9uZSJ9fX1dLFsxNyw1LCIiLDEseyJsZXZlbCI6MX1dLFswLDE4LCIiLDEseyJsZXZlbCI6MX1dXQ==
\begin{tikzcd}[row sep=small, column sep=tiny,  every label/.append style = {font=\normalsize}]
	& {} &&& {x'} \\
	\\
	\\
	& \textcolor{white}{O} & \textcolor{white}{O} & v && {} \\
	\\
	\\
	{} && {}
	\arrow[""{name=0, anchor=center, inner sep=0}, "x"{description}, draw=none, from=1-2, to=4-3]
	\arrow["\cdots"{description}, sloped,,  ,  from=1-5, to=4-4]
	\arrow[""{name=1, anchor=center, inner sep=0}, "\cdots"{description}, sloped,,  ,  draw=none, from=1-5, to=4-6]
	\arrow[""{name=2, anchor=center, inner sep=0}, draw=none, from=4-2, to=7-1]
	\arrow[""{name=3, anchor=center, inner sep=0}, "v"{description}, draw=none, from=4-4, to=7-3]
	\arrow[""{name=4, anchor=center, inner sep=0}, draw=none, from=0, to=4-2]
	\arrow["\cdots"{description}, sloped,,  ,  draw=none, from=0, to=2]
	\arrow["\cdots"{description}, sloped,,  ,  draw=none, from=0, to=3]
	\arrow[""{name=5, anchor=center, inner sep=0}, draw=none, from=1-5, to=1]
	\arrow[""{name=6, anchor=center, inner sep=0}, draw=none, from=4-3, to=0]
	\arrow[""{name=7, anchor=center, inner sep=0}, draw=none, from=4-3, to=3]
	\arrow[""{name=8, anchor=center, inner sep=0}, draw=none, from=4-4, to=3]
	\arrow[from=4, to=4-2]
	\arrow[from=1-5, to=5]
	\arrow[no head, from=6, to=4-3]
	\arrow[from=4-3, to=7]
	\arrow[from=4-4, to=8]
\end{tikzcd}. \\
        Case 2.8.
% https://q.uiver.app/#q=WzAsNyxbMiw2LCJ2Il0sWzMsMywiTyIsWzAsMCwxMDAsMV1dLFs0LDBdLFsxLDksIngiXSxbMCw2XSxbNCwzLCJPIixbMCwwLDEwMCwxXV0sWzUsNl0sWzAsMSwidiIsMSx7InN0eWxlIjp7ImJvZHkiOnsibmFtZSI6Im5vbmUifSwiaGVhZCI6eyJuYW1lIjoibm9uZSJ9fX1dLFsxLDIsIngnIiwxLHsic3R5bGUiOnsiYm9keSI6eyJuYW1lIjoibm9uZSJ9LCJoZWFkIjp7Im5hbWUiOiJub25lIn19fV0sWzAsMywiXFxjZG90cyIsMV0sWzMsNCwiXFxjZG90cyIsMSx7InN0eWxlIjp7ImJvZHkiOnsibmFtZSI6Im5vbmUifSwiaGVhZCI6eyJuYW1lIjoibm9uZSJ9fX1dLFs1LDYsIiIsMSx7InN0eWxlIjp7ImJvZHkiOnsibmFtZSI6Im5vbmUifSwiaGVhZCI6eyJuYW1lIjoibm9uZSJ9fX1dLFsxLDgsIiIsMSx7InNob3J0ZW4iOnsidGFyZ2V0IjoyMH0sInN0eWxlIjp7ImJvZHkiOnsibmFtZSI6Im5vbmUifSwiaGVhZCI6eyJuYW1lIjoibm9uZSJ9fX1dLFsxLDcsIiIsMSx7InNob3J0ZW4iOnsidGFyZ2V0IjoyMH0sInN0eWxlIjp7ImJvZHkiOnsibmFtZSI6Im5vbmUifSwiaGVhZCI6eyJuYW1lIjoibm9uZSJ9fX1dLFs3LDgsIlxcY2RvdHMiLDEseyJzaG9ydGVuIjp7InNvdXJjZSI6MjAsInRhcmdldCI6MjB9LCJzdHlsZSI6eyJib2R5Ijp7Im5hbWUiOiJub25lIn0sImhlYWQiOnsibmFtZSI6Im5vbmUifX19XSxbMCw3LCIiLDEseyJzaG9ydGVuIjp7InRhcmdldCI6MjB9LCJzdHlsZSI6eyJib2R5Ijp7Im5hbWUiOiJub25lIn0sImhlYWQiOnsibmFtZSI6Im5vbmUifX19XSxbMywxMCwiIiwxLHsic2hvcnRlbiI6eyJ0YXJnZXQiOjIwfSwic3R5bGUiOnsiYm9keSI6eyJuYW1lIjoibm9uZSJ9LCJoZWFkIjp7Im5hbWUiOiJub25lIn19fV0sWzgsMTEsIlxcY2RvdHMiLDEseyJzaG9ydGVuIjp7InNvdXJjZSI6MjAsInRhcmdldCI6MjB9LCJzdHlsZSI6eyJib2R5Ijp7Im5hbWUiOiJub25lIn0sImhlYWQiOnsibmFtZSI6Im5vbmUifX19XSxbOCw1LCIiLDEseyJzaG9ydGVuIjp7InNvdXJjZSI6MjB9LCJzdHlsZSI6eyJib2R5Ijp7Im5hbWUiOiJub25lIn0sImhlYWQiOnsibmFtZSI6Im5vbmUifX19XSxbMTIsMSwiIiwxLHsibGV2ZWwiOjEsInN0eWxlIjp7ImhlYWQiOnsibmFtZSI6Im5vbmUifX19XSxbMSwxMywiIiwxLHsibGV2ZWwiOjF9XSxbMTUsMCwiIiwxLHsibGV2ZWwiOjF9XSxbMTYsMywiIiwxLHsibGV2ZWwiOjF9XSxbMTgsNSwiIiwxLHsibGV2ZWwiOjF9XV0=
\begin{tikzcd}[row sep=small, column sep=tiny,  every label/.append style = {font=\normalsize}]
	&&&& {} \\
	\\
	\\
	&&& \textcolor{white}{O} & \textcolor{white}{O} \\
	\\
	\\
	{} && v &&& {} \\
	\\
	\\
	& x
	\arrow[""{name=0, anchor=center, inner sep=0}, "{x'}"{description}, draw=none, from=4-4, to=1-5]
	\arrow[""{name=1, anchor=center, inner sep=0}, draw=none, from=4-5, to=7-6]
	\arrow[""{name=2, anchor=center, inner sep=0}, "v"{description}, draw=none, from=7-3, to=4-4]
	\arrow["\cdots"{description}, sloped,,  ,  from=7-3, to=10-2]
	\arrow[""{name=3, anchor=center, inner sep=0}, "\cdots"{description}, sloped,,  ,  draw=none, from=10-2, to=7-1]
	\arrow[""{name=4, anchor=center, inner sep=0}, draw=none, from=4-4, to=0]
	\arrow[""{name=5, anchor=center, inner sep=0}, draw=none, from=4-4, to=2]
	\arrow[""{name=6, anchor=center, inner sep=0}, draw=none, from=0, to=4-5]
	\arrow["\cdots"{description}, sloped,,  ,  draw=none, from=0, to=1]
	\arrow["\cdots"{description}, sloped,,  ,  draw=none, from=2, to=0]
	\arrow[""{name=7, anchor=center, inner sep=0}, draw=none, from=7-3, to=2]
	\arrow[""{name=8, anchor=center, inner sep=0}, draw=none, from=10-2, to=3]
	\arrow[no head, from=4, to=4-4]
	\arrow[from=4-4, to=5]
	\arrow[from=6, to=4-5]
	\arrow[from=7, to=7-3]
	\arrow[from=8, to=10-2]
\end{tikzcd}. \quad
        Case 2.9.
% https://q.uiver.app/#q=WzAsNyxbMywwLCJ2Il0sWzQsMywieCciXSxbMiwzLCJPIixbMCwwLDEwMCwxXV0sWzEsNl0sWzEsMywiTyIsWzAsMCwxMDAsMV1dLFswLDBdLFs1LDBdLFswLDEsIlxcY2RvdHMiLDFdLFswLDIsInYiLDEseyJzdHlsZSI6eyJib2R5Ijp7Im5hbWUiOiJub25lIn0sImhlYWQiOnsibmFtZSI6Im5vbmUifX19XSxbMiwzLCJ4IiwxLHsic3R5bGUiOnsiYm9keSI6eyJuYW1lIjoibm9uZSJ9LCJoZWFkIjp7Im5hbWUiOiJub25lIn19fV0sWzUsNCwiIiwxLHsic3R5bGUiOnsiYm9keSI6eyJuYW1lIjoibm9uZSJ9LCJoZWFkIjp7Im5hbWUiOiJub25lIn19fV0sWzYsMSwiXFxjZG90cyIsMSx7InN0eWxlIjp7ImJvZHkiOnsibmFtZSI6Im5vbmUifSwiaGVhZCI6eyJuYW1lIjoibm9uZSJ9fX1dLFswLDgsIiIsMSx7InNob3J0ZW4iOnsidGFyZ2V0IjoyMH0sInN0eWxlIjp7ImJvZHkiOnsibmFtZSI6Im5vbmUifSwiaGVhZCI6eyJuYW1lIjoibm9uZSJ9fX1dLFsyLDksIiIsMSx7InNob3J0ZW4iOnsidGFyZ2V0IjoyMH0sInN0eWxlIjp7ImJvZHkiOnsibmFtZSI6Im5vbmUifSwiaGVhZCI6eyJuYW1lIjoibm9uZSJ9fX1dLFsyLDgsIiIsMSx7InNob3J0ZW4iOnsidGFyZ2V0IjoyMH0sInN0eWxlIjp7ImJvZHkiOnsibmFtZSI6Im5vbmUifSwiaGVhZCI6eyJuYW1lIjoibm9uZSJ9fX1dLFs0LDksIiIsMSx7InNob3J0ZW4iOnsidGFyZ2V0IjoyMH0sInN0eWxlIjp7ImJvZHkiOnsibmFtZSI6Im5vbmUifSwiaGVhZCI6eyJuYW1lIjoibm9uZSJ9fX1dLFsxMCw5LCJcXGNkb3RzIiwxLHsic2hvcnRlbiI6eyJzb3VyY2UiOjIwLCJ0YXJnZXQiOjIwfSwic3R5bGUiOnsiYm9keSI6eyJuYW1lIjoibm9uZSJ9LCJoZWFkIjp7Im5hbWUiOiJub25lIn19fV0sWzExLDEsIiIsMSx7InNob3J0ZW4iOnsic291cmNlIjoyMH0sInN0eWxlIjp7ImJvZHkiOnsibmFtZSI6Im5vbmUifSwiaGVhZCI6eyJuYW1lIjoibm9uZSJ9fX1dLFswLDEyLCIiLDEseyJsZXZlbCI6MX1dLFsyLDEzLCIiLDEseyJsZXZlbCI6MX1dLFsyLDE0LCIiLDEseyJsZXZlbCI6MSwic3R5bGUiOnsiaGVhZCI6eyJuYW1lIjoibm9uZSJ9fX1dLFsxNCwxMywiXFxjZG90cyIsMSx7InNob3J0ZW4iOnsic291cmNlIjoyMCwidGFyZ2V0IjoyMH0sInN0eWxlIjp7ImJvZHkiOnsibmFtZSI6Im5vbmUifSwiaGVhZCI6eyJuYW1lIjoibm9uZSJ9fX1dLFs0LDE1LCIiLDEseyJsZXZlbCI6MX1dLFsxNywxLCIiLDEseyJsZXZlbCI6MX1dXQ==
\begin{tikzcd}[row sep=small, column sep=tiny,  every label/.append style = {font=\normalsize}]
	{} &&& v && {} \\
	\\
	\\
	& \textcolor{white}{O} & \textcolor{white}{O} && {x'} \\
	\\
	\\
	& {}
	\arrow[""{name=0, anchor=center, inner sep=0}, draw=none, from=1-1, to=4-2]
	\arrow[""{name=1, anchor=center, inner sep=0}, "v"{description}, draw=none, from=1-4, to=4-3]
	\arrow["\cdots"{description}, sloped,,  ,  from=1-4, to=4-5]
	\arrow[""{name=2, anchor=center, inner sep=0}, "\cdots"{description}, sloped,,  ,  draw=none, from=1-6, to=4-5]
	\arrow[""{name=3, anchor=center, inner sep=0}, "x"{description}, draw=none, from=4-3, to=7-2]
	\arrow["\cdots"{description}, sloped,,  ,  draw=none, from=0, to=3]
	\arrow[""{name=4, anchor=center, inner sep=0}, draw=none, from=1-4, to=1]
	\arrow[""{name=5, anchor=center, inner sep=0}, draw=none, from=2, to=4-5]
	\arrow[""{name=6, anchor=center, inner sep=0}, draw=none, from=4-2, to=3]
	\arrow[""{name=7, anchor=center, inner sep=0}, draw=none, from=4-3, to=1]
	\arrow[""{name=8, anchor=center, inner sep=0}, draw=none, from=4-3, to=3]
	\arrow[from=1-4, to=4]
	\arrow[from=5, to=4-5]
	\arrow[from=4-2, to=6]
	\arrow[no head, from=4-3, to=7]
	\arrow[from=4-3, to=8]
	\arrow["\cdots"{description}, sloped,,  ,  draw=none, from=7, to=8]
\end{tikzcd}. \\
\ \\
        Case 2.10.
% https://q.uiver.app/#q=WzAsOCxbMiwzLCJ2Il0sWzEsMCwieCJdLFszLDBdLFszLDZdLFs0LDMsIk8iLFswLDAsMTAwLDFdXSxbMywzLCJPIixbMCwwLDEwMCwxXV0sWzUsMF0sWzAsM10sWzEsMCwiXFxjZG90cyIsMV0sWzIsMCwidiIsMSx7InN0eWxlIjp7ImJvZHkiOnsibmFtZSI6Im5vbmUifSwiaGVhZCI6eyJuYW1lIjoibm9uZSJ9fX1dLFszLDQsIngnIiwxLHsic3R5bGUiOnsiYm9keSI6eyJuYW1lIjoibm9uZSJ9LCJoZWFkIjp7Im5hbWUiOiJub25lIn19fV0sWzYsNCwiIiwxLHsic3R5bGUiOnsiYm9keSI6eyJuYW1lIjoibm9uZSJ9LCJoZWFkIjp7Im5hbWUiOiJub25lIn19fV0sWzEsNywiXFxjZG90cyIsMSx7InN0eWxlIjp7ImJvZHkiOnsibmFtZSI6Im5vbmUifSwiaGVhZCI6eyJuYW1lIjoibm9uZSJ9fX1dLFs5LDAsIiIsMSx7InN0eWxlIjp7ImJvZHkiOnsibmFtZSI6Im5vbmUifSwiaGVhZCI6eyJuYW1lIjoibm9uZSJ9fX1dLFs1LDksIiIsMSx7InNob3J0ZW4iOnsidGFyZ2V0IjoyMH0sInN0eWxlIjp7ImJvZHkiOnsibmFtZSI6Im5vbmUifSwiaGVhZCI6eyJuYW1lIjoibm9uZSJ9fX1dLFs1LDEwLCIiLDEseyJzaG9ydGVuIjp7InRhcmdldCI6MjB9LCJzdHlsZSI6eyJib2R5Ijp7Im5hbWUiOiJub25lIn0sImhlYWQiOnsibmFtZSI6Im5vbmUifX19XSxbOSwxMCwiXFxjZG90cyIsMSx7InNob3J0ZW4iOnsic291cmNlIjoyMCwidGFyZ2V0IjoyMH0sInN0eWxlIjp7ImJvZHkiOnsibmFtZSI6Im5vbmUifSwiaGVhZCI6eyJuYW1lIjoibm9uZSJ9fX1dLFsxMSwxMCwiXFxjZG90cyIsMSx7InNob3J0ZW4iOnsic291cmNlIjoyMCwidGFyZ2V0IjoyMH0sInN0eWxlIjp7ImJvZHkiOnsibmFtZSI6Im5vbmUifSwiaGVhZCI6eyJuYW1lIjoibm9uZSJ9fX1dLFs0LDEwLCIiLDEseyJzaG9ydGVuIjp7InRhcmdldCI6MjB9LCJzdHlsZSI6eyJib2R5Ijp7Im5hbWUiOiJub25lIn0sImhlYWQiOnsibmFtZSI6Im5vbmUifX19XSxbMSwxMiwiIiwxLHsic2hvcnRlbiI6eyJ0YXJnZXQiOjIwfSwic3R5bGUiOnsiYm9keSI6eyJuYW1lIjoibm9uZSJ9LCJoZWFkIjp7Im5hbWUiOiJub25lIn19fV0sWzEzLDAsIiIsMSx7ImxldmVsIjoxfV0sWzE0LDUsIiIsMSx7ImxldmVsIjoxLCJzdHlsZSI6eyJoZWFkIjp7Im5hbWUiOiJub25lIn19fV0sWzUsMTUsIiIsMSx7ImxldmVsIjoxfV0sWzQsMTgsIiIsMSx7ImxldmVsIjoxfV0sWzEsMTksIiIsMSx7ImxldmVsIjoxfV1d
\begin{tikzcd}[row sep=small, column sep=tiny,  every label/.append style = {font=\normalsize}]
	& x && {} && {} \\
	\\
	\\
	{} && v & \textcolor{white}{O} & \textcolor{white}{O} \\
	\\
	\\
	&&& {}
	\arrow[""{name=0, anchor=center, inner sep=0}, "\cdots"{description}, sloped,,  ,  draw=none, from=1-2, to=4-1]
	\arrow["\cdots"{description}, sloped,,  ,  from=1-2, to=4-3]
	\arrow[""{name=1, anchor=center, inner sep=0}, "v"{description}, draw=none, from=1-4, to=4-3]
	\arrow[""{name=2, anchor=center, inner sep=0}, draw=none, from=1-6, to=4-5]
	\arrow[""{name=3, anchor=center, inner sep=0}, "{x'}"{description}, draw=none, from=7-4, to=4-5]
	\arrow[""{name=4, anchor=center, inner sep=0}, draw=none, from=1-2, to=0]
	\arrow[""{name=5, anchor=center, inner sep=0}, draw=none, from=1, to=4-3]
	\arrow["\cdots"{description}, sloped,,  ,  draw=none, from=1, to=3]
	\arrow["\cdots"{description}, sloped,,  ,  draw=none, from=2, to=3]
	\arrow[""{name=6, anchor=center, inner sep=0}, draw=none, from=4-4, to=1]
	\arrow[""{name=7, anchor=center, inner sep=0}, draw=none, from=4-4, to=3]
	\arrow[""{name=8, anchor=center, inner sep=0}, draw=none, from=4-5, to=3]
	\arrow[from=1-2, to=4]
	\arrow[from=5, to=4-3]
	\arrow[no head, from=6, to=4-4]
	\arrow[from=4-4, to=7]
	\arrow[from=4-5, to=8]
\end{tikzcd}.

		{\bf Case 3.} $N$ contains three vertices that are peaks or deeps of $Z$, that is, $N=$
% https://q.uiver.app/#q=WzAsMyxbMSwwLCJ2Il0sWzIsMywidiJdLFswLDNdLFswLDEsIlxcY2RvdHMiLDFdLFswLDIsInYiLDEseyJzdHlsZSI6eyJib2R5Ijp7Im5hbWUiOiJub25lIn0sImhlYWQiOnsibmFtZSI6Im5vbmUifX19XSxbMCw0LCIiLDEseyJzaG9ydGVuIjp7InRhcmdldCI6MjB9LCJzdHlsZSI6eyJib2R5Ijp7Im5hbWUiOiJub25lIn0sImhlYWQiOnsibmFtZSI6Im5vbmUifX19XSxbMCw1LCIiLDEseyJsZXZlbCI6MX1dXQ==
\begin{tikzcd}[row sep=small, column sep=tiny,  every label/.append style = {font=\normalsize}]
	& v \\
	\\
	\\
	{} && v
	\arrow[""{name=0, anchor=center, inner sep=0}, "v"{description}, draw=none, from=1-2, to=4-1]
	\arrow["\cdots"{description}, sloped,,  ,  from=1-2, to=4-3]
	\arrow[""{name=1, anchor=center, inner sep=0}, draw=none, from=1-2, to=0]
	\arrow[from=1-2, to=1]
\end{tikzcd} or
% https://q.uiver.app/#q=WzAsMyxbMCwwLCJ2Il0sWzEsMywidiJdLFsyLDBdLFswLDEsIlxcY2RvdHMiLDFdLFsyLDEsInYiLDEseyJzdHlsZSI6eyJib2R5Ijp7Im5hbWUiOiJub25lIn0sImhlYWQiOnsibmFtZSI6Im5vbmUifX19XSxbMSw0LCIiLDEseyJzaG9ydGVuIjp7InRhcmdldCI6MjB9LCJzdHlsZSI6eyJib2R5Ijp7Im5hbWUiOiJub25lIn0sImhlYWQiOnsibmFtZSI6Im5vbmUifX19XSxbNSwxLCIiLDEseyJsZXZlbCI6MX1dXQ==
\begin{tikzcd}[row sep=small, column sep=tiny,  every label/.append style = {font=\normalsize}]
	v && {} \\
	\\
	\\
	& v
	\arrow["\cdots"{description}, sloped,,  ,  from=1-1, to=4-2]
	\arrow[""{name=0, anchor=center, inner sep=0}, "v"{description}, draw=none, from=1-3, to=4-2]
	\arrow[""{name=1, anchor=center, inner sep=0}, draw=none, from=4-2, to=0]
	\arrow[from=1, to=4-2]
\end{tikzcd}. After possibly taking the flip of $Y$, the local configuration of $Y$ containing $N$ as a subdiagram can be classified into the following cases: \\
		Case 3.1.
\begin{tikzcd}[row sep=small, column sep=tiny,  every label/.append style = {font=\normalsize}]
	& v \\
	\\
	\\
	{} && v
	\arrow[""{name=0, anchor=center, inner sep=0}, "v"{description}, draw=none, from=1-2, to=4-1]
	\arrow["\cdots"{description}, sloped,,  ,  from=1-2, to=4-3]
	\arrow[""{name=1, anchor=center, inner sep=0}, draw=none, from=1-2, to=0]
	\arrow[from=1-2, to=1]
\end{tikzcd}. \quad
		Case 3.2.
% https://q.uiver.app/#q=WzAsNyxbMSwwLCJ4Il0sWzIsMywidiJdLFswLDAsIk8iLFswLDAsMTAwLDFdXSxbMCwzXSxbMywwXSxbMywzLCJPIixbMCwwLDEwMCwxXV0sWzQsNl0sWzAsMSwiXFxjZG90cyIsMV0sWzAsMywiXFxjZG90cyIsMSx7InN0eWxlIjp7ImJvZHkiOnsibmFtZSI6Im5vbmUifSwiaGVhZCI6eyJuYW1lIjoibm9uZSJ9fX1dLFs0LDEsInYiLDEseyJzdHlsZSI6eyJib2R5Ijp7Im5hbWUiOiJub25lIn0sImhlYWQiOnsibmFtZSI6Im5vbmUifX19XSxbNSw2LCJ2IiwxLHsic3R5bGUiOnsiYm9keSI6eyJuYW1lIjoibm9uZSJ9LCJoZWFkIjp7Im5hbWUiOiJub25lIn19fV0sWzAsOCwiIiwxLHsic2hvcnRlbiI6eyJ0YXJnZXQiOjIwfSwic3R5bGUiOnsiYm9keSI6eyJuYW1lIjoibm9uZSJ9LCJoZWFkIjp7Im5hbWUiOiJub25lIn19fV0sWzEsOSwiIiwxLHsic2hvcnRlbiI6eyJ0YXJnZXQiOjIwfSwic3R5bGUiOnsiYm9keSI6eyJuYW1lIjoibm9uZSJ9LCJoZWFkIjp7Im5hbWUiOiJub25lIn19fV0sWzksMTAsIiIsMSx7InN0eWxlIjp7ImJvZHkiOnsibmFtZSI6Im5vbmUifSwiaGVhZCI6eyJuYW1lIjoibm9uZSJ9fX1dLFswLDExLCIiLDEseyJsZXZlbCI6MX1dLFsxMiwxLCIiLDEseyJsZXZlbCI6MX1dLFsxMCwxMywiIiwxLHsic2hvcnRlbiI6eyJzb3VyY2UiOjIwLCJ0YXJnZXQiOjIwfSwic3R5bGUiOnsiYm9keSI6eyJuYW1lIjoibm9uZSJ9LCJoZWFkIjp7Im5hbWUiOiJub25lIn19fV0sWzksMTMsIiIsMSx7InNob3J0ZW4iOnsic291cmNlIjoyMCwidGFyZ2V0IjoyMH0sInN0eWxlIjp7ImJvZHkiOnsibmFtZSI6Im5vbmUifSwiaGVhZCI6eyJuYW1lIjoibm9uZSJ9fX1dLFsxNywxNiwiXFxjZG90cyIsMSx7ImxldmVsIjoxfV1d
\begin{tikzcd}[row sep=small, column sep=tiny,  every label/.append style = {font=\normalsize}]
	\textcolor{white}{O} & x && {} \\
	\\
	\\
	{} && v & \textcolor{white}{O} \\
	\\
	\\
	&&&& {}
	\arrow[""{name=0, anchor=center, inner sep=0}, "\cdots"{description}, sloped,,  ,  draw=none, from=1-2, to=4-1]
	\arrow["\cdots"{description}, sloped,,  ,  from=1-2, to=4-3]
	\arrow[""{name=1, anchor=center, inner sep=0}, "v"{description}, draw=none, from=1-4, to=4-3]
	\arrow[""{name=2, anchor=center, inner sep=0}, "v"{description}, draw=none, from=4-4, to=7-5]
	\arrow[""{name=3, anchor=center, inner sep=0}, draw=none, from=1-2, to=0]
	\arrow[""{name=4, anchor=center, inner sep=0}, draw=none, from=1, to=2]
	\arrow[""{name=5, anchor=center, inner sep=0}, draw=none, from=4-3, to=1]
	\arrow[from=1-2, to=3]
	\arrow[""{name=6, anchor=center, inner sep=0}, draw=none, from=1, to=4]
	\arrow[from=5, to=4-3]
	\arrow[""{name=7, anchor=center, inner sep=0}, draw=none, from=2, to=4]
	\arrow["\cdots"{description}, sloped,,  ,  from=6, to=7]
\end{tikzcd}. \quad
		Case 3.3.
% https://q.uiver.app/#q=WzAsOCxbMSw2LCJ4Il0sWzIsMywidiJdLFswLDAsIk8iLFswLDAsMTAwLDFdXSxbMCwzXSxbMywzLCJPIixbMCwwLDEwMCwxXV0sWzQsNl0sWzMsMF0sWzAsOF0sWzEsMCwiXFxjZG90cyIsMV0sWzAsMywiXFxjZG90cyIsMSx7InN0eWxlIjp7ImJvZHkiOnsibmFtZSI6Im5vbmUifSwiaGVhZCI6eyJuYW1lIjoibm9uZSJ9fX1dLFs0LDUsInYiLDEseyJzdHlsZSI6eyJib2R5Ijp7Im5hbWUiOiJub25lIn0sImhlYWQiOnsibmFtZSI6Im5vbmUifX19XSxbNiwxLCJ2IiwxLHsic3R5bGUiOnsiYm9keSI6eyJuYW1lIjoibm9uZSJ9LCJoZWFkIjp7Im5hbWUiOiJub25lIn19fV0sWzAsOSwiIiwxLHsic2hvcnRlbiI6eyJ0YXJnZXQiOjIwfSwic3R5bGUiOnsiYm9keSI6eyJuYW1lIjoibm9uZSJ9LCJoZWFkIjp7Im5hbWUiOiJub25lIn19fV0sWzEsMTEsIiIsMSx7InNob3J0ZW4iOnsidGFyZ2V0IjoyMH0sInN0eWxlIjp7ImJvZHkiOnsibmFtZSI6Im5vbmUifSwiaGVhZCI6eyJuYW1lIjoibm9uZSJ9fX1dLFsxMSwxMCwiIiwxLHsic3R5bGUiOnsiYm9keSI6eyJuYW1lIjoibm9uZSJ9LCJoZWFkIjp7Im5hbWUiOiJub25lIn19fV0sWzEyLDAsIiIsMSx7ImxldmVsIjoxfV0sWzExLDE0LCIiLDEseyJzaG9ydGVuIjp7InNvdXJjZSI6MjAsInRhcmdldCI6MjB9LCJzdHlsZSI6eyJib2R5Ijp7Im5hbWUiOiJub25lIn0sImhlYWQiOnsibmFtZSI6Im5vbmUifX19XSxbMTMsMSwiIiwxLHsibGV2ZWwiOjF9XSxbMTAsMTQsIiIsMSx7InNob3J0ZW4iOnsic291cmNlIjoyMCwidGFyZ2V0IjoyMH0sInN0eWxlIjp7ImJvZHkiOnsibmFtZSI6Im5vbmUifSwiaGVhZCI6eyJuYW1lIjoibm9uZSJ9fX1dLFsxNiwxOCwiXFxjZG90cyIsMSx7ImxldmVsIjoxfV1d
\begin{tikzcd}[row sep=small, column sep=tiny,  every label/.append style = {font=\normalsize}]
	\textcolor{white}{O} &&& {} \\
	\\
	\\
	{} && v & \textcolor{white}{O} \\
	\\
	\\
	& x &&& {} \\
	\\
	{}
	\arrow[""{name=0, anchor=center, inner sep=0}, "v"{description}, draw=none, from=1-4, to=4-3]
	\arrow["\cdots"{description}, sloped,,  ,  from=4-3, to=7-2]
	\arrow[""{name=1, anchor=center, inner sep=0}, "v"{description}, draw=none, from=4-4, to=7-5]
	\arrow[""{name=2, anchor=center, inner sep=0}, "\cdots"{description}, sloped,,  ,  draw=none, from=7-2, to=4-1]
	\arrow[""{name=3, anchor=center, inner sep=0}, draw=none, from=0, to=1]
	\arrow[""{name=4, anchor=center, inner sep=0}, draw=none, from=4-3, to=0]
	\arrow[""{name=5, anchor=center, inner sep=0}, draw=none, from=7-2, to=2]
	\arrow[""{name=6, anchor=center, inner sep=0}, draw=none, from=0, to=3]
	\arrow[from=4, to=4-3]
	\arrow[""{name=7, anchor=center, inner sep=0}, draw=none, from=1, to=3]
	\arrow[from=5, to=7-2]
	\arrow["\cdots"{description}, sloped,,  ,  from=6, to=7]
\end{tikzcd}.
\\
		Case 3.4.
\begin{tikzcd}[row sep=small, column sep=tiny,  every label/.append style = {font=\normalsize}]
	v && {} \\
	\\
	\\
	& v
	\arrow["\cdots"{description}, sloped,,  ,  from=1-1, to=4-2]
	\arrow[""{name=0, anchor=center, inner sep=0}, "v"{description}, draw=none, from=1-3, to=4-2]
	\arrow[""{name=1, anchor=center, inner sep=0}, draw=none, from=4-2, to=0]
	\arrow[from=1, to=4-2]
\end{tikzcd}. \quad
		Case 3.5.
% https://q.uiver.app/#q=WzAsNyxbMCwzLCJ2Il0sWzEsNiwidiJdLFsyLDMsIk8iLFswLDAsMTAwLDFdXSxbMywwXSxbMywzLCJPIixbMCwwLDEwMCwxXV0sWzQsNl0sWzEsN10sWzAsMSwiXFxjZG90cyIsMV0sWzIsMSwidiIsMSx7InN0eWxlIjp7ImJvZHkiOnsibmFtZSI6Im5vbmUifSwiaGVhZCI6eyJuYW1lIjoibm9uZSJ9fX1dLFszLDIsIngiLDEseyJzdHlsZSI6eyJib2R5Ijp7Im5hbWUiOiJub25lIn0sImhlYWQiOnsibmFtZSI6Im5vbmUifX19XSxbNCw1LCIiLDEseyJzdHlsZSI6eyJib2R5Ijp7Im5hbWUiOiJub25lIn0sImhlYWQiOnsibmFtZSI6Im5vbmUifX19XSxbMSw4LCIiLDEseyJzaG9ydGVuIjp7InRhcmdldCI6MjB9LCJzdHlsZSI6eyJib2R5Ijp7Im5hbWUiOiJub25lIn0sImhlYWQiOnsibmFtZSI6Im5vbmUifX19XSxbOSwyLCIiLDEseyJzaG9ydGVuIjp7InNvdXJjZSI6MjB9LCJzdHlsZSI6eyJib2R5Ijp7Im5hbWUiOiJub25lIn0sImhlYWQiOnsibmFtZSI6Im5vbmUifX19XSxbOCwyLCIiLDEseyJzaG9ydGVuIjp7InNvdXJjZSI6MjB9LCJzdHlsZSI6eyJib2R5Ijp7Im5hbWUiOiJub25lIn0sImhlYWQiOnsibmFtZSI6Im5vbmUifX19XSxbOSwxMCwiXFxjZG90cyIsMSx7InNob3J0ZW4iOnsic291cmNlIjoyMCwidGFyZ2V0IjoyMH0sInN0eWxlIjp7ImJvZHkiOnsibmFtZSI6Im5vbmUifSwiaGVhZCI6eyJuYW1lIjoibm9uZSJ9fX1dLFs0LDksIiIsMSx7InNob3J0ZW4iOnsidGFyZ2V0IjoyMH0sInN0eWxlIjp7ImJvZHkiOnsibmFtZSI6Im5vbmUifSwiaGVhZCI6eyJuYW1lIjoibm9uZSJ9fX1dLFsxMSwxLCIiLDEseyJsZXZlbCI6MX1dLFsxMiwxMywiXFxjZG90cyIsMSx7ImxldmVsIjoxfV0sWzE1LDQsIiIsMSx7ImxldmVsIjoxfV1d
\begin{tikzcd}[row sep=small, column sep=tiny,  every label/.append style = {font=\normalsize}]
	&&& {} \\
	\\
	\\
	v && \textcolor{white}{O} & \textcolor{white}{O} \\
	\\
	\\
	& v &&& {} \\
	& {}
	\arrow[""{name=0, anchor=center, inner sep=0}, "x"{description}, draw=none, from=1-4, to=4-3]
	\arrow["\cdots"{description}, sloped,,  ,  from=4-1, to=7-2]
	\arrow[""{name=1, anchor=center, inner sep=0}, "v"{description}, draw=none, from=4-3, to=7-2]
	\arrow[""{name=2, anchor=center, inner sep=0}, draw=none, from=4-4, to=7-5]
	\arrow[""{name=3, anchor=center, inner sep=0}, draw=none, from=0, to=4-3]
	\arrow["\cdots"{description}, sloped,,  ,  draw=none, from=0, to=2]
	\arrow[""{name=4, anchor=center, inner sep=0}, draw=none, from=1, to=4-3]
	\arrow[""{name=5, anchor=center, inner sep=0}, draw=none, from=4-4, to=0]
	\arrow[""{name=6, anchor=center, inner sep=0}, draw=none, from=7-2, to=1]
	\arrow["\cdots"{description}, sloped,,  ,  from=3, to=4]
	\arrow[from=5, to=4-4]
	\arrow[from=6, to=7-2]
\end{tikzcd}. \quad
		Case 3.6.
% https://q.uiver.app/#q=WzAsOCxbMCwxLCJ2Il0sWzEsNCwidiJdLFsyLDQsIk8iLFswLDAsMTAwLDFdXSxbMyw3XSxbMyw0LCJPIixbMCwwLDEwMCwxXV0sWzQsMV0sWzIsMV0sWzAsMF0sWzAsMSwiXFxjZG90cyIsMV0sWzMsMiwieCIsMSx7InN0eWxlIjp7ImJvZHkiOnsibmFtZSI6Im5vbmUifSwiaGVhZCI6eyJuYW1lIjoibm9uZSJ9fX1dLFs0LDUsIiIsMSx7InN0eWxlIjp7ImJvZHkiOnsibmFtZSI6Im5vbmUifSwiaGVhZCI6eyJuYW1lIjoibm9uZSJ9fX1dLFs2LDEsInYiLDEseyJzdHlsZSI6eyJib2R5Ijp7Im5hbWUiOiJub25lIn0sImhlYWQiOnsibmFtZSI6Im5vbmUifX19XSxbOSwxMCwiXFxjZG90cyIsMSx7InNob3J0ZW4iOnsic291cmNlIjoyMCwidGFyZ2V0IjoyMH0sInN0eWxlIjp7ImJvZHkiOnsibmFtZSI6Im5vbmUifSwiaGVhZCI6eyJuYW1lIjoibm9uZSJ9fX1dLFs0LDksIiIsMSx7InNob3J0ZW4iOnsidGFyZ2V0IjoyMH0sInN0eWxlIjp7ImJvZHkiOnsibmFtZSI6Im5vbmUifSwiaGVhZCI6eyJuYW1lIjoibm9uZSJ9fX1dLFsxLDExLCIiLDEseyJzaG9ydGVuIjp7InRhcmdldCI6MjB9LCJzdHlsZSI6eyJib2R5Ijp7Im5hbWUiOiJub25lIn0sImhlYWQiOnsibmFtZSI6Im5vbmUifX19XSxbOSwyLCIiLDEseyJzaG9ydGVuIjp7InNvdXJjZSI6MjB9LCJzdHlsZSI6eyJib2R5Ijp7Im5hbWUiOiJub25lIn0sImhlYWQiOnsibmFtZSI6Im5vbmUifX19XSxbMTEsMiwiIiwxLHsic2hvcnRlbiI6eyJzb3VyY2UiOjIwfSwic3R5bGUiOnsiYm9keSI6eyJuYW1lIjoibm9uZSJ9LCJoZWFkIjp7Im5hbWUiOiJub25lIn19fV0sWzQsMTMsIiIsMSx7ImxldmVsIjoxfV0sWzE0LDEsIiIsMSx7ImxldmVsIjoxfV0sWzE2LDE1LCJcXGNkb3RzIiwxLHsibGV2ZWwiOjF9XV0=
\begin{tikzcd}[row sep=small, column sep=tiny,  every label/.append style = {font=\normalsize}]
	{} \\
	v && {} && {} \\
	\\
	\\
	& v & \textcolor{white}{O} & \textcolor{white}{O} \\
	\\
	\\
	&&& {}
	\arrow["\cdots"{description}, sloped,,  ,  from=2-1, to=5-2]
	\arrow[""{name=0, anchor=center, inner sep=0}, "v"{description}, draw=none, from=2-3, to=5-2]
	\arrow[""{name=1, anchor=center, inner sep=0}, draw=none, from=5-4, to=2-5]
	\arrow[""{name=2, anchor=center, inner sep=0}, "x"{description}, draw=none, from=8-4, to=5-3]
	\arrow[""{name=3, anchor=center, inner sep=0}, draw=none, from=0, to=5-3]
	\arrow[""{name=4, anchor=center, inner sep=0}, draw=none, from=5-2, to=0]
	\arrow[""{name=5, anchor=center, inner sep=0}, draw=none, from=5-4, to=2]
	\arrow["\cdots"{description}, sloped,,  ,  draw=none, from=2, to=1]
	\arrow[""{name=6, anchor=center, inner sep=0}, draw=none, from=2, to=5-3]
	\arrow["\cdots"{description}, sloped,,  ,  from=3, to=6]
	\arrow[from=4, to=5-2]
	\arrow[from=5-4, to=5]
\end{tikzcd}.

{\bf \dots\dots}

\

{\bf Case 2n-1.} $N$ contains an odd number (say, $2n-1$ for $n>1$) vertices that are peaks or deeps of $Z$, that is, $N=$

\[% https://q.uiver.app/#q=WzAsMTEsWzAsMywidiJdLFsxLDBdLFsyLDMsInYiXSxbMSw2XSxbMSwzLCJPIixbMCwwLDEwMCwxXV0sWzUsNiwidiJdLFszLDYsIk8iLFswLDAsMTAwLDFdXSxbNCw2XSxbNCwzXSxbNiw5LCJ2Il0sWzUsOV0sWzEsMCwidiIsMSx7InN0eWxlIjp7ImJvZHkiOnsibmFtZSI6Im5vbmUifSwiaGVhZCI6eyJuYW1lIjoibm9uZSJ9fX1dLFsyLDMsInYiLDEseyJzdHlsZSI6eyJib2R5Ijp7Im5hbWUiOiJub25lIn0sImhlYWQiOnsibmFtZSI6Im5vbmUifX19XSxbMiw2LCJcXGNkb3RzIiwxXSxbOCw2LCIiLDEseyJzdHlsZSI6eyJib2R5Ijp7Im5hbWUiOiJub25lIn0sImhlYWQiOnsibmFtZSI6Im5vbmUifX19XSxbNSw5LCJcXGNkb3RzIiwxXSxbNywxMCwidiIsMSx7InN0eWxlIjp7ImJvZHkiOnsibmFtZSI6Im5vbmUifSwiaGVhZCI6eyJuYW1lIjoibm9uZSJ9fX1dLFs0LDExLCIiLDEseyJzaG9ydGVuIjp7InRhcmdldCI6MjB9LCJzdHlsZSI6eyJib2R5Ijp7Im5hbWUiOiJub25lIn0sImhlYWQiOnsibmFtZSI6Im5vbmUifX19XSxbNCwxMiwiIiwxLHsic2hvcnRlbiI6eyJ0YXJnZXQiOjIwfSwic3R5bGUiOnsiYm9keSI6eyJuYW1lIjoibm9uZSJ9LCJoZWFkIjp7Im5hbWUiOiJub25lIn19fV0sWzAsMTEsIiIsMSx7InNob3J0ZW4iOnsidGFyZ2V0IjoyMH0sInN0eWxlIjp7ImJvZHkiOnsibmFtZSI6Im5vbmUifSwiaGVhZCI6eyJuYW1lIjoibm9uZSJ9fX1dLFsyLDEyLCIiLDEseyJzaG9ydGVuIjp7InRhcmdldCI6MjB9LCJzdHlsZSI6eyJib2R5Ijp7Im5hbWUiOiJub25lIn0sImhlYWQiOnsibmFtZSI6Im5vbmUifX19XSxbNiwxNCwiIiwxLHsic2hvcnRlbiI6eyJ0YXJnZXQiOjIwfSwic3R5bGUiOnsiYm9keSI6eyJuYW1lIjoibm9uZSJ9LCJoZWFkIjp7Im5hbWUiOiJub25lIn19fV0sWzUsMTYsIiIsMSx7InNob3J0ZW4iOnsidGFyZ2V0IjoyMH0sInN0eWxlIjp7ImJvZHkiOnsibmFtZSI6Im5vbmUifSwiaGVhZCI6eyJuYW1lIjoibm9uZSJ9fX1dLFs3LDE2LCIiLDEseyJzaG9ydGVuIjp7InRhcmdldCI6MjB9LCJzdHlsZSI6eyJib2R5Ijp7Im5hbWUiOiJub25lIn0sImhlYWQiOnsibmFtZSI6Im5vbmUifX19XSxbMTQsNywiIiwxLHsic2hvcnRlbiI6eyJzb3VyY2UiOjIwfSwic3R5bGUiOnsiYm9keSI6eyJuYW1lIjoibm9uZSJ9LCJoZWFkIjp7Im5hbWUiOiJub25lIn19fV0sWzE3LDE4LCJcXGNkb3RzIiwxLHsibGV2ZWwiOjF9XSxbMTksMCwiIiwxLHsibGV2ZWwiOjF9XSxbMiwyMCwiIiwxLHsibGV2ZWwiOjF9XSxbNSwyMiwiIiwxLHsibGV2ZWwiOjF9XSxbMjEsNiwiXFxjZG90cyIsMSx7ImxldmVsIjoxLCJzdHlsZSI6eyJib2R5Ijp7Im5hbWUiOiJub25lIn0sImhlYWQiOnsibmFtZSI6Im5vbmUifX19XSxbMjQsMjMsIlxcY2RvdHMiLDEseyJsZXZlbCI6MX1dXQ==
\begin{tikzcd}[row sep=small, column sep=tiny,  every label/.append style = {font=\normalsize}]
	& {} \\
	\\
	\\
	v & \textcolor{white}{O} & v && {} \\
	\\
	\\
	& {} && \textcolor{white}{O} & {} & v \\
	\\
	\\
	&&&&& {} & v
	\arrow[""{name=0, anchor=center, inner sep=0}, "v"{description}, draw=none, from=1-2, to=4-1]
	\arrow[""{name=1, anchor=center, inner sep=0}, "v"{description}, draw=none, from=4-3, to=7-2]
	\arrow["\cdots"{description}, sloped,from=4-3, to=7-4]
	\arrow[""{name=2, anchor=center, inner sep=0}, draw=none, from=4-5, to=7-4]
	\arrow[""{name=3, anchor=center, inner sep=0}, "v"{description}, draw=none, from=7-5, to=10-6]
	\arrow["\cdots"{description}, sloped,from=7-6, to=10-7]
	\arrow[""{name=4, anchor=center, inner sep=0}, draw=none, from=4-1, to=0]
	\arrow[""{name=5, anchor=center, inner sep=0}, draw=none, from=4-2, to=0]
	\arrow[""{name=6, anchor=center, inner sep=0}, draw=none, from=4-2, to=1]
	\arrow[""{name=7, anchor=center, inner sep=0}, draw=none, from=4-3, to=1]
	\arrow[""{name=8, anchor=center, inner sep=0}, draw=none, from=2, to=7-5]
	\arrow[""{name=9, anchor=center, inner sep=0}, draw=none, from=7-4, to=2]
	\arrow[""{name=10, anchor=center, inner sep=0}, draw=none, from=7-5, to=3]
	\arrow[""{name=11, anchor=center, inner sep=0}, draw=none, from=7-6, to=3]
	\arrow[from=4, to=4-1]
	\arrow["\cdots"{description}, sloped,from=5, to=6]
	\arrow[from=4-3, to=7]
	\arrow["\cdots"{description}, sloped,from=8, to=10]
	\arrow["\cdots"{description}, sloped,draw=none, from=9, to=7-4]
	\arrow[from=7-6, to=11]
\end{tikzcd} \text{\quad or \quad}
% https://q.uiver.app/#q=WzAsMTEsWzIsMywiTyIsWzAsMCwxMDAsMV1dLFswLDAsInYiXSxbMSwzLCJ2Il0sWzIsNl0sWzMsM10sWzIsMF0sWzUsNiwiTyIsWzAsMCwxMDAsMV1dLFs0LDYsInYiXSxbNSw5XSxbNiw2LCJ2Il0sWzUsM10sWzQsMywiIiwxLHsic3R5bGUiOnsiYm9keSI6eyJuYW1lIjoibm9uZSJ9LCJoZWFkIjp7Im5hbWUiOiJub25lIn19fV0sWzUsMiwidiIsMSx7InN0eWxlIjp7ImJvZHkiOnsibmFtZSI6Im5vbmUifSwiaGVhZCI6eyJuYW1lIjoibm9uZSJ9fX1dLFsxLDIsIlxcY2RvdHMiLDFdLFs5LDgsInYiLDEseyJzdHlsZSI6eyJib2R5Ijp7Im5hbWUiOiJub25lIn0sImhlYWQiOnsibmFtZSI6Im5vbmUifX19XSxbMTAsNywidiIsMSx7InN0eWxlIjp7ImJvZHkiOnsibmFtZSI6Im5vbmUifSwiaGVhZCI6eyJuYW1lIjoibm9uZSJ9fX1dLFs0LDcsIlxcY2RvdHMiLDFdLFs0LDExLCJcXGNkb3RzIiwxLHsic2hvcnRlbiI6eyJ0YXJnZXQiOjIwfSwic3R5bGUiOnsiYm9keSI6eyJuYW1lIjoibm9uZSJ9LCJoZWFkIjp7Im5hbWUiOiJub25lIn19fV0sWzAsMTEsIiIsMSx7InNob3J0ZW4iOnsidGFyZ2V0IjoyMH0sInN0eWxlIjp7ImJvZHkiOnsibmFtZSI6Im5vbmUifSwiaGVhZCI6eyJuYW1lIjoibm9uZSJ9fX1dLFswLDEyLCIiLDEseyJzaG9ydGVuIjp7InRhcmdldCI6MjB9LCJzdHlsZSI6eyJib2R5Ijp7Im5hbWUiOiJub25lIn0sImhlYWQiOnsibmFtZSI6Im5vbmUifX19XSxbMiwxMiwiIiwxLHsic2hvcnRlbiI6eyJ0YXJnZXQiOjIwfSwic3R5bGUiOnsiYm9keSI6eyJuYW1lIjoibm9uZSJ9LCJoZWFkIjp7Im5hbWUiOiJub25lIn19fV0sWzksMTQsIiIsMSx7InNob3J0ZW4iOnsidGFyZ2V0IjoyMH0sInN0eWxlIjp7ImJvZHkiOnsibmFtZSI6Im5vbmUifSwiaGVhZCI6eyJuYW1lIjoibm9uZSJ9fX1dLFs2LDE0LCIiLDEseyJzaG9ydGVuIjp7InRhcmdldCI6MjB9LCJzdHlsZSI6eyJib2R5Ijp7Im5hbWUiOiJub25lIn0sImhlYWQiOnsibmFtZSI6Im5vbmUifX19XSxbNiwxNSwiIiwxLHsic2hvcnRlbiI6eyJ0YXJnZXQiOjIwfSwic3R5bGUiOnsiYm9keSI6eyJuYW1lIjoibm9uZSJ9LCJoZWFkIjp7Im5hbWUiOiJub25lIn19fV0sWzcsMTUsIiIsMSx7InNob3J0ZW4iOnsidGFyZ2V0IjoyMH0sInN0eWxlIjp7ImJvZHkiOnsibmFtZSI6Im5vbmUifSwiaGVhZCI6eyJuYW1lIjoibm9uZSJ9fX1dLFsxOSwxOCwiXFxjZG90cyIsMSx7ImxldmVsIjoxfV0sWzIwLDIsIiIsMSx7ImxldmVsIjoxfV0sWzIzLDIyLCJcXGNkb3RzIiwxLHsibGV2ZWwiOjF9XSxbMjQsNywiIiwxLHsibGV2ZWwiOjF9XSxbOSwyMSwiIiwxLHsibGV2ZWwiOjF9XV0=
\begin{tikzcd}[row sep=small, column sep=tiny,  every label/.append style = {font=\normalsize}]
	v && {} \\
	\\
	\\
	& v & \textcolor{white}{O} & {} && {} \\
	\\
	\\
	&& {} && v & \textcolor{white}{O} & v \\
	\\
	\\
	&&&&& {}
	\arrow["\cdots"{description}, sloped,,  ,  from=1-1, to=4-2]
	\arrow[""{name=0, anchor=center, inner sep=0}, "v"{description}, draw=none, from=1-3, to=4-2]
	\arrow[""{name=1, anchor=center, inner sep=0}, draw=none, from=4-4, to=7-3]
	\arrow["\cdots"{description}, sloped,,  ,  from=4-4, to=7-5]
	\arrow[""{name=2, anchor=center, inner sep=0}, "v"{description}, draw=none, from=4-6, to=7-5]
	\arrow[""{name=3, anchor=center, inner sep=0}, "v"{description}, draw=none, from=7-7, to=10-6]
	\arrow[""{name=4, anchor=center, inner sep=0}, draw=none, from=4-2, to=0]
	\arrow[""{name=5, anchor=center, inner sep=0}, draw=none, from=4-3, to=0]
	\arrow[""{name=6, anchor=center, inner sep=0}, draw=none, from=4-3, to=1]
	\arrow["\cdots"{description}, sloped,,  ,  draw=none, from=4-4, to=1]
	\arrow[""{name=7, anchor=center, inner sep=0}, draw=none, from=7-5, to=2]
	\arrow[""{name=8, anchor=center, inner sep=0}, draw=none, from=7-6, to=2]
	\arrow[""{name=9, anchor=center, inner sep=0}, draw=none, from=7-6, to=3]
	\arrow[""{name=10, anchor=center, inner sep=0}, draw=none, from=7-7, to=3]
	\arrow[from=4, to=4-2]
	\arrow["\cdots"{description}, sloped,,  ,  from=5, to=6]
	\arrow[from=7, to=7-5]
	\arrow["\cdots"{description}, sloped,,  ,  from=8, to=9]
	\arrow[from=7-7, to=10]
\end{tikzcd}.\] After possibly taking the flip of $Y$, the local configuration of $Y$ containing $N$ as a subdiagram can be classified into the following cases:\\
Case 2n-1.1.
% https://q.uiver.app/#q=WzAsMTEsWzAsMywidiJdLFsxLDBdLFsyLDMsInYiXSxbMSw2XSxbMSwzLCJPIixbMCwwLDEwMCwxXV0sWzUsNiwidiJdLFszLDYsIk8iLFswLDAsMTAwLDFdXSxbNCw2XSxbNCwzXSxbNiw5LCJ2Il0sWzUsOV0sWzEsMCwidiIsMSx7InN0eWxlIjp7ImJvZHkiOnsibmFtZSI6Im5vbmUifSwiaGVhZCI6eyJuYW1lIjoibm9uZSJ9fX1dLFsyLDMsInYiLDEseyJzdHlsZSI6eyJib2R5Ijp7Im5hbWUiOiJub25lIn0sImhlYWQiOnsibmFtZSI6Im5vbmUifX19XSxbMiw2LCJcXGNkb3RzIiwxXSxbOCw2LCIiLDEseyJzdHlsZSI6eyJib2R5Ijp7Im5hbWUiOiJub25lIn0sImhlYWQiOnsibmFtZSI6Im5vbmUifX19XSxbNSw5LCJcXGNkb3RzIiwxXSxbNywxMCwidiIsMSx7InN0eWxlIjp7ImJvZHkiOnsibmFtZSI6Im5vbmUifSwiaGVhZCI6eyJuYW1lIjoibm9uZSJ9fX1dLFs0LDExLCIiLDEseyJzaG9ydGVuIjp7InRhcmdldCI6MjB9LCJzdHlsZSI6eyJib2R5Ijp7Im5hbWUiOiJub25lIn0sImhlYWQiOnsibmFtZSI6Im5vbmUifX19XSxbNCwxMiwiIiwxLHsic2hvcnRlbiI6eyJ0YXJnZXQiOjIwfSwic3R5bGUiOnsiYm9keSI6eyJuYW1lIjoibm9uZSJ9LCJoZWFkIjp7Im5hbWUiOiJub25lIn19fV0sWzAsMTEsIiIsMSx7InNob3J0ZW4iOnsidGFyZ2V0IjoyMH0sInN0eWxlIjp7ImJvZHkiOnsibmFtZSI6Im5vbmUifSwiaGVhZCI6eyJuYW1lIjoibm9uZSJ9fX1dLFsyLDEyLCIiLDEseyJzaG9ydGVuIjp7InRhcmdldCI6MjB9LCJzdHlsZSI6eyJib2R5Ijp7Im5hbWUiOiJub25lIn0sImhlYWQiOnsibmFtZSI6Im5vbmUifX19XSxbNiwxNCwiIiwxLHsic2hvcnRlbiI6eyJ0YXJnZXQiOjIwfSwic3R5bGUiOnsiYm9keSI6eyJuYW1lIjoibm9uZSJ9LCJoZWFkIjp7Im5hbWUiOiJub25lIn19fV0sWzUsMTYsIiIsMSx7InNob3J0ZW4iOnsidGFyZ2V0IjoyMH0sInN0eWxlIjp7ImJvZHkiOnsibmFtZSI6Im5vbmUifSwiaGVhZCI6eyJuYW1lIjoibm9uZSJ9fX1dLFs3LDE2LCIiLDEseyJzaG9ydGVuIjp7InRhcmdldCI6MjB9LCJzdHlsZSI6eyJib2R5Ijp7Im5hbWUiOiJub25lIn0sImhlYWQiOnsibmFtZSI6Im5vbmUifX19XSxbMTQsNywiIiwxLHsic2hvcnRlbiI6eyJzb3VyY2UiOjIwfSwic3R5bGUiOnsiYm9keSI6eyJuYW1lIjoibm9uZSJ9LCJoZWFkIjp7Im5hbWUiOiJub25lIn19fV0sWzE3LDE4LCJcXGNkb3RzIiwxLHsibGV2ZWwiOjF9XSxbMTksMCwiIiwxLHsibGV2ZWwiOjF9XSxbMiwyMCwiIiwxLHsibGV2ZWwiOjF9XSxbNSwyMiwiIiwxLHsibGV2ZWwiOjF9XSxbMjEsNiwiXFxjZG90cyIsMSx7ImxldmVsIjoxLCJzdHlsZSI6eyJib2R5Ijp7Im5hbWUiOiJub25lIn0sImhlYWQiOnsibmFtZSI6Im5vbmUifX19XSxbMjQsMjMsIlxcY2RvdHMiLDEseyJsZXZlbCI6MX1dXQ==
\begin{tikzcd}[row sep=small, column sep=tiny,  every label/.append style = {font=\normalsize}]
	& {} \\
	\\
	\\
	v & \textcolor{white}{O} & v && {} \\
	\\
	\\
	& {} && \textcolor{white}{O} & {} & v \\
	\\
	\\
	&&&&& {} & v
	\arrow[""{name=0, anchor=center, inner sep=0}, "v"{description}, draw=none, from=1-2, to=4-1]
	\arrow[""{name=1, anchor=center, inner sep=0}, "v"{description}, draw=none, from=4-3, to=7-2]
	\arrow["\cdots"{description}, sloped,from=4-3, to=7-4]
	\arrow[""{name=2, anchor=center, inner sep=0}, draw=none, from=4-5, to=7-4]
	\arrow[""{name=3, anchor=center, inner sep=0}, "v"{description}, draw=none, from=7-5, to=10-6]
	\arrow["\cdots"{description}, sloped,from=7-6, to=10-7]
	\arrow[""{name=4, anchor=center, inner sep=0}, draw=none, from=4-1, to=0]
	\arrow[""{name=5, anchor=center, inner sep=0}, draw=none, from=4-2, to=0]
	\arrow[""{name=6, anchor=center, inner sep=0}, draw=none, from=4-2, to=1]
	\arrow[""{name=7, anchor=center, inner sep=0}, draw=none, from=4-3, to=1]
	\arrow[""{name=8, anchor=center, inner sep=0}, draw=none, from=2, to=7-5]
	\arrow[""{name=9, anchor=center, inner sep=0}, draw=none, from=7-4, to=2]
	\arrow[""{name=10, anchor=center, inner sep=0}, draw=none, from=7-5, to=3]
	\arrow[""{name=11, anchor=center, inner sep=0}, draw=none, from=7-6, to=3]
	\arrow[from=4, to=4-1]
	\arrow["\cdots"{description}, sloped,from=5, to=6]
	\arrow[from=4-3, to=7]
	\arrow["\cdots"{description}, sloped,from=8, to=10]
	\arrow["\cdots"{description}, sloped,draw=none, from=9, to=7-4]
	\arrow[from=7-6, to=11]
\end{tikzcd}.
		Case 2n-1.2.
% https://q.uiver.app/#q=WzAsMTMsWzIsMywidiJdLFszLDBdLFs0LDMsInYiXSxbMyw2XSxbMywzLCJPIixbMCwwLDEwMCwxXV0sWzcsNiwidiJdLFs1LDYsIk8iLFswLDAsMTAwLDFdXSxbNiw2XSxbNiwzXSxbOCw5LCJ2Il0sWzcsOV0sWzEsMCwieCJdLFswLDNdLFsxLDAsInYiLDEseyJzdHlsZSI6eyJib2R5Ijp7Im5hbWUiOiJub25lIn0sImhlYWQiOnsibmFtZSI6Im5vbmUifX19XSxbMiwzLCJ2IiwxLHsic3R5bGUiOnsiYm9keSI6eyJuYW1lIjoibm9uZSJ9LCJoZWFkIjp7Im5hbWUiOiJub25lIn19fV0sWzIsNiwiXFxjZG90cyIsMV0sWzgsNiwiIiwxLHsic3R5bGUiOnsiYm9keSI6eyJuYW1lIjoibm9uZSJ9LCJoZWFkIjp7Im5hbWUiOiJub25lIn19fV0sWzUsOSwiXFxjZG90cyIsMV0sWzcsMTAsInYiLDEseyJzdHlsZSI6eyJib2R5Ijp7Im5hbWUiOiJub25lIn0sImhlYWQiOnsibmFtZSI6Im5vbmUifX19XSxbMTEsMCwiXFxjZG90cyIsMV0sWzExLDEyLCJcXGNkb3RzIiwxLHsic3R5bGUiOnsiYm9keSI6eyJuYW1lIjoibm9uZSJ9LCJoZWFkIjp7Im5hbWUiOiJub25lIn19fV0sWzQsMTMsIiIsMSx7InNob3J0ZW4iOnsidGFyZ2V0IjoyMH0sInN0eWxlIjp7ImJvZHkiOnsibmFtZSI6Im5vbmUifSwiaGVhZCI6eyJuYW1lIjoibm9uZSJ9fX1dLFs0LDE0LCIiLDEseyJzaG9ydGVuIjp7InRhcmdldCI6MjB9LCJzdHlsZSI6eyJib2R5Ijp7Im5hbWUiOiJub25lIn0sImhlYWQiOnsibmFtZSI6Im5vbmUifX19XSxbMCwxMywiIiwxLHsic2hvcnRlbiI6eyJ0YXJnZXQiOjIwfSwic3R5bGUiOnsiYm9keSI6eyJuYW1lIjoibm9uZSJ9LCJoZWFkIjp7Im5hbWUiOiJub25lIn19fV0sWzIsMTQsIiIsMSx7InNob3J0ZW4iOnsidGFyZ2V0IjoyMH0sInN0eWxlIjp7ImJvZHkiOnsibmFtZSI6Im5vbmUifSwiaGVhZCI6eyJuYW1lIjoibm9uZSJ9fX1dLFs2LDE2LCIiLDEseyJzaG9ydGVuIjp7InRhcmdldCI6MjB9LCJzdHlsZSI6eyJib2R5Ijp7Im5hbWUiOiJub25lIn0sImhlYWQiOnsibmFtZSI6Im5vbmUifX19XSxbNSwxOCwiIiwxLHsic2hvcnRlbiI6eyJ0YXJnZXQiOjIwfSwic3R5bGUiOnsiYm9keSI6eyJuYW1lIjoibm9uZSJ9LCJoZWFkIjp7Im5hbWUiOiJub25lIn19fV0sWzcsMTgsIiIsMSx7InNob3J0ZW4iOnsidGFyZ2V0IjoyMH0sInN0eWxlIjp7ImJvZHkiOnsibmFtZSI6Im5vbmUifSwiaGVhZCI6eyJuYW1lIjoibm9uZSJ9fX1dLFsxMSwyMCwiIiwxLHsic2hvcnRlbiI6eyJ0YXJnZXQiOjIwfSwic3R5bGUiOnsiYm9keSI6eyJuYW1lIjoibm9uZSJ9LCJoZWFkIjp7Im5hbWUiOiJub25lIn19fV0sWzE2LDcsIiIsMSx7InNob3J0ZW4iOnsic291cmNlIjoyMH0sInN0eWxlIjp7ImJvZHkiOnsibmFtZSI6Im5vbmUifSwiaGVhZCI6eyJuYW1lIjoibm9uZSJ9fX1dLFsyMSwyMiwiXFxjZG90cyIsMSx7ImxldmVsIjoxfV0sWzIzLDAsIiIsMSx7ImxldmVsIjoxfV0sWzIsMjQsIiIsMSx7ImxldmVsIjoxfV0sWzUsMjYsIiIsMSx7ImxldmVsIjoxfV0sWzI1LDYsIlxcY2RvdHMiLDEseyJsZXZlbCI6MSwic3R5bGUiOnsiYm9keSI6eyJuYW1lIjoibm9uZSJ9LCJoZWFkIjp7Im5hbWUiOiJub25lIn19fV0sWzExLDI4LCIiLDEseyJsZXZlbCI6MX1dLFsyOSwyNywiXFxjZG90cyIsMSx7ImxldmVsIjoxfV1d
\begin{tikzcd}[row sep=small, column sep=tiny,  every label/.append style = {font=\normalsize}]
	& x && {} \\
	\\
	\\
	{} && v & \textcolor{white}{O} & v && {} \\
	\\
	\\
	&&& {} && \textcolor{white}{O} & {} & v \\
	\\
	\\
	&&&&&&& {} & v
	\arrow[""{name=0, anchor=center, inner sep=0}, "\cdots"{description}, sloped,draw=none, from=1-2, to=4-1]
	\arrow["\cdots"{description}, sloped,from=1-2, to=4-3]
	\arrow[""{name=1, anchor=center, inner sep=0}, "v"{description}, draw=none, from=1-4, to=4-3]
	\arrow[""{name=2, anchor=center, inner sep=0}, "v"{description}, draw=none, from=4-5, to=7-4]
	\arrow["\cdots"{description}, sloped,from=4-5, to=7-6]
	\arrow[""{name=3, anchor=center, inner sep=0}, draw=none, from=4-7, to=7-6]
	\arrow[""{name=4, anchor=center, inner sep=0}, "v"{description}, draw=none, from=7-7, to=10-8]
	\arrow["\cdots"{description}, sloped,from=7-8, to=10-9]
	\arrow[""{name=5, anchor=center, inner sep=0}, draw=none, from=1-2, to=0]
	\arrow[""{name=6, anchor=center, inner sep=0}, draw=none, from=4-3, to=1]
	\arrow[""{name=7, anchor=center, inner sep=0}, draw=none, from=4-4, to=1]
	\arrow[""{name=8, anchor=center, inner sep=0}, draw=none, from=4-4, to=2]
	\arrow[""{name=9, anchor=center, inner sep=0}, draw=none, from=4-5, to=2]
	\arrow[""{name=10, anchor=center, inner sep=0}, draw=none, from=3, to=7-7]
	\arrow[""{name=11, anchor=center, inner sep=0}, draw=none, from=7-6, to=3]
	\arrow[""{name=12, anchor=center, inner sep=0}, draw=none, from=7-7, to=4]
	\arrow[""{name=13, anchor=center, inner sep=0}, draw=none, from=7-8, to=4]
	\arrow[from=1-2, to=5]
	\arrow[from=6, to=4-3]
	\arrow["\cdots"{description}, sloped,from=7, to=8]
	\arrow[from=4-5, to=9]
	\arrow["\cdots"{description}, sloped,from=10, to=12]
	\arrow["\cdots"{description}, sloped,draw=none, from=11, to=7-6]
	\arrow[from=7-8, to=13]
\end{tikzcd}. \\
Case 2n-1.3.
% https://q.uiver.app/#q=WzAsMTMsWzIsMywidiJdLFszLDBdLFs0LDMsInYiXSxbMyw2XSxbMywzLCJPIixbMCwwLDEwMCwxXV0sWzcsNiwidiJdLFs1LDYsIk8iLFswLDAsMTAwLDFdXSxbNiw2XSxbNiwzXSxbOCw5LCJ2Il0sWzcsOV0sWzEsNiwieCJdLFswLDNdLFsxLDAsInYiLDEseyJzdHlsZSI6eyJib2R5Ijp7Im5hbWUiOiJub25lIn0sImhlYWQiOnsibmFtZSI6Im5vbmUifX19XSxbMiwzLCJ2IiwxLHsic3R5bGUiOnsiYm9keSI6eyJuYW1lIjoibm9uZSJ9LCJoZWFkIjp7Im5hbWUiOiJub25lIn19fV0sWzIsNiwiXFxjZG90cyIsMV0sWzgsNiwiIiwxLHsic3R5bGUiOnsiYm9keSI6eyJuYW1lIjoibm9uZSJ9LCJoZWFkIjp7Im5hbWUiOiJub25lIn19fV0sWzUsOSwiXFxjZG90cyIsMV0sWzcsMTAsInYiLDEseyJzdHlsZSI6eyJib2R5Ijp7Im5hbWUiOiJub25lIn0sImhlYWQiOnsibmFtZSI6Im5vbmUifX19XSxbMCwxMSwiXFxjZG90cyIsMV0sWzExLDEyLCJcXGNkb3RzIiwxLHsic3R5bGUiOnsiYm9keSI6eyJuYW1lIjoibm9uZSJ9LCJoZWFkIjp7Im5hbWUiOiJub25lIn19fV0sWzQsMTMsIiIsMSx7InNob3J0ZW4iOnsidGFyZ2V0IjoyMH0sInN0eWxlIjp7ImJvZHkiOnsibmFtZSI6Im5vbmUifSwiaGVhZCI6eyJuYW1lIjoibm9uZSJ9fX1dLFs0LDE0LCIiLDEseyJzaG9ydGVuIjp7InRhcmdldCI6MjB9LCJzdHlsZSI6eyJib2R5Ijp7Im5hbWUiOiJub25lIn0sImhlYWQiOnsibmFtZSI6Im5vbmUifX19XSxbMCwxMywiIiwxLHsic2hvcnRlbiI6eyJ0YXJnZXQiOjIwfSwic3R5bGUiOnsiYm9keSI6eyJuYW1lIjoibm9uZSJ9LCJoZWFkIjp7Im5hbWUiOiJub25lIn19fV0sWzIsMTQsIiIsMSx7InNob3J0ZW4iOnsidGFyZ2V0IjoyMH0sInN0eWxlIjp7ImJvZHkiOnsibmFtZSI6Im5vbmUifSwiaGVhZCI6eyJuYW1lIjoibm9uZSJ9fX1dLFs2LDE2LCIiLDEseyJzaG9ydGVuIjp7InRhcmdldCI6MjB9LCJzdHlsZSI6eyJib2R5Ijp7Im5hbWUiOiJub25lIn0sImhlYWQiOnsibmFtZSI6Im5vbmUifX19XSxbNSwxOCwiIiwxLHsic2hvcnRlbiI6eyJ0YXJnZXQiOjIwfSwic3R5bGUiOnsiYm9keSI6eyJuYW1lIjoibm9uZSJ9LCJoZWFkIjp7Im5hbWUiOiJub25lIn19fV0sWzcsMTgsIiIsMSx7InNob3J0ZW4iOnsidGFyZ2V0IjoyMH0sInN0eWxlIjp7ImJvZHkiOnsibmFtZSI6Im5vbmUifSwiaGVhZCI6eyJuYW1lIjoibm9uZSJ9fX1dLFsxMSwyMCwiIiwxLHsic2hvcnRlbiI6eyJ0YXJnZXQiOjIwfSwic3R5bGUiOnsiYm9keSI6eyJuYW1lIjoibm9uZSJ9LCJoZWFkIjp7Im5hbWUiOiJub25lIn19fV0sWzE2LDcsIiIsMSx7InNob3J0ZW4iOnsic291cmNlIjoyMH0sInN0eWxlIjp7ImJvZHkiOnsibmFtZSI6Im5vbmUifSwiaGVhZCI6eyJuYW1lIjoibm9uZSJ9fX1dLFsyMSwyMiwiXFxjZG90cyIsMSx7ImxldmVsIjoxfV0sWzIzLDAsIiIsMSx7ImxldmVsIjoxfV0sWzIsMjQsIiIsMSx7ImxldmVsIjoxfV0sWzUsMjYsIiIsMSx7ImxldmVsIjoxfV0sWzI1LDYsIlxcY2RvdHMiLDEseyJsZXZlbCI6MSwic3R5bGUiOnsiYm9keSI6eyJuYW1lIjoibm9uZSJ9LCJoZWFkIjp7Im5hbWUiOiJub25lIn19fV0sWzI4LDExLCIiLDEseyJsZXZlbCI6MX1dLFsyOSwyNywiXFxjZG90cyIsMSx7ImxldmVsIjoxfV1d
\begin{tikzcd}[row sep=small, column sep=tiny,  every label/.append style = {font=\normalsize}]
	&&& {} \\
	\\
	\\
	{} && v & \textcolor{white}{O} & v && {} \\
	\\
	\\
	& x && {} && \textcolor{white}{O} & {} & v \\
	\\
	\\
	&&&&&&& {} & v
	\arrow[""{name=0, anchor=center, inner sep=0}, "v"{description}, draw=none, from=1-4, to=4-3]
	\arrow["\cdots"{description}, sloped,from=4-3, to=7-2]
	\arrow[""{name=1, anchor=center, inner sep=0}, "v"{description}, draw=none, from=4-5, to=7-4]
	\arrow["\cdots"{description}, sloped,from=4-5, to=7-6]
	\arrow[""{name=2, anchor=center, inner sep=0}, draw=none, from=4-7, to=7-6]
	\arrow[""{name=3, anchor=center, inner sep=0}, "\cdots"{description}, sloped,draw=none, from=7-2, to=4-1]
	\arrow[""{name=4, anchor=center, inner sep=0}, "v"{description}, draw=none, from=7-7, to=10-8]
	\arrow["\cdots"{description}, sloped,from=7-8, to=10-9]
	\arrow[""{name=5, anchor=center, inner sep=0}, draw=none, from=4-3, to=0]
	\arrow[""{name=6, anchor=center, inner sep=0}, draw=none, from=4-4, to=0]
	\arrow[""{name=7, anchor=center, inner sep=0}, draw=none, from=4-4, to=1]
	\arrow[""{name=8, anchor=center, inner sep=0}, draw=none, from=4-5, to=1]
	\arrow[""{name=9, anchor=center, inner sep=0}, draw=none, from=2, to=7-7]
	\arrow[""{name=10, anchor=center, inner sep=0}, draw=none, from=7-2, to=3]
	\arrow[""{name=11, anchor=center, inner sep=0}, draw=none, from=7-6, to=2]
	\arrow[""{name=12, anchor=center, inner sep=0}, draw=none, from=7-7, to=4]
	\arrow[""{name=13, anchor=center, inner sep=0}, draw=none, from=7-8, to=4]
	\arrow[from=5, to=4-3]
	\arrow["\cdots"{description}, sloped,from=6, to=7]
	\arrow[from=4-5, to=8]
	\arrow["\cdots"{description}, sloped,from=9, to=12]
	\arrow[from=10, to=7-2]
	\arrow["\cdots"{description}, sloped,draw=none, from=11, to=7-6]
	\arrow[from=7-8, to=13]
\end{tikzcd}. \\
Case 2n-1.4.
% https://q.uiver.app/#q=WzAsMTEsWzIsMywiTyIsWzAsMCwxMDAsMV1dLFswLDAsInYiXSxbMSwzLCJ2Il0sWzIsNl0sWzMsMywiTyIsWzAsMCwxMDAsMV1dLFsyLDBdLFs1LDYsIk8iLFswLDAsMTAwLDFdXSxbNCw2LCJ2Il0sWzUsOV0sWzYsNiwidiJdLFs1LDNdLFs0LDMsIiIsMSx7InN0eWxlIjp7ImJvZHkiOnsibmFtZSI6Im5vbmUifSwiaGVhZCI6eyJuYW1lIjoibm9uZSJ9fX1dLFs1LDIsInYiLDEseyJzdHlsZSI6eyJib2R5Ijp7Im5hbWUiOiJub25lIn0sImhlYWQiOnsibmFtZSI6Im5vbmUifX19XSxbMSwyLCJcXGNkb3RzIiwxXSxbOSw4LCJ2IiwxLHsic3R5bGUiOnsiYm9keSI6eyJuYW1lIjoibm9uZSJ9LCJoZWFkIjp7Im5hbWUiOiJub25lIn19fV0sWzEwLDcsInYiLDEseyJzdHlsZSI6eyJib2R5Ijp7Im5hbWUiOiJub25lIn0sImhlYWQiOnsibmFtZSI6Im5vbmUifX19XSxbNCw3LCJcXGNkb3RzIiwxXSxbNCwxMSwiXFxjZG90cyIsMSx7InNob3J0ZW4iOnsidGFyZ2V0IjoyMH0sInN0eWxlIjp7ImJvZHkiOnsibmFtZSI6Im5vbmUifSwiaGVhZCI6eyJuYW1lIjoibm9uZSJ9fX1dLFswLDExLCIiLDEseyJzaG9ydGVuIjp7InRhcmdldCI6MjB9LCJzdHlsZSI6eyJib2R5Ijp7Im5hbWUiOiJub25lIn0sImhlYWQiOnsibmFtZSI6Im5vbmUifX19XSxbMCwxMiwiIiwxLHsic2hvcnRlbiI6eyJ0YXJnZXQiOjIwfSwic3R5bGUiOnsiYm9keSI6eyJuYW1lIjoibm9uZSJ9LCJoZWFkIjp7Im5hbWUiOiJub25lIn19fV0sWzIsMTIsIiIsMSx7InNob3J0ZW4iOnsidGFyZ2V0IjoyMH0sInN0eWxlIjp7ImJvZHkiOnsibmFtZSI6Im5vbmUifSwiaGVhZCI6eyJuYW1lIjoibm9uZSJ9fX1dLFs5LDE0LCIiLDEseyJzaG9ydGVuIjp7InRhcmdldCI6MjB9LCJzdHlsZSI6eyJib2R5Ijp7Im5hbWUiOiJub25lIn0sImhlYWQiOnsibmFtZSI6Im5vbmUifX19XSxbNiwxNCwiIiwxLHsic2hvcnRlbiI6eyJ0YXJnZXQiOjIwfSwic3R5bGUiOnsiYm9keSI6eyJuYW1lIjoibm9uZSJ9LCJoZWFkIjp7Im5hbWUiOiJub25lIn19fV0sWzYsMTUsIiIsMSx7InNob3J0ZW4iOnsidGFyZ2V0IjoyMH0sInN0eWxlIjp7ImJvZHkiOnsibmFtZSI6Im5vbmUifSwiaGVhZCI6eyJuYW1lIjoibm9uZSJ9fX1dLFs3LDE1LCIiLDEseyJzaG9ydGVuIjp7InRhcmdldCI6MjB9LCJzdHlsZSI6eyJib2R5Ijp7Im5hbWUiOiJub25lIn0sImhlYWQiOnsibmFtZSI6Im5vbmUifX19XSxbMTksMTgsIlxcY2RvdHMiLDEseyJsZXZlbCI6MX1dLFsyMCwyLCIiLDEseyJsZXZlbCI6MX1dLFsyMywyMiwiXFxjZG90cyIsMSx7ImxldmVsIjoxfV0sWzI0LDcsIiIsMSx7ImxldmVsIjoxfV0sWzksMjEsIiIsMSx7ImxldmVsIjoxfV1d
\begin{tikzcd}[row sep=small, column sep=tiny,  every label/.append style = {font=\normalsize}]
	v && {} \\
	\\
	\\
	& v & \textcolor{white}{O} & \textcolor{white}{O} && {} \\
	\\
	\\
	&& {} && v & \textcolor{white}{O} & v \\
	\\
	\\
	&&&&& {}
	\arrow["\cdots"{description}, sloped,from=1-1, to=4-2]
	\arrow[""{name=0, anchor=center, inner sep=0}, "v"{description}, draw=none, from=1-3, to=4-2]
	\arrow[""{name=1, anchor=center, inner sep=0}, draw=none, from=4-4, to=7-3]
	\arrow["\cdots"{description}, sloped,from=4-4, to=7-5]
	\arrow[""{name=2, anchor=center, inner sep=0}, "v"{description}, draw=none, from=4-6, to=7-5]
	\arrow[""{name=3, anchor=center, inner sep=0}, "v"{description}, draw=none, from=7-7, to=10-6]
	\arrow[""{name=4, anchor=center, inner sep=0}, draw=none, from=4-2, to=0]
	\arrow[""{name=5, anchor=center, inner sep=0}, draw=none, from=4-3, to=0]
	\arrow[""{name=6, anchor=center, inner sep=0}, draw=none, from=4-3, to=1]
	\arrow["\cdots"{description}, sloped,draw=none, from=4-4, to=1]
	\arrow[""{name=7, anchor=center, inner sep=0}, draw=none, from=7-5, to=2]
	\arrow[""{name=8, anchor=center, inner sep=0}, draw=none, from=7-6, to=2]
	\arrow[""{name=9, anchor=center, inner sep=0}, draw=none, from=7-6, to=3]
	\arrow[""{name=10, anchor=center, inner sep=0}, draw=none, from=7-7, to=3]
	\arrow[from=4, to=4-2]
	\arrow["\cdots"{description}, sloped,from=5, to=6]
	\arrow[from=7, to=7-5]
	\arrow["\cdots"{description}, sloped,from=8, to=9]
	\arrow[from=7-7, to=10]
\end{tikzcd}.
Case 2n-1.5.
% https://q.uiver.app/#q=WzAsMTMsWzIsMywiTyIsWzAsMCwxMDAsMV1dLFswLDAsInYiXSxbMSwzLCJ2Il0sWzIsNl0sWzMsMywiTyIsWzAsMCwxMDAsMV1dLFsyLDBdLFs1LDYsIk8iLFswLDAsMTAwLDFdXSxbNCw2LCJ2Il0sWzUsOV0sWzYsNiwidiJdLFs1LDNdLFs3LDMsIngiXSxbOCw2XSxbNCwzLCIiLDEseyJzdHlsZSI6eyJib2R5Ijp7Im5hbWUiOiJub25lIn0sImhlYWQiOnsibmFtZSI6Im5vbmUifX19XSxbNSwyLCJ2IiwxLHsic3R5bGUiOnsiYm9keSI6eyJuYW1lIjoibm9uZSJ9LCJoZWFkIjp7Im5hbWUiOiJub25lIn19fV0sWzEsMiwiXFxjZG90cyIsMV0sWzksOCwidiIsMSx7InN0eWxlIjp7ImJvZHkiOnsibmFtZSI6Im5vbmUifSwiaGVhZCI6eyJuYW1lIjoibm9uZSJ9fX1dLFsxMCw3LCJ2IiwxLHsic3R5bGUiOnsiYm9keSI6eyJuYW1lIjoibm9uZSJ9LCJoZWFkIjp7Im5hbWUiOiJub25lIn19fV0sWzQsNywiXFxjZG90cyIsMV0sWzExLDksIlxcY2RvdHMiLDFdLFsxMSwxMiwiXFxjZG90cyIsMSx7InN0eWxlIjp7ImJvZHkiOnsibmFtZSI6Im5vbmUifSwiaGVhZCI6eyJuYW1lIjoibm9uZSJ9fX1dLFs0LDEzLCJcXGNkb3RzIiwxLHsic2hvcnRlbiI6eyJ0YXJnZXQiOjIwfSwic3R5bGUiOnsiYm9keSI6eyJuYW1lIjoibm9uZSJ9LCJoZWFkIjp7Im5hbWUiOiJub25lIn19fV0sWzAsMTMsIiIsMSx7InNob3J0ZW4iOnsidGFyZ2V0IjoyMH0sInN0eWxlIjp7ImJvZHkiOnsibmFtZSI6Im5vbmUifSwiaGVhZCI6eyJuYW1lIjoibm9uZSJ9fX1dLFswLDE0LCIiLDEseyJzaG9ydGVuIjp7InRhcmdldCI6MjB9LCJzdHlsZSI6eyJib2R5Ijp7Im5hbWUiOiJub25lIn0sImhlYWQiOnsibmFtZSI6Im5vbmUifX19XSxbMiwxNCwiIiwxLHsic2hvcnRlbiI6eyJ0YXJnZXQiOjIwfSwic3R5bGUiOnsiYm9keSI6eyJuYW1lIjoibm9uZSJ9LCJoZWFkIjp7Im5hbWUiOiJub25lIn19fV0sWzksMTYsIiIsMSx7InNob3J0ZW4iOnsidGFyZ2V0IjoyMH0sInN0eWxlIjp7ImJvZHkiOnsibmFtZSI6Im5vbmUifSwiaGVhZCI6eyJuYW1lIjoibm9uZSJ9fX1dLFs2LDE2LCIiLDEseyJzaG9ydGVuIjp7InRhcmdldCI6MjB9LCJzdHlsZSI6eyJib2R5Ijp7Im5hbWUiOiJub25lIn0sImhlYWQiOnsibmFtZSI6Im5vbmUifX19XSxbNiwxNywiIiwxLHsic2hvcnRlbiI6eyJ0YXJnZXQiOjIwfSwic3R5bGUiOnsiYm9keSI6eyJuYW1lIjoibm9uZSJ9LCJoZWFkIjp7Im5hbWUiOiJub25lIn19fV0sWzcsMTcsIiIsMSx7InNob3J0ZW4iOnsidGFyZ2V0IjoyMH0sInN0eWxlIjp7ImJvZHkiOnsibmFtZSI6Im5vbmUifSwiaGVhZCI6eyJuYW1lIjoibm9uZSJ9fX1dLFsxMSwyMCwiIiwxLHsic2hvcnRlbiI6eyJ0YXJnZXQiOjIwfSwic3R5bGUiOnsiYm9keSI6eyJuYW1lIjoibm9uZSJ9LCJoZWFkIjp7Im5hbWUiOiJub25lIn19fV0sWzIzLDIyLCJcXGNkb3RzIiwxLHsibGV2ZWwiOjF9XSxbMjQsMiwiIiwxLHsibGV2ZWwiOjF9XSxbMjcsMjYsIlxcY2RvdHMiLDEseyJsZXZlbCI6MX1dLFsyOCw3LCIiLDEseyJsZXZlbCI6MX1dLFs5LDI1LCIiLDEseyJsZXZlbCI6MX1dLFsxMSwyOSwiIiwxLHsibGV2ZWwiOjF9XV0=
\begin{tikzcd}[row sep=small, column sep=tiny,  every label/.append style = {font=\normalsize}]
	v && {} \\
	\\
	\\
	& v & \textcolor{white}{O} & \textcolor{white}{O} && {} && x \\
	\\
	\\
	&& {} && v & \textcolor{white}{O} & v && {} \\
	\\
	\\
	&&&&& {}
	\arrow["\cdots"{description}, sloped,from=1-1, to=4-2]
	\arrow[""{name=0, anchor=center, inner sep=0}, "v"{description}, draw=none, from=1-3, to=4-2]
	\arrow[""{name=1, anchor=center, inner sep=0}, draw=none, from=4-4, to=7-3]
	\arrow["\cdots"{description}, sloped,from=4-4, to=7-5]
	\arrow[""{name=2, anchor=center, inner sep=0}, "v"{description}, draw=none, from=4-6, to=7-5]
	\arrow["\cdots"{description}, sloped,from=4-8, to=7-7]
	\arrow[""{name=3, anchor=center, inner sep=0}, "\cdots"{description}, sloped,draw=none, from=4-8, to=7-9]
	\arrow[""{name=4, anchor=center, inner sep=0}, "v"{description}, draw=none, from=7-7, to=10-6]
	\arrow[""{name=5, anchor=center, inner sep=0}, draw=none, from=4-2, to=0]
	\arrow[""{name=6, anchor=center, inner sep=0}, draw=none, from=4-3, to=0]
	\arrow[""{name=7, anchor=center, inner sep=0}, draw=none, from=4-3, to=1]
	\arrow["\cdots"{description}, sloped,draw=none, from=4-4, to=1]
	\arrow[""{name=8, anchor=center, inner sep=0}, draw=none, from=4-8, to=3]
	\arrow[""{name=9, anchor=center, inner sep=0}, draw=none, from=7-5, to=2]
	\arrow[""{name=10, anchor=center, inner sep=0}, draw=none, from=7-6, to=2]
	\arrow[""{name=11, anchor=center, inner sep=0}, draw=none, from=7-6, to=4]
	\arrow[""{name=12, anchor=center, inner sep=0}, draw=none, from=7-7, to=4]
	\arrow[from=5, to=4-2]
	\arrow["\cdots"{description}, sloped,from=6, to=7]
	\arrow[from=4-8, to=8]
	\arrow[from=9, to=7-5]
	\arrow["\cdots"{description}, sloped,from=10, to=11]
	\arrow[from=7-7, to=12]
\end{tikzcd}. \\
Case 2n-1.6.
% https://q.uiver.app/#q=WzAsMTMsWzIsMywiTyIsWzAsMCwxMDAsMV1dLFswLDAsInYiXSxbMSwzLCJ2Il0sWzIsNl0sWzMsMywiTyIsWzAsMCwxMDAsMV1dLFsyLDBdLFs1LDYsIk8iLFswLDAsMTAwLDFdXSxbNCw2LCJ2Il0sWzUsOV0sWzYsNiwidiJdLFs1LDNdLFs3LDksIngiXSxbOCw2XSxbNCwzLCIiLDEseyJzdHlsZSI6eyJib2R5Ijp7Im5hbWUiOiJub25lIn0sImhlYWQiOnsibmFtZSI6Im5vbmUifX19XSxbNSwyLCJ2IiwxLHsic3R5bGUiOnsiYm9keSI6eyJuYW1lIjoibm9uZSJ9LCJoZWFkIjp7Im5hbWUiOiJub25lIn19fV0sWzEsMiwiXFxjZG90cyIsMV0sWzksOCwidiIsMSx7InN0eWxlIjp7ImJvZHkiOnsibmFtZSI6Im5vbmUifSwiaGVhZCI6eyJuYW1lIjoibm9uZSJ9fX1dLFsxMCw3LCJ2IiwxLHsic3R5bGUiOnsiYm9keSI6eyJuYW1lIjoibm9uZSJ9LCJoZWFkIjp7Im5hbWUiOiJub25lIn19fV0sWzksMTEsIlxcY2RvdHMiLDFdLFsxMSwxMiwiXFxjZG90cyIsMSx7InN0eWxlIjp7ImJvZHkiOnsibmFtZSI6Im5vbmUifSwiaGVhZCI6eyJuYW1lIjoibm9uZSJ9fX1dLFs0LDcsIlxcY2RvdHMiLDFdLFs0LDEzLCJcXGNkb3RzIiwxLHsic2hvcnRlbiI6eyJ0YXJnZXQiOjIwfSwic3R5bGUiOnsiYm9keSI6eyJuYW1lIjoibm9uZSJ9LCJoZWFkIjp7Im5hbWUiOiJub25lIn19fV0sWzAsMTMsIiIsMSx7InNob3J0ZW4iOnsidGFyZ2V0IjoyMH0sInN0eWxlIjp7ImJvZHkiOnsibmFtZSI6Im5vbmUifSwiaGVhZCI6eyJuYW1lIjoibm9uZSJ9fX1dLFswLDE0LCIiLDEseyJzaG9ydGVuIjp7InRhcmdldCI6MjB9LCJzdHlsZSI6eyJib2R5Ijp7Im5hbWUiOiJub25lIn0sImhlYWQiOnsibmFtZSI6Im5vbmUifX19XSxbMiwxNCwiIiwxLHsic2hvcnRlbiI6eyJ0YXJnZXQiOjIwfSwic3R5bGUiOnsiYm9keSI6eyJuYW1lIjoibm9uZSJ9LCJoZWFkIjp7Im5hbWUiOiJub25lIn19fV0sWzksMTYsIiIsMSx7InNob3J0ZW4iOnsidGFyZ2V0IjoyMH0sInN0eWxlIjp7ImJvZHkiOnsibmFtZSI6Im5vbmUifSwiaGVhZCI6eyJuYW1lIjoibm9uZSJ9fX1dLFs2LDE2LCIiLDEseyJzaG9ydGVuIjp7InRhcmdldCI6MjB9LCJzdHlsZSI6eyJib2R5Ijp7Im5hbWUiOiJub25lIn0sImhlYWQiOnsibmFtZSI6Im5vbmUifX19XSxbNiwxNywiIiwxLHsic2hvcnRlbiI6eyJ0YXJnZXQiOjIwfSwic3R5bGUiOnsiYm9keSI6eyJuYW1lIjoibm9uZSJ9LCJoZWFkIjp7Im5hbWUiOiJub25lIn19fV0sWzcsMTcsIiIsMSx7InNob3J0ZW4iOnsidGFyZ2V0IjoyMH0sInN0eWxlIjp7ImJvZHkiOnsibmFtZSI6Im5vbmUifSwiaGVhZCI6eyJuYW1lIjoibm9uZSJ9fX1dLFsxMSwxOSwiIiwxLHsic2hvcnRlbiI6eyJ0YXJnZXQiOjIwfSwic3R5bGUiOnsiYm9keSI6eyJuYW1lIjoibm9uZSJ9LCJoZWFkIjp7Im5hbWUiOiJub25lIn19fV0sWzIzLDIyLCJcXGNkb3RzIiwxLHsibGV2ZWwiOjF9XSxbMjQsMiwiIiwxLHsibGV2ZWwiOjF9XSxbMjcsMjYsIlxcY2RvdHMiLDEseyJsZXZlbCI6MX1dLFsyOCw3LCIiLDEseyJsZXZlbCI6MX1dLFs5LDI1LCIiLDEseyJsZXZlbCI6MX1dLFsyOSwxMSwiIiwxLHsibGV2ZWwiOjF9XV0=
\begin{tikzcd}[row sep=small, column sep=tiny,  every label/.append style = {font=\normalsize}]
	v && {} \\
	\\
	\\
	& v & \textcolor{white}{O} & \textcolor{white}{O} && {} \\
	\\
	\\
	&& {} && v & \textcolor{white}{O} & v && {} \\
	\\
	\\
	&&&&& {} && x
	\arrow["\cdots"{description}, sloped,from=1-1, to=4-2]
	\arrow[""{name=0, anchor=center, inner sep=0}, "v"{description}, draw=none, from=1-3, to=4-2]
	\arrow[""{name=1, anchor=center, inner sep=0}, draw=none, from=4-4, to=7-3]
	\arrow["\cdots"{description}, sloped,from=4-4, to=7-5]
	\arrow[""{name=2, anchor=center, inner sep=0}, "v"{description}, draw=none, from=4-6, to=7-5]
	\arrow[""{name=3, anchor=center, inner sep=0}, "v"{description}, draw=none, from=7-7, to=10-6]
	\arrow["\cdots"{description}, sloped,from=7-7, to=10-8]
	\arrow[""{name=4, anchor=center, inner sep=0}, "\cdots"{description}, sloped,draw=none, from=10-8, to=7-9]
	\arrow[""{name=5, anchor=center, inner sep=0}, draw=none, from=4-2, to=0]
	\arrow[""{name=6, anchor=center, inner sep=0}, draw=none, from=4-3, to=0]
	\arrow[""{name=7, anchor=center, inner sep=0}, draw=none, from=4-3, to=1]
	\arrow["\cdots"{description}, sloped,draw=none, from=4-4, to=1]
	\arrow[""{name=8, anchor=center, inner sep=0}, draw=none, from=7-5, to=2]
	\arrow[""{name=9, anchor=center, inner sep=0}, draw=none, from=7-6, to=2]
	\arrow[""{name=10, anchor=center, inner sep=0}, draw=none, from=7-6, to=3]
	\arrow[""{name=11, anchor=center, inner sep=0}, draw=none, from=7-7, to=3]
	\arrow[""{name=12, anchor=center, inner sep=0}, draw=none, from=10-8, to=4]
	\arrow[from=5, to=4-2]
	\arrow["\cdots"{description}, sloped,from=6, to=7]
	\arrow[from=8, to=7-5]
	\arrow["\cdots"{description}, sloped,from=9, to=10]
	\arrow[from=7-7, to=11]
	\arrow[from=12, to=10-8]
\end{tikzcd}.

		{\bf Case 2n.} $N$ contains an even number (say, $2n$ for $n\geq 1$) vertices that are peaks or deeps of $Z$, that is, $N=$

\[% https://q.uiver.app/#q=WzAsMTEsWzIsMywiTyIsWzAsMCwxMDAsMV1dLFswLDAsInYiXSxbMSwzLCJ2Il0sWzIsNl0sWzMsMywiTyIsWzAsMCwxMDAsMV1dLFsyLDBdLFs1LDYsIk8iLFswLDAsMTAwLDFdXSxbNCw2LCJ2Il0sWzUsOV0sWzYsNl0sWzUsM10sWzQsMywiIiwxLHsic3R5bGUiOnsiYm9keSI6eyJuYW1lIjoibm9uZSJ9LCJoZWFkIjp7Im5hbWUiOiJub25lIn19fV0sWzUsMiwidiIsMSx7InN0eWxlIjp7ImJvZHkiOnsibmFtZSI6Im5vbmUifSwiaGVhZCI6eyJuYW1lIjoibm9uZSJ9fX1dLFsxLDIsIlxcY2RvdHMiLDFdLFs5LDgsInYiLDEseyJzdHlsZSI6eyJib2R5Ijp7Im5hbWUiOiJub25lIn0sImhlYWQiOnsibmFtZSI6Im5vbmUifX19XSxbMTAsNywidiIsMSx7InN0eWxlIjp7ImJvZHkiOnsibmFtZSI6Im5vbmUifSwiaGVhZCI6eyJuYW1lIjoibm9uZSJ9fX1dLFs0LDcsIlxcY2RvdHMiLDFdLFs0LDExLCJcXGNkb3RzIiwxLHsic2hvcnRlbiI6eyJ0YXJnZXQiOjIwfSwic3R5bGUiOnsiYm9keSI6eyJuYW1lIjoibm9uZSJ9LCJoZWFkIjp7Im5hbWUiOiJub25lIn19fV0sWzAsMTEsIiIsMSx7InNob3J0ZW4iOnsidGFyZ2V0IjoyMH0sInN0eWxlIjp7ImJvZHkiOnsibmFtZSI6Im5vbmUifSwiaGVhZCI6eyJuYW1lIjoibm9uZSJ9fX1dLFswLDEyLCIiLDEseyJzaG9ydGVuIjp7InRhcmdldCI6MjB9LCJzdHlsZSI6eyJib2R5Ijp7Im5hbWUiOiJub25lIn0sImhlYWQiOnsibmFtZSI6Im5vbmUifX19XSxbMiwxMiwiIiwxLHsic2hvcnRlbiI6eyJ0YXJnZXQiOjIwfSwic3R5bGUiOnsiYm9keSI6eyJuYW1lIjoibm9uZSJ9LCJoZWFkIjp7Im5hbWUiOiJub25lIn19fV0sWzksMTQsIiIsMSx7InNob3J0ZW4iOnsidGFyZ2V0IjoyMH0sInN0eWxlIjp7ImJvZHkiOnsibmFtZSI6Im5vbmUifSwiaGVhZCI6eyJuYW1lIjoibm9uZSJ9fX1dLFs2LDE0LCIiLDEseyJzaG9ydGVuIjp7InRhcmdldCI6MjB9LCJzdHlsZSI6eyJib2R5Ijp7Im5hbWUiOiJub25lIn0sImhlYWQiOnsibmFtZSI6Im5vbmUifX19XSxbNiwxNSwiIiwxLHsic2hvcnRlbiI6eyJ0YXJnZXQiOjIwfSwic3R5bGUiOnsiYm9keSI6eyJuYW1lIjoibm9uZSJ9LCJoZWFkIjp7Im5hbWUiOiJub25lIn19fV0sWzcsMTUsIiIsMSx7InNob3J0ZW4iOnsidGFyZ2V0IjoyMH0sInN0eWxlIjp7ImJvZHkiOnsibmFtZSI6Im5vbmUifSwiaGVhZCI6eyJuYW1lIjoibm9uZSJ9fX1dLFsxOSwxOCwiXFxjZG90cyIsMSx7ImxldmVsIjoxfV0sWzIwLDIsIiIsMSx7ImxldmVsIjoxfV0sWzIzLDIyLCJcXGNkb3RzIiwxLHsibGV2ZWwiOjF9XSxbMjQsNywiIiwxLHsibGV2ZWwiOjF9XV0=
\begin{tikzcd}[row sep=small, column sep=tiny,  every label/.append style = {font=\normalsize}]
	v && {} \\
	\\
	\\
	& v & \textcolor{white}{O} & \textcolor{white}{O} && {} \\
	\\
	\\
	&& {} && v & \textcolor{white}{O} & {} \\
	\\
	\\
	&&&&& {}
	\arrow["\cdots"{description}, sloped,from=1-1, to=4-2]
	\arrow[""{name=0, anchor=center, inner sep=0}, "v"{description}, draw=none, from=1-3, to=4-2]
	\arrow[""{name=1, anchor=center, inner sep=0}, draw=none, from=4-4, to=7-3]
	\arrow["\cdots"{description}, sloped,from=4-4, to=7-5]
	\arrow[""{name=2, anchor=center, inner sep=0}, "v"{description}, draw=none, from=4-6, to=7-5]
	\arrow[""{name=3, anchor=center, inner sep=0}, "v"{description}, draw=none, from=7-7, to=10-6]
	\arrow[""{name=4, anchor=center, inner sep=0}, draw=none, from=4-2, to=0]
	\arrow[""{name=5, anchor=center, inner sep=0}, draw=none, from=4-3, to=0]
	\arrow[""{name=6, anchor=center, inner sep=0}, draw=none, from=4-3, to=1]
	\arrow["\cdots"{description}, sloped,draw=none, from=4-4, to=1]
	\arrow[""{name=7, anchor=center, inner sep=0}, draw=none, from=7-5, to=2]
	\arrow[""{name=8, anchor=center, inner sep=0}, draw=none, from=7-6, to=2]
	\arrow[""{name=9, anchor=center, inner sep=0}, draw=none, from=7-6, to=3]
	\arrow[draw=none, from=7-7, to=3]
	\arrow[from=4, to=4-2]
	\arrow["\cdots"{description}, sloped,from=5, to=6]
	\arrow[from=7, to=7-5]
	\arrow["\cdots"{description}, sloped,from=8, to=9]
\end{tikzcd} \text{\qquad or \quad\qquad}
% https://q.uiver.app/#q=WzAsMTAsWzEsMywiTyIsWzAsMCwxMDAsMV1dLFswLDMsInYiXSxbMSw2XSxbMiwzLCJPIixbMCwwLDEwMCwxXV0sWzEsMF0sWzQsNiwiTyIsWzAsMCwxMDAsMV1dLFszLDYsInYiXSxbNCw5XSxbNSw2LCJ2Il0sWzQsM10sWzMsMiwiIiwxLHsic3R5bGUiOnsiYm9keSI6eyJuYW1lIjoibm9uZSJ9LCJoZWFkIjp7Im5hbWUiOiJub25lIn19fV0sWzQsMSwidiIsMSx7InN0eWxlIjp7ImJvZHkiOnsibmFtZSI6Im5vbmUifSwiaGVhZCI6eyJuYW1lIjoibm9uZSJ9fX1dLFs4LDcsInYiLDEseyJzdHlsZSI6eyJib2R5Ijp7Im5hbWUiOiJub25lIn0sImhlYWQiOnsibmFtZSI6Im5vbmUifX19XSxbOSw2LCJ2IiwxLHsic3R5bGUiOnsiYm9keSI6eyJuYW1lIjoibm9uZSJ9LCJoZWFkIjp7Im5hbWUiOiJub25lIn19fV0sWzMsNiwiXFxjZG90cyIsMV0sWzMsMTAsIlxcY2RvdHMiLDEseyJzaG9ydGVuIjp7InRhcmdldCI6MjB9LCJzdHlsZSI6eyJib2R5Ijp7Im5hbWUiOiJub25lIn0sImhlYWQiOnsibmFtZSI6Im5vbmUifX19XSxbMCwxMCwiIiwxLHsic2hvcnRlbiI6eyJ0YXJnZXQiOjIwfSwic3R5bGUiOnsiYm9keSI6eyJuYW1lIjoibm9uZSJ9LCJoZWFkIjp7Im5hbWUiOiJub25lIn19fV0sWzAsMTEsIiIsMSx7InNob3J0ZW4iOnsidGFyZ2V0IjoyMH0sInN0eWxlIjp7ImJvZHkiOnsibmFtZSI6Im5vbmUifSwiaGVhZCI6eyJuYW1lIjoibm9uZSJ9fX1dLFsxLDExLCIiLDEseyJzaG9ydGVuIjp7InRhcmdldCI6MjB9LCJzdHlsZSI6eyJib2R5Ijp7Im5hbWUiOiJub25lIn0sImhlYWQiOnsibmFtZSI6Im5vbmUifX19XSxbOCwxMiwiIiwxLHsic2hvcnRlbiI6eyJ0YXJnZXQiOjIwfSwic3R5bGUiOnsiYm9keSI6eyJuYW1lIjoibm9uZSJ9LCJoZWFkIjp7Im5hbWUiOiJub25lIn19fV0sWzUsMTIsIiIsMSx7InNob3J0ZW4iOnsidGFyZ2V0IjoyMH0sInN0eWxlIjp7ImJvZHkiOnsibmFtZSI6Im5vbmUifSwiaGVhZCI6eyJuYW1lIjoibm9uZSJ9fX1dLFs1LDEzLCIiLDEseyJzaG9ydGVuIjp7InRhcmdldCI6MjB9LCJzdHlsZSI6eyJib2R5Ijp7Im5hbWUiOiJub25lIn0sImhlYWQiOnsibmFtZSI6Im5vbmUifX19XSxbNiwxMywiIiwxLHsic2hvcnRlbiI6eyJ0YXJnZXQiOjIwfSwic3R5bGUiOnsiYm9keSI6eyJuYW1lIjoibm9uZSJ9LCJoZWFkIjp7Im5hbWUiOiJub25lIn19fV0sWzE3LDE2LCJcXGNkb3RzIiwxLHsibGV2ZWwiOjF9XSxbMTgsMSwiIiwxLHsibGV2ZWwiOjF9XSxbMjEsMjAsIlxcY2RvdHMiLDEseyJsZXZlbCI6MX1dLFsyMiw2LCIiLDEseyJsZXZlbCI6MX1dLFs4LDE5LCIiLDEseyJsZXZlbCI6MX1dXQ==
\begin{tikzcd}[row sep=small, column sep=tiny,  every label/.append style = {font=\normalsize}]
	& {} \\
	\\
	\\
	v & \textcolor{white}{O} & \textcolor{white}{O} && {} \\
	\\
	\\
	& {} && v & \textcolor{white}{O} & v \\
	\\
	\\
	&&&& {}
	\arrow[""{name=0, anchor=center, inner sep=0}, "v"{description}, draw=none, from=1-2, to=4-1]
	\arrow[""{name=1, anchor=center, inner sep=0}, draw=none, from=4-3, to=7-2]
	\arrow["\cdots"{description}, sloped,from=4-3, to=7-4]
	\arrow[""{name=2, anchor=center, inner sep=0}, "v"{description}, draw=none, from=4-5, to=7-4]
	\arrow[""{name=3, anchor=center, inner sep=0}, "v"{description}, draw=none, from=7-6, to=10-5]
	\arrow[""{name=4, anchor=center, inner sep=0}, draw=none, from=4-1, to=0]
	\arrow[""{name=5, anchor=center, inner sep=0}, draw=none, from=4-2, to=0]
	\arrow[""{name=6, anchor=center, inner sep=0}, draw=none, from=4-2, to=1]
	\arrow["\cdots"{description}, sloped,draw=none, from=4-3, to=1]
	\arrow[""{name=7, anchor=center, inner sep=0}, draw=none, from=7-4, to=2]
	\arrow[""{name=8, anchor=center, inner sep=0}, draw=none, from=7-5, to=2]
	\arrow[""{name=9, anchor=center, inner sep=0}, draw=none, from=7-5, to=3]
	\arrow[""{name=10, anchor=center, inner sep=0}, draw=none, from=7-6, to=3]
	\arrow[from=4, to=4-1]
	\arrow["\cdots"{description}, sloped,from=5, to=6]
	\arrow[from=7, to=7-4]
	\arrow["\cdots"{description}, sloped,from=8, to=9]
	\arrow[from=7-6, to=10]
\end{tikzcd}. \]
After possibly taking the flip of $Y$, the local configuration of $Y$ containing $N$ as a subdiagram can be classified into the following cases:\\
		Case 2n.1.
% https://q.uiver.app/#q=WzAsMTEsWzIsMywiTyIsWzAsMCwxMDAsMV1dLFswLDAsInYiXSxbMSwzLCJ2Il0sWzIsNl0sWzMsM10sWzIsMF0sWzUsNiwiTyIsWzAsMCwxMDAsMV1dLFs0LDYsInYiXSxbNSw5XSxbNiw2XSxbNSwzXSxbNCwzLCIiLDEseyJzdHlsZSI6eyJib2R5Ijp7Im5hbWUiOiJub25lIn0sImhlYWQiOnsibmFtZSI6Im5vbmUifX19XSxbNSwyLCJ2IiwxLHsic3R5bGUiOnsiYm9keSI6eyJuYW1lIjoibm9uZSJ9LCJoZWFkIjp7Im5hbWUiOiJub25lIn19fV0sWzEsMiwiXFxjZG90cyIsMV0sWzksOCwidiIsMSx7InN0eWxlIjp7ImJvZHkiOnsibmFtZSI6Im5vbmUifSwiaGVhZCI6eyJuYW1lIjoibm9uZSJ9fX1dLFsxMCw3LCJ2IiwxLHsic3R5bGUiOnsiYm9keSI6eyJuYW1lIjoibm9uZSJ9LCJoZWFkIjp7Im5hbWUiOiJub25lIn19fV0sWzQsNywiXFxjZG90cyIsMV0sWzQsMTEsIlxcY2RvdHMiLDEseyJzaG9ydGVuIjp7InRhcmdldCI6MjB9LCJzdHlsZSI6eyJib2R5Ijp7Im5hbWUiOiJub25lIn0sImhlYWQiOnsibmFtZSI6Im5vbmUifX19XSxbMCwxMSwiIiwxLHsic2hvcnRlbiI6eyJ0YXJnZXQiOjIwfSwic3R5bGUiOnsiYm9keSI6eyJuYW1lIjoibm9uZSJ9LCJoZWFkIjp7Im5hbWUiOiJub25lIn19fV0sWzAsMTIsIiIsMSx7InNob3J0ZW4iOnsidGFyZ2V0IjoyMH0sInN0eWxlIjp7ImJvZHkiOnsibmFtZSI6Im5vbmUifSwiaGVhZCI6eyJuYW1lIjoibm9uZSJ9fX1dLFsyLDEyLCIiLDEseyJzaG9ydGVuIjp7InRhcmdldCI6MjB9LCJzdHlsZSI6eyJib2R5Ijp7Im5hbWUiOiJub25lIn0sImhlYWQiOnsibmFtZSI6Im5vbmUifX19XSxbOSwxNCwiIiwxLHsic2hvcnRlbiI6eyJ0YXJnZXQiOjIwfSwic3R5bGUiOnsiYm9keSI6eyJuYW1lIjoibm9uZSJ9LCJoZWFkIjp7Im5hbWUiOiJub25lIn19fV0sWzYsMTQsIiIsMSx7InNob3J0ZW4iOnsidGFyZ2V0IjoyMH0sInN0eWxlIjp7ImJvZHkiOnsibmFtZSI6Im5vbmUifSwiaGVhZCI6eyJuYW1lIjoibm9uZSJ9fX1dLFs2LDE1LCIiLDEseyJzaG9ydGVuIjp7InRhcmdldCI6MjB9LCJzdHlsZSI6eyJib2R5Ijp7Im5hbWUiOiJub25lIn0sImhlYWQiOnsibmFtZSI6Im5vbmUifX19XSxbNywxNSwiIiwxLHsic2hvcnRlbiI6eyJ0YXJnZXQiOjIwfSwic3R5bGUiOnsiYm9keSI6eyJuYW1lIjoibm9uZSJ9LCJoZWFkIjp7Im5hbWUiOiJub25lIn19fV0sWzE5LDE4LCJcXGNkb3RzIiwxLHsibGV2ZWwiOjF9XSxbMjAsMiwiIiwxLHsibGV2ZWwiOjF9XSxbMjMsMjIsIlxcY2RvdHMiLDEseyJsZXZlbCI6MX1dLFsyNCw3LCIiLDEseyJsZXZlbCI6MX1dXQ==
\begin{tikzcd}[row sep=small, column sep=tiny,  every label/.append style = {font=\normalsize}]
	v && {} \\
	\\
	\\
	& v & \textcolor{white}{O} & \textcolor{white}{O} && {} \\
	\\
	\\
	&& {} && v & \textcolor{white}{O} & {} \\
	\\
	\\
	&&&&& {}
	\arrow["\cdots"{description}, sloped,,  ,  from=1-1, to=4-2]
	\arrow[""{name=0, anchor=center, inner sep=0}, "v"{description}, draw=none, from=1-3, to=4-2]
	\arrow[""{name=1, anchor=center, inner sep=0}, draw=none, from=4-4, to=7-3]
	\arrow["\cdots"{description}, sloped,,  ,  from=4-4, to=7-5]
	\arrow[""{name=2, anchor=center, inner sep=0}, "v"{description}, draw=none, from=4-6, to=7-5]
	\arrow[""{name=3, anchor=center, inner sep=0}, "v"{description}, draw=none, from=7-7, to=10-6]
	\arrow[""{name=4, anchor=center, inner sep=0}, draw=none, from=4-2, to=0]
	\arrow[""{name=5, anchor=center, inner sep=0}, draw=none, from=4-3, to=0]
	\arrow[""{name=6, anchor=center, inner sep=0}, draw=none, from=4-3, to=1]
	\arrow["\cdots"{description}, sloped,,  ,  draw=none, from=4-4, to=1]
	\arrow[""{name=7, anchor=center, inner sep=0}, draw=none, from=7-5, to=2]
	\arrow[""{name=8, anchor=center, inner sep=0}, draw=none, from=7-6, to=2]
	\arrow[""{name=9, anchor=center, inner sep=0}, draw=none, from=7-6, to=3]
	\arrow[draw=none, from=7-7, to=3]
	\arrow[from=4, to=4-2]
	\arrow["\cdots"{description}, sloped,,  ,  from=5, to=6]
	\arrow[from=7, to=7-5]
	\arrow["\cdots"{description}, sloped,,  ,  from=8, to=9]
\end{tikzcd}. \quad
		Case 2n.2.
% https://q.uiver.app/#q=WzAsMTAsWzEsMywiTyIsWzAsMCwxMDAsMV1dLFswLDMsInYiXSxbMSw2XSxbMiwzXSxbMSwwXSxbNCw2LCJPIixbMCwwLDEwMCwxXV0sWzMsNiwidiJdLFs0LDldLFs1LDYsInYiXSxbNCwzXSxbMywyLCIiLDEseyJzdHlsZSI6eyJib2R5Ijp7Im5hbWUiOiJub25lIn0sImhlYWQiOnsibmFtZSI6Im5vbmUifX19XSxbNCwxLCJ2IiwxLHsic3R5bGUiOnsiYm9keSI6eyJuYW1lIjoibm9uZSJ9LCJoZWFkIjp7Im5hbWUiOiJub25lIn19fV0sWzgsNywidiIsMSx7InN0eWxlIjp7ImJvZHkiOnsibmFtZSI6Im5vbmUifSwiaGVhZCI6eyJuYW1lIjoibm9uZSJ9fX1dLFs5LDYsInYiLDEseyJzdHlsZSI6eyJib2R5Ijp7Im5hbWUiOiJub25lIn0sImhlYWQiOnsibmFtZSI6Im5vbmUifX19XSxbMyw2LCJcXGNkb3RzIiwxXSxbMywxMCwiXFxjZG90cyIsMSx7InNob3J0ZW4iOnsidGFyZ2V0IjoyMH0sInN0eWxlIjp7ImJvZHkiOnsibmFtZSI6Im5vbmUifSwiaGVhZCI6eyJuYW1lIjoibm9uZSJ9fX1dLFswLDEwLCIiLDEseyJzaG9ydGVuIjp7InRhcmdldCI6MjB9LCJzdHlsZSI6eyJib2R5Ijp7Im5hbWUiOiJub25lIn0sImhlYWQiOnsibmFtZSI6Im5vbmUifX19XSxbMCwxMSwiIiwxLHsic2hvcnRlbiI6eyJ0YXJnZXQiOjIwfSwic3R5bGUiOnsiYm9keSI6eyJuYW1lIjoibm9uZSJ9LCJoZWFkIjp7Im5hbWUiOiJub25lIn19fV0sWzEsMTEsIiIsMSx7InNob3J0ZW4iOnsidGFyZ2V0IjoyMH0sInN0eWxlIjp7ImJvZHkiOnsibmFtZSI6Im5vbmUifSwiaGVhZCI6eyJuYW1lIjoibm9uZSJ9fX1dLFs4LDEyLCIiLDEseyJzaG9ydGVuIjp7InRhcmdldCI6MjB9LCJzdHlsZSI6eyJib2R5Ijp7Im5hbWUiOiJub25lIn0sImhlYWQiOnsibmFtZSI6Im5vbmUifX19XSxbNSwxMiwiIiwxLHsic2hvcnRlbiI6eyJ0YXJnZXQiOjIwfSwic3R5bGUiOnsiYm9keSI6eyJuYW1lIjoibm9uZSJ9LCJoZWFkIjp7Im5hbWUiOiJub25lIn19fV0sWzUsMTMsIiIsMSx7InNob3J0ZW4iOnsidGFyZ2V0IjoyMH0sInN0eWxlIjp7ImJvZHkiOnsibmFtZSI6Im5vbmUifSwiaGVhZCI6eyJuYW1lIjoibm9uZSJ9fX1dLFs2LDEzLCIiLDEseyJzaG9ydGVuIjp7InRhcmdldCI6MjB9LCJzdHlsZSI6eyJib2R5Ijp7Im5hbWUiOiJub25lIn0sImhlYWQiOnsibmFtZSI6Im5vbmUifX19XSxbMTcsMTYsIlxcY2RvdHMiLDEseyJsZXZlbCI6MX1dLFsxOCwxLCIiLDEseyJsZXZlbCI6MX1dLFsyMSwyMCwiXFxjZG90cyIsMSx7ImxldmVsIjoxfV0sWzIyLDYsIiIsMSx7ImxldmVsIjoxfV0sWzgsMTksIiIsMSx7ImxldmVsIjoxfV1d
\begin{tikzcd}[row sep=small, column sep=tiny,  every label/.append style = {font=\normalsize}]
	& {} \\
	\\
	\\
	v & \textcolor{white}{O} & \textcolor{white}{O} && {} \\
	\\
	\\
	& {} && v & \textcolor{white}{O} & v \\
	\\
	\\
	&&&& {}
	\arrow[""{name=0, anchor=center, inner sep=0}, "v"{description}, draw=none, from=1-2, to=4-1]
	\arrow[""{name=1, anchor=center, inner sep=0}, draw=none, from=4-3, to=7-2]
	\arrow["\cdots"{description}, sloped,,  ,  from=4-3, to=7-4]
	\arrow[""{name=2, anchor=center, inner sep=0}, "v"{description}, draw=none, from=4-5, to=7-4]
	\arrow[""{name=3, anchor=center, inner sep=0}, "v"{description}, draw=none, from=7-6, to=10-5]
	\arrow[""{name=4, anchor=center, inner sep=0}, draw=none, from=4-1, to=0]
	\arrow[""{name=5, anchor=center, inner sep=0}, draw=none, from=4-2, to=0]
	\arrow[""{name=6, anchor=center, inner sep=0}, draw=none, from=4-2, to=1]
	\arrow["\cdots"{description}, sloped,,  ,  draw=none, from=4-3, to=1]
	\arrow[""{name=7, anchor=center, inner sep=0}, draw=none, from=7-4, to=2]
	\arrow[""{name=8, anchor=center, inner sep=0}, draw=none, from=7-5, to=2]
	\arrow[""{name=9, anchor=center, inner sep=0}, draw=none, from=7-5, to=3]
	\arrow[""{name=10, anchor=center, inner sep=0}, draw=none, from=7-6, to=3]
	\arrow[from=4, to=4-1]
	\arrow["\cdots"{description}, sloped,,  ,  from=5, to=6]
	\arrow[from=7, to=7-4]
	\arrow["\cdots"{description}, sloped,,  ,  from=8, to=9]
	\arrow[from=7-6, to=10]
\end{tikzcd}. \\
		Case 2n.3.
% https://q.uiver.app/#q=WzAsMTIsWzMsMywiTyIsWzAsMCwxMDAsMV1dLFsyLDMsInYiXSxbMyw2XSxbNCwzXSxbMywwXSxbNiw2LCJPIixbMCwwLDEwMCwxXV0sWzUsNiwidiJdLFs2LDldLFs3LDYsInYiXSxbNiwzXSxbMSwwLCJ4Il0sWzAsM10sWzMsMiwiIiwxLHsic3R5bGUiOnsiYm9keSI6eyJuYW1lIjoibm9uZSJ9LCJoZWFkIjp7Im5hbWUiOiJub25lIn19fV0sWzQsMSwidiIsMSx7InN0eWxlIjp7ImJvZHkiOnsibmFtZSI6Im5vbmUifSwiaGVhZCI6eyJuYW1lIjoibm9uZSJ9fX1dLFs4LDcsInYiLDEseyJzdHlsZSI6eyJib2R5Ijp7Im5hbWUiOiJub25lIn0sImhlYWQiOnsibmFtZSI6Im5vbmUifX19XSxbOSw2LCJ2IiwxLHsic3R5bGUiOnsiYm9keSI6eyJuYW1lIjoibm9uZSJ9LCJoZWFkIjp7Im5hbWUiOiJub25lIn19fV0sWzMsNiwiXFxjZG90cyIsMV0sWzEwLDEsIlxcY2RvdHMiLDFdLFsxMCwxMSwiXFxjZG90cyIsMSx7InN0eWxlIjp7ImJvZHkiOnsibmFtZSI6Im5vbmUifSwiaGVhZCI6eyJuYW1lIjoibm9uZSJ9fX1dLFszLDEyLCJcXGNkb3RzIiwxLHsic2hvcnRlbiI6eyJ0YXJnZXQiOjIwfSwic3R5bGUiOnsiYm9keSI6eyJuYW1lIjoibm9uZSJ9LCJoZWFkIjp7Im5hbWUiOiJub25lIn19fV0sWzAsMTIsIiIsMSx7InNob3J0ZW4iOnsidGFyZ2V0IjoyMH0sInN0eWxlIjp7ImJvZHkiOnsibmFtZSI6Im5vbmUifSwiaGVhZCI6eyJuYW1lIjoibm9uZSJ9fX1dLFswLDEzLCIiLDEseyJzaG9ydGVuIjp7InRhcmdldCI6MjB9LCJzdHlsZSI6eyJib2R5Ijp7Im5hbWUiOiJub25lIn0sImhlYWQiOnsibmFtZSI6Im5vbmUifX19XSxbMSwxMywiIiwxLHsic2hvcnRlbiI6eyJ0YXJnZXQiOjIwfSwic3R5bGUiOnsiYm9keSI6eyJuYW1lIjoibm9uZSJ9LCJoZWFkIjp7Im5hbWUiOiJub25lIn19fV0sWzgsMTQsIiIsMSx7InNob3J0ZW4iOnsidGFyZ2V0IjoyMH0sInN0eWxlIjp7ImJvZHkiOnsibmFtZSI6Im5vbmUifSwiaGVhZCI6eyJuYW1lIjoibm9uZSJ9fX1dLFs1LDE0LCIiLDEseyJzaG9ydGVuIjp7InRhcmdldCI6MjB9LCJzdHlsZSI6eyJib2R5Ijp7Im5hbWUiOiJub25lIn0sImhlYWQiOnsibmFtZSI6Im5vbmUifX19XSxbNSwxNSwiIiwxLHsic2hvcnRlbiI6eyJ0YXJnZXQiOjIwfSwic3R5bGUiOnsiYm9keSI6eyJuYW1lIjoibm9uZSJ9LCJoZWFkIjp7Im5hbWUiOiJub25lIn19fV0sWzYsMTUsIiIsMSx7InNob3J0ZW4iOnsidGFyZ2V0IjoyMH0sInN0eWxlIjp7ImJvZHkiOnsibmFtZSI6Im5vbmUifSwiaGVhZCI6eyJuYW1lIjoibm9uZSJ9fX1dLFsxMCwxOCwiIiwxLHsic2hvcnRlbiI6eyJ0YXJnZXQiOjIwfSwic3R5bGUiOnsiYm9keSI6eyJuYW1lIjoibm9uZSJ9LCJoZWFkIjp7Im5hbWUiOiJub25lIn19fV0sWzIxLDIwLCJcXGNkb3RzIiwxLHsibGV2ZWwiOjF9XSxbMjIsMSwiIiwxLHsibGV2ZWwiOjF9XSxbMjUsMjQsIlxcY2RvdHMiLDEseyJsZXZlbCI6MX1dLFsyNiw2LCIiLDEseyJsZXZlbCI6MX1dLFs4LDIzLCIiLDEseyJsZXZlbCI6MX1dLFsxMCwyNywiIiwxLHsibGV2ZWwiOjF9XV0=
\begin{tikzcd}[row sep=small, column sep=tiny,  every label/.append style = {font=\normalsize}]
	& x && {} \\
	\\
	\\
	{} && v & \textcolor{white}{O} & \textcolor{white}{O} && {} \\
	\\
	\\
	&&& {} && v & \textcolor{white}{O} & v \\
	\\
	\\
	&&&&&& {}
	\arrow[""{name=0, anchor=center, inner sep=0}, "\cdots"{description}, sloped,,  ,  draw=none, from=1-2, to=4-1]
	\arrow["\cdots"{description}, sloped,,  ,  from=1-2, to=4-3]
	\arrow[""{name=1, anchor=center, inner sep=0}, "v"{description}, draw=none, from=1-4, to=4-3]
	\arrow[""{name=2, anchor=center, inner sep=0}, draw=none, from=4-5, to=7-4]
	\arrow["\cdots"{description}, sloped,,  ,  from=4-5, to=7-6]
	\arrow[""{name=3, anchor=center, inner sep=0}, "v"{description}, draw=none, from=4-7, to=7-6]
	\arrow[""{name=4, anchor=center, inner sep=0}, "v"{description}, draw=none, from=7-8, to=10-7]
	\arrow[""{name=5, anchor=center, inner sep=0}, draw=none, from=1-2, to=0]
	\arrow[""{name=6, anchor=center, inner sep=0}, draw=none, from=4-3, to=1]
	\arrow[""{name=7, anchor=center, inner sep=0}, draw=none, from=4-4, to=1]
	\arrow[""{name=8, anchor=center, inner sep=0}, draw=none, from=4-4, to=2]
	\arrow["\cdots"{description}, sloped,,  ,  draw=none, from=4-5, to=2]
	\arrow[""{name=9, anchor=center, inner sep=0}, draw=none, from=7-6, to=3]
	\arrow[""{name=10, anchor=center, inner sep=0}, draw=none, from=7-7, to=3]
	\arrow[""{name=11, anchor=center, inner sep=0}, draw=none, from=7-7, to=4]
	\arrow[""{name=12, anchor=center, inner sep=0}, draw=none, from=7-8, to=4]
	\arrow[from=1-2, to=5]
	\arrow[from=6, to=4-3]
	\arrow["\cdots"{description}, sloped,,  ,  from=7, to=8]
	\arrow[from=9, to=7-6]
	\arrow["\cdots"{description}, sloped,,  ,  from=10, to=11]
	\arrow[from=7-8, to=12]
\end{tikzcd}. \quad
		Case 2n.4.
% https://q.uiver.app/#q=WzAsMTIsWzEsMywiTyIsWzAsMCwxMDAsMV1dLFswLDMsInYiXSxbMSw2XSxbMiwzXSxbMSwwXSxbNCw2LCJPIixbMCwwLDEwMCwxXV0sWzMsNiwidiJdLFs0LDldLFs1LDYsInYiXSxbNCwzXSxbNiwzLCJ4Il0sWzcsNl0sWzMsMiwiIiwxLHsic3R5bGUiOnsiYm9keSI6eyJuYW1lIjoibm9uZSJ9LCJoZWFkIjp7Im5hbWUiOiJub25lIn19fV0sWzQsMSwidiIsMSx7InN0eWxlIjp7ImJvZHkiOnsibmFtZSI6Im5vbmUifSwiaGVhZCI6eyJuYW1lIjoibm9uZSJ9fX1dLFs4LDcsInYiLDEseyJzdHlsZSI6eyJib2R5Ijp7Im5hbWUiOiJub25lIn0sImhlYWQiOnsibmFtZSI6Im5vbmUifX19XSxbOSw2LCJ2IiwxLHsic3R5bGUiOnsiYm9keSI6eyJuYW1lIjoibm9uZSJ9LCJoZWFkIjp7Im5hbWUiOiJub25lIn19fV0sWzMsNiwiXFxjZG90cyIsMV0sWzEwLDgsIlxcY2RvdHMiLDFdLFsxMCwxMSwiXFxjZG90cyIsMSx7InN0eWxlIjp7ImJvZHkiOnsibmFtZSI6Im5vbmUifSwiaGVhZCI6eyJuYW1lIjoibm9uZSJ9fX1dLFszLDEyLCJcXGNkb3RzIiwxLHsic2hvcnRlbiI6eyJ0YXJnZXQiOjIwfSwic3R5bGUiOnsiYm9keSI6eyJuYW1lIjoibm9uZSJ9LCJoZWFkIjp7Im5hbWUiOiJub25lIn19fV0sWzAsMTIsIiIsMSx7InNob3J0ZW4iOnsidGFyZ2V0IjoyMH0sInN0eWxlIjp7ImJvZHkiOnsibmFtZSI6Im5vbmUifSwiaGVhZCI6eyJuYW1lIjoibm9uZSJ9fX1dLFswLDEzLCIiLDEseyJzaG9ydGVuIjp7InRhcmdldCI6MjB9LCJzdHlsZSI6eyJib2R5Ijp7Im5hbWUiOiJub25lIn0sImhlYWQiOnsibmFtZSI6Im5vbmUifX19XSxbMSwxMywiIiwxLHsic2hvcnRlbiI6eyJ0YXJnZXQiOjIwfSwic3R5bGUiOnsiYm9keSI6eyJuYW1lIjoibm9uZSJ9LCJoZWFkIjp7Im5hbWUiOiJub25lIn19fV0sWzgsMTQsIiIsMSx7InNob3J0ZW4iOnsidGFyZ2V0IjoyMH0sInN0eWxlIjp7ImJvZHkiOnsibmFtZSI6Im5vbmUifSwiaGVhZCI6eyJuYW1lIjoibm9uZSJ9fX1dLFs1LDE0LCIiLDEseyJzaG9ydGVuIjp7InRhcmdldCI6MjB9LCJzdHlsZSI6eyJib2R5Ijp7Im5hbWUiOiJub25lIn0sImhlYWQiOnsibmFtZSI6Im5vbmUifX19XSxbNSwxNSwiIiwxLHsic2hvcnRlbiI6eyJ0YXJnZXQiOjIwfSwic3R5bGUiOnsiYm9keSI6eyJuYW1lIjoibm9uZSJ9LCJoZWFkIjp7Im5hbWUiOiJub25lIn19fV0sWzYsMTUsIiIsMSx7InNob3J0ZW4iOnsidGFyZ2V0IjoyMH0sInN0eWxlIjp7ImJvZHkiOnsibmFtZSI6Im5vbmUifSwiaGVhZCI6eyJuYW1lIjoibm9uZSJ9fX1dLFsxMCwxOCwiIiwxLHsic2hvcnRlbiI6eyJ0YXJnZXQiOjIwfSwic3R5bGUiOnsiYm9keSI6eyJuYW1lIjoibm9uZSJ9LCJoZWFkIjp7Im5hbWUiOiJub25lIn19fV0sWzIxLDIwLCJcXGNkb3RzIiwxLHsibGV2ZWwiOjF9XSxbMjIsMSwiIiwxLHsibGV2ZWwiOjF9XSxbMjUsMjQsIlxcY2RvdHMiLDEseyJsZXZlbCI6MX1dLFsyNiw2LCIiLDEseyJsZXZlbCI6MX1dLFs4LDIzLCIiLDEseyJsZXZlbCI6MX1dLFsxMCwyNywiIiwxLHsibGV2ZWwiOjF9XV0=
\begin{tikzcd}[row sep=small, column sep=tiny,  every label/.append style = {font=\normalsize}]
	& {} \\
	\\
	\\
	v & \textcolor{white}{O} & \textcolor{white}{O} && {} && x \\
	\\
	\\
	& {} && v & \textcolor{white}{O} & v && {} \\
	\\
	\\
	&&&& {}
	\arrow[""{name=0, anchor=center, inner sep=0}, "v"{description}, draw=none, from=1-2, to=4-1]
	\arrow[""{name=1, anchor=center, inner sep=0}, draw=none, from=4-3, to=7-2]
	\arrow["\cdots"{description}, sloped,,  ,  from=4-3, to=7-4]
	\arrow[""{name=2, anchor=center, inner sep=0}, "v"{description}, draw=none, from=4-5, to=7-4]
	\arrow["\cdots"{description}, sloped,,  ,  from=4-7, to=7-6]
	\arrow[""{name=3, anchor=center, inner sep=0}, "\cdots"{description}, sloped,,  ,  draw=none, from=4-7, to=7-8]
	\arrow[""{name=4, anchor=center, inner sep=0}, "v"{description}, draw=none, from=7-6, to=10-5]
	\arrow[""{name=5, anchor=center, inner sep=0}, draw=none, from=4-1, to=0]
	\arrow[""{name=6, anchor=center, inner sep=0}, draw=none, from=4-2, to=0]
	\arrow[""{name=7, anchor=center, inner sep=0}, draw=none, from=4-2, to=1]
	\arrow["\cdots"{description}, sloped,,  ,  draw=none, from=4-3, to=1]
	\arrow[""{name=8, anchor=center, inner sep=0}, draw=none, from=4-7, to=3]
	\arrow[""{name=9, anchor=center, inner sep=0}, draw=none, from=7-4, to=2]
	\arrow[""{name=10, anchor=center, inner sep=0}, draw=none, from=7-5, to=2]
	\arrow[""{name=11, anchor=center, inner sep=0}, draw=none, from=7-5, to=4]
	\arrow[""{name=12, anchor=center, inner sep=0}, draw=none, from=7-6, to=4]
	\arrow[from=5, to=4-1]
	\arrow["\cdots"{description}, sloped,,  ,  from=6, to=7]
	\arrow[from=4-7, to=8]
	\arrow[from=9, to=7-4]
	\arrow["\cdots"{description}, sloped,,  ,  from=10, to=11]
	\arrow[from=7-6, to=12]
\end{tikzcd}. \\
		Case 2n.5.
% https://q.uiver.app/#q=WzAsMTIsWzMsMywiTyIsWzAsMCwxMDAsMV1dLFsyLDMsInYiXSxbMyw2XSxbNCwzXSxbMywwXSxbNiw2LCJPIixbMCwwLDEwMCwxXV0sWzUsNiwidiJdLFs2LDldLFs3LDYsInYiXSxbNiwzXSxbMSw2LCJ4Il0sWzAsM10sWzMsMiwiIiwxLHsic3R5bGUiOnsiYm9keSI6eyJuYW1lIjoibm9uZSJ9LCJoZWFkIjp7Im5hbWUiOiJub25lIn19fV0sWzQsMSwidiIsMSx7InN0eWxlIjp7ImJvZHkiOnsibmFtZSI6Im5vbmUifSwiaGVhZCI6eyJuYW1lIjoibm9uZSJ9fX1dLFs4LDcsInYiLDEseyJzdHlsZSI6eyJib2R5Ijp7Im5hbWUiOiJub25lIn0sImhlYWQiOnsibmFtZSI6Im5vbmUifX19XSxbOSw2LCJ2IiwxLHsic3R5bGUiOnsiYm9keSI6eyJuYW1lIjoibm9uZSJ9LCJoZWFkIjp7Im5hbWUiOiJub25lIn19fV0sWzMsNiwiXFxjZG90cyIsMV0sWzEsMTAsIlxcY2RvdHMiLDFdLFsxMCwxMSwiXFxjZG90cyIsMSx7InN0eWxlIjp7ImJvZHkiOnsibmFtZSI6Im5vbmUifSwiaGVhZCI6eyJuYW1lIjoibm9uZSJ9fX1dLFszLDEyLCJcXGNkb3RzIiwxLHsic2hvcnRlbiI6eyJ0YXJnZXQiOjIwfSwic3R5bGUiOnsiYm9keSI6eyJuYW1lIjoibm9uZSJ9LCJoZWFkIjp7Im5hbWUiOiJub25lIn19fV0sWzAsMTIsIiIsMSx7InNob3J0ZW4iOnsidGFyZ2V0IjoyMH0sInN0eWxlIjp7ImJvZHkiOnsibmFtZSI6Im5vbmUifSwiaGVhZCI6eyJuYW1lIjoibm9uZSJ9fX1dLFswLDEzLCIiLDEseyJzaG9ydGVuIjp7InRhcmdldCI6MjB9LCJzdHlsZSI6eyJib2R5Ijp7Im5hbWUiOiJub25lIn0sImhlYWQiOnsibmFtZSI6Im5vbmUifX19XSxbMSwxMywiIiwxLHsic2hvcnRlbiI6eyJ0YXJnZXQiOjIwfSwic3R5bGUiOnsiYm9keSI6eyJuYW1lIjoibm9uZSJ9LCJoZWFkIjp7Im5hbWUiOiJub25lIn19fV0sWzgsMTQsIiIsMSx7InNob3J0ZW4iOnsidGFyZ2V0IjoyMH0sInN0eWxlIjp7ImJvZHkiOnsibmFtZSI6Im5vbmUifSwiaGVhZCI6eyJuYW1lIjoibm9uZSJ9fX1dLFs1LDE0LCIiLDEseyJzaG9ydGVuIjp7InRhcmdldCI6MjB9LCJzdHlsZSI6eyJib2R5Ijp7Im5hbWUiOiJub25lIn0sImhlYWQiOnsibmFtZSI6Im5vbmUifX19XSxbNSwxNSwiIiwxLHsic2hvcnRlbiI6eyJ0YXJnZXQiOjIwfSwic3R5bGUiOnsiYm9keSI6eyJuYW1lIjoibm9uZSJ9LCJoZWFkIjp7Im5hbWUiOiJub25lIn19fV0sWzYsMTUsIiIsMSx7InNob3J0ZW4iOnsidGFyZ2V0IjoyMH0sInN0eWxlIjp7ImJvZHkiOnsibmFtZSI6Im5vbmUifSwiaGVhZCI6eyJuYW1lIjoibm9uZSJ9fX1dLFsxMCwxOCwiIiwxLHsic2hvcnRlbiI6eyJ0YXJnZXQiOjIwfSwic3R5bGUiOnsiYm9keSI6eyJuYW1lIjoibm9uZSJ9LCJoZWFkIjp7Im5hbWUiOiJub25lIn19fV0sWzIxLDIwLCJcXGNkb3RzIiwxLHsibGV2ZWwiOjF9XSxbMjIsMSwiIiwxLHsibGV2ZWwiOjF9XSxbMjUsMjQsIlxcY2RvdHMiLDEseyJsZXZlbCI6MX1dLFsyNiw2LCIiLDEseyJsZXZlbCI6MX1dLFs4LDIzLCIiLDEseyJsZXZlbCI6MX1dLFsyNywxMCwiIiwxLHsibGV2ZWwiOjF9XV0=
\begin{tikzcd}[row sep=small, column sep=tiny,  every label/.append style = {font=\normalsize}]
	&&& {} \\
	\\
	\\
	{} && v & \textcolor{white}{O} & \textcolor{white}{O} && {} \\
	\\
	\\
	& x && {} && v & \textcolor{white}{O} & v \\
	\\
	\\
	&&&&&& {}
	\arrow[""{name=0, anchor=center, inner sep=0}, "v"{description}, draw=none, from=1-4, to=4-3]
	\arrow["\cdots"{description}, sloped,,  ,  from=4-3, to=7-2]
	\arrow[""{name=1, anchor=center, inner sep=0}, draw=none, from=4-5, to=7-4]
	\arrow["\cdots"{description}, sloped,,  ,  from=4-5, to=7-6]
	\arrow[""{name=2, anchor=center, inner sep=0}, "v"{description}, draw=none, from=4-7, to=7-6]
	\arrow[""{name=3, anchor=center, inner sep=0}, "\cdots"{description}, sloped,,  ,  draw=none, from=7-2, to=4-1]
	\arrow[""{name=4, anchor=center, inner sep=0}, "v"{description}, draw=none, from=7-8, to=10-7]
	\arrow[""{name=5, anchor=center, inner sep=0}, draw=none, from=4-3, to=0]
	\arrow[""{name=6, anchor=center, inner sep=0}, draw=none, from=4-4, to=0]
	\arrow[""{name=7, anchor=center, inner sep=0}, draw=none, from=4-4, to=1]
	\arrow["\cdots"{description}, sloped,,  ,  draw=none, from=4-5, to=1]
	\arrow[""{name=8, anchor=center, inner sep=0}, draw=none, from=7-2, to=3]
	\arrow[""{name=9, anchor=center, inner sep=0}, draw=none, from=7-6, to=2]
	\arrow[""{name=10, anchor=center, inner sep=0}, draw=none, from=7-7, to=2]
	\arrow[""{name=11, anchor=center, inner sep=0}, draw=none, from=7-7, to=4]
	\arrow[""{name=12, anchor=center, inner sep=0}, draw=none, from=7-8, to=4]
	\arrow[from=5, to=4-3]
	\arrow["\cdots"{description}, sloped,,  ,  from=6, to=7]
	\arrow[from=8, to=7-2]
	\arrow[from=9, to=7-6]
	\arrow["\cdots"{description}, sloped,,  ,  from=10, to=11]
	\arrow[from=7-8, to=12]
\end{tikzcd}. \quad
		Case 2n.6.
% https://q.uiver.app/#q=WzAsMTIsWzEsMywiTyIsWzAsMCwxMDAsMV1dLFswLDMsInYiXSxbMSw2XSxbMiwzXSxbMSwwXSxbNCw2LCJPIixbMCwwLDEwMCwxXV0sWzMsNiwidiJdLFs0LDldLFs1LDYsInYiXSxbNCwzXSxbNiw5LCJ4Il0sWzcsNl0sWzMsMiwiIiwxLHsic3R5bGUiOnsiYm9keSI6eyJuYW1lIjoibm9uZSJ9LCJoZWFkIjp7Im5hbWUiOiJub25lIn19fV0sWzQsMSwidiIsMSx7InN0eWxlIjp7ImJvZHkiOnsibmFtZSI6Im5vbmUifSwiaGVhZCI6eyJuYW1lIjoibm9uZSJ9fX1dLFs4LDcsInYiLDEseyJzdHlsZSI6eyJib2R5Ijp7Im5hbWUiOiJub25lIn0sImhlYWQiOnsibmFtZSI6Im5vbmUifX19XSxbOSw2LCJ2IiwxLHsic3R5bGUiOnsiYm9keSI6eyJuYW1lIjoibm9uZSJ9LCJoZWFkIjp7Im5hbWUiOiJub25lIn19fV0sWzMsNiwiXFxjZG90cyIsMV0sWzgsMTAsIlxcY2RvdHMiLDFdLFsxMCwxMSwiXFxjZG90cyIsMSx7InN0eWxlIjp7ImJvZHkiOnsibmFtZSI6Im5vbmUifSwiaGVhZCI6eyJuYW1lIjoibm9uZSJ9fX1dLFszLDEyLCJcXGNkb3RzIiwxLHsic2hvcnRlbiI6eyJ0YXJnZXQiOjIwfSwic3R5bGUiOnsiYm9keSI6eyJuYW1lIjoibm9uZSJ9LCJoZWFkIjp7Im5hbWUiOiJub25lIn19fV0sWzAsMTIsIiIsMSx7InNob3J0ZW4iOnsidGFyZ2V0IjoyMH0sInN0eWxlIjp7ImJvZHkiOnsibmFtZSI6Im5vbmUifSwiaGVhZCI6eyJuYW1lIjoibm9uZSJ9fX1dLFswLDEzLCIiLDEseyJzaG9ydGVuIjp7InRhcmdldCI6MjB9LCJzdHlsZSI6eyJib2R5Ijp7Im5hbWUiOiJub25lIn0sImhlYWQiOnsibmFtZSI6Im5vbmUifX19XSxbMSwxMywiIiwxLHsic2hvcnRlbiI6eyJ0YXJnZXQiOjIwfSwic3R5bGUiOnsiYm9keSI6eyJuYW1lIjoibm9uZSJ9LCJoZWFkIjp7Im5hbWUiOiJub25lIn19fV0sWzgsMTQsIiIsMSx7InNob3J0ZW4iOnsidGFyZ2V0IjoyMH0sInN0eWxlIjp7ImJvZHkiOnsibmFtZSI6Im5vbmUifSwiaGVhZCI6eyJuYW1lIjoibm9uZSJ9fX1dLFs1LDE0LCIiLDEseyJzaG9ydGVuIjp7InRhcmdldCI6MjB9LCJzdHlsZSI6eyJib2R5Ijp7Im5hbWUiOiJub25lIn0sImhlYWQiOnsibmFtZSI6Im5vbmUifX19XSxbNSwxNSwiIiwxLHsic2hvcnRlbiI6eyJ0YXJnZXQiOjIwfSwic3R5bGUiOnsiYm9keSI6eyJuYW1lIjoibm9uZSJ9LCJoZWFkIjp7Im5hbWUiOiJub25lIn19fV0sWzYsMTUsIiIsMSx7InNob3J0ZW4iOnsidGFyZ2V0IjoyMH0sInN0eWxlIjp7ImJvZHkiOnsibmFtZSI6Im5vbmUifSwiaGVhZCI6eyJuYW1lIjoibm9uZSJ9fX1dLFsxMCwxOCwiIiwxLHsic2hvcnRlbiI6eyJ0YXJnZXQiOjIwfSwic3R5bGUiOnsiYm9keSI6eyJuYW1lIjoibm9uZSJ9LCJoZWFkIjp7Im5hbWUiOiJub25lIn19fV0sWzIxLDIwLCJcXGNkb3RzIiwxLHsibGV2ZWwiOjF9XSxbMjIsMSwiIiwxLHsibGV2ZWwiOjF9XSxbMjUsMjQsIlxcY2RvdHMiLDEseyJsZXZlbCI6MX1dLFsyNiw2LCIiLDEseyJsZXZlbCI6MX1dLFs4LDIzLCIiLDEseyJsZXZlbCI6MX1dLFsyNywxMCwiIiwxLHsibGV2ZWwiOjF9XV0=
\begin{tikzcd}[row sep=small, column sep=tiny,  every label/.append style = {font=\normalsize}]
	& {} \\
	\\
	\\
	v & \textcolor{white}{O} & \textcolor{white}{O} && {} \\
	\\
	\\
	& {} && v & \textcolor{white}{O} & v && {} \\
	\\
	\\
	&&&& {} && x
	\arrow[""{name=0, anchor=center, inner sep=0}, "v"{description}, draw=none, from=1-2, to=4-1]
	\arrow[""{name=1, anchor=center, inner sep=0}, draw=none, from=4-3, to=7-2]
	\arrow["\cdots"{description}, sloped,,  ,  from=4-3, to=7-4]
	\arrow[""{name=2, anchor=center, inner sep=0}, "v"{description}, draw=none, from=4-5, to=7-4]
	\arrow[""{name=3, anchor=center, inner sep=0}, "v"{description}, draw=none, from=7-6, to=10-5]
	\arrow["\cdots"{description}, sloped,,  ,  from=7-6, to=10-7]
	\arrow[""{name=4, anchor=center, inner sep=0}, "\cdots"{description}, sloped,,  ,  draw=none, from=10-7, to=7-8]
	\arrow[""{name=5, anchor=center, inner sep=0}, draw=none, from=4-1, to=0]
	\arrow[""{name=6, anchor=center, inner sep=0}, draw=none, from=4-2, to=0]
	\arrow[""{name=7, anchor=center, inner sep=0}, draw=none, from=4-2, to=1]
	\arrow["\cdots"{description}, sloped,,  ,  draw=none, from=4-3, to=1]
	\arrow[""{name=8, anchor=center, inner sep=0}, draw=none, from=7-4, to=2]
	\arrow[""{name=9, anchor=center, inner sep=0}, draw=none, from=7-5, to=2]
	\arrow[""{name=10, anchor=center, inner sep=0}, draw=none, from=7-5, to=3]
	\arrow[""{name=11, anchor=center, inner sep=0}, draw=none, from=7-6, to=3]
	\arrow[""{name=12, anchor=center, inner sep=0}, draw=none, from=10-7, to=4]
	\arrow[from=5, to=4-1]
	\arrow["\cdots"{description}, sloped,,  ,  from=6, to=7]
	\arrow[from=8, to=7-4]
	\arrow["\cdots"{description}, sloped,,  ,  from=9, to=10]
	\arrow[from=7-6, to=11]
	\arrow[from=12, to=10-7]
\end{tikzcd}. \\
\\
		Case 2n.7.
% https://q.uiver.app/#q=WzAsMTQsWzMsMywiTyIsWzAsMCwxMDAsMV1dLFsyLDMsInYiXSxbMyw2XSxbNCwzXSxbMywwXSxbNiw2LCJPIixbMCwwLDEwMCwxXV0sWzUsNiwidiJdLFs2LDldLFs3LDYsInYiXSxbNiwzXSxbMSwwLCJ4Il0sWzAsM10sWzgsMywieCciXSxbOSw2XSxbMywyLCIiLDEseyJzdHlsZSI6eyJib2R5Ijp7Im5hbWUiOiJub25lIn0sImhlYWQiOnsibmFtZSI6Im5vbmUifX19XSxbNCwxLCJ2IiwxLHsic3R5bGUiOnsiYm9keSI6eyJuYW1lIjoibm9uZSJ9LCJoZWFkIjp7Im5hbWUiOiJub25lIn19fV0sWzgsNywidiIsMSx7InN0eWxlIjp7ImJvZHkiOnsibmFtZSI6Im5vbmUifSwiaGVhZCI6eyJuYW1lIjoibm9uZSJ9fX1dLFs5LDYsInYiLDEseyJzdHlsZSI6eyJib2R5Ijp7Im5hbWUiOiJub25lIn0sImhlYWQiOnsibmFtZSI6Im5vbmUifX19XSxbMyw2LCJcXGNkb3RzIiwxXSxbMTAsMSwiXFxjZG90cyIsMV0sWzEwLDExLCJcXGNkb3RzIiwxLHsic3R5bGUiOnsiYm9keSI6eyJuYW1lIjoibm9uZSJ9LCJoZWFkIjp7Im5hbWUiOiJub25lIn19fV0sWzEyLDgsIlxcY2RvdHMiLDFdLFsxMiwxMywiXFxjZG90cyIsMSx7InN0eWxlIjp7ImJvZHkiOnsibmFtZSI6Im5vbmUifSwiaGVhZCI6eyJuYW1lIjoibm9uZSJ9fX1dLFszLDE0LCJcXGNkb3RzIiwxLHsic2hvcnRlbiI6eyJ0YXJnZXQiOjIwfSwic3R5bGUiOnsiYm9keSI6eyJuYW1lIjoibm9uZSJ9LCJoZWFkIjp7Im5hbWUiOiJub25lIn19fV0sWzAsMTQsIiIsMSx7InNob3J0ZW4iOnsidGFyZ2V0IjoyMH0sInN0eWxlIjp7ImJvZHkiOnsibmFtZSI6Im5vbmUifSwiaGVhZCI6eyJuYW1lIjoibm9uZSJ9fX1dLFswLDE1LCIiLDEseyJzaG9ydGVuIjp7InRhcmdldCI6MjB9LCJzdHlsZSI6eyJib2R5Ijp7Im5hbWUiOiJub25lIn0sImhlYWQiOnsibmFtZSI6Im5vbmUifX19XSxbMSwxNSwiIiwxLHsic2hvcnRlbiI6eyJ0YXJnZXQiOjIwfSwic3R5bGUiOnsiYm9keSI6eyJuYW1lIjoibm9uZSJ9LCJoZWFkIjp7Im5hbWUiOiJub25lIn19fV0sWzgsMTYsIiIsMSx7InNob3J0ZW4iOnsidGFyZ2V0IjoyMH0sInN0eWxlIjp7ImJvZHkiOnsibmFtZSI6Im5vbmUifSwiaGVhZCI6eyJuYW1lIjoibm9uZSJ9fX1dLFs1LDE2LCIiLDEseyJzaG9ydGVuIjp7InRhcmdldCI6MjB9LCJzdHlsZSI6eyJib2R5Ijp7Im5hbWUiOiJub25lIn0sImhlYWQiOnsibmFtZSI6Im5vbmUifX19XSxbNSwxNywiIiwxLHsic2hvcnRlbiI6eyJ0YXJnZXQiOjIwfSwic3R5bGUiOnsiYm9keSI6eyJuYW1lIjoibm9uZSJ9LCJoZWFkIjp7Im5hbWUiOiJub25lIn19fV0sWzYsMTcsIiIsMSx7InNob3J0ZW4iOnsidGFyZ2V0IjoyMH0sInN0eWxlIjp7ImJvZHkiOnsibmFtZSI6Im5vbmUifSwiaGVhZCI6eyJuYW1lIjoibm9uZSJ9fX1dLFsxMCwyMCwiIiwxLHsic2hvcnRlbiI6eyJ0YXJnZXQiOjIwfSwic3R5bGUiOnsiYm9keSI6eyJuYW1lIjoibm9uZSJ9LCJoZWFkIjp7Im5hbWUiOiJub25lIn19fV0sWzEyLDIyLCIiLDEseyJzaG9ydGVuIjp7InRhcmdldCI6MjB9LCJzdHlsZSI6eyJib2R5Ijp7Im5hbWUiOiJub25lIn0sImhlYWQiOnsibmFtZSI6Im5vbmUifX19XSxbMjUsMjQsIlxcY2RvdHMiLDEseyJsZXZlbCI6MX1dLFsyNiwxLCIiLDEseyJsZXZlbCI6MX1dLFsyOSwyOCwiXFxjZG90cyIsMSx7ImxldmVsIjoxfV0sWzMwLDYsIiIsMSx7ImxldmVsIjoxfV0sWzgsMjcsIiIsMSx7ImxldmVsIjoxfV0sWzEwLDMxLCIiLDEseyJsZXZlbCI6MX1dLFsxMiwzMiwiIiwxLHsibGV2ZWwiOjF9XV0=
\begin{tikzcd}[row sep=small, column sep=tiny,  every label/.append style = {font=\normalsize}]
	& x && {} \\
	\\
	\\
	{} && v & \textcolor{white}{O} & \textcolor{white}{O} && {} && {x'} \\
	\\
	\\
	&&& {} && v & \textcolor{white}{O} & v && {} \\
	\\
	\\
	&&&&&& {}
	\arrow[""{name=0, anchor=center, inner sep=0}, "\cdots"{description}, sloped,,  ,  draw=none, from=1-2, to=4-1]
	\arrow["\cdots"{description}, sloped,,  ,  from=1-2, to=4-3]
	\arrow[""{name=1, anchor=center, inner sep=0}, "v"{description}, draw=none, from=1-4, to=4-3]
	\arrow[""{name=2, anchor=center, inner sep=0}, draw=none, from=4-5, to=7-4]
	\arrow["\cdots"{description}, sloped,,  ,  from=4-5, to=7-6]
	\arrow[""{name=3, anchor=center, inner sep=0}, "v"{description}, draw=none, from=4-7, to=7-6]
	\arrow["\cdots"{description}, sloped,,  ,  from=4-9, to=7-8]
	\arrow[""{name=4, anchor=center, inner sep=0}, "\cdots"{description}, sloped,,  ,  draw=none, from=4-9, to=7-10]
	\arrow[""{name=5, anchor=center, inner sep=0}, "v"{description}, draw=none, from=7-8, to=10-7]
	\arrow[""{name=6, anchor=center, inner sep=0}, draw=none, from=1-2, to=0]
	\arrow[""{name=7, anchor=center, inner sep=0}, draw=none, from=4-3, to=1]
	\arrow[""{name=8, anchor=center, inner sep=0}, draw=none, from=4-4, to=1]
	\arrow[""{name=9, anchor=center, inner sep=0}, draw=none, from=4-4, to=2]
	\arrow["\cdots"{description}, sloped,,  ,  draw=none, from=4-5, to=2]
	\arrow[""{name=10, anchor=center, inner sep=0}, draw=none, from=4-9, to=4]
	\arrow[""{name=11, anchor=center, inner sep=0}, draw=none, from=7-6, to=3]
	\arrow[""{name=12, anchor=center, inner sep=0}, draw=none, from=7-7, to=3]
	\arrow[""{name=13, anchor=center, inner sep=0}, draw=none, from=7-7, to=5]
	\arrow[""{name=14, anchor=center, inner sep=0}, draw=none, from=7-8, to=5]
	\arrow[from=1-2, to=6]
	\arrow[from=7, to=4-3]
	\arrow["\cdots"{description}, sloped,,  ,  from=8, to=9]
	\arrow[from=4-9, to=10]
	\arrow[from=11, to=7-6]
	\arrow["\cdots"{description}, sloped,,  ,  from=12, to=13]
	\arrow[from=7-8, to=14]
\end{tikzcd}.
\\
		Case 2n.8.
% https://q.uiver.app/#q=WzAsMTQsWzMsMywiTyIsWzAsMCwxMDAsMV1dLFsyLDMsInYiXSxbMyw2XSxbNCwzXSxbMywwXSxbNiw2LCJPIixbMCwwLDEwMCwxXV0sWzUsNiwidiJdLFs2LDldLFs3LDYsInYiXSxbNiwzXSxbMSw2LCJ4Il0sWzAsM10sWzgsMywieCciXSxbOSw2XSxbMywyLCIiLDEseyJzdHlsZSI6eyJib2R5Ijp7Im5hbWUiOiJub25lIn0sImhlYWQiOnsibmFtZSI6Im5vbmUifX19XSxbNCwxLCJ2IiwxLHsic3R5bGUiOnsiYm9keSI6eyJuYW1lIjoibm9uZSJ9LCJoZWFkIjp7Im5hbWUiOiJub25lIn19fV0sWzgsNywidiIsMSx7InN0eWxlIjp7ImJvZHkiOnsibmFtZSI6Im5vbmUifSwiaGVhZCI6eyJuYW1lIjoibm9uZSJ9fX1dLFs5LDYsInYiLDEseyJzdHlsZSI6eyJib2R5Ijp7Im5hbWUiOiJub25lIn0sImhlYWQiOnsibmFtZSI6Im5vbmUifX19XSxbMyw2LCJcXGNkb3RzIiwxXSxbMSwxMCwiXFxjZG90cyIsMV0sWzEwLDExLCJcXGNkb3RzIiwxLHsic3R5bGUiOnsiYm9keSI6eyJuYW1lIjoibm9uZSJ9LCJoZWFkIjp7Im5hbWUiOiJub25lIn19fV0sWzEyLDgsIlxcY2RvdHMiLDFdLFsxMiwxMywiXFxjZG90cyIsMSx7InN0eWxlIjp7ImJvZHkiOnsibmFtZSI6Im5vbmUifSwiaGVhZCI6eyJuYW1lIjoibm9uZSJ9fX1dLFszLDE0LCJcXGNkb3RzIiwxLHsic2hvcnRlbiI6eyJ0YXJnZXQiOjIwfSwic3R5bGUiOnsiYm9keSI6eyJuYW1lIjoibm9uZSJ9LCJoZWFkIjp7Im5hbWUiOiJub25lIn19fV0sWzAsMTQsIiIsMSx7InNob3J0ZW4iOnsidGFyZ2V0IjoyMH0sInN0eWxlIjp7ImJvZHkiOnsibmFtZSI6Im5vbmUifSwiaGVhZCI6eyJuYW1lIjoibm9uZSJ9fX1dLFswLDE1LCIiLDEseyJzaG9ydGVuIjp7InRhcmdldCI6MjB9LCJzdHlsZSI6eyJib2R5Ijp7Im5hbWUiOiJub25lIn0sImhlYWQiOnsibmFtZSI6Im5vbmUifX19XSxbMSwxNSwiIiwxLHsic2hvcnRlbiI6eyJ0YXJnZXQiOjIwfSwic3R5bGUiOnsiYm9keSI6eyJuYW1lIjoibm9uZSJ9LCJoZWFkIjp7Im5hbWUiOiJub25lIn19fV0sWzgsMTYsIiIsMSx7InNob3J0ZW4iOnsidGFyZ2V0IjoyMH0sInN0eWxlIjp7ImJvZHkiOnsibmFtZSI6Im5vbmUifSwiaGVhZCI6eyJuYW1lIjoibm9uZSJ9fX1dLFs1LDE2LCIiLDEseyJzaG9ydGVuIjp7InRhcmdldCI6MjB9LCJzdHlsZSI6eyJib2R5Ijp7Im5hbWUiOiJub25lIn0sImhlYWQiOnsibmFtZSI6Im5vbmUifX19XSxbNSwxNywiIiwxLHsic2hvcnRlbiI6eyJ0YXJnZXQiOjIwfSwic3R5bGUiOnsiYm9keSI6eyJuYW1lIjoibm9uZSJ9LCJoZWFkIjp7Im5hbWUiOiJub25lIn19fV0sWzYsMTcsIiIsMSx7InNob3J0ZW4iOnsidGFyZ2V0IjoyMH0sInN0eWxlIjp7ImJvZHkiOnsibmFtZSI6Im5vbmUifSwiaGVhZCI6eyJuYW1lIjoibm9uZSJ9fX1dLFsxMCwyMCwiIiwxLHsic2hvcnRlbiI6eyJ0YXJnZXQiOjIwfSwic3R5bGUiOnsiYm9keSI6eyJuYW1lIjoibm9uZSJ9LCJoZWFkIjp7Im5hbWUiOiJub25lIn19fV0sWzEyLDIyLCIiLDEseyJzaG9ydGVuIjp7InRhcmdldCI6MjB9LCJzdHlsZSI6eyJib2R5Ijp7Im5hbWUiOiJub25lIn0sImhlYWQiOnsibmFtZSI6Im5vbmUifX19XSxbMjUsMjQsIlxcY2RvdHMiLDEseyJsZXZlbCI6MX1dLFsyNiwxLCIiLDEseyJsZXZlbCI6MX1dLFsyOSwyOCwiXFxjZG90cyIsMSx7ImxldmVsIjoxfV0sWzMwLDYsIiIsMSx7ImxldmVsIjoxfV0sWzgsMjcsIiIsMSx7ImxldmVsIjoxfV0sWzMxLDEwLCIiLDEseyJsZXZlbCI6MX1dLFsxMiwzMiwiIiwxLHsibGV2ZWwiOjF9XV0=
\begin{tikzcd}[row sep=small, column sep=tiny,  every label/.append style = {font=\normalsize}]
	&&& {} \\
	\\
	\\
	{} && v & \textcolor{white}{O} & \textcolor{white}{O} && {} && {x'} \\
	\\
	\\
	& x && {} && v & \textcolor{white}{O} & v && {} \\
	\\
	\\
	&&&&&& {}
	\arrow[""{name=0, anchor=center, inner sep=0}, "v"{description}, draw=none, from=1-4, to=4-3]
	\arrow["\cdots"{description}, sloped,,  ,  from=4-3, to=7-2]
	\arrow[""{name=1, anchor=center, inner sep=0}, draw=none, from=4-5, to=7-4]
	\arrow["\cdots"{description}, sloped,,  ,  from=4-5, to=7-6]
	\arrow[""{name=2, anchor=center, inner sep=0}, "v"{description}, draw=none, from=4-7, to=7-6]
	\arrow["\cdots"{description}, sloped,,  ,  from=4-9, to=7-8]
	\arrow[""{name=3, anchor=center, inner sep=0}, "\cdots"{description}, sloped,,  ,  draw=none, from=4-9, to=7-10]
	\arrow[""{name=4, anchor=center, inner sep=0}, "\cdots"{description}, sloped,,  ,  draw=none, from=7-2, to=4-1]
	\arrow[""{name=5, anchor=center, inner sep=0}, "v"{description}, draw=none, from=7-8, to=10-7]
	\arrow[""{name=6, anchor=center, inner sep=0}, draw=none, from=4-3, to=0]
	\arrow[""{name=7, anchor=center, inner sep=0}, draw=none, from=4-4, to=0]
	\arrow[""{name=8, anchor=center, inner sep=0}, draw=none, from=4-4, to=1]
	\arrow["\cdots"{description}, sloped,,  ,  draw=none, from=4-5, to=1]
	\arrow[""{name=9, anchor=center, inner sep=0}, draw=none, from=4-9, to=3]
	\arrow[""{name=10, anchor=center, inner sep=0}, draw=none, from=7-2, to=4]
	\arrow[""{name=11, anchor=center, inner sep=0}, draw=none, from=7-6, to=2]
	\arrow[""{name=12, anchor=center, inner sep=0}, draw=none, from=7-7, to=2]
	\arrow[""{name=13, anchor=center, inner sep=0}, draw=none, from=7-7, to=5]
	\arrow[""{name=14, anchor=center, inner sep=0}, draw=none, from=7-8, to=5]
	\arrow[from=6, to=4-3]
	\arrow["\cdots"{description}, sloped,,  ,  from=7, to=8]
	\arrow[from=4-9, to=9]
	\arrow[from=10, to=7-2]
	\arrow[from=11, to=7-6]
	\arrow["\cdots"{description}, sloped,,  ,  from=12, to=13]
	\arrow[from=7-8, to=14]
\end{tikzcd}.
\\
		Case 2n.9.
% https://q.uiver.app/#q=WzAsMTQsWzMsMywiTyIsWzAsMCwxMDAsMV1dLFsyLDMsInYiXSxbMyw2XSxbNCwzXSxbMywwXSxbNiw2LCJPIixbMCwwLDEwMCwxXV0sWzUsNiwidiJdLFs2LDldLFs3LDYsInYiXSxbNiwzXSxbMSw2LCJ4Il0sWzAsM10sWzgsOSwieCciXSxbOSw2XSxbMywyLCIiLDEseyJzdHlsZSI6eyJib2R5Ijp7Im5hbWUiOiJub25lIn0sImhlYWQiOnsibmFtZSI6Im5vbmUifX19XSxbNCwxLCJ2IiwxLHsic3R5bGUiOnsiYm9keSI6eyJuYW1lIjoibm9uZSJ9LCJoZWFkIjp7Im5hbWUiOiJub25lIn19fV0sWzgsNywidiIsMSx7InN0eWxlIjp7ImJvZHkiOnsibmFtZSI6Im5vbmUifSwiaGVhZCI6eyJuYW1lIjoibm9uZSJ9fX1dLFs5LDYsInYiLDEseyJzdHlsZSI6eyJib2R5Ijp7Im5hbWUiOiJub25lIn0sImhlYWQiOnsibmFtZSI6Im5vbmUifX19XSxbMyw2LCJcXGNkb3RzIiwxXSxbMSwxMCwiXFxjZG90cyIsMV0sWzEwLDExLCJcXGNkb3RzIiwxLHsic3R5bGUiOnsiYm9keSI6eyJuYW1lIjoibm9uZSJ9LCJoZWFkIjp7Im5hbWUiOiJub25lIn19fV0sWzgsMTIsIlxcY2RvdHMiLDFdLFsxMiwxMywiXFxjZG90cyIsMSx7InN0eWxlIjp7ImJvZHkiOnsibmFtZSI6Im5vbmUifSwiaGVhZCI6eyJuYW1lIjoibm9uZSJ9fX1dLFszLDE0LCJcXGNkb3RzIiwxLHsic2hvcnRlbiI6eyJ0YXJnZXQiOjIwfSwic3R5bGUiOnsiYm9keSI6eyJuYW1lIjoibm9uZSJ9LCJoZWFkIjp7Im5hbWUiOiJub25lIn19fV0sWzAsMTQsIiIsMSx7InNob3J0ZW4iOnsidGFyZ2V0IjoyMH0sInN0eWxlIjp7ImJvZHkiOnsibmFtZSI6Im5vbmUifSwiaGVhZCI6eyJuYW1lIjoibm9uZSJ9fX1dLFswLDE1LCIiLDEseyJzaG9ydGVuIjp7InRhcmdldCI6MjB9LCJzdHlsZSI6eyJib2R5Ijp7Im5hbWUiOiJub25lIn0sImhlYWQiOnsibmFtZSI6Im5vbmUifX19XSxbMSwxNSwiIiwxLHsic2hvcnRlbiI6eyJ0YXJnZXQiOjIwfSwic3R5bGUiOnsiYm9keSI6eyJuYW1lIjoibm9uZSJ9LCJoZWFkIjp7Im5hbWUiOiJub25lIn19fV0sWzgsMTYsIiIsMSx7InNob3J0ZW4iOnsidGFyZ2V0IjoyMH0sInN0eWxlIjp7ImJvZHkiOnsibmFtZSI6Im5vbmUifSwiaGVhZCI6eyJuYW1lIjoibm9uZSJ9fX1dLFs1LDE2LCIiLDEseyJzaG9ydGVuIjp7InRhcmdldCI6MjB9LCJzdHlsZSI6eyJib2R5Ijp7Im5hbWUiOiJub25lIn0sImhlYWQiOnsibmFtZSI6Im5vbmUifX19XSxbNSwxNywiIiwxLHsic2hvcnRlbiI6eyJ0YXJnZXQiOjIwfSwic3R5bGUiOnsiYm9keSI6eyJuYW1lIjoibm9uZSJ9LCJoZWFkIjp7Im5hbWUiOiJub25lIn19fV0sWzYsMTcsIiIsMSx7InNob3J0ZW4iOnsidGFyZ2V0IjoyMH0sInN0eWxlIjp7ImJvZHkiOnsibmFtZSI6Im5vbmUifSwiaGVhZCI6eyJuYW1lIjoibm9uZSJ9fX1dLFsxMCwyMCwiIiwxLHsic2hvcnRlbiI6eyJ0YXJnZXQiOjIwfSwic3R5bGUiOnsiYm9keSI6eyJuYW1lIjoibm9uZSJ9LCJoZWFkIjp7Im5hbWUiOiJub25lIn19fV0sWzEyLDIyLCIiLDEseyJzaG9ydGVuIjp7InRhcmdldCI6MjB9LCJzdHlsZSI6eyJib2R5Ijp7Im5hbWUiOiJub25lIn0sImhlYWQiOnsibmFtZSI6Im5vbmUifX19XSxbMjUsMjQsIlxcY2RvdHMiLDEseyJsZXZlbCI6MX1dLFsyNiwxLCIiLDEseyJsZXZlbCI6MX1dLFsyOSwyOCwiXFxjZG90cyIsMSx7ImxldmVsIjoxfV0sWzMwLDYsIiIsMSx7ImxldmVsIjoxfV0sWzgsMjcsIiIsMSx7ImxldmVsIjoxfV0sWzMxLDEwLCIiLDEseyJsZXZlbCI6MX1dLFszMiwxMiwiIiwxLHsibGV2ZWwiOjF9XV0=
\begin{tikzcd}[row sep=small, column sep=tiny,  every label/.append style = {font=\normalsize}]
	&&& {} \\
	\\
	\\
	{} && v & \textcolor{white}{O} & \textcolor{white}{O} && {} \\
	\\
	\\
	& x && {} && v & \textcolor{white}{O} & v && {} \\
	\\
	\\
	&&&&&& {} && {x'}
	\arrow[""{name=0, anchor=center, inner sep=0}, "v"{description}, draw=none, from=1-4, to=4-3]
	\arrow["\cdots"{description}, sloped,,  ,  from=4-3, to=7-2]
	\arrow[""{name=1, anchor=center, inner sep=0}, draw=none, from=4-5, to=7-4]
	\arrow["\cdots"{description}, sloped,,  ,  from=4-5, to=7-6]
	\arrow[""{name=2, anchor=center, inner sep=0}, "v"{description}, draw=none, from=4-7, to=7-6]
	\arrow[""{name=3, anchor=center, inner sep=0}, "\cdots"{description}, sloped,,  ,  draw=none, from=7-2, to=4-1]
	\arrow[""{name=4, anchor=center, inner sep=0}, "v"{description}, draw=none, from=7-8, to=10-7]
	\arrow["\cdots"{description}, sloped,,  ,  from=7-8, to=10-9]
	\arrow[""{name=5, anchor=center, inner sep=0}, "\cdots"{description}, sloped,,  ,  draw=none, from=10-9, to=7-10]
	\arrow[""{name=6, anchor=center, inner sep=0}, draw=none, from=4-3, to=0]
	\arrow[""{name=7, anchor=center, inner sep=0}, draw=none, from=4-4, to=0]
	\arrow[""{name=8, anchor=center, inner sep=0}, draw=none, from=4-4, to=1]
	\arrow["\cdots"{description}, sloped,,  ,  draw=none, from=4-5, to=1]
	\arrow[""{name=9, anchor=center, inner sep=0}, draw=none, from=7-2, to=3]
	\arrow[""{name=10, anchor=center, inner sep=0}, draw=none, from=7-6, to=2]
	\arrow[""{name=11, anchor=center, inner sep=0}, draw=none, from=7-7, to=2]
	\arrow[""{name=12, anchor=center, inner sep=0}, draw=none, from=7-7, to=4]
	\arrow[""{name=13, anchor=center, inner sep=0}, draw=none, from=7-8, to=4]
	\arrow[""{name=14, anchor=center, inner sep=0}, draw=none, from=10-9, to=5]
	\arrow[from=6, to=4-3]
	\arrow["\cdots"{description}, sloped,,  ,  from=7, to=8]
	\arrow[from=9, to=7-2]
	\arrow[from=10, to=7-6]
	\arrow["\cdots"{description}, sloped,,  ,  from=11, to=12]
	\arrow[from=7-8, to=13]
	\arrow[from=14, to=10-9]
\end{tikzcd}.
\\
		Case 2n.10.
% https://q.uiver.app/#q=WzAsMTQsWzMsMywiTyIsWzAsMCwxMDAsMV1dLFsyLDMsInYiXSxbMyw2XSxbNCwzXSxbMywwXSxbNiw2LCJPIixbMCwwLDEwMCwxXV0sWzUsNiwidiJdLFs2LDldLFs3LDYsInYiXSxbNiwzXSxbMSwwLCJ4Il0sWzAsM10sWzgsOSwieCciXSxbOSw2XSxbMywyLCIiLDEseyJzdHlsZSI6eyJib2R5Ijp7Im5hbWUiOiJub25lIn0sImhlYWQiOnsibmFtZSI6Im5vbmUifX19XSxbNCwxLCJ2IiwxLHsic3R5bGUiOnsiYm9keSI6eyJuYW1lIjoibm9uZSJ9LCJoZWFkIjp7Im5hbWUiOiJub25lIn19fV0sWzgsNywidiIsMSx7InN0eWxlIjp7ImJvZHkiOnsibmFtZSI6Im5vbmUifSwiaGVhZCI6eyJuYW1lIjoibm9uZSJ9fX1dLFs5LDYsInYiLDEseyJzdHlsZSI6eyJib2R5Ijp7Im5hbWUiOiJub25lIn0sImhlYWQiOnsibmFtZSI6Im5vbmUifX19XSxbMyw2LCJcXGNkb3RzIiwxXSxbMTAsMSwiXFxjZG90cyIsMV0sWzEwLDExLCJcXGNkb3RzIiwxLHsic3R5bGUiOnsiYm9keSI6eyJuYW1lIjoibm9uZSJ9LCJoZWFkIjp7Im5hbWUiOiJub25lIn19fV0sWzgsMTIsIlxcY2RvdHMiLDFdLFsxMiwxMywiXFxjZG90cyIsMSx7InN0eWxlIjp7ImJvZHkiOnsibmFtZSI6Im5vbmUifSwiaGVhZCI6eyJuYW1lIjoibm9uZSJ9fX1dLFszLDE0LCJcXGNkb3RzIiwxLHsic2hvcnRlbiI6eyJ0YXJnZXQiOjIwfSwic3R5bGUiOnsiYm9keSI6eyJuYW1lIjoibm9uZSJ9LCJoZWFkIjp7Im5hbWUiOiJub25lIn19fV0sWzAsMTQsIiIsMSx7InNob3J0ZW4iOnsidGFyZ2V0IjoyMH0sInN0eWxlIjp7ImJvZHkiOnsibmFtZSI6Im5vbmUifSwiaGVhZCI6eyJuYW1lIjoibm9uZSJ9fX1dLFswLDE1LCIiLDEseyJzaG9ydGVuIjp7InRhcmdldCI6MjB9LCJzdHlsZSI6eyJib2R5Ijp7Im5hbWUiOiJub25lIn0sImhlYWQiOnsibmFtZSI6Im5vbmUifX19XSxbMSwxNSwiIiwxLHsic2hvcnRlbiI6eyJ0YXJnZXQiOjIwfSwic3R5bGUiOnsiYm9keSI6eyJuYW1lIjoibm9uZSJ9LCJoZWFkIjp7Im5hbWUiOiJub25lIn19fV0sWzgsMTYsIiIsMSx7InNob3J0ZW4iOnsidGFyZ2V0IjoyMH0sInN0eWxlIjp7ImJvZHkiOnsibmFtZSI6Im5vbmUifSwiaGVhZCI6eyJuYW1lIjoibm9uZSJ9fX1dLFs1LDE2LCIiLDEseyJzaG9ydGVuIjp7InRhcmdldCI6MjB9LCJzdHlsZSI6eyJib2R5Ijp7Im5hbWUiOiJub25lIn0sImhlYWQiOnsibmFtZSI6Im5vbmUifX19XSxbNSwxNywiIiwxLHsic2hvcnRlbiI6eyJ0YXJnZXQiOjIwfSwic3R5bGUiOnsiYm9keSI6eyJuYW1lIjoibm9uZSJ9LCJoZWFkIjp7Im5hbWUiOiJub25lIn19fV0sWzYsMTcsIiIsMSx7InNob3J0ZW4iOnsidGFyZ2V0IjoyMH0sInN0eWxlIjp7ImJvZHkiOnsibmFtZSI6Im5vbmUifSwiaGVhZCI6eyJuYW1lIjoibm9uZSJ9fX1dLFsxMCwyMCwiIiwxLHsic2hvcnRlbiI6eyJ0YXJnZXQiOjIwfSwic3R5bGUiOnsiYm9keSI6eyJuYW1lIjoibm9uZSJ9LCJoZWFkIjp7Im5hbWUiOiJub25lIn19fV0sWzEyLDIyLCIiLDEseyJzaG9ydGVuIjp7InRhcmdldCI6MjB9LCJzdHlsZSI6eyJib2R5Ijp7Im5hbWUiOiJub25lIn0sImhlYWQiOnsibmFtZSI6Im5vbmUifX19XSxbMjUsMjQsIlxcY2RvdHMiLDEseyJsZXZlbCI6MX1dLFsyNiwxLCIiLDEseyJsZXZlbCI6MX1dLFsyOSwyOCwiXFxjZG90cyIsMSx7ImxldmVsIjoxfV0sWzMwLDYsIiIsMSx7ImxldmVsIjoxfV0sWzgsMjcsIiIsMSx7ImxldmVsIjoxfV0sWzEwLDMxLCIiLDEseyJsZXZlbCI6MX1dLFszMiwxMiwiIiwxLHsibGV2ZWwiOjF9XV0=
\begin{tikzcd}[row sep=small, column sep=tiny,  every label/.append style = {font=\normalsize}]
	& x && {} \\
	\\
	\\
	{} && v & \textcolor{white}{O} & \textcolor{white}{O} && {} \\
	\\
	\\
	&&& {} && v & \textcolor{white}{O} & v && {} \\
	\\
	\\
	&&&&&& {} && {x'}
	\arrow[""{name=0, anchor=center, inner sep=0}, "\cdots"{description}, sloped,,  ,  draw=none, from=1-2, to=4-1]
	\arrow["\cdots"{description}, sloped,,  ,  from=1-2, to=4-3]
	\arrow[""{name=1, anchor=center, inner sep=0}, "v"{description}, draw=none, from=1-4, to=4-3]
	\arrow[""{name=2, anchor=center, inner sep=0}, draw=none, from=4-5, to=7-4]
	\arrow["\cdots"{description}, sloped,,  ,  from=4-5, to=7-6]
	\arrow[""{name=3, anchor=center, inner sep=0}, "v"{description}, draw=none, from=4-7, to=7-6]
	\arrow[""{name=4, anchor=center, inner sep=0}, "v"{description}, draw=none, from=7-8, to=10-7]
	\arrow["\cdots"{description}, sloped,,  ,  from=7-8, to=10-9]
	\arrow[""{name=5, anchor=center, inner sep=0}, "\cdots"{description}, sloped,,  ,  draw=none, from=10-9, to=7-10]
	\arrow[""{name=6, anchor=center, inner sep=0}, draw=none, from=1-2, to=0]
	\arrow[""{name=7, anchor=center, inner sep=0}, draw=none, from=4-3, to=1]
	\arrow[""{name=8, anchor=center, inner sep=0}, draw=none, from=4-4, to=1]
	\arrow[""{name=9, anchor=center, inner sep=0}, draw=none, from=4-4, to=2]
	\arrow["\cdots"{description}, sloped,,  ,  draw=none, from=4-5, to=2]
	\arrow[""{name=10, anchor=center, inner sep=0}, draw=none, from=7-6, to=3]
	\arrow[""{name=11, anchor=center, inner sep=0}, draw=none, from=7-7, to=3]
	\arrow[""{name=12, anchor=center, inner sep=0}, draw=none, from=7-7, to=4]
	\arrow[""{name=13, anchor=center, inner sep=0}, draw=none, from=7-8, to=4]
	\arrow[""{name=14, anchor=center, inner sep=0}, draw=none, from=10-9, to=5]
	\arrow[from=1-2, to=6]
	\arrow[from=7, to=4-3]
	\arrow["\cdots"{description}, sloped,,  ,  from=8, to=9]
	\arrow[from=10, to=7-6]
	\arrow["\cdots"{description}, sloped,,  ,  from=11, to=12]
	\arrow[from=7-8, to=13]
	\arrow[from=14, to=10-9]
\end{tikzcd}.

		We define an equivalence relation $\sim$ on the set $\mathscr{B}$ of diagrammatic morphisms from $X$ to $Y$ as follows. For two diagrammatic morphisms $f,f':X\rightarrow Y$, denote by $N_f,N_{f'}$ the corresponding subdiagrams of $Y$ (which are defined at the beginning of the proof). Then $f\sim f'$ if and only if $N_f$ and $N_{f'}$ are the same subdiagram of $Y$. We denote by $\mathscr{B}/\sim$ the quotient set of $\mathscr{B}$ by $\sim$.

        {For each $\mathscr{C}\in\mathscr{B}/\sim$, choose a diagrammatic morphism $f\in\mathscr{C}$, and denote the subdiagram of $Y$ corresponding to $f$ by $N$ (now $N$ is independent to the choice of $f$). We will define a subset $\mathscr{C}'$ of $\Hom(\Omega(X),Y)$ by discussing the local configuration of $Y$ containing $N$ as a subdiagram} case by case.

        {In the following discussion, we {denote by} $\underline{\mathscr{C}}$ (resp. $\underline{\mathscr{C}'}$) a basis of the $k$-vector space spanned by the image of $\mathscr{C}$ (resp. $\mathscr{C}'$) in the stable category.} Note that some of the above sets may be empty.

		For convenience, we denote the vertex corresponding to $\top(X)$ and $\top(\Omega(X))$ by $v_X$, and we also denote the generator of $\top(X)$ and $\top(\Omega(X))$ by $v_X$.

        {Case 1.1. $Y=v$. \\
        $\mathscr{C}=\{f_1\colon v_X\mapsto v\}$, $\underline{\mathscr{C}}=\{\underline{f_1}\}$. We take $\mathscr{C}'=\{g_1\colon v_X\mapsto v\}$, which implies that $\underline{\mathscr{C}'}=\{\underline{g_1}\}$.}

        {Case 1.2. $Y=$}
        \begin{tikzcd}[row sep=small, column sep=tiny,  every label/.append style = {font=\normalsize}]
	& x \\
	\\
	\\
	{} && v
	\arrow[""{name=0, anchor=center, inner sep=0}, "\cdots"{description}, sloped,,  ,  draw=none, from=1-2, to=4-1]
	\arrow["\cdots"{description}, sloped,,  ,  from=1-2, to=4-3]
	\arrow[""{name=1, anchor=center, inner sep=0}, draw=none, from=1-2, to=0]
	\arrow[from=1-2, to=1]
\end{tikzcd}. \\
{$\mathscr{C}=\{f_1\colon v_X\mapsto v\}$, $\underline{\mathscr{C}}=\{\underline{f_1}\}$. We take $\mathscr{C}'=\{g_1\colon v_X\mapsto v\}$, which implies that $\underline{\mathscr{C}'}=\{\underline{g_1}\}$.}

         {Case 1.3. $Y=$}
\begin{tikzcd}[row sep=small, column sep=tiny,  every label/.append style = {font=\normalsize}]
	v && {} \\
	\\
	\\
	& x
	\arrow["\cdots"{description}, sloped,,  ,  from=1-1, to=4-2]
	\arrow[""{name=0, anchor=center, inner sep=0}, "\cdots"{description}, sloped,,  ,  draw=none, from=4-2, to=1-3]
	\arrow[""{name=1, anchor=center, inner sep=0}, draw=none, from=4-2, to=0]
	\arrow[from=1, to=4-2]
\end{tikzcd}. \\
{$\mathscr{C}=\{f_1\colon v_X\mapsto v\}$, $\underline{\mathscr{C}}=\{\underline{f_1}\}$. We take $\mathscr{C}'=\{g_1\colon v_X\mapsto v\}$, which implies that $\underline{\mathscr{C}'}=\{\underline{g_1}\}$.}

		In the following cases, we relabel the vertices originally denoted by $v$ as $v_i$ and \(v_i{\urcorner}\), so that different morphisms in $\mathscr{C}$ can be distinguished explicitly.

		Case 2n-1.1. $Y=$
% https://q.uiver.app/#q=WzAsMTEsWzAsMywidl9uIl0sWzEsMF0sWzIsMywidl97bi0xfVxcdXJjb3JuZXIiXSxbMSw2XSxbMSwzLCJPIixbMCwwLDEwMCwxXV0sWzUsNiwidl8yXFx1cmNvcm5lciJdLFszLDYsIk8iLFswLDAsMTAwLDFdXSxbNCw2XSxbNCwzLCJPIixbMCwwLDEwMCwxXV0sWzYsOSwidl8xIl0sWzUsOV0sWzEsMCwidl9uXFx1cmNvcm5lciIsMSx7InN0eWxlIjp7ImJvZHkiOnsibmFtZSI6Im5vbmUifSwiaGVhZCI6eyJuYW1lIjoibm9uZSJ9fX1dLFsyLDMsInZfe24tMX0iLDEseyJzdHlsZSI6eyJib2R5Ijp7Im5hbWUiOiJub25lIn0sImhlYWQiOnsibmFtZSI6Im5vbmUifX19XSxbMiw2LCJcXGNkb3RzIiwxXSxbOCw2LCIiLDEseyJzdHlsZSI6eyJib2R5Ijp7Im5hbWUiOiJub25lIn0sImhlYWQiOnsibmFtZSI6Im5vbmUifX19XSxbNSw5LCJcXGNkb3RzIiwxXSxbNywxMCwidl8yIiwxLHsic3R5bGUiOnsiYm9keSI6eyJuYW1lIjoibm9uZSJ9LCJoZWFkIjp7Im5hbWUiOiJub25lIn19fV0sWzQsMTEsIiIsMSx7InNob3J0ZW4iOnsidGFyZ2V0IjoyMH0sInN0eWxlIjp7ImJvZHkiOnsibmFtZSI6Im5vbmUifSwiaGVhZCI6eyJuYW1lIjoibm9uZSJ9fX1dLFs0LDEyLCIiLDEseyJzaG9ydGVuIjp7InRhcmdldCI6MjB9LCJzdHlsZSI6eyJib2R5Ijp7Im5hbWUiOiJub25lIn0sImhlYWQiOnsibmFtZSI6Im5vbmUifX19XSxbMCwxMSwiIiwxLHsic2hvcnRlbiI6eyJ0YXJnZXQiOjIwfSwic3R5bGUiOnsiYm9keSI6eyJuYW1lIjoibm9uZSJ9LCJoZWFkIjp7Im5hbWUiOiJub25lIn19fV0sWzIsMTIsIiIsMSx7InNob3J0ZW4iOnsidGFyZ2V0IjoyMH0sInN0eWxlIjp7ImJvZHkiOnsibmFtZSI6Im5vbmUifSwiaGVhZCI6eyJuYW1lIjoibm9uZSJ9fX1dLFs2LDE0LCIiLDEseyJzaG9ydGVuIjp7InRhcmdldCI6MjB9LCJzdHlsZSI6eyJib2R5Ijp7Im5hbWUiOiJub25lIn0sImhlYWQiOnsibmFtZSI6Im5vbmUifX19XSxbNSwxNiwiIiwxLHsic2hvcnRlbiI6eyJ0YXJnZXQiOjIwfSwic3R5bGUiOnsiYm9keSI6eyJuYW1lIjoibm9uZSJ9LCJoZWFkIjp7Im5hbWUiOiJub25lIn19fV0sWzcsMTYsIiIsMSx7InNob3J0ZW4iOnsidGFyZ2V0IjoyMH0sInN0eWxlIjp7ImJvZHkiOnsibmFtZSI6Im5vbmUifSwiaGVhZCI6eyJuYW1lIjoibm9uZSJ9fX1dLFs3LDE0LCIiLDEseyJzaG9ydGVuIjp7InRhcmdldCI6MjB9LCJzdHlsZSI6eyJib2R5Ijp7Im5hbWUiOiJub25lIn0sImhlYWQiOnsibmFtZSI6Im5vbmUifX19XSxbMTcsMTgsIlxcY2RvdHMiLDEseyJsZXZlbCI6MX1dLFsxOSwwLCIiLDEseyJsZXZlbCI6MX1dLFsyLDIwLCIiLDEseyJsZXZlbCI6MX1dLFs1LDIyLCIiLDEseyJsZXZlbCI6MX1dLFsyMSw2LCJcXGNkb3RzIiwxLHsibGV2ZWwiOjEsInN0eWxlIjp7ImJvZHkiOnsibmFtZSI6Im5vbmUifSwiaGVhZCI6eyJuYW1lIjoibm9uZSJ9fX1dLFsyNCwyMywiXFxjZG90cyIsMSx7ImxldmVsIjoxfV1d
\begin{tikzcd}[row sep=small, column sep=tiny,  every label/.append style = {font=\normalsize}]
	& {} \\
	\\
	\\
	{v_n} & \textcolor{white}{O} & {v_{n-1}\urcorner} && \textcolor{white}{O} \\
	\\
	\\
	& {} && \textcolor{white}{O} & {} & {v_2\urcorner} \\
	\\
	\\
	&&&&& {} & {v_1}
	\arrow[""{name=0, anchor=center, inner sep=0}, "{v_n\urcorner}"{description}, draw=none, from=1-2, to=4-1]
	\arrow[""{name=1, anchor=center, inner sep=0}, "{v_{n-1}}"{description}, draw=none, from=4-3, to=7-2]
	\arrow["\cdots"{description}, sloped,from=4-3, to=7-4]
	\arrow[""{name=2, anchor=center, inner sep=0}, draw=none, from=4-5, to=7-4]
	\arrow[""{name=3, anchor=center, inner sep=0}, "{v_2}"{description}, draw=none, from=7-5, to=10-6]
	\arrow["\cdots"{description}, sloped,from=7-6, to=10-7]
	\arrow[""{name=4, anchor=center, inner sep=0}, draw=none, from=4-1, to=0]
	\arrow[""{name=5, anchor=center, inner sep=0}, draw=none, from=4-2, to=0]
	\arrow[""{name=6, anchor=center, inner sep=0}, draw=none, from=4-2, to=1]
	\arrow[""{name=7, anchor=center, inner sep=0}, draw=none, from=4-3, to=1]
	\arrow[""{name=8, anchor=center, inner sep=0}, draw=none, from=7-4, to=2]
	\arrow[""{name=9, anchor=center, inner sep=0}, draw=none, from=7-5, to=2]
	\arrow[""{name=10, anchor=center, inner sep=0}, draw=none, from=7-5, to=3]
	\arrow[""{name=11, anchor=center, inner sep=0}, draw=none, from=7-6, to=3]
	\arrow[from=4, to=4-1]
	\arrow["\cdots"{description}, sloped,from=5, to=6]
	\arrow[from=4-3, to=7]
	\arrow["\cdots"{description}, sloped,draw=none, from=8, to=7-4]
	\arrow["\cdots"{description}, sloped,from=9, to=10]
	\arrow[from=7-6, to=11]
\end{tikzcd}. \\
${\mathscr{C}}=\{f_i\colon v_X\mapsto v_i\mid i=1,2,\cdots,n\}$,
where $\underline{f_{i+1}-tf_i}=0$ for $i=1,2,\cdots,n-1$, which implies that ${\underline{\mathscr{C}}}=\{\underline{f_n}\}$. And we take ${\mathscr{C}'}=\{g_j\colon v_X\mapsto v_j\mid j=1,2,\cdots,n\}$, where $\underline{g_j}=0$ for $j\neq 1$, which implies ${\underline{\mathscr{C}'}}=\{\underline{g_1}\}$.

		 Case 2n-1.2. $Y=$
% https://q.uiver.app/#q=WzAsMTMsWzIsMywidl9uIl0sWzMsMF0sWzQsMywidl97bi0xfVxcdXJjb3JuZXIiXSxbMyw2XSxbMywzLCJPIixbMCwwLDEwMCwxXV0sWzcsNiwidl8yXFx1cmNvcm5lciJdLFs1LDYsIk8iLFswLDAsMTAwLDFdXSxbNiw2XSxbNiwzLCJPIixbMCwwLDEwMCwxXV0sWzgsOSwidl8xIl0sWzcsOV0sWzEsMCwieCJdLFswLDNdLFsxLDAsInZfblxcdXJjb3JuZXIiLDEseyJzdHlsZSI6eyJib2R5Ijp7Im5hbWUiOiJub25lIn0sImhlYWQiOnsibmFtZSI6Im5vbmUifX19XSxbMiwzLCJ2X3tuLTF9IiwxLHsic3R5bGUiOnsiYm9keSI6eyJuYW1lIjoibm9uZSJ9LCJoZWFkIjp7Im5hbWUiOiJub25lIn19fV0sWzIsNiwiXFxjZG90cyIsMV0sWzgsNiwiIiwxLHsic3R5bGUiOnsiYm9keSI6eyJuYW1lIjoibm9uZSJ9LCJoZWFkIjp7Im5hbWUiOiJub25lIn19fV0sWzUsOSwiXFxjZG90cyIsMV0sWzcsMTAsInZfMiIsMSx7InN0eWxlIjp7ImJvZHkiOnsibmFtZSI6Im5vbmUifSwiaGVhZCI6eyJuYW1lIjoibm9uZSJ9fX1dLFsxMSwwLCJcXGNkb3RzIiwxXSxbMTEsMTIsIlxcY2RvdHMiLDEseyJzdHlsZSI6eyJib2R5Ijp7Im5hbWUiOiJub25lIn0sImhlYWQiOnsibmFtZSI6Im5vbmUifX19XSxbNCwxMywiIiwxLHsic2hvcnRlbiI6eyJ0YXJnZXQiOjIwfSwic3R5bGUiOnsiYm9keSI6eyJuYW1lIjoibm9uZSJ9LCJoZWFkIjp7Im5hbWUiOiJub25lIn19fV0sWzQsMTQsIiIsMSx7InNob3J0ZW4iOnsidGFyZ2V0IjoyMH0sInN0eWxlIjp7ImJvZHkiOnsibmFtZSI6Im5vbmUifSwiaGVhZCI6eyJuYW1lIjoibm9uZSJ9fX1dLFswLDEzLCIiLDEseyJzaG9ydGVuIjp7InRhcmdldCI6MjB9LCJzdHlsZSI6eyJib2R5Ijp7Im5hbWUiOiJub25lIn0sImhlYWQiOnsibmFtZSI6Im5vbmUifX19XSxbMiwxNCwiIiwxLHsic2hvcnRlbiI6eyJ0YXJnZXQiOjIwfSwic3R5bGUiOnsiYm9keSI6eyJuYW1lIjoibm9uZSJ9LCJoZWFkIjp7Im5hbWUiOiJub25lIn19fV0sWzYsMTYsIiIsMSx7InNob3J0ZW4iOnsidGFyZ2V0IjoyMH0sInN0eWxlIjp7ImJvZHkiOnsibmFtZSI6Im5vbmUifSwiaGVhZCI6eyJuYW1lIjoibm9uZSJ9fX1dLFs1LDE4LCIiLDEseyJzaG9ydGVuIjp7InRhcmdldCI6MjB9LCJzdHlsZSI6eyJib2R5Ijp7Im5hbWUiOiJub25lIn0sImhlYWQiOnsibmFtZSI6Im5vbmUifX19XSxbNywxOCwiIiwxLHsic2hvcnRlbiI6eyJ0YXJnZXQiOjIwfSwic3R5bGUiOnsiYm9keSI6eyJuYW1lIjoibm9uZSJ9LCJoZWFkIjp7Im5hbWUiOiJub25lIn19fV0sWzExLDIwLCIiLDEseyJzaG9ydGVuIjp7InRhcmdldCI6MjB9LCJzdHlsZSI6eyJib2R5Ijp7Im5hbWUiOiJub25lIn0sImhlYWQiOnsibmFtZSI6Im5vbmUifX19XSxbMTYsNywiIiwxLHsic2hvcnRlbiI6eyJzb3VyY2UiOjIwfSwic3R5bGUiOnsiYm9keSI6eyJuYW1lIjoibm9uZSJ9LCJoZWFkIjp7Im5hbWUiOiJub25lIn19fV0sWzIxLDIyLCJcXGNkb3RzIiwxLHsibGV2ZWwiOjF9XSxbMjMsMCwiIiwxLHsibGV2ZWwiOjF9XSxbMiwyNCwiIiwxLHsibGV2ZWwiOjF9XSxbNSwyNiwiIiwxLHsibGV2ZWwiOjF9XSxbMjUsNiwiXFxjZG90cyIsMSx7ImxldmVsIjoxLCJzdHlsZSI6eyJib2R5Ijp7Im5hbWUiOiJub25lIn0sImhlYWQiOnsibmFtZSI6Im5vbmUifX19XSxbMTEsMjgsIiIsMSx7ImxldmVsIjoxfV0sWzI5LDI3LCJcXGNkb3RzIiwxLHsibGV2ZWwiOjF9XV0=
\begin{tikzcd}[row sep=small, column sep=tiny,  every label/.append style = {font=\normalsize}]
	& x && {} \\
	\\
	\\
	{} && {v_n} & \textcolor{white}{O} & {v_{n-1}\urcorner} && \textcolor{white}{O} \\
	\\
	\\
	&&& {} && \textcolor{white}{O} & {} & {v_2\urcorner} \\
	\\
	\\
	&&&&&&& {} & {v_1}
	\arrow[""{name=0, anchor=center, inner sep=0}, "\cdots"{description}, sloped,draw=none, from=1-2, to=4-1]
	\arrow["\cdots"{description}, sloped,from=1-2, to=4-3]
	\arrow[""{name=1, anchor=center, inner sep=0}, "{v_n\urcorner}"{description}, draw=none, from=1-4, to=4-3]
	\arrow[""{name=2, anchor=center, inner sep=0}, "{v_{n-1}}"{description}, draw=none, from=4-5, to=7-4]
	\arrow["\cdots"{description}, sloped,from=4-5, to=7-6]
	\arrow[""{name=3, anchor=center, inner sep=0}, draw=none, from=4-7, to=7-6]
	\arrow[""{name=4, anchor=center, inner sep=0}, "{v_2}"{description}, draw=none, from=7-7, to=10-8]
	\arrow["\cdots"{description}, sloped,from=7-8, to=10-9]
	\arrow[""{name=5, anchor=center, inner sep=0}, draw=none, from=1-2, to=0]
	\arrow[""{name=6, anchor=center, inner sep=0}, draw=none, from=4-3, to=1]
	\arrow[""{name=7, anchor=center, inner sep=0}, draw=none, from=4-4, to=1]
	\arrow[""{name=8, anchor=center, inner sep=0}, draw=none, from=4-4, to=2]
	\arrow[""{name=9, anchor=center, inner sep=0}, draw=none, from=4-5, to=2]
	\arrow[""{name=10, anchor=center, inner sep=0}, draw=none, from=3, to=7-7]
	\arrow[""{name=11, anchor=center, inner sep=0}, draw=none, from=7-6, to=3]
	\arrow[""{name=12, anchor=center, inner sep=0}, draw=none, from=7-7, to=4]
	\arrow[""{name=13, anchor=center, inner sep=0}, draw=none, from=7-8, to=4]
	\arrow[from=1-2, to=5]
	\arrow[from=6, to=4-3]
	\arrow["\cdots"{description}, sloped,from=7, to=8]
	\arrow[from=4-5, to=9]
	\arrow["\cdots"{description}, sloped,from=10, to=12]
	\arrow["\cdots"{description}, sloped,draw=none, from=11, to=7-6]
	\arrow[from=7-8, to=13]
\end{tikzcd}. \\
Similarly,
${\mathscr{C}}=\{f_i\colon v_X\mapsto v_i\mid i=1,2,\cdots,n\}$, ${\underline{\mathscr{C}}}=\{\underline{f_n}\}$. \\ {We define} ${\mathscr{C}'}=\{g_j\colon v_X\mapsto v_j\mid j=1,2,\cdots,n\}$, and we have ${\underline{\mathscr{C}'}}=\{\underline{g_1}\}$.

		Case 2n-1.3. $Y=$
% https://q.uiver.app/#q=WzAsMTMsWzIsMywidl9uIl0sWzMsMF0sWzQsMywidl97bi0xfVxcdXJjb3JuZXIiXSxbMyw2XSxbMywzLCJPIixbMCwwLDEwMCwxXV0sWzcsNiwidl8yXFx1cmNvcm5lciJdLFs1LDYsIk8iLFswLDAsMTAwLDFdXSxbNiw2LCJPIixbMCwwLDEwMCwxXV0sWzYsM10sWzgsOSwidl8xIl0sWzcsOV0sWzEsNiwieCJdLFswLDNdLFsxLDAsInZfblxcdXJjb3JuZXIiLDEseyJzdHlsZSI6eyJib2R5Ijp7Im5hbWUiOiJub25lIn0sImhlYWQiOnsibmFtZSI6Im5vbmUifX19XSxbMiwzLCJ2X3tuLTF9IiwxLHsic3R5bGUiOnsiYm9keSI6eyJuYW1lIjoibm9uZSJ9LCJoZWFkIjp7Im5hbWUiOiJub25lIn19fV0sWzIsNiwiXFxjZG90cyIsMV0sWzgsNiwiIiwxLHsic3R5bGUiOnsiYm9keSI6eyJuYW1lIjoibm9uZSJ9LCJoZWFkIjp7Im5hbWUiOiJub25lIn19fV0sWzUsOSwiXFxjZG90cyIsMV0sWzcsMTAsInZfMiIsMSx7InN0eWxlIjp7ImJvZHkiOnsibmFtZSI6Im5vbmUifSwiaGVhZCI6eyJuYW1lIjoibm9uZSJ9fX1dLFswLDExLCJcXGNkb3RzIiwxXSxbMTEsMTIsIlxcY2RvdHMiLDEseyJzdHlsZSI6eyJib2R5Ijp7Im5hbWUiOiJub25lIn0sImhlYWQiOnsibmFtZSI6Im5vbmUifX19XSxbNCwxMywiIiwxLHsic2hvcnRlbiI6eyJ0YXJnZXQiOjIwfSwic3R5bGUiOnsiYm9keSI6eyJuYW1lIjoibm9uZSJ9LCJoZWFkIjp7Im5hbWUiOiJub25lIn19fV0sWzQsMTQsIiIsMSx7InNob3J0ZW4iOnsidGFyZ2V0IjoyMH0sInN0eWxlIjp7ImJvZHkiOnsibmFtZSI6Im5vbmUifSwiaGVhZCI6eyJuYW1lIjoibm9uZSJ9fX1dLFswLDEzLCIiLDEseyJzaG9ydGVuIjp7InRhcmdldCI6MjB9LCJzdHlsZSI6eyJib2R5Ijp7Im5hbWUiOiJub25lIn0sImhlYWQiOnsibmFtZSI6Im5vbmUifX19XSxbMiwxNCwiIiwxLHsic2hvcnRlbiI6eyJ0YXJnZXQiOjIwfSwic3R5bGUiOnsiYm9keSI6eyJuYW1lIjoibm9uZSJ9LCJoZWFkIjp7Im5hbWUiOiJub25lIn19fV0sWzYsMTYsIiIsMSx7InNob3J0ZW4iOnsidGFyZ2V0IjoyMH0sInN0eWxlIjp7ImJvZHkiOnsibmFtZSI6Im5vbmUifSwiaGVhZCI6eyJuYW1lIjoibm9uZSJ9fX1dLFs1LDE4LCIiLDEseyJzaG9ydGVuIjp7InRhcmdldCI6MjB9LCJzdHlsZSI6eyJib2R5Ijp7Im5hbWUiOiJub25lIn0sImhlYWQiOnsibmFtZSI6Im5vbmUifX19XSxbNywxOCwiIiwxLHsic2hvcnRlbiI6eyJ0YXJnZXQiOjIwfSwic3R5bGUiOnsiYm9keSI6eyJuYW1lIjoibm9uZSJ9LCJoZWFkIjp7Im5hbWUiOiJub25lIn19fV0sWzExLDIwLCIiLDEseyJzaG9ydGVuIjp7InRhcmdldCI6MjB9LCJzdHlsZSI6eyJib2R5Ijp7Im5hbWUiOiJub25lIn0sImhlYWQiOnsibmFtZSI6Im5vbmUifX19XSxbMTYsNywiIiwxLHsic2hvcnRlbiI6eyJzb3VyY2UiOjIwfSwic3R5bGUiOnsiYm9keSI6eyJuYW1lIjoibm9uZSJ9LCJoZWFkIjp7Im5hbWUiOiJub25lIn19fV0sWzIxLDIyLCJcXGNkb3RzIiwxLHsibGV2ZWwiOjF9XSxbMjMsMCwiIiwxLHsibGV2ZWwiOjF9XSxbMiwyNCwiIiwxLHsibGV2ZWwiOjF9XSxbNSwyNiwiIiwxLHsibGV2ZWwiOjF9XSxbMjUsNiwiXFxjZG90cyIsMSx7ImxldmVsIjoxLCJzdHlsZSI6eyJib2R5Ijp7Im5hbWUiOiJub25lIn0sImhlYWQiOnsibmFtZSI6Im5vbmUifX19XSxbMjgsMTEsIiIsMSx7ImxldmVsIjoxfV0sWzI5LDI3LCJcXGNkb3RzIiwxLHsibGV2ZWwiOjF9XV0=
\begin{tikzcd}[row sep=small, column sep=tiny,  every label/.append style = {font=\normalsize}]
	&&& {} \\
	\\
	\\
	{} && {v_n} & \textcolor{white}{O} & {v_{n-1}\urcorner} && {} \\
	\\
	\\
	& x && {} && \textcolor{white}{O} & \textcolor{white}{O} & {v_2\urcorner} \\
	\\
	\\
	&&&&&&& {} & {v_1}
	\arrow[""{name=0, anchor=center, inner sep=0}, "{v_n\urcorner}"{description}, draw=none, from=1-4, to=4-3]
	\arrow["\cdots"{description}, sloped,from=4-3, to=7-2]
	\arrow[""{name=1, anchor=center, inner sep=0}, "{v_{n-1}}"{description}, draw=none, from=4-5, to=7-4]
	\arrow["\cdots"{description}, sloped,from=4-5, to=7-6]
	\arrow[""{name=2, anchor=center, inner sep=0}, draw=none, from=4-7, to=7-6]
	\arrow[""{name=3, anchor=center, inner sep=0}, "\cdots"{description}, sloped,draw=none, from=7-2, to=4-1]
	\arrow[""{name=4, anchor=center, inner sep=0}, "{v_2}"{description}, draw=none, from=7-7, to=10-8]
	\arrow["\cdots"{description}, sloped,from=7-8, to=10-9]
	\arrow[""{name=5, anchor=center, inner sep=0}, draw=none, from=4-3, to=0]
	\arrow[""{name=6, anchor=center, inner sep=0}, draw=none, from=4-4, to=0]
	\arrow[""{name=7, anchor=center, inner sep=0}, draw=none, from=4-4, to=1]
	\arrow[""{name=8, anchor=center, inner sep=0}, draw=none, from=4-5, to=1]
	\arrow[""{name=9, anchor=center, inner sep=0}, draw=none, from=2, to=7-7]
	\arrow[""{name=10, anchor=center, inner sep=0}, draw=none, from=7-2, to=3]
	\arrow[""{name=11, anchor=center, inner sep=0}, draw=none, from=7-6, to=2]
	\arrow[""{name=12, anchor=center, inner sep=0}, draw=none, from=7-7, to=4]
	\arrow[""{name=13, anchor=center, inner sep=0}, draw=none, from=7-8, to=4]
	\arrow[from=5, to=4-3]
	\arrow["\cdots"{description}, sloped,from=6, to=7]
	\arrow[from=4-5, to=8]
	\arrow["\cdots"{description}, sloped,from=9, to=12]
	\arrow[from=10, to=7-2]
	\arrow["\cdots"{description}, sloped,draw=none, from=11, to=7-6]
	\arrow[from=7-8, to=13]
\end{tikzcd}. \\
		  ${\mathscr{C}}=\{f_i\colon v_X\mapsto v_i\mid i=1,2,\cdots,n\}$, ${\underline{\mathscr{C}}}=\{\underline{f_n}\}$. \\ {We define} ${\mathscr{C}'}=\{g_j\colon v_X\mapsto v_j\mid j=1,2,\cdots,n\}$, and we have ${\underline{\mathscr{C}'}}=\{\underline{g_1}\}$.

		  Case 2n-1.4. $Y=$
% https://q.uiver.app/#q=WzAsMTEsWzIsMywiTyIsWzAsMCwxMDAsMV1dLFswLDAsInZfblxcdXJjb3JuZXIiXSxbMSwzLCJ2X3tuLTF9Il0sWzIsNl0sWzMsMywiTyIsWzAsMCwxMDAsMV1dLFsyLDBdLFs1LDYsIk8iLFswLDAsMTAwLDFdXSxbNCw2LCJ2XzIiXSxbNSw5XSxbNiw2LCJ2XzFcXHVyY29ybmVyIl0sWzUsM10sWzQsMywiIiwxLHsic3R5bGUiOnsiYm9keSI6eyJuYW1lIjoibm9uZSJ9LCJoZWFkIjp7Im5hbWUiOiJub25lIn19fV0sWzUsMiwidl97bi0xfVxcdXJjb3JuZXIiLDEseyJzdHlsZSI6eyJib2R5Ijp7Im5hbWUiOiJub25lIn0sImhlYWQiOnsibmFtZSI6Im5vbmUifX19XSxbMSwyLCJcXGNkb3RzIiwxXSxbOSw4LCJ2XzEiLDEseyJzdHlsZSI6eyJib2R5Ijp7Im5hbWUiOiJub25lIn0sImhlYWQiOnsibmFtZSI6Im5vbmUifX19XSxbMTAsNywidl8yXFx1cmNvcm5lciIsMSx7InN0eWxlIjp7ImJvZHkiOnsibmFtZSI6Im5vbmUifSwiaGVhZCI6eyJuYW1lIjoibm9uZSJ9fX1dLFs0LDcsIlxcY2RvdHMiLDFdLFs0LDExLCJcXGNkb3RzIiwxLHsic2hvcnRlbiI6eyJ0YXJnZXQiOjIwfSwic3R5bGUiOnsiYm9keSI6eyJuYW1lIjoibm9uZSJ9LCJoZWFkIjp7Im5hbWUiOiJub25lIn19fV0sWzAsMTEsIiIsMSx7InNob3J0ZW4iOnsidGFyZ2V0IjoyMH0sInN0eWxlIjp7ImJvZHkiOnsibmFtZSI6Im5vbmUifSwiaGVhZCI6eyJuYW1lIjoibm9uZSJ9fX1dLFswLDEyLCIiLDEseyJzaG9ydGVuIjp7InRhcmdldCI6MjB9LCJzdHlsZSI6eyJib2R5Ijp7Im5hbWUiOiJub25lIn0sImhlYWQiOnsibmFtZSI6Im5vbmUifX19XSxbMiwxMiwiIiwxLHsic2hvcnRlbiI6eyJ0YXJnZXQiOjIwfSwic3R5bGUiOnsiYm9keSI6eyJuYW1lIjoibm9uZSJ9LCJoZWFkIjp7Im5hbWUiOiJub25lIn19fV0sWzksMTQsIiIsMSx7InNob3J0ZW4iOnsidGFyZ2V0IjoyMH0sInN0eWxlIjp7ImJvZHkiOnsibmFtZSI6Im5vbmUifSwiaGVhZCI6eyJuYW1lIjoibm9uZSJ9fX1dLFs2LDE0LCIiLDEseyJzaG9ydGVuIjp7InRhcmdldCI6MjB9LCJzdHlsZSI6eyJib2R5Ijp7Im5hbWUiOiJub25lIn0sImhlYWQiOnsibmFtZSI6Im5vbmUifX19XSxbNiwxNSwiIiwxLHsic2hvcnRlbiI6eyJ0YXJnZXQiOjIwfSwic3R5bGUiOnsiYm9keSI6eyJuYW1lIjoibm9uZSJ9LCJoZWFkIjp7Im5hbWUiOiJub25lIn19fV0sWzcsMTUsIiIsMSx7InNob3J0ZW4iOnsidGFyZ2V0IjoyMH0sInN0eWxlIjp7ImJvZHkiOnsibmFtZSI6Im5vbmUifSwiaGVhZCI6eyJuYW1lIjoibm9uZSJ9fX1dLFsxOSwxOCwiXFxjZG90cyIsMSx7ImxldmVsIjoxfV0sWzIwLDIsIiIsMSx7ImxldmVsIjoxfV0sWzIzLDIyLCJcXGNkb3RzIiwxLHsibGV2ZWwiOjF9XSxbMjQsNywiIiwxLHsibGV2ZWwiOjF9XSxbOSwyMSwiIiwxLHsibGV2ZWwiOjF9XV0=
\begin{tikzcd}[row sep=small, column sep=tiny,  every label/.append style = {font=\normalsize}]
	{v_n\urcorner} && {} \\
	\\
	\\
	& {v_{n-1}} & \textcolor{white}{O} & \textcolor{white}{O} && {} \\
	\\
	\\
	&& {} && {v_2} & \textcolor{white}{O} & {v_1\urcorner} \\
	\\
	\\
	&&&&& {}
	\arrow["\cdots"{description}, sloped,from=1-1, to=4-2]
	\arrow[""{name=0, anchor=center, inner sep=0}, "{v_{n-1}\urcorner}"{description}, draw=none, from=1-3, to=4-2]
	\arrow[""{name=1, anchor=center, inner sep=0}, draw=none, from=4-4, to=7-3]
	\arrow["\cdots"{description}, sloped,from=4-4, to=7-5]
	\arrow[""{name=2, anchor=center, inner sep=0}, "{v_2\urcorner}"{description}, draw=none, from=4-6, to=7-5]
	\arrow[""{name=3, anchor=center, inner sep=0}, "{v_1}"{description}, draw=none, from=7-7, to=10-6]
	\arrow[""{name=4, anchor=center, inner sep=0}, draw=none, from=4-2, to=0]
	\arrow[""{name=5, anchor=center, inner sep=0}, draw=none, from=4-3, to=0]
	\arrow[""{name=6, anchor=center, inner sep=0}, draw=none, from=4-3, to=1]
	\arrow["\cdots"{description}, sloped,draw=none, from=4-4, to=1]
	\arrow[""{name=7, anchor=center, inner sep=0}, draw=none, from=7-5, to=2]
	\arrow[""{name=8, anchor=center, inner sep=0}, draw=none, from=7-6, to=2]
	\arrow[""{name=9, anchor=center, inner sep=0}, draw=none, from=7-6, to=3]
	\arrow[""{name=10, anchor=center, inner sep=0}, draw=none, from=7-7, to=3]
	\arrow[from=4, to=4-2]
	\arrow["\cdots"{description}, sloped,from=5, to=6]
	\arrow[from=7, to=7-5]
	\arrow["\cdots"{description}, sloped,from=8, to=9]
	\arrow[from=7-7, to=10]
\end{tikzcd}. \\
		  ${\mathscr{C}}=\{f_i\colon v_X\mapsto v_i\mid i=1,2,\cdots,n-1\}\cup\{f_n\colon v_X\mapsto v_n\urcorner\}$, where $\underline{f_{i+1}-tf_i}=0$ for \\ $i=1,2,\cdots,n-2$, and $\underline{f_1}=0$. Then ${\underline{\mathscr{C}}}=\{\underline{f_n}\}$. And we take \[{\mathscr{C}'}=\{g_j\colon v_X\mapsto v_j\mid j=1,2,\cdots,n-1\}\cup\{g_*\colon v_X\mapsto t^{n-1}v_1\urcorner+t^{n-2}v_2\urcorner+\cdots+tv_{n-1}\urcorner+v_n\urcorner\}.\] Then ${\underline{\mathscr{C}'}}=\{\underline{g_*}\}$.

		  Case 2n-1.5. $Y=$
% https://q.uiver.app/#q=WzAsMTMsWzIsMywiTyIsWzAsMCwxMDAsMV1dLFswLDAsInZfblxcdXJjb3JuZXIiXSxbMSwzLCJ2X3tuLTF9Il0sWzIsNl0sWzMsMywiTyIsWzAsMCwxMDAsMV1dLFsyLDBdLFs1LDYsIk8iLFswLDAsMTAwLDFdXSxbNCw2LCJ2XzIiXSxbNSw5XSxbNiw2LCJ2XzFcXHVyY29ybmVyIl0sWzUsMywiTyIsWzAsMCwxMDAsMV1dLFs3LDMsIngiXSxbOCw2XSxbNCwzLCIiLDEseyJzdHlsZSI6eyJib2R5Ijp7Im5hbWUiOiJub25lIn0sImhlYWQiOnsibmFtZSI6Im5vbmUifX19XSxbNSwyLCJ2X3tuLTF9XFx1cmNvcm5lciIsMSx7InN0eWxlIjp7ImJvZHkiOnsibmFtZSI6Im5vbmUifSwiaGVhZCI6eyJuYW1lIjoibm9uZSJ9fX1dLFsxLDIsIlxcY2RvdHMiLDFdLFs5LDgsInZfMSIsMSx7InN0eWxlIjp7ImJvZHkiOnsibmFtZSI6Im5vbmUifSwiaGVhZCI6eyJuYW1lIjoibm9uZSJ9fX1dLFsxMCw3LCJ2XzJcXHVyY29ybmVyIiwxLHsic3R5bGUiOnsiYm9keSI6eyJuYW1lIjoibm9uZSJ9LCJoZWFkIjp7Im5hbWUiOiJub25lIn19fV0sWzQsNywiXFxjZG90cyIsMV0sWzExLDksIlxcY2RvdHMiLDFdLFsxMSwxMiwiXFxjZG90cyIsMSx7InN0eWxlIjp7ImJvZHkiOnsibmFtZSI6Im5vbmUifSwiaGVhZCI6eyJuYW1lIjoibm9uZSJ9fX1dLFs0LDEzLCJcXGNkb3RzIiwxLHsic2hvcnRlbiI6eyJ0YXJnZXQiOjIwfSwic3R5bGUiOnsiYm9keSI6eyJuYW1lIjoibm9uZSJ9LCJoZWFkIjp7Im5hbWUiOiJub25lIn19fV0sWzAsMTMsIiIsMSx7InNob3J0ZW4iOnsidGFyZ2V0IjoyMH0sInN0eWxlIjp7ImJvZHkiOnsibmFtZSI6Im5vbmUifSwiaGVhZCI6eyJuYW1lIjoibm9uZSJ9fX1dLFswLDE0LCIiLDEseyJzaG9ydGVuIjp7InRhcmdldCI6MjB9LCJzdHlsZSI6eyJib2R5Ijp7Im5hbWUiOiJub25lIn0sImhlYWQiOnsibmFtZSI6Im5vbmUifX19XSxbMiwxNCwiIiwxLHsic2hvcnRlbiI6eyJ0YXJnZXQiOjIwfSwic3R5bGUiOnsiYm9keSI6eyJuYW1lIjoibm9uZSJ9LCJoZWFkIjp7Im5hbWUiOiJub25lIn19fV0sWzksMTYsIiIsMSx7InNob3J0ZW4iOnsidGFyZ2V0IjoyMH0sInN0eWxlIjp7ImJvZHkiOnsibmFtZSI6Im5vbmUifSwiaGVhZCI6eyJuYW1lIjoibm9uZSJ9fX1dLFs2LDE2LCIiLDEseyJzaG9ydGVuIjp7InRhcmdldCI6MjB9LCJzdHlsZSI6eyJib2R5Ijp7Im5hbWUiOiJub25lIn0sImhlYWQiOnsibmFtZSI6Im5vbmUifX19XSxbNiwxNywiIiwxLHsic2hvcnRlbiI6eyJ0YXJnZXQiOjIwfSwic3R5bGUiOnsiYm9keSI6eyJuYW1lIjoibm9uZSJ9LCJoZWFkIjp7Im5hbWUiOiJub25lIn19fV0sWzcsMTcsIiIsMSx7InNob3J0ZW4iOnsidGFyZ2V0IjoyMH0sInN0eWxlIjp7ImJvZHkiOnsibmFtZSI6Im5vbmUifSwiaGVhZCI6eyJuYW1lIjoibm9uZSJ9fX1dLFsxMSwyMCwiIiwxLHsic2hvcnRlbiI6eyJ0YXJnZXQiOjIwfSwic3R5bGUiOnsiYm9keSI6eyJuYW1lIjoibm9uZSJ9LCJoZWFkIjp7Im5hbWUiOiJub25lIn19fV0sWzIzLDIyLCJcXGNkb3RzIiwxLHsibGV2ZWwiOjF9XSxbMjQsMiwiIiwxLHsibGV2ZWwiOjF9XSxbMjcsMjYsIlxcY2RvdHMiLDEseyJsZXZlbCI6MX1dLFsyOCw3LCIiLDEseyJsZXZlbCI6MX1dLFs5LDI1LCIiLDEseyJsZXZlbCI6MX1dLFsxMSwyOSwiIiwxLHsibGV2ZWwiOjF9XV0=
\begin{tikzcd}[row sep=small, column sep=tiny,  every label/.append style = {font=\normalsize}]
	{v_n\urcorner} && {} \\
	\\
	\\
	& {v_{n-1}} & \textcolor{white}{O} & \textcolor{white}{O} && \textcolor{white}{O} && x \\
	\\
	\\
	&& {} && {v_2} & \textcolor{white}{O} & {v_1\urcorner} && {} \\
	\\
	\\
	&&&&& {}
	\arrow["\cdots"{description}, sloped,from=1-1, to=4-2]
	\arrow[""{name=0, anchor=center, inner sep=0}, "{v_{n-1}\urcorner}"{description}, draw=none, from=1-3, to=4-2]
	\arrow[""{name=1, anchor=center, inner sep=0}, draw=none, from=4-4, to=7-3]
	\arrow["\cdots"{description}, sloped,from=4-4, to=7-5]
	\arrow[""{name=2, anchor=center, inner sep=0}, "{v_2\urcorner}"{description}, draw=none, from=4-6, to=7-5]
	\arrow["\cdots"{description}, sloped,from=4-8, to=7-7]
	\arrow[""{name=3, anchor=center, inner sep=0}, "\cdots"{description}, sloped,draw=none, from=4-8, to=7-9]
	\arrow[""{name=4, anchor=center, inner sep=0}, "{v_1}"{description}, draw=none, from=7-7, to=10-6]
	\arrow[""{name=5, anchor=center, inner sep=0}, draw=none, from=4-2, to=0]
	\arrow[""{name=6, anchor=center, inner sep=0}, draw=none, from=4-3, to=0]
	\arrow[""{name=7, anchor=center, inner sep=0}, draw=none, from=4-3, to=1]
	\arrow["\cdots"{description}, sloped,draw=none, from=4-4, to=1]
	\arrow[""{name=8, anchor=center, inner sep=0}, draw=none, from=4-8, to=3]
	\arrow[""{name=9, anchor=center, inner sep=0}, draw=none, from=7-5, to=2]
	\arrow[""{name=10, anchor=center, inner sep=0}, draw=none, from=7-6, to=2]
	\arrow[""{name=11, anchor=center, inner sep=0}, draw=none, from=7-6, to=4]
	\arrow[""{name=12, anchor=center, inner sep=0}, draw=none, from=7-7, to=4]
	\arrow[from=5, to=4-2]
	\arrow["\cdots"{description}, sloped,from=6, to=7]
	\arrow[from=4-8, to=8]
	\arrow[from=9, to=7-5]
	\arrow["\cdots"{description}, sloped,from=10, to=11]
	\arrow[from=7-7, to=12]
\end{tikzcd}. \\
		${\mathscr{C}}=\{f_i\colon v_X\mapsto v_i\mid i=1,2,\cdots,n-1\}\cup\{f_n\colon v_X\mapsto v_n\urcorner\}$ and ${\underline{\mathscr{C}}}=\{\underline{f_n}\}$. {We take} \[{\mathscr{C}'}=\{g_j\colon v_X\mapsto v_j\mid j=1,2,\cdots,n-1\}\cup\{g_*\colon v_X\mapsto t^{n-1}v_1\urcorner+t^{n-2}v_2\urcorner+\cdots+tv_{n-1}\urcorner+v_n\urcorner\}.\] Then ${\underline{\mathscr{C}'}}=\{\underline{g_*}\}$.

		Case 2n-1.6. $Y=$
% https://q.uiver.app/#q=WzAsMTMsWzIsMywiTyIsWzAsMCwxMDAsMV1dLFswLDAsInZfblxcdXJjb3JuZXIiXSxbMSwzLCJ2X3tuLTF9Il0sWzIsNl0sWzMsMywiTyIsWzAsMCwxMDAsMV1dLFsyLDBdLFs1LDYsIk8iLFswLDAsMTAwLDFdXSxbNCw2LCJ2XzIiXSxbNSw5XSxbNiw2LCJ2XzFcXHVyY29ybmVyIl0sWzUsM10sWzcsOSwieCJdLFs4LDZdLFs0LDMsIiIsMSx7InN0eWxlIjp7ImJvZHkiOnsibmFtZSI6Im5vbmUifSwiaGVhZCI6eyJuYW1lIjoibm9uZSJ9fX1dLFs1LDIsInZfe24tMX1cXHVyY29ybmVyIiwxLHsic3R5bGUiOnsiYm9keSI6eyJuYW1lIjoibm9uZSJ9LCJoZWFkIjp7Im5hbWUiOiJub25lIn19fV0sWzEsMiwiXFxjZG90cyIsMV0sWzksOCwidl8xIiwxLHsic3R5bGUiOnsiYm9keSI6eyJuYW1lIjoibm9uZSJ9LCJoZWFkIjp7Im5hbWUiOiJub25lIn19fV0sWzEwLDcsInZfMlxcdXJjb3JuZXIiLDEseyJzdHlsZSI6eyJib2R5Ijp7Im5hbWUiOiJub25lIn0sImhlYWQiOnsibmFtZSI6Im5vbmUifX19XSxbOSwxMSwiXFxjZG90cyIsMV0sWzExLDEyLCJcXGNkb3RzIiwxLHsic3R5bGUiOnsiYm9keSI6eyJuYW1lIjoibm9uZSJ9LCJoZWFkIjp7Im5hbWUiOiJub25lIn19fV0sWzQsNywiXFxjZG90cyIsMV0sWzQsMTMsIlxcY2RvdHMiLDEseyJzaG9ydGVuIjp7InRhcmdldCI6MjB9LCJzdHlsZSI6eyJib2R5Ijp7Im5hbWUiOiJub25lIn0sImhlYWQiOnsibmFtZSI6Im5vbmUifX19XSxbMCwxMywiIiwxLHsic2hvcnRlbiI6eyJ0YXJnZXQiOjIwfSwic3R5bGUiOnsiYm9keSI6eyJuYW1lIjoibm9uZSJ9LCJoZWFkIjp7Im5hbWUiOiJub25lIn19fV0sWzAsMTQsIiIsMSx7InNob3J0ZW4iOnsidGFyZ2V0IjoyMH0sInN0eWxlIjp7ImJvZHkiOnsibmFtZSI6Im5vbmUifSwiaGVhZCI6eyJuYW1lIjoibm9uZSJ9fX1dLFsyLDE0LCIiLDEseyJzaG9ydGVuIjp7InRhcmdldCI6MjB9LCJzdHlsZSI6eyJib2R5Ijp7Im5hbWUiOiJub25lIn0sImhlYWQiOnsibmFtZSI6Im5vbmUifX19XSxbOSwxNiwiIiwxLHsic2hvcnRlbiI6eyJ0YXJnZXQiOjIwfSwic3R5bGUiOnsiYm9keSI6eyJuYW1lIjoibm9uZSJ9LCJoZWFkIjp7Im5hbWUiOiJub25lIn19fV0sWzYsMTYsIiIsMSx7InNob3J0ZW4iOnsidGFyZ2V0IjoyMH0sInN0eWxlIjp7ImJvZHkiOnsibmFtZSI6Im5vbmUifSwiaGVhZCI6eyJuYW1lIjoibm9uZSJ9fX1dLFs2LDE3LCIiLDEseyJzaG9ydGVuIjp7InRhcmdldCI6MjB9LCJzdHlsZSI6eyJib2R5Ijp7Im5hbWUiOiJub25lIn0sImhlYWQiOnsibmFtZSI6Im5vbmUifX19XSxbNywxNywiIiwxLHsic2hvcnRlbiI6eyJ0YXJnZXQiOjIwfSwic3R5bGUiOnsiYm9keSI6eyJuYW1lIjoibm9uZSJ9LCJoZWFkIjp7Im5hbWUiOiJub25lIn19fV0sWzExLDE5LCIiLDEseyJzaG9ydGVuIjp7InRhcmdldCI6MjB9LCJzdHlsZSI6eyJib2R5Ijp7Im5hbWUiOiJub25lIn0sImhlYWQiOnsibmFtZSI6Im5vbmUifX19XSxbMjMsMjIsIlxcY2RvdHMiLDEseyJsZXZlbCI6MX1dLFsyNCwyLCIiLDEseyJsZXZlbCI6MX1dLFsyNywyNiwiXFxjZG90cyIsMSx7ImxldmVsIjoxfV0sWzI4LDcsIiIsMSx7ImxldmVsIjoxfV0sWzksMjUsIiIsMSx7ImxldmVsIjoxfV0sWzI5LDExLCIiLDEseyJsZXZlbCI6MX1dXQ==
\begin{tikzcd}[row sep=small, column sep=tiny,  every label/.append style = {font=\normalsize}]
	{v_n\urcorner} && {} \\
	\\
	\\
	& {v_{n-1}} & \textcolor{white}{O} & \textcolor{white}{O} && {} \\
	\\
	\\
	&& {} && {v_2} & \textcolor{white}{O} & {v_1\urcorner} && {} \\
	\\
	\\
	&&&&& {} && x
	\arrow["\cdots"{description}, sloped,from=1-1, to=4-2]
	\arrow[""{name=0, anchor=center, inner sep=0}, "{v_{n-1}\urcorner}"{description}, draw=none, from=1-3, to=4-2]
	\arrow[""{name=1, anchor=center, inner sep=0}, draw=none, from=4-4, to=7-3]
	\arrow["\cdots"{description}, sloped,from=4-4, to=7-5]
	\arrow[""{name=2, anchor=center, inner sep=0}, "{v_2\urcorner}"{description}, draw=none, from=4-6, to=7-5]
	\arrow[""{name=3, anchor=center, inner sep=0}, "{v_1}"{description}, draw=none, from=7-7, to=10-6]
	\arrow["\cdots"{description}, sloped,from=7-7, to=10-8]
	\arrow[""{name=4, anchor=center, inner sep=0}, "\cdots"{description}, sloped,draw=none, from=10-8, to=7-9]
	\arrow[""{name=5, anchor=center, inner sep=0}, draw=none, from=4-2, to=0]
	\arrow[""{name=6, anchor=center, inner sep=0}, draw=none, from=4-3, to=0]
	\arrow[""{name=7, anchor=center, inner sep=0}, draw=none, from=4-3, to=1]
	\arrow["\cdots"{description}, sloped,draw=none, from=4-4, to=1]
	\arrow[""{name=8, anchor=center, inner sep=0}, draw=none, from=7-5, to=2]
	\arrow[""{name=9, anchor=center, inner sep=0}, draw=none, from=7-6, to=2]
	\arrow[""{name=10, anchor=center, inner sep=0}, draw=none, from=7-6, to=3]
	\arrow[""{name=11, anchor=center, inner sep=0}, draw=none, from=7-7, to=3]
	\arrow[""{name=12, anchor=center, inner sep=0}, draw=none, from=10-8, to=4]
	\arrow[from=5, to=4-2]
	\arrow["\cdots"{description}, sloped,from=6, to=7]
	\arrow[from=8, to=7-5]
	\arrow["\cdots"{description}, sloped,from=9, to=10]
	\arrow[from=7-7, to=11]
	\arrow[from=12, to=10-8]
\end{tikzcd}. \\
		${\mathscr{C}}=\{f_i\colon v_X\mapsto v_i\mid i=1,2,\cdots,n-1\}\cup\{f_n\colon v_X\mapsto v_n\urcorner\}$, ${\underline{\mathscr{C}}}=\{\underline{f_n}\}$. {We take} \[{\mathscr{C}'}=\{g_j\colon v_X\mapsto v_j\mid j=1,2,\cdots,n-1\}\cup\{g_*\colon v_X\mapsto t^{n-1}v_1\urcorner+t^{n-2}v_2\urcorner+\cdots+tv_{n-1}\urcorner+v_n\urcorner\}.\] Then ${\underline{\mathscr{C}'}}=\{\underline{g_*}\}$.

		Case 2n.1. $Y=$
% https://q.uiver.app/#q=WzAsMTEsWzIsMywiTyIsWzAsMCwxMDAsMV1dLFswLDAsInZfblxcdXJjb3JuZXIiXSxbMSwzLCJ2X3tuLTF9Il0sWzIsNl0sWzMsMywiTyIsWzAsMCwxMDAsMV1dLFsyLDBdLFs1LDYsIk8iLFswLDAsMTAwLDFdXSxbNCw2LCJ2XzEiXSxbNSw5XSxbNiw2XSxbNSwzXSxbNCwzLCIiLDEseyJzdHlsZSI6eyJib2R5Ijp7Im5hbWUiOiJub25lIn0sImhlYWQiOnsibmFtZSI6Im5vbmUifX19XSxbNSwyLCJ2X3tuLTF9XFx1cmNvcm5lciIsMSx7InN0eWxlIjp7ImJvZHkiOnsibmFtZSI6Im5vbmUifSwiaGVhZCI6eyJuYW1lIjoibm9uZSJ9fX1dLFsxLDIsIlxcY2RvdHMiLDFdLFs5LDgsInZfMFxcdXJjb3JuZXIiLDEseyJzdHlsZSI6eyJib2R5Ijp7Im5hbWUiOiJub25lIn0sImhlYWQiOnsibmFtZSI6Im5vbmUifX19XSxbMTAsNywidl8xXFx1cmNvcm5lciIsMSx7InN0eWxlIjp7ImJvZHkiOnsibmFtZSI6Im5vbmUifSwiaGVhZCI6eyJuYW1lIjoibm9uZSJ9fX1dLFs0LDcsIlxcY2RvdHMiLDFdLFs0LDExLCJcXGNkb3RzIiwxLHsic2hvcnRlbiI6eyJ0YXJnZXQiOjIwfSwic3R5bGUiOnsiYm9keSI6eyJuYW1lIjoibm9uZSJ9LCJoZWFkIjp7Im5hbWUiOiJub25lIn19fV0sWzAsMTEsIiIsMSx7InNob3J0ZW4iOnsidGFyZ2V0IjoyMH0sInN0eWxlIjp7ImJvZHkiOnsibmFtZSI6Im5vbmUifSwiaGVhZCI6eyJuYW1lIjoibm9uZSJ9fX1dLFswLDEyLCIiLDEseyJzaG9ydGVuIjp7InRhcmdldCI6MjB9LCJzdHlsZSI6eyJib2R5Ijp7Im5hbWUiOiJub25lIn0sImhlYWQiOnsibmFtZSI6Im5vbmUifX19XSxbMiwxMiwiIiwxLHsic2hvcnRlbiI6eyJ0YXJnZXQiOjIwfSwic3R5bGUiOnsiYm9keSI6eyJuYW1lIjoibm9uZSJ9LCJoZWFkIjp7Im5hbWUiOiJub25lIn19fV0sWzksMTQsIiIsMSx7InNob3J0ZW4iOnsidGFyZ2V0IjoyMH0sInN0eWxlIjp7ImJvZHkiOnsibmFtZSI6Im5vbmUifSwiaGVhZCI6eyJuYW1lIjoibm9uZSJ9fX1dLFs2LDE0LCIiLDEseyJzaG9ydGVuIjp7InRhcmdldCI6MjB9LCJzdHlsZSI6eyJib2R5Ijp7Im5hbWUiOiJub25lIn0sImhlYWQiOnsibmFtZSI6Im5vbmUifX19XSxbNiwxNSwiIiwxLHsic2hvcnRlbiI6eyJ0YXJnZXQiOjIwfSwic3R5bGUiOnsiYm9keSI6eyJuYW1lIjoibm9uZSJ9LCJoZWFkIjp7Im5hbWUiOiJub25lIn19fV0sWzcsMTUsIiIsMSx7InNob3J0ZW4iOnsidGFyZ2V0IjoyMH0sInN0eWxlIjp7ImJvZHkiOnsibmFtZSI6Im5vbmUifSwiaGVhZCI6eyJuYW1lIjoibm9uZSJ9fX1dLFsxOSwxOCwiXFxjZG90cyIsMSx7ImxldmVsIjoxfV0sWzIwLDIsIiIsMSx7ImxldmVsIjoxfV0sWzIzLDIyLCJcXGNkb3RzIiwxLHsibGV2ZWwiOjF9XSxbMjQsNywiIiwxLHsibGV2ZWwiOjF9XV0=
\begin{tikzcd}[row sep=small, column sep=tiny,  every label/.append style = {font=\normalsize}]
	{v_n\urcorner} && {} \\
	\\
	\\
	& {v_{n-1}} & \textcolor{white}{O} & \textcolor{white}{O} && {} \\
	\\
	\\
	&& {} && {v_1} & \textcolor{white}{O} & {} \\
	\\
	\\
	&&&&& {}
	\arrow["\cdots"{description}, sloped,from=1-1, to=4-2]
	\arrow[""{name=0, anchor=center, inner sep=0}, "{v_{n-1}\urcorner}"{description}, draw=none, from=1-3, to=4-2]
	\arrow[""{name=1, anchor=center, inner sep=0}, draw=none, from=4-4, to=7-3]
	\arrow["\cdots"{description}, sloped,from=4-4, to=7-5]
	\arrow[""{name=2, anchor=center, inner sep=0}, "{v_1\urcorner}"{description}, draw=none, from=4-6, to=7-5]
	\arrow[""{name=3, anchor=center, inner sep=0}, "{v_0}"{description}, draw=none, from=7-7, to=10-6]
	\arrow[""{name=4, anchor=center, inner sep=0}, draw=none, from=4-2, to=0]
	\arrow[""{name=5, anchor=center, inner sep=0}, draw=none, from=4-3, to=0]
	\arrow[""{name=6, anchor=center, inner sep=0}, draw=none, from=4-3, to=1]
	\arrow["\cdots"{description}, sloped,draw=none, from=4-4, to=1]
	\arrow[""{name=7, anchor=center, inner sep=0}, draw=none, from=7-5, to=2]
	\arrow[""{name=8, anchor=center, inner sep=0}, draw=none, from=7-6, to=2]
	\arrow[""{name=9, anchor=center, inner sep=0}, draw=none, from=7-6, to=3]
	\arrow[draw=none, from=7-7, to=3]
	\arrow[from=4, to=4-2]
	\arrow["\cdots"{description}, sloped,from=5, to=6]
	\arrow[from=7, to=7-5]
	\arrow["\cdots"{description}, sloped,from=8, to=9]
\end{tikzcd}. \\
		${\mathscr{C}}=\{f_i\colon v_X\mapsto v_i\mid i=0,1,\cdots,n-1\}\cup\{f_n\colon v_X\mapsto v_n\urcorner\}$, where $\underline{f_{i+1}-tf_i}=0$ for $i=0,1,\cdots,n-2$, and $\underline{f_{n-1}}=0$. Hence ${\underline{\mathscr{C}}}=\{\underline{f_n}\}$. And we take \\${\mathscr{C}'}=\{g_j\colon v_X\mapsto v_j\mid j=0,1,\cdots,n-1\}$. Then ${\underline{\mathscr{C}'}}=\{\underline{g_0}\}$.

		Case 2n.2-2n.10. $Y$ has a local configuration containing the same subdiagram \\ \[N=
% https://q.uiver.app/#q=WzAsMTAsWzEsMywiTyIsWzAsMCwxMDAsMV1dLFswLDMsInZfbiJdLFsxLDZdLFsyLDMsIk8iLFswLDAsMTAwLDFdXSxbMSwwXSxbNCw2LCJPIixbMCwwLDEwMCwxXV0sWzMsNiwidl8yIl0sWzQsOV0sWzUsNiwidl8xXFx1cmNvcm5lciJdLFs0LDNdLFszLDIsIiIsMSx7InN0eWxlIjp7ImJvZHkiOnsibmFtZSI6Im5vbmUifSwiaGVhZCI6eyJuYW1lIjoibm9uZSJ9fX1dLFs0LDEsInZfblxcdXJjb3JuZXIiLDEseyJzdHlsZSI6eyJib2R5Ijp7Im5hbWUiOiJub25lIn0sImhlYWQiOnsibmFtZSI6Im5vbmUifX19XSxbOCw3LCJ2XzEiLDEseyJzdHlsZSI6eyJib2R5Ijp7Im5hbWUiOiJub25lIn0sImhlYWQiOnsibmFtZSI6Im5vbmUifX19XSxbOSw2LCJ2XzJcXHVyY29ybmVyIiwxLHsic3R5bGUiOnsiYm9keSI6eyJuYW1lIjoibm9uZSJ9LCJoZWFkIjp7Im5hbWUiOiJub25lIn19fV0sWzMsNiwiXFxjZG90cyIsMV0sWzMsMTAsIlxcY2RvdHMiLDEseyJzaG9ydGVuIjp7InRhcmdldCI6MjB9LCJzdHlsZSI6eyJib2R5Ijp7Im5hbWUiOiJub25lIn0sImhlYWQiOnsibmFtZSI6Im5vbmUifX19XSxbMCwxMCwiIiwxLHsic2hvcnRlbiI6eyJ0YXJnZXQiOjIwfSwic3R5bGUiOnsiYm9keSI6eyJuYW1lIjoibm9uZSJ9LCJoZWFkIjp7Im5hbWUiOiJub25lIn19fV0sWzAsMTEsIiIsMSx7InNob3J0ZW4iOnsidGFyZ2V0IjoyMH0sInN0eWxlIjp7ImJvZHkiOnsibmFtZSI6Im5vbmUifSwiaGVhZCI6eyJuYW1lIjoibm9uZSJ9fX1dLFsxLDExLCIiLDEseyJzaG9ydGVuIjp7InRhcmdldCI6MjB9LCJzdHlsZSI6eyJib2R5Ijp7Im5hbWUiOiJub25lIn0sImhlYWQiOnsibmFtZSI6Im5vbmUifX19XSxbOCwxMiwiIiwxLHsic2hvcnRlbiI6eyJ0YXJnZXQiOjIwfSwic3R5bGUiOnsiYm9keSI6eyJuYW1lIjoibm9uZSJ9LCJoZWFkIjp7Im5hbWUiOiJub25lIn19fV0sWzUsMTIsIiIsMSx7InNob3J0ZW4iOnsidGFyZ2V0IjoyMH0sInN0eWxlIjp7ImJvZHkiOnsibmFtZSI6Im5vbmUifSwiaGVhZCI6eyJuYW1lIjoibm9uZSJ9fX1dLFs1LDEzLCIiLDEseyJzaG9ydGVuIjp7InRhcmdldCI6MjB9LCJzdHlsZSI6eyJib2R5Ijp7Im5hbWUiOiJub25lIn0sImhlYWQiOnsibmFtZSI6Im5vbmUifX19XSxbNiwxMywiIiwxLHsic2hvcnRlbiI6eyJ0YXJnZXQiOjIwfSwic3R5bGUiOnsiYm9keSI6eyJuYW1lIjoibm9uZSJ9LCJoZWFkIjp7Im5hbWUiOiJub25lIn19fV0sWzE3LDE2LCJcXGNkb3RzIiwxLHsibGV2ZWwiOjF9XSxbMTgsMSwiIiwxLHsibGV2ZWwiOjF9XSxbMjEsMjAsIlxcY2RvdHMiLDEseyJsZXZlbCI6MX1dLFsyMiw2LCIiLDEseyJsZXZlbCI6MX1dLFs4LDE5LCIiLDEseyJsZXZlbCI6MX1dXQ==
\begin{tikzcd}[row sep=small, column sep=tiny,  every label/.append style = {font=\normalsize}]
	& {} \\
	\\
	\\
	{v_n} & \textcolor{white}{O} & \textcolor{white}{O} && {} \\
	\\
	\\
	& {} && {v_2} & \textcolor{white}{O} & {v_1\urcorner} \\
	\\
	\\
	&&&& {}
	\arrow[""{name=0, anchor=center, inner sep=0}, "{v_n\urcorner}"{description}, draw=none, from=1-2, to=4-1]
	\arrow[""{name=1, anchor=center, inner sep=0}, draw=none, from=4-3, to=7-2]
	\arrow["\cdots"{description}, sloped,from=4-3, to=7-4]
	\arrow[""{name=2, anchor=center, inner sep=0}, "{v_2\urcorner}"{description}, draw=none, from=4-5, to=7-4]
	\arrow[""{name=3, anchor=center, inner sep=0}, "{v_1}"{description}, draw=none, from=7-6, to=10-5]
	\arrow[""{name=4, anchor=center, inner sep=0}, draw=none, from=4-1, to=0]
	\arrow[""{name=5, anchor=center, inner sep=0}, draw=none, from=4-2, to=0]
	\arrow[""{name=6, anchor=center, inner sep=0}, draw=none, from=4-2, to=1]
	\arrow["\cdots"{description}, sloped,draw=none, from=4-3, to=1]
	\arrow[""{name=7, anchor=center, inner sep=0}, draw=none, from=7-4, to=2]
	\arrow[""{name=8, anchor=center, inner sep=0}, draw=none, from=7-5, to=2]
	\arrow[""{name=9, anchor=center, inner sep=0}, draw=none, from=7-5, to=3]
	\arrow[""{name=10, anchor=center, inner sep=0}, draw=none, from=7-6, to=3]
	\arrow[from=4, to=4-1]
	\arrow["\cdots"{description}, sloped,from=5, to=6]
	\arrow[from=7, to=7-4]
	\arrow["\cdots"{description}, sloped,from=8, to=9]
	\arrow[from=7-6, to=10]
\end{tikzcd}.\] \\
		${\mathscr{C}}=\{f_i\colon v_X\mapsto v_i\mid i=1,2,\cdots,n\}$, ${\underline{\mathscr{C}}}=\varnothing$. {We take} \\${\mathscr{C}'}=\{g_j\colon v_X\mapsto v_j\mid j=1,2,\cdots,n\}$, {and then} ${\underline{\mathscr{C}'}}=\varnothing$.

		Notice that in each case, we have the same cardinality:
		\[|{\underline{\mathscr{C}}}|=|{\underline{\mathscr{C}'}}|.\]
		And only in {Case 1.1-1.3, Case 2n-1.1-Case 2n-1.6 and Case 2n.1}, ${\underline{\mathscr{C}}}\neq\varnothing$, and {Case 1.1, Case 2n-1.1, Case 2n-1.4 and Case 2n.1} lead to $Y=N$, while in all other cases $N$ contains an endpoint of $Y$. Therefore, {there are at most two equivalence classes $\mathscr{C}_1,\mathscr{C}_2\in\mathscr{B}/\sim$ with $\underline{\mathscr{C}_1},\underline{\mathscr{C}_2}\neq\varnothing$.} By the calculation in the proof, the cardinality of $\underline{\mathscr{C}}$ is $1$ for each $\mathscr{C}\in\mathscr{B}/\sim$ with $\underline{\mathscr{C}}\neq\varnothing$. Thus we have \[\dim\nolimits_k \sHom_{\varLambda}(X,Y)\leq 2.\]
			 Moreover, we can verify that {if $\mathscr{C}_1$ and $\mathscr{C}_2$ are two different equivalence classes with $\underline{\mathscr{C}_1}, \underline{\mathscr{C}_2}\neq\varnothing$, then $\underline{\mathscr{C}_1}$ and $\underline{\mathscr{C}_2}$ are linearly independent in $\sHom(X,Y)$. Therefore
$$\bigcup\limits_{\mathscr{C}\in\mathscr{B}/\sim}\underline{\mathscr{C}}$$
forms a basis of $\sHom(X,Y)$.} Similarly, we can verify that the morphism set
{$$\bigcup\limits_{\mathscr{C}\in\mathscr{B}/\sim}\underline{\mathscr{C'}}$$} is linearly independent in $\sHom(\Omega(X),Y)$, and hence
		\[\begin{aligned}
			\dim\nolimits_k \sHom(X,Y)&={\sum_{\mathscr{C}\in\mathscr{B}/\sim}|\underline{\mathscr{C}}|}\\
			&={\sum_{\mathscr{C}\in\mathscr{B}/\sim}|\underline{\mathscr{C}'}|}\leq\dim\nolimits_k \sHom(\Omega(X),Y).
		\end{aligned}\]

		Now it remains to prove that $\dim\nolimits_k \sHom_{\varLambda}(\Omega(X),Y)\leq \dim\nolimits_k \sHom(X,Y)$.

        Since $\Omega(X)/\soc(\Omega(X))$ is isomorphic to $X/\soc(X)$, we can view $\Hom(\Omega(X)/\soc(\Omega(X)),Y)$ as a common subspace of $\Hom(\Omega(X),Y)$ and $\Hom(X,Y)$. As before we {denote by} $\mathscr{B}$ the set of diagrammatic morphisms from $X$ to $Y$, which is a basis of $\Hom(X,Y)$. Then the subset $\mathscr{B}'$ of $\mathscr{B}$ consisting of diagrammatic morphisms $f$ with $f(\soc(X))=0$ forms a basis of $\Hom(\Omega(X)/\soc(\Omega(X)),Y)$. Recall that for each $\mathscr{C}\in\mathscr{B}/\sim$, we have constructed a subset $\mathscr{C}'$ of $\Hom(\Omega(X),Y)$. Denote by $\mathscr{E}$ the union of all $\mathscr{C}'$. By the definition of each subset $\mathscr{C}'$ of $\Hom(\Omega(X),Y)$, it is straightforward to show that $\mathscr{B}'$ is contained in $\mathscr{E}$.

        {Next we will show that $\mathscr{E}$ generates $\Hom(\Omega(X),Y)$ as a $k$-vector space.}

        {Let $g$ be a morphism in $\Hom(\Omega(X),Y)$. We denote the vertex corresponding to $\top(\Omega(X))$ (resp. $\soc(\Omega(X))$) by $v_X$ (resp. $v'_X$). The morphism $g$ maps $v_X$ to a linear combination of vertices $v_i$ in $Y$, where each $v_i$ corresponds to the edge $v$ in the Brauer graph of $\varLambda$. Each $v_i$ gives rise to a subdiagram $N_i$ of $Y$, which is defined similarly as the subdiagram $N$ of $Y$ at the beginning of the proof. Then $v_i$ appears either at a peak or a deep in $N_i$. If the vertex $v_i$ appears at a deep in $N_i$, or if $N_i$ contains only one vertex, then there exists a morphism $f\in\mathscr{B}'$ such that $f$ maps $v_X$ to $v_i$. Then after subtracting a linear combination of morphisms in $\mathscr{B}'$, we may assume that $g$ maps $v_X$ to a linear combination of vertices $v_i$, where each $N_i$ contains at least two vertices and each $v_i$ appears at a peak in $N_i$.}

        {We will show that $N_i$ can only appears as a subdiagram of $Y$ as in Case 2n-1.4-2n-1.6.}

        Case 2n-1.1-2n-1.3. $Y$ has a local configuration containing the same subdiagram \\ \[N_i=
% https://q.uiver.app/#q=WzAsMTEsWzAsMywidl9uIl0sWzEsMF0sWzIsMywidl97bi0xfVxcdXJjb3JuZXIiXSxbMSw2XSxbMSwzLCJPIixbMCwwLDEwMCwxXV0sWzUsNiwidl8yXFx1cmNvcm5lciJdLFszLDYsIk8iLFswLDAsMTAwLDFdXSxbNCw2XSxbNCwzLCJPIixbMCwwLDEwMCwxXV0sWzYsOSwidl8xIl0sWzUsOV0sWzEsMCwidl9uXFx1cmNvcm5lciIsMSx7InN0eWxlIjp7ImJvZHkiOnsibmFtZSI6Im5vbmUifSwiaGVhZCI6eyJuYW1lIjoibm9uZSJ9fX1dLFsyLDMsInZfe24tMX0iLDEseyJzdHlsZSI6eyJib2R5Ijp7Im5hbWUiOiJub25lIn0sImhlYWQiOnsibmFtZSI6Im5vbmUifX19XSxbMiw2LCJcXGNkb3RzIiwxXSxbOCw2LCIiLDEseyJzdHlsZSI6eyJib2R5Ijp7Im5hbWUiOiJub25lIn0sImhlYWQiOnsibmFtZSI6Im5vbmUifX19XSxbNSw5LCJcXGNkb3RzIiwxXSxbNywxMCwidl8yIiwxLHsic3R5bGUiOnsiYm9keSI6eyJuYW1lIjoibm9uZSJ9LCJoZWFkIjp7Im5hbWUiOiJub25lIn19fV0sWzQsMTEsIiIsMSx7InNob3J0ZW4iOnsidGFyZ2V0IjoyMH0sInN0eWxlIjp7ImJvZHkiOnsibmFtZSI6Im5vbmUifSwiaGVhZCI6eyJuYW1lIjoibm9uZSJ9fX1dLFs0LDEyLCIiLDEseyJzaG9ydGVuIjp7InRhcmdldCI6MjB9LCJzdHlsZSI6eyJib2R5Ijp7Im5hbWUiOiJub25lIn0sImhlYWQiOnsibmFtZSI6Im5vbmUifX19XSxbMCwxMSwiIiwxLHsic2hvcnRlbiI6eyJ0YXJnZXQiOjIwfSwic3R5bGUiOnsiYm9keSI6eyJuYW1lIjoibm9uZSJ9LCJoZWFkIjp7Im5hbWUiOiJub25lIn19fV0sWzIsMTIsIiIsMSx7InNob3J0ZW4iOnsidGFyZ2V0IjoyMH0sInN0eWxlIjp7ImJvZHkiOnsibmFtZSI6Im5vbmUifSwiaGVhZCI6eyJuYW1lIjoibm9uZSJ9fX1dLFs2LDE0LCIiLDEseyJzaG9ydGVuIjp7InRhcmdldCI6MjB9LCJzdHlsZSI6eyJib2R5Ijp7Im5hbWUiOiJub25lIn0sImhlYWQiOnsibmFtZSI6Im5vbmUifX19XSxbNSwxNiwiIiwxLHsic2hvcnRlbiI6eyJ0YXJnZXQiOjIwfSwic3R5bGUiOnsiYm9keSI6eyJuYW1lIjoibm9uZSJ9LCJoZWFkIjp7Im5hbWUiOiJub25lIn19fV0sWzcsMTYsIiIsMSx7InNob3J0ZW4iOnsidGFyZ2V0IjoyMH0sInN0eWxlIjp7ImJvZHkiOnsibmFtZSI6Im5vbmUifSwiaGVhZCI6eyJuYW1lIjoibm9uZSJ9fX1dLFs3LDE0LCIiLDEseyJzaG9ydGVuIjp7InRhcmdldCI6MjB9LCJzdHlsZSI6eyJib2R5Ijp7Im5hbWUiOiJub25lIn0sImhlYWQiOnsibmFtZSI6Im5vbmUifX19XSxbMTcsMTgsIlxcY2RvdHMiLDEseyJsZXZlbCI6MX1dLFsxOSwwLCIiLDEseyJsZXZlbCI6MX1dLFsyLDIwLCIiLDEseyJsZXZlbCI6MX1dLFs1LDIyLCIiLDEseyJsZXZlbCI6MX1dLFsyMSw2LCJcXGNkb3RzIiwxLHsibGV2ZWwiOjEsInN0eWxlIjp7ImJvZHkiOnsibmFtZSI6Im5vbmUifSwiaGVhZCI6eyJuYW1lIjoibm9uZSJ9fX1dLFsyNCwyMywiXFxjZG90cyIsMSx7ImxldmVsIjoxfV1d
\begin{tikzcd}[row sep=small, column sep=tiny,  every label/.append style = {font=\normalsize}]
	& {} \\
	\\
	\\
	{v_n} & \textcolor{white}{O} & {v_{n-1}\urcorner} && \textcolor{white}{O} \\
	\\
	\\
	& {} && \textcolor{white}{O} & {} & {v_2\urcorner} \\
	\\
	\\
	&&&&& {} & {v_1}
	\arrow[""{name=0, anchor=center, inner sep=0}, "{v_n\urcorner}"{description}, draw=none, from=1-2, to=4-1]
	\arrow[""{name=1, anchor=center, inner sep=0}, "{v_{n-1}}"{description}, draw=none, from=4-3, to=7-2]
	\arrow["\cdots"{description}, sloped,from=4-3, to=7-4]
	\arrow[""{name=2, anchor=center, inner sep=0}, draw=none, from=4-5, to=7-4]
	\arrow[""{name=3, anchor=center, inner sep=0}, "{v_2}"{description}, draw=none, from=7-5, to=10-6]
	\arrow["\cdots"{description}, sloped,from=7-6, to=10-7]
	\arrow[""{name=4, anchor=center, inner sep=0}, draw=none, from=4-1, to=0]
	\arrow[""{name=5, anchor=center, inner sep=0}, draw=none, from=4-2, to=0]
	\arrow[""{name=6, anchor=center, inner sep=0}, draw=none, from=4-2, to=1]
	\arrow[""{name=7, anchor=center, inner sep=0}, draw=none, from=4-3, to=1]
	\arrow[""{name=8, anchor=center, inner sep=0}, draw=none, from=7-4, to=2]
	\arrow[""{name=9, anchor=center, inner sep=0}, draw=none, from=7-5, to=2]
	\arrow[""{name=10, anchor=center, inner sep=0}, draw=none, from=7-5, to=3]
	\arrow[""{name=11, anchor=center, inner sep=0}, draw=none, from=7-6, to=3]
	\arrow[from=4, to=4-1]
	\arrow["\cdots"{description}, sloped,from=5, to=6]
	\arrow[from=4-3, to=7]
	\arrow["\cdots"{description}, sloped,draw=none, from=8, to=7-4]
	\arrow["\cdots"{description}, sloped,from=9, to=10]
	\arrow[from=7-6, to=11]
\end{tikzcd}. \]\\
{Suppose that the coefficient of $v_j\urcorner$ in $g(v_X)$ is $\lambda_j$ ($2\leq j\leq n$). Let $\eta'$ be the string $v\ra v$. Since $\eta\cdot v_X=v'_X$ and $\eta'\cdot v_X=t v'_X$, we have $t \eta\cdot v_X=\eta'\cdot v_X$ and $t \eta\cdot g(v_X)=\eta'\cdot g(v_X)$. A calculation shows that
\begin{equation*}
\text{the coefficient of } v_j \text{ in } \eta\cdot g(v_X)=\begin{cases}
\lambda_{j+1}, & \text{ if } 1\leq j\leq n-1; \\
0, & \text{ if } j=n,
\end{cases}
\end{equation*} and
\begin{equation*}
\text{the coefficient of } v_j \text{ in } \eta'\cdot g(v_X)=\begin{cases}
\lambda_{j}, & \text{ if } 2\leq j\leq n; \\
0, & \text{ if } j=1.
\end{cases}
\end{equation*}
Then $\lambda_n=0$ and $\lambda_j=t\lambda_{j+1}$ for $2\leq j\leq n-1$. Thus $\lambda_j=0$ for $2\leq j\leq n$.}

{Case 2n.1. $Y=$}
% https://q.uiver.app/#q=WzAsMTEsWzIsMywiTyIsWzAsMCwxMDAsMV1dLFswLDAsInZfblxcdXJjb3JuZXIiXSxbMSwzLCJ2X3tuLTF9Il0sWzIsNl0sWzMsMywiTyIsWzAsMCwxMDAsMV1dLFsyLDBdLFs1LDYsIk8iLFswLDAsMTAwLDFdXSxbNCw2LCJ2XzEiXSxbNSw5XSxbNiw2XSxbNSwzXSxbNCwzLCIiLDEseyJzdHlsZSI6eyJib2R5Ijp7Im5hbWUiOiJub25lIn0sImhlYWQiOnsibmFtZSI6Im5vbmUifX19XSxbNSwyLCJ2X3tuLTF9XFx1cmNvcm5lciIsMSx7InN0eWxlIjp7ImJvZHkiOnsibmFtZSI6Im5vbmUifSwiaGVhZCI6eyJuYW1lIjoibm9uZSJ9fX1dLFsxLDIsIlxcY2RvdHMiLDFdLFs5LDgsInZfMFxcdXJjb3JuZXIiLDEseyJzdHlsZSI6eyJib2R5Ijp7Im5hbWUiOiJub25lIn0sImhlYWQiOnsibmFtZSI6Im5vbmUifX19XSxbMTAsNywidl8xXFx1cmNvcm5lciIsMSx7InN0eWxlIjp7ImJvZHkiOnsibmFtZSI6Im5vbmUifSwiaGVhZCI6eyJuYW1lIjoibm9uZSJ9fX1dLFs0LDcsIlxcY2RvdHMiLDFdLFs0LDExLCJcXGNkb3RzIiwxLHsic2hvcnRlbiI6eyJ0YXJnZXQiOjIwfSwic3R5bGUiOnsiYm9keSI6eyJuYW1lIjoibm9uZSJ9LCJoZWFkIjp7Im5hbWUiOiJub25lIn19fV0sWzAsMTEsIiIsMSx7InNob3J0ZW4iOnsidGFyZ2V0IjoyMH0sInN0eWxlIjp7ImJvZHkiOnsibmFtZSI6Im5vbmUifSwiaGVhZCI6eyJuYW1lIjoibm9uZSJ9fX1dLFswLDEyLCIiLDEseyJzaG9ydGVuIjp7InRhcmdldCI6MjB9LCJzdHlsZSI6eyJib2R5Ijp7Im5hbWUiOiJub25lIn0sImhlYWQiOnsibmFtZSI6Im5vbmUifX19XSxbMiwxMiwiIiwxLHsic2hvcnRlbiI6eyJ0YXJnZXQiOjIwfSwic3R5bGUiOnsiYm9keSI6eyJuYW1lIjoibm9uZSJ9LCJoZWFkIjp7Im5hbWUiOiJub25lIn19fV0sWzksMTQsIiIsMSx7InNob3J0ZW4iOnsidGFyZ2V0IjoyMH0sInN0eWxlIjp7ImJvZHkiOnsibmFtZSI6Im5vbmUifSwiaGVhZCI6eyJuYW1lIjoibm9uZSJ9fX1dLFs2LDE0LCIiLDEseyJzaG9ydGVuIjp7InRhcmdldCI6MjB9LCJzdHlsZSI6eyJib2R5Ijp7Im5hbWUiOiJub25lIn0sImhlYWQiOnsibmFtZSI6Im5vbmUifX19XSxbNiwxNSwiIiwxLHsic2hvcnRlbiI6eyJ0YXJnZXQiOjIwfSwic3R5bGUiOnsiYm9keSI6eyJuYW1lIjoibm9uZSJ9LCJoZWFkIjp7Im5hbWUiOiJub25lIn19fV0sWzcsMTUsIiIsMSx7InNob3J0ZW4iOnsidGFyZ2V0IjoyMH0sInN0eWxlIjp7ImJvZHkiOnsibmFtZSI6Im5vbmUifSwiaGVhZCI6eyJuYW1lIjoibm9uZSJ9fX1dLFsxOSwxOCwiXFxjZG90cyIsMSx7ImxldmVsIjoxfV0sWzIwLDIsIiIsMSx7ImxldmVsIjoxfV0sWzIzLDIyLCJcXGNkb3RzIiwxLHsibGV2ZWwiOjF9XSxbMjQsNywiIiwxLHsibGV2ZWwiOjF9XV0=
\begin{tikzcd}[row sep=small, column sep=tiny,  every label/.append style = {font=\normalsize}]
	{v_n\urcorner} && {} \\
	\\
	\\
	& {v_{n-1}} & \textcolor{white}{O} & \textcolor{white}{O} && {} \\
	\\
	\\
	&& {} && {v_1} & \textcolor{white}{O} & {} \\
	\\
	\\
	&&&&& {}
	\arrow["\cdots"{description}, sloped,from=1-1, to=4-2]
	\arrow[""{name=0, anchor=center, inner sep=0}, "{v_{n-1}\urcorner}"{description}, draw=none, from=1-3, to=4-2]
	\arrow[""{name=1, anchor=center, inner sep=0}, draw=none, from=4-4, to=7-3]
	\arrow["\cdots"{description}, sloped,from=4-4, to=7-5]
	\arrow[""{name=2, anchor=center, inner sep=0}, "{v_1\urcorner}"{description}, draw=none, from=4-6, to=7-5]
	\arrow[""{name=3, anchor=center, inner sep=0}, "{v_0}"{description}, draw=none, from=7-7, to=10-6]
	\arrow[""{name=4, anchor=center, inner sep=0}, draw=none, from=4-2, to=0]
	\arrow[""{name=5, anchor=center, inner sep=0}, draw=none, from=4-3, to=0]
	\arrow[""{name=6, anchor=center, inner sep=0}, draw=none, from=4-3, to=1]
	\arrow["\cdots"{description}, sloped,draw=none, from=4-4, to=1]
	\arrow[""{name=7, anchor=center, inner sep=0}, draw=none, from=7-5, to=2]
	\arrow[""{name=8, anchor=center, inner sep=0}, draw=none, from=7-6, to=2]
	\arrow[""{name=9, anchor=center, inner sep=0}, draw=none, from=7-6, to=3]
	\arrow[draw=none, from=7-7, to=3]
	\arrow[from=4, to=4-2]
	\arrow["\cdots"{description}, sloped,from=5, to=6]
	\arrow[from=7, to=7-5]
	\arrow["\cdots"{description}, sloped,from=8, to=9]
\end{tikzcd}. \\
{Suppose that the coefficient of $v_j\urcorner$ in $g(v_X)$ is $\lambda_j$ ($1\leq j\leq n$). Let $\eta'$ be the diagram $v\ra v$. We have $t \eta\cdot g(v_X)=\eta'\cdot g(v_X)$. A calculation shows that
the coefficient of $v_j$ in $\eta\cdot g(v_X)$ is $\lambda_{j+1}$ for $0\leq j\leq n-1$, and
\begin{equation*}
\text{the coefficient of } v_j \text{ in } \eta'\cdot g(v_X)=\begin{cases}
\lambda_{j}, & \text{ if } 1\leq j\leq n-1; \\
0, & \text{ if } j=0.
\end{cases}
\end{equation*}
Then $t\lambda_1=0$ and $\lambda_j=t\lambda_{j+1}$ for $1\leq j\leq n-1$. Thus $\lambda_j=0$ for $1\leq j\leq n$.}

Case 2n.2-2n.10. $Y$ has a local configuration containing the same subdiagram\\ \[N_i=
% https://q.uiver.app/#q=WzAsMTAsWzEsMywiTyIsWzAsMCwxMDAsMV1dLFswLDMsInZfbiJdLFsxLDZdLFsyLDMsIk8iLFswLDAsMTAwLDFdXSxbMSwwXSxbNCw2LCJPIixbMCwwLDEwMCwxXV0sWzMsNiwidl8yIl0sWzQsOV0sWzUsNiwidl8xXFx1cmNvcm5lciJdLFs0LDNdLFszLDIsIiIsMSx7InN0eWxlIjp7ImJvZHkiOnsibmFtZSI6Im5vbmUifSwiaGVhZCI6eyJuYW1lIjoibm9uZSJ9fX1dLFs0LDEsInZfblxcdXJjb3JuZXIiLDEseyJzdHlsZSI6eyJib2R5Ijp7Im5hbWUiOiJub25lIn0sImhlYWQiOnsibmFtZSI6Im5vbmUifX19XSxbOCw3LCJ2XzEiLDEseyJzdHlsZSI6eyJib2R5Ijp7Im5hbWUiOiJub25lIn0sImhlYWQiOnsibmFtZSI6Im5vbmUifX19XSxbOSw2LCJ2XzJcXHVyY29ybmVyIiwxLHsic3R5bGUiOnsiYm9keSI6eyJuYW1lIjoibm9uZSJ9LCJoZWFkIjp7Im5hbWUiOiJub25lIn19fV0sWzMsNiwiXFxjZG90cyIsMV0sWzMsMTAsIlxcY2RvdHMiLDEseyJzaG9ydGVuIjp7InRhcmdldCI6MjB9LCJzdHlsZSI6eyJib2R5Ijp7Im5hbWUiOiJub25lIn0sImhlYWQiOnsibmFtZSI6Im5vbmUifX19XSxbMCwxMCwiIiwxLHsic2hvcnRlbiI6eyJ0YXJnZXQiOjIwfSwic3R5bGUiOnsiYm9keSI6eyJuYW1lIjoibm9uZSJ9LCJoZWFkIjp7Im5hbWUiOiJub25lIn19fV0sWzAsMTEsIiIsMSx7InNob3J0ZW4iOnsidGFyZ2V0IjoyMH0sInN0eWxlIjp7ImJvZHkiOnsibmFtZSI6Im5vbmUifSwiaGVhZCI6eyJuYW1lIjoibm9uZSJ9fX1dLFsxLDExLCIiLDEseyJzaG9ydGVuIjp7InRhcmdldCI6MjB9LCJzdHlsZSI6eyJib2R5Ijp7Im5hbWUiOiJub25lIn0sImhlYWQiOnsibmFtZSI6Im5vbmUifX19XSxbOCwxMiwiIiwxLHsic2hvcnRlbiI6eyJ0YXJnZXQiOjIwfSwic3R5bGUiOnsiYm9keSI6eyJuYW1lIjoibm9uZSJ9LCJoZWFkIjp7Im5hbWUiOiJub25lIn19fV0sWzUsMTIsIiIsMSx7InNob3J0ZW4iOnsidGFyZ2V0IjoyMH0sInN0eWxlIjp7ImJvZHkiOnsibmFtZSI6Im5vbmUifSwiaGVhZCI6eyJuYW1lIjoibm9uZSJ9fX1dLFs1LDEzLCIiLDEseyJzaG9ydGVuIjp7InRhcmdldCI6MjB9LCJzdHlsZSI6eyJib2R5Ijp7Im5hbWUiOiJub25lIn0sImhlYWQiOnsibmFtZSI6Im5vbmUifX19XSxbNiwxMywiIiwxLHsic2hvcnRlbiI6eyJ0YXJnZXQiOjIwfSwic3R5bGUiOnsiYm9keSI6eyJuYW1lIjoibm9uZSJ9LCJoZWFkIjp7Im5hbWUiOiJub25lIn19fV0sWzE3LDE2LCJcXGNkb3RzIiwxLHsibGV2ZWwiOjF9XSxbMTgsMSwiIiwxLHsibGV2ZWwiOjF9XSxbMjEsMjAsIlxcY2RvdHMiLDEseyJsZXZlbCI6MX1dLFsyMiw2LCIiLDEseyJsZXZlbCI6MX1dLFs4LDE5LCIiLDEseyJsZXZlbCI6MX1dXQ==
\begin{tikzcd}[row sep=small, column sep=tiny,  every label/.append style = {font=\normalsize}]
	& {} \\
	\\
	\\
	{v_n} & \textcolor{white}{O} & \textcolor{white}{O} && {} \\
	\\
	\\
	& {} && {v_2} & \textcolor{white}{O} & {v_1\urcorner} \\
	\\
	\\
	&&&& {}
	\arrow[""{name=0, anchor=center, inner sep=0}, "{v_n\urcorner}"{description}, draw=none, from=1-2, to=4-1]
	\arrow[""{name=1, anchor=center, inner sep=0}, draw=none, from=4-3, to=7-2]
	\arrow["\cdots"{description}, sloped,from=4-3, to=7-4]
	\arrow[""{name=2, anchor=center, inner sep=0}, "{v_2\urcorner}"{description}, draw=none, from=4-5, to=7-4]
	\arrow[""{name=3, anchor=center, inner sep=0}, "{v_1}"{description}, draw=none, from=7-6, to=10-5]
	\arrow[""{name=4, anchor=center, inner sep=0}, draw=none, from=4-1, to=0]
	\arrow[""{name=5, anchor=center, inner sep=0}, draw=none, from=4-2, to=0]
	\arrow[""{name=6, anchor=center, inner sep=0}, draw=none, from=4-2, to=1]
	\arrow["\cdots"{description}, sloped,draw=none, from=4-3, to=1]
	\arrow[""{name=7, anchor=center, inner sep=0}, draw=none, from=7-4, to=2]
	\arrow[""{name=8, anchor=center, inner sep=0}, draw=none, from=7-5, to=2]
	\arrow[""{name=9, anchor=center, inner sep=0}, draw=none, from=7-5, to=3]
	\arrow[""{name=10, anchor=center, inner sep=0}, draw=none, from=7-6, to=3]
	\arrow[from=4, to=4-1]
	\arrow["\cdots"{description}, sloped,from=5, to=6]
	\arrow[from=7, to=7-4]
	\arrow["\cdots"{description}, sloped,from=8, to=9]
	\arrow[from=7-6, to=10]
\end{tikzcd}. \]\\
{Suppose that the coefficient of $v_j\urcorner$ in $g(v_X)$ is $\lambda_j$ ($1\leq j\leq n$). Let $\eta'$ be the diagram $v\ra v$. We have $t \eta\cdot g(v_X)=\eta'\cdot g(v_X)$. A calculation shows that
\begin{equation*}
\text{the coefficient of } v_j \text{ in } \eta\cdot g(v_X)=\begin{cases}
\lambda_{j+1}, & \text{ if } 1\leq j\leq n-1; \\
0, & \text{ if } j=n,
\end{cases}
\end{equation*}
and the coefficient of $v_j$ in $\eta'\cdot g(v_X)$ is $\lambda_j$ for $1\leq j\leq n$.
Then $t\lambda_{j+1}=\lambda_j$ for $1\leq j\leq n-1$ and $\lambda_n=0$. Thus $\lambda_j=0$ for $1\leq j\leq n$.}

{Therefore $N_i$ can only appears as a subdiagram of $Y$ as in Case 2n-1.4-2n-1.6. It is straightforward to show that $g$ is a linear combination of the morphisms $g_*$ given by Case 2n-1.4-2n-1.6. Then we imply that $\mathscr{E}$ generates $\Hom(\Omega(X),Y)$ as a $k$-vector space.}

{By definition,
$$\mathscr{E}=\bigcup\limits_{\mathscr{C}\in\mathscr{B}/\sim}\mathscr{C}'.$$
Then $$\bigcup\limits_{\mathscr{C}\in\mathscr{B}/\sim}\underline{\mathscr{C}'}$$
generates $\sHom(\Omega(X),Y)$.
Since $$\bigcup\limits_{\mathscr{C}\in\mathscr{B}/\sim}\underline{\mathscr{C}}$$
forms a basis of $\sHom(X,Y)$, and since $|\underline{\mathscr{C}}|=|\underline{\mathscr{C}'}|$ for each $\mathscr{C}\in\mathscr{B}/\sim$, we have
\[\begin{aligned}
			\dim\nolimits_k \sHom(\Omega(X),Y)&\leq\sum_{\mathscr{C}\in\mathscr{B}/\sim}|\underline{\mathscr{C}'}|\\
			&=\sum_{\mathscr{C}\in\mathscr{B}/\sim}|\underline{\mathscr{C}}|=\dim\nolimits_k \sHom(X,Y).
		\end{aligned}\]}

		The lemma follows.
    \end{proof}
	
	\begin{remark}
		From the proof of Lemma \ref{lem:equal-dimension} we have the following observations:
		\begin{enumerate}
			\item In Case 2n.1 (see also Case 2.1 in Example \ref{example: illustrate}), the subdiagram $\eta$ of $N$ is closed under taking successors, and hence can arise as the image of a diagrammatic monomorphism from $X$ to $Y$. However, since $X \not\cong \Omega (X)$, $\eta$ cannot be the image of a diagrammatic monomorphism from $\Omega (X)$ to $Y$. Therefore, if Case 2n.1 occurs (in this case $Y\cong N$), then the dimensions of $\Hom(X,Y)$ and $\Hom(\Omega(X),Y)$ differ by $1$.  
			\item Case 2.8 may occur multiple times if the multiplicity $m>1$ and $Y$ contains a sufficiently long subdiagram.
			\item The diagrammatic morphism $g_*$ from $\Omega(X)$ to $Y$ appearing in Case 2n-1.4-2n-1.6 arises from the module structure of $\Omega(X)$. It can be obtained from a diagrammatic morphism from a universal cover of $\Omega(X)$ to $Y$. We recommend the reader refer to \cite{Kra}, where basis morphisms of Hom-spaces are described in a more general setting for special biserial algebras, covering not only string modules but also band modules.
		\end{enumerate}
	\end{remark}

	We give {a concrete} example to illustrate {the proof of Lemma \ref{lem:equal-dimension}}.

	\begin{example}\label{example: illustrate}
		Let $\varLambda$ be a symmetric stably biserial algebra given by the Brauer graph
		\[\begin{tikzpicture}
            \draw (-0.5,0) circle (0.5);
            \fill (0,0) circle (0.5ex);
            \fill (1,1) circle (0.5ex);
            \fill (1,-1) circle (0.5ex);
            \node at(-1.2,0) {$v$};
            \node at(-0.3,0) {$3$};
            \draw[-] (0,0) -- (1,1);
            \draw[-] (0,0) -- (1,-1);
            \draw (-0.45,-0.15) rectangle (-0.15,0.15);
            \node at(0.5,0.75) {$x_1$};
            \node at(0.5,-0.75) {$x_2$};
        \end{tikzpicture}\]
		where $3$ is the multiplicity of the central vertex, and the cyclic order of half-edges is given by clockwise orientation. We can denote the projective indecomposable module $P_v$ corresponding to the vertex $v$ by the following diagram:
		\[\begin{tikzcd}[row sep=-4, column sep=0]
			% https://q.uiver.app/#q=WzAsMjQsWzEsMCwidiJdLFswLDEsInYiXSxbMCwyLCJ4XzEiXSxbMCwzLCJ4XzIiXSxbMCw0LCJ2Il0sWzAsNSwidiJdLFswLDYsInhfMSJdLFswLDcsInhfMiJdLFswLDgsInYiXSxbMCw5LCJ2Il0sWzAsMTAsInhfMSJdLFswLDExLCJ4XzIiXSxbMSwxMiwidiJdLFsyLDEsInhfMSJdLFsyLDIsInhfMiJdLFsyLDMsInYiXSxbMiw0LCJ2Il0sWzIsNSwieF8xIl0sWzIsNiwieF8yIl0sWzIsNywidiJdLFsyLDgsInYiXSxbMiw5LCJ4XzEiXSxbMiwxMCwieF8yIl0sWzIsMTEsInYiXSxbMSwxMiwiIiwwLHsic3R5bGUiOnsiYm9keSI6eyJuYW1lIjoiZGFzaGVkIn0sImhlYWQiOnsibmFtZSI6Im5vbmUifX19XV0=
				& v \\
				v && {x_1} \\
				{x_1} && {x_2} \\
				{x_2} && v \\
				v && v \\
				v && {x_1} \\
				{x_1} && {x_2} \\
				{x_2} && v \\
				v && v \\
				v && {x_1} \\
				{x_1} && {x_2} \\
				{x_2} && v \\
				& v
				\arrow[dashed, no head, from=2-1, to=13-2]
        \end{tikzcd},\]
		where the dashed line means that $\alpha^2-t C_\alpha^3$ acts as zero on $\top(P_v)$, $t\in k^*$, with $\alpha\colon v\ra v$ being the deformed loop and $C_\alpha$ being the cycle given by the permutation, see Theorem \ref{thm:sym-StBA}. Let $X$ be the $\varLambda$-module given by the diagram
		\[v\ra x_1 \ra x_2 \ra v \ra v \ra x_1 \ra x_2 \ra v \ra v \ra x_1 \ra x_2 \ra v.\]
We illustrate the proof of Lemma \ref{lem:equal-dimension} for the string $\varLambda$-module $Y=X$. For $Y=X$ and for a diagrammatic morphism $f:X\rightarrow Y$, $N_f$ can only appears as a subdiagram in $Y$ as Case 2.1 or Case 2.8 (note that there are two subdiagrams $N_{f}$ and $N_{f'}$ of $Y$ which belong to Case 2.8). Denote by $\mathscr{B}$ the set of diagrammatic morphisms from $X$ to $Y$. Recall that we have defined an equivalence relation $\sim$ on $\mathscr{B}$: for $f,f'\in\mathscr{B}$, $f\sim f'$ if $N_{f}$ and $N_{f'}$ is the same subdiagram of $Y$. We {denote by} $\mathscr{C}$ the equivalence class of a fixed diagrammatic morphism $f:X\rightarrow Y$ and $\mathscr{C}'$ the corresponding subset of $\sHom(\Omega(X),Y)$.

		In Case 2.1, {$N_f =Y$.} We label the vertices in the diagrams of $X,Y$ and $P_v$ as follows:
		\[X=v_X \ra x_1 \ra x_2 \ra v \ra v \ra x_1 \ra x_2 \ra v \ra v \ra x_1 \ra x_2 \ra v,\]
		\[Y=v_1\urcorner\ra x_1 \ra x_2 \ra v' \ra v' \ra x_1 \ra x_2 \ra v' \ra v' \ra x_1 \ra x_2 \ra v_0,\]
		\[{P_v=} % https://q.uiver.app/#q=WzAsMjUsWzMsMCwidl9ZIl0sWzIsMSwidl9hIl0sWzIsMiwieF8xIl0sWzIsMywieF8yIl0sWzIsNCwidiJdLFsyLDUsInYiXSxbMiw2LCJ4XzEiXSxbMiw3LCJ4XzIiXSxbMiw4LCJ2Il0sWzIsOSwidiJdLFsyLDEwLCJ4XzEiXSxbMiwxMSwieF8yIl0sWzMsMTIsInYiXSxbNCwxLCJ4XzEiXSxbNCwyLCJ4XzIiXSxbNCwzLCJ2Il0sWzQsNCwidiJdLFs0LDUsInhfMSJdLFs0LDYsInhfMiJdLFs0LDcsInYiXSxbNCw4LCJ2Il0sWzQsOSwieF8xIl0sWzQsMTAsInhfMiJdLFs0LDExLCJ2X2IiXSxbMCw2LCJQX3Y9Il0sWzEsMTIsIiIsMCx7InN0eWxlIjp7ImJvZHkiOnsibmFtZSI6ImRhc2hlZCJ9LCJoZWFkIjp7Im5hbWUiOiJub25lIn19fV1d
		\begin{tikzcd}[row sep=-4, column sep=0]
			&&& {v_Y} \\
			&& {v_a} && {x_1} \\
			&& {x_1} && {x_2} \\
			&& {x_2} && v \\
			&& v && v \\
			&& v && {x_1} \\
			&& {x_1} && {x_2} \\
			&& {x_2} && v \\
			&& v && v \\
			&& v && {x_1} \\
			&& {x_1} && {x_2} \\
			&& {x_2} && {v_b} \\
			&&& v
			\arrow[dashed, no head, from=2-3, to=13-4]
		\end{tikzcd}.\]
		Then ${\mathscr{C}}=\{f_0\colon v_X\mapsto v_0\}\cup\{f_1\colon v_X\mapsto v_1\urcorner\}$, where $\underline{f_0}=0$, since $f_0$ factors through the projective module $P_v$:
		\[
		\begin{array}{cccccc}
		f_0\colon & X   & \lra     & P_v       & \lra     & Y \\
				& v_X & \longmapsto & {-\frac{1}{t}v_a+v_b} &         &   \\
				&     &         & v_Y       & \longmapsto & v_1\urcorner \\
				&     &         & v_a       & \longmapsto & 0 \\
				&     &         & v_b       & \longmapsto & v_0
		\end{array}
		,\]
		and $\underline{f_1}\neq 0$. Hence the basis of $\span \{\underline{f_0},\underline{f_1}\}$ is ${\underline{\mathscr{C}}}=\{\underline{f_1}\}$. We denote the $k$-basis element of $\Omega(X)$ corresponding to its top vertex also by $v_X$. In this case, ${\mathscr{C}'}$ has only one morphism $g_0\colon v_X\mapsto v_0$ given by the diagrammatic morphism from the string module $\Omega(X)/\soc (\Omega(X))$ to $Y$, and we have ${\underline{\mathscr{C}'}}=\{\underline{g_0}\}$.

		In the first occurrence of Case 2.8, $N_f$ is the higher subdiagram $v\rightarrow v$ of $Y$. After taking the flip of $Y$, we label the vertices in the diagrams of $X,Y$ and $P_v$ as follows:
		\[X=v_X \ra x_1 \ra x_2 \ra v \ra v \ra x_1 \ra x_2 \ra v \ra v \ra x_1 \ra x_2 \ra v,\]
		\[Y=v' \la x_2 \la x_1 \la v' \la v' \la x_2 \la x_1 \la v_1 \la v_1\urcorner \la x_2 \la x_1 \la v'_Y,\]
		\[{P_v=} % https://q.uiver.app/#q=WzAsMjUsWzMsMCwidl9ZIl0sWzIsMSwidl9hIl0sWzIsMiwieF8xIl0sWzIsMywieF8yIl0sWzIsNCwidiJdLFsyLDUsInYiXSxbMiw2LCJ4XzEiXSxbMiw3LCJ4XzIiXSxbMiw4LCJ2Il0sWzIsOSwidiJdLFsyLDEwLCJ4XzEiXSxbMiwxMSwieF8yIl0sWzMsMTIsInYiXSxbNCwxLCJ4XzEiXSxbNCwyLCJ4XzIiXSxbNCwzLCJ2Il0sWzQsNCwidl9iIl0sWzQsNSwieF8xIl0sWzQsNiwieF8yIl0sWzQsNywidiJdLFs0LDgsInYiXSxbNCw5LCJ4XzEiXSxbNCwxMCwieF8yIl0sWzQsMTEsInYiXSxbMCw2LCJQX3Y9Il0sWzEsMTIsIiIsMCx7InN0eWxlIjp7ImJvZHkiOnsibmFtZSI6ImRhc2hlZCJ9LCJoZWFkIjp7Im5hbWUiOiJub25lIn19fV1d
		\begin{tikzcd}[row sep=-4, column sep=0]
			&&& {v_Y} \\
			&& {v_a} && {x_1} \\
			&& {x_1} && {x_2} \\
			&& {x_2} && v \\
			&& v && {v_b} \\
			&& v && {x_1} \\
			&& {x_1} && {x_2} \\
			&& {x_2} && v \\
			&& v && v \\
			&& v && {x_1} \\
			&& {x_1} && {x_2} \\
			&& {x_2} && v \\
			&&& v
			\arrow[dashed, no head, from=2-3, to=13-4]
		\end{tikzcd}.\]
		Then ${\mathscr{C}}=\{f_1\colon v_X\mapsto v_1\}$, where $\underline{f_1}=0$ since $f_1$ factors through the projective module $P_v$:
		\[
		\begin{array}{cccccl}
		f_1\colon & X   & \lra     & P_v       & \lra     & Y \\
				& v_X & \longmapsto & v_b &         &   \\
				&     &         & v_Y       & \longmapsto & v'_Y  \\
				&     &         & v_a       & \longmapsto & 0 \\
				&     &         & v_b       & \longmapsto & v_1
		\end{array}
		,\] which implies that ${\underline{\mathscr{C}}}=\varnothing$. In this case, ${\mathscr{C}'}=\{g_1\colon v_X\mapsto v_1\}$, and we have ${\underline{\mathscr{C}'}}=\varnothing$ since $g_1$ factors through the projective module $P_v$ as well:
		\[
		\begin{array}{cccccl}
		g_1\colon & \Omega(X)   & \lra     & P_v       & \lra     & Y \\
				& v_X & \longmapsto & v_b &         &   \\
				&     &         & v_Y       & \longmapsto & v'_Y \\
				&     &         & v_a       & \longmapsto & 0 \\
				&     &         & v_b       & \longmapsto & v_1
		\end{array}
		.\]

		In the second occurrence of Case 2.8, {$N_f$ is the lower subdiagram $v\rightarrow v$ of $Y$.} After taking the flip of $Y$, we label the vertices in the diagrams of $X,Y$ and $P_v$ as follows:
		\[X=v_X \ra x_1 \ra x_2 \ra v \ra v \ra x_1 \ra x_2 \ra v \ra v \ra x_1 \ra x_2 \ra v,\]
		\[Y=v' \la x_2 \la x_1 \la v_1 \la v_1\urcorner \la x_2 \la x_1 \la v' \la v' \la x_2 \la x_1 \la v'_Y,\]
		\[{P_v=}
		\begin{tikzcd}[row sep=-4, column sep=0]
			&&& {v_Y} \\
			&& {v_a} && {x_1} \\
			&& {x_1} && {x_2} \\
			&& {x_2} && v \\
			&& v && v \\
			&& v && {x_1} \\
			&& {x_1} && {x_2} \\
			&& {x_2} && v \\
			&& v && {v_b} \\
			&& v && {x_1} \\
			&& {x_1} && {x_2} \\
			&& {x_2} && v \\
			&&& v
			\arrow[dashed, no head, from=2-3, to=13-4]
		\end{tikzcd}.\]
		Then ${\mathscr{C}}=\{f_1\colon v_X\mapsto v_1\}$, where $\underline{f_1}=0$ since $f_1$ factors through the projective module $P_v$:
		\[
		\begin{array}{cccccl}
		f_1\colon & X   & \lra     & P_v       & \lra     & Y \\
				& v_X & \longmapsto & v_b &         &   \\
				&     &         & v_Y       & \longmapsto & v'_Y \\
				&     &         & v_a       & \longmapsto & 0 \\
				&     &         & v_b       & \longmapsto & v_1
		\end{array}
		,\] which implies that ${\underline{\mathscr{C}}}=\varnothing$. In this case, ${\mathscr{C}'}=\{g_1\colon v_X\mapsto v_1\}$, and we have ${\underline{\mathscr{C}'}}=\varnothing$ since $g_1$ factors through the projective module $P_v$ as well:
		\[
		\begin{array}{cccccl}
		g_1\colon & \Omega(X)   & \lra     & P_v       & \lra     & Y \\
				& v_X & \longmapsto & v_b &         &   \\
				&     &         & v_Y       & \longmapsto & v'_Y \\
				&     &         & v_a       & \longmapsto & 0 \\
				&     &         & v_b       & \longmapsto & v_1
		\end{array}
		.\]

		We can verify that 
		\[\Hom(X,Y)=\span{\bigcup\limits_{\mathscr{C}\in\mathscr{B}/\sim}{\mathscr{C}}}=\span\{f_0^{2.1},f_1^{2.1},f_1^{2.8.1},f_1^{2.8.2}\}\]
		and 
		\[\Hom(\Omega(X),Y)=\span{\bigcup\limits_{\mathscr{C}\in\mathscr{B}/\sim}{\mathscr{C}'}}=\span\{g_0^{2.1},g_1^{2.8.1},g_1^{2.8.2}\},\]
		where the superscripts indicate the cases in which the morphisms appear. Thus, by the above computation, we conclude that
		\[\sHom(X,Y)=\span {\bigcup\limits_{\mathscr{C}\in\mathscr{B}/\sim}\underline{\mathscr{C}}}= \span \{\underline{f_1^{2.1}}\}\text{\quad and\quad }\sHom(\Omega(X),Y)=\span {\bigcup\limits_{\mathscr{C}\in\mathscr{B}/\sim}\underline{\mathscr{C}'}}= \span \{\underline{g_0^{2.1}}\}.\]
		Therefore,
		\[\dim\nolimits_k \sHom(X,Y)={\sum_{\mathscr{C}\in\mathscr{B}/\sim}|\underline{\mathscr{C}}|}=1={\sum_{\mathscr{C}\in\mathscr{B}/\sim}|\underline{\mathscr{C}'}|}=\dim\nolimits_k \sHom(\Omega(X),Y).\]
		
	\end{example}


\section{Brauer graph algebras are closed under stable equivalence of Morita type}\label{sec:stb-closed}

Before proceeding to the proof of the main theorems, we establish two more lemmas.

\begin{Lem}\label{lem:stable-tree=BGA}
Let $A$ be a Brauer graph algebra associated with the Brauer graph $(\Gamma, m)$, and let $B$ be a symmetric stably biserial algebra defined by the Brauer graph $(\Gamma', m')$ and a set $\mathcal{L}$ of deformed loops.
\begin{enumerate}
	\item If
\begin{enumerate}
	\item $A$ and $B$ are stably equivalent;
	\item $A$ is not local;
	\item$\Gamma$ has only one face,
\end{enumerate} then $B$ is also a Brauer graph algebra.
\item If
\begin{enumerate}
	\item $A$ and $B$ are stably equivalent of Morita type;
	\item$\Gamma$ has only one face,
\end{enumerate} then $B$ is also a Brauer graph algebra.
\end{enumerate}

\end{Lem}

\begin{proof}
We first prove (1). We note that both algebras $A$ and $B$ are indecomposable because we assume that the corresponding Brauer graphs are connected. According to Theorem \ref{thm:sym-StBA}, the conclusion holds when the field characteristic is not $2$, since in that case every symmetric stably biserial algebra is isomorphic to a Brauer graph algebra. Thus, we may restrict to the case where the field $k$ has characteristic $2$. We may {also} assume that $A$ is representation-infinite. Let $n$ be the number of isomorphism classes of simple $A$-modules. Then the unique face of $\Gamma$ has perimeter $2n$.

Since $A$ is not local, $n>1$. By Proposition \ref{prop:stbA-tube-and-face}, the number of faces of perimeter $2n$ in $\Gamma$ (resp.~$\Gamma'$) coincides with the number of tubes of rank $n$ in the Auslander-Reiten quiver of the algebra $A$ (resp.~$B$) that are not stable under the syzygy $\Omega$, divided by $2$. Since $A$ and $B$ are stably equivalent and $\Gamma$ has a face of perimeter $2n$, it follows that $\Gamma'$ must also contain a face of perimeter $2n$.

As $B$ is stably equivalent to $A$, it follows from Theorem \ref{thm:AR-conj} that the number of isomorphism classes of simple $B$-modules is also $n$. Hence, the ribbon graph $\Gamma'$ has $n$ edges, and the sum of the perimeters of all its faces is $2n$. This forces $\Gamma'$ to have exactly one face.

Now, suppose that $B$ is not a Brauer graph algebra. Then the quiver of $B$ must contain deformed loops, implying that $\Gamma'$ contains loops. In that case, $(\Gamma', m')$ would have a face of perimeter $1$, which is a contradiction. We conclude that $\mathcal{L} = \varnothing$ and therefore $B$ is a Brauer graph algebra.

We now prove (2). In fact, it suffices to consider the case where $A$ is a local algebra, that is, $\Gamma$ is an edge connected with two different vertices with multiplicities $m_1$ and $m_2$ (both $m_1$ and $m_2$ are bigger than $1$, since $A$ is representation-infinite). Then by Proposition \ref{prop:center-of-stBA}, we have $$Z(A)/R(A)\cong k[x_1,x_2]/\langle x_1^{m_1},x_2^{m_2},x_1x_2\rangle.$$ As $B$ is stably equivalent to $A$, it follows from Theorem \ref{thm:AR-conj} that the number of isomorphism classes of simple $B$-modules is also $1$. Therefore, $\Gamma'$ is a loop or an edge connected with two different vertices. Suppose that $B$ is not a Brauer graph algebra. Then $\Gamma'$ is a loop around a vertex with multiplicity $m$ and $\mathcal{L} \neq \varnothing$. By Theorem~\ref{thm:sym-StBA}, the local algebra $B$ has the form
$$B\cong k\langle x,y\rangle/\big\langle
(xy)^m-(yx)^m,\ x^2-\lambda_1(yx)^m,\ y^2-\lambda_2(yx)^m,
\ (yx)^m x,\ (yx)^m y\big\rangle,$$
where $\lambda_1,\lambda_2\in k$ and $(\lambda_1,\lambda_2)\neq(0,0)$. Put
$$s=(yx)^m=(xy)^m,
\qquad
q_x=y(xy)^{m-1},
\qquad
q_y=x(yx)^{m-1},
\qquad
p_{1,t}=(yx)^t+(xy)^t.$$
Thus $R(B)=ks$. Set $x_1'=q_x+R(B)$, $x_2'=q_y+R(B)$ and $x_3'=p_{1,1}+R(B)$
in $Z(B)/R(B)$ (when $m=1$, the last element is zero). Since $(x_3')^t=p_{1,t}+R(B)$, Proposition~\ref{prop:center-of-stBA} gives
$$Z(B)/R(B)\cong k[x_1',x_2',x_3']\big/\big\langle(x_1')^2,(x_2')^2,(x_3')^m,x_1'x_2',x_1'x_3',x_2'x_3'\big\rangle.$$
By Proposition \ref{prop:ZZ}, we have $m=1$ and $m_1=m_2=2$. Thus, in this case, the algebras $A$ and $B$ are precisely those appearing in Lemma~\ref{lem:stb-local-iso}. By Lemma~\ref{lem:stb-local-iso}, if $B$ is symmetric, then $A\cong B$. Hence $B$ is also a Brauer graph algebra, which contradicts the initial assumption. 
\end{proof}

We prove the following lemma, which will be needed in the proof of our main result.

\begin{Lem}\label{lem:exist-string-then-tree}
	Let $A=kQ/I$ be a representation-infinite Brauer graph algebra with associated Brauer graph $\Gamma$. {Suppose that}
 there exist non-isomorphic string modules $M_1$ and $M_2$ such that
\begin{enumerate}
    \item $\Omega_A(M_1) = M_2$ and $\Omega_A(M_2) = M_1$;
    \item $\mathrm{top}(M_1) \cong \mathrm{top}(M_2)$;
    \item $M_1$ and $M_2$ lie on the mouths of distinct tubes of rank~$1$ in the Auslander-Reiten quiver of $A$.
\end{enumerate}
{Then} the ribbon graph $\Gamma$ consists of a single edge connecting two distinct vertices.
\end{Lem}

\begin{proof}
	Since the string modules $M_1$ and $M_2$ lie on the mouths of distinct tubes of rank~$1$, they are maximal uniserial (cf. \cite[II.6.2(1)]{E}).
Let $P$ be the projective cover of $M_1$.
Since $A$ is symmetric and both $M_1$ and $M_2$ are string modules, we have {the isomorphisms of simple modules}
\[
\soc(M_2) \cong \soc(P) \cong \top(P) \cong \top(M_1).
\]
Denote this simple module by $S_v$, corresponding to a vertex $v$ in the quiver $Q$.
By symmetry and the assumption $\mathrm{top}(M_1) \cong \mathrm{top}(M_2)$, it follows that
\[
S_v \cong \top(M_1) \cong \soc(M_1) \cong \top(M_2) \cong \soc(M_2).
\]

Since $M_1 \not\cong M_2$, the starting arrow $\alpha$ of the string defining $M_1$ must differ from the starting arrow $\beta$ of the string defining $M_2$.
Consequently, the quotient $P / \soc(P)$ is a string module whose corresponding string contains a subword of the form
\[
\cdots \xleftarrow{\beta} v \xleftarrow{\beta} v \xrightarrow{\alpha} \cdots .
\]
Therefore, $\beta$ is a loop.
By symmetry, $\alpha$ is also a loop.
Since $\beta^2 \neq 0$ and $A$ is special biserial, we must have $\beta \alpha = \alpha \beta = 0$.
Hence, the ribbon graph $\Gamma$ consists of a single edge connecting two distinct vertices.
\end{proof}

We now prove the main results of this section.

\begin{Thm}\label{thm:st-BGA}
 Let $A$ be a non-local Brauer graph algebra. Then for any basic symmetric algebra $B$ without semisimple summands, if $B$ and $A$ are stably equivalent, then $B$ is also a Brauer graph algebra.
\end{Thm}

\begin{proof}
	Without loss of generality, we may assume that $A$ is indecomposable representation-infinite and the field characteristic is $2$. By the assumptions and Theorem~\ref{thm:sta-to-BGA=StB}, $B$ is a symmetric stably biserial algebra. Assume that the Brauer graph associated with $A$ is $(\Gamma, m)$, and let $B$ be a symmetric stably biserial algebra defined by the Brauer graph $(\Gamma', m')$ and a set $\mathcal{L}$ of deformed loops.

{If $B$ is not special biserial, then the set $\mathcal{L}$ is nonempty. So} $\Gamma'$ contains a loop, which corresponds to a deformed loop in $\mathcal{L}$. Let
$$\begin{tikzpicture}
\draw (-0.5,0) circle (0.5);
\fill (0,0) circle (0.5ex);
\node at(-1.2,0) {$1$};
\node at(-0.3,0) {$m$};
\draw[-] (0,0) -- (1,1);
\draw[-] (0,0) -- (1,-1);
\draw (-0.45,-0.15) rectangle (-0.15,0.15);
\node at(0.5,0.75) {$2$};
\node at(0.5,-0.75) {$n$};
\draw[dotted] (0.866,-0.5) arc (-30:30:1);
\end{tikzpicture}$$
be a subgraph of $(\Gamma',m')$, where the edge marked by $1$ gives a deformed loop of $B$. Let
$$\begin{tikzpicture}
\node at(0,0) {$1$};
\node at(0,-0.3) {$2$};
\draw[dotted] (0,-0.5) -- (0,-0.9);
\node at(0,-1.1) {$n$};
\node at(0,-1.4) {$1$};
\node at(0,-1.7) {$1$};
\draw[dotted] (0,-1.9) -- (0,-2.3);
\node at(0,-2.5) {$1$};
\node at(0,-2.8) {$2$};
\draw[dotted] (0,-3) -- (0,-3.4);
\node at(0,-3.6) {$n$};
\node at(0,-3.9) {$1$};
\node at(-1,-2) {$X$};
\node at(-0.5,-2) {$=$};
\end{tikzpicture}$$
be a maximal uniserial $B$-module, then
$$\begin{tikzpicture}
\node at(0,0) {$1$};
\node at(0,-0.3) {$2$};
\draw[dotted] (0,-0.5) -- (0,-0.9);
\node at(0,-1.1) {$n$};
\node at(0,-1.4) {$1$};
\node at(0,-1.7) {$1$};
\draw[dotted] (0,-1.9) -- (0,-2.3);
\node at(0,-2.5) {$1$};
\node at(0,-2.8) {$2$};
\draw[dotted] (0,-3) -- (0,-3.4);
\node at(0,-3.6) {$n$};
\node at(0,-3.9) {$1$};
\node at(-1.5,-2) {$\Omega_{B}(X)$};
\node at(-0.5,-2) {$=$};
\node at(0.6,-2) {$t$};
\draw[dashed] (0.15,-3.9) arc (-15:15:7.5);
\end{tikzpicture},$$
where $t\in k^{*}$.
{Note that $X$ is a string module and $\Omega_B(X)$ is an exceptional band module, and these two modules lie on the mouths of distinct tubes of rank~$1$ in the Auslander-Reiten quiver of $B$ (cf. Subsection 2.3).}
For any string module $Y \neq X$ of $B$, Lemma~\ref{lem:equal-dimension} implies that
\[
\mathrm{dim}_{k} \underline{\mathrm{Hom}}_{B}(X, Y) = \mathrm{dim}_{k} \underline{\mathrm{Hom}}_{B}(\Omega_{B}(X), Y).
\]
Let $M$ be the $A$-module corresponding to $X$ under {a} stable equivalence
$\alpha\colon B\text{-}\underline{\mathrm{mod}} \to A\text{-}\underline{\mathrm{mod}}$. The above identity implies that for any simple $A$-module $S$,
\[
\mathrm{dim}_{k} \underline{\mathrm{Hom}}_{A}(M, S) = \mathrm{dim}_{k} \underline{\mathrm{Hom}}_{A}(\Omega_{A}(M), S).
\]
This equality of dimensions means that the multiplicity of each simple module $S$ in the top of $M$ coincides with its multiplicity in the top of $\Omega_A(M)$. Therefore,
\[
\mathrm{top}(M) \cong \mathrm{top}(\Omega_{A}(M)).
\]

%Since soc$(M)\cong$top$(\Omega_{A}(M))$ and soc$(\Omega_{A}(M))\cong$top$(M)$, we have top$(M)\cong$top$(\Omega_{A}(M))\cong$soc$(M)\cong$soc$(\Omega_{A}(M))$. Since $X$ and $\Omega_{B}(X)$ are exceptional $B$-modules, $M$ and $\Omega_{A}(M)$ are exceptional $A$-modules.

Since $X$ and $\Omega_B(X)$ lie on the mouths of distinct tubes of rank~$1$ in the Auslander-Reiten quiver of $B$, their images under the stable equivalence, say $M$ and $\Omega_A(M)$, will also lie on the mouths of distinct tubes of rank~$1$ in the Auslander-Reiten quiver of $A$.

If $M$ is a string module, then so is $\Omega_A(M)$ since $A$ is a Brauer graph algebra.
By Lemma~\ref{lem:exist-string-then-tree}, the Brauer graph $\Gamma$ is a single edge, which implies that $A$ is a local algebra. This contradicts the assumption that $A$ is non-local. Therefore, $M$ cannot be a string module.

If $M$ is a non-exceptional band module, then by Proposition~\ref{prop:non-exceptional-tube}, for some integer $n > 2$, there exists a string module $Y$ {which does not lie in a tube of rank $1$ in the Auslander-Reiten quiver of $A$} such that
\[
\mathrm{dim}_k\, \underline{\Hom}_A(M, Y) > n.
\]
Consequently,
\[
\mathrm{dim}_k\, \underline{\Hom}_B(X, \alpha^{-1}(Y)) > n.
\]
Since $Y$ does not lie in a tube of rank $1$ of the Auslander-Reiten quiver of $A$, the same holds for $\alpha^{-1}(Y)$. Therefore, $\alpha^{-1}(Y)$ is a string module.
However, by {Lemma~\ref{lem:equal-dimension}}, we have
\[
\mathrm{dim}_k\, \underline{\Hom}_B(X, \alpha^{-1}(Y)) \le 2,
\]
which gives a contradiction.


If $M$ is an exceptional band module, then suppose $M$ is given by the exceptional band
\[
p_{2k}^{-1} p_{2k-1} \cdots p_{2}^{-1} p_{1}
\]
of $A$, with constant $\lambda \in k^*$, where $p_1, p_3, \dots, p_{2k-1}$ are minimal band lines and $p_2, p_4, \dots, p_{2k}$ are maximal band lines.
By Proposition~\ref{prop:double-faces and exceptional bands}, this exceptional band corresponds to a double-face $F' = (h_1, h_2, \dots, h_k)$ of $\widetilde{\Gamma}$. Suppose $F'$ belongs to the face $F$ of $\widetilde{\Gamma}$. Then the length of $F$ is either $k$ or $2k$.

\textbf{Case 1:} If the length of $F$ is $k$, then $F$ contains exactly the same half-edges as $F'$, and the exceptional band of $\Omega_A(M)$ is also given by the double-face $F'$ of $\widetilde{\Gamma}$. A direct calculation shows that the constant of the band module \(\Omega_A(M)\) is $-\lambda$, which equals $\lambda$ since $\mathrm{char}(k) = 2$. Hence $M \cong \Omega_A(M)$, which implies $X \cong \Omega_B(X)$, a contradiction.

\textbf{Case 2:} If the length of $F$ is $2k$. In this case, $F$ splits into two double-faces $F'$ and $F''$ of $\widetilde{\Gamma}$, and the exceptional band of $\Omega_A(M)$ is given by $F''$.
Let $R$ be the quiver {defined} by the ribbon graph $\widetilde{\Gamma}$. Then the face $F$ of $\widetilde{\Gamma}$ induces an oriented cycle
\[
c = \alpha_{2k} \cdots \alpha_2 \alpha_1
\]
in $R$. By the definition of a face, all $\alpha_i$ are distinct. Let $v_i = s(\alpha_i)$ for $1 \leq i \leq 2k$ (note that $v_i$ may equal $v_j$ for different $i, j$). Then:
\[
\mathrm{top}(M) = S_{v_1} \oplus S_{v_3} \oplus \cdots \oplus S_{v_{2k-1}}, \quad
\mathrm{top}(\Omega_A(M)) = S_{v_2} \oplus S_{v_4} \oplus \cdots \oplus S_{v_{2k}}.
\]
Since $\mathrm{top}(M) \cong \mathrm{top}(\Omega_A(M))$, for each vertex $v$ in the oriented cycle $c$, there exist an odd index $2i-1$ and an even index $2j$ such that $v = v_{2i-1} = v_{2j}$. This implies that there are two arrows in $c$ starting at $v$ (namely $\alpha_{2i-1}$ and $\alpha_{2j}$) and two arrows terminating at $v$ (namely $\alpha_{2i-2}$ and $\alpha_{2j-1}$).
Since $R$ is connected, and each vertex in $R$ has at most two incoming arrows and at most two outgoing arrows, the oriented cycle $c$ passes through every vertex of $R$, and every arrow of $R$ occurs in $c$ exactly once. Equivalently, the sequence of arrows appearing in $c$ exhausts the arrow set of $R$. Therefore, $F$ is the only face of $\widetilde{\Gamma}$, which means that $\Gamma$ also has only one face (cf. Subsection \ref{subsec:excep-band}). By Lemma~\ref{lem:stable-tree=BGA}, $B$ is a Brauer graph algebra, which contradicts the initial assumption.
\end{proof}

\begin{Thm}\label{thm:st.M-BGA}
 Let $A$ be a Brauer graph algebra. Then for any basic algebra $B$ without semisimple summands, if $B$ and $A$ are stably equivalent of Morita type, then $B$ is also a Brauer graph algebra.
\end{Thm}

\begin{proof}
	Without loss of generality, we may assume that $A$ is indecomposable representation-infinite and the field characteristic is $2$. By \cite[Corollary 2.4 and Proposition 2.1]{L} and Theorem \ref{thm:sta-to-BGA=StB}, $B$ is indecomposable and symmetric stably biserial. We retain the notation from the proof of Theorem~\ref{thm:st-BGA}. In fact, the only case that remains to be checked is the case where $M=\alpha(X)$ is a string module.

If $M$ is a string module, then so is $\Omega_A(M)$ since $A=kQ_\Gamma/I_\Gamma$ is the Brauer graph algebra associated with the Brauer graph $(\Gamma,m)$.
By Lemma~\ref{lem:exist-string-then-tree}, the Brauer graph $\Gamma$ is a single edge connecting two distinct vertices, which has only one face.
Hence, by Lemma~\ref{lem:stable-tree=BGA}, the algebra $B$ is also a Brauer graph algebra.
\end{proof}

The theorem above allows us to give an alternative proof of the main result in \cite{AZ2}.

\begin{Cor} \label{closed-under-derived-equivalence}
	Let $A$ be a Brauer graph algebra. Then for any basic algebra $B$, if $B$ and $A$ are derived equivalent, then $B$ is also a Brauer graph algebra.
\end{Cor}

\begin{proof}
Since $A$ is symmetric, it follows from \cite[Corollary 5.3]{Ric} that $B$ is also symmetric. Moreover, it follows from \cite[Corollary 5.5]{Ric} that $A$ and $B$ are stably equivalent of Morita type. Since two derived equivalent algebras have isomorphic centers and (as an algebra) $A$ is indecomposable, $B$ is also indecomposable. Moreover, since $A$ is not a simple algebra, $B$ is also not a simple algebra. Then, by Theorem \ref{thm:st.M-BGA}, $B$ is also a Brauer graph algebra.
\end{proof}


\begin{thebibliography}{88}

\bibitem{A}
{{\sc M. Antipov,} Invariants of the stable equivalence of symmetric special biserial algebras. 
Journal of Mathematical Sciences \textbf{140}(5) (2007), 611--621.}

\bibitem{AZ1}
{{\sc M. Antipov and A. Zvonareva,} On stably biserial algebras and the Auslander-Reiten conjecture for special biserial algebras. 
Journal of Mathematical Sciences \textbf{240}(4) (2019), 375--394.}

\bibitem{AZ2}
{{\sc M. Antipov and A. Zvonareva,} Brauer graph algebras are closed under derived equivalence. 
Mathematische Zeitschrift \textbf{301}(2) (2022), 1963--1981.}

\bibitem{AIP}
{{\sc S. Ariki, K. Iijima and E. Park,} Representation type of finite quiver Hecke algebras of type $A^{(1)}_{l}$ for arbitrary parameters. 
International Mathematics Research Notices \textbf{2015}(15) (2015), 6070--6135.}

\bibitem{Asa2}
{{\sc H. Asashiba,}
On a lift of an individual stable equivalence to a standard derived equivalence for representation-finite self-injective algebras.
Algebras and Representation Theory \textbf{6}(4) (2003), 427--447.}

\bibitem{ARS}
{{\sc M. Auslander, I. Reiten and S. O. Smal{\o},}
{\it Representation theory of Artin algebras}.
Cambridge Studies in Advanced Mathematics, vol. 36.
Cambridge University Press, 1995.}

\bibitem{BC}
{{\sc K. Baur and R. Coelho Sim{\~o}es,}
A geometric model for the module category of a string algebra.
arXiv preprint arXiv:2403.07810, 2024.}

\bibitem{Ben}
{{\sc D. J. Benson,}
{\it Modular representation theory: new trends and methods}.
Lecture Notes in Mathematics, vol. 1081.
Springer, 1984.}

\bibitem{BHS}
{{\sc R. Bocian, T. Holm and A. Skowro{\'n}ski,}
Derived equivalence classification of weakly symmetric algebras of Euclidean type.
Journal of Pure and Applied Algebra \textbf{191}(1--2) (2004), 43--74.}

\bibitem{BR}
{{\sc M. C. R. Butler and C. M. Ringel,}
Auslander--Reiten sequences with few middle terms and applications to string algebras.
Communications in Algebra \textbf{15}(1--2) (1987), 145--179.}

\bibitem{Cra89}
{{\sc W. W. Crawley-Boevey,}
Maps between representations of zero-relation algebras.
Journal of Algebra \textbf{126}(2) (1989), 259--263.}

\bibitem{DF}
{{\sc P. W. Donovan and M.-R. Freislich,}
The indecomposable modular representations of certain groups with dihedral Sylow subgroup.
Mathematische Annalen \textbf{238}(3) (1978), 207--216.}

\bibitem{D}
{{\sc D. Duffield,}
Auslander--Reiten components of symmetric special biserial algebras.
Journal of Algebra \textbf{508} (2018), 475--511.}

\bibitem{E}
{{\sc K. Erdmann,}
{\it Blocks of tame representation type and related algebras}.
Lecture Notes in Mathematics, vol. 1428.
Springer, 1990.}

\bibitem{GreenJA}
{{\sc J. A. Green,}
Walking around the Brauer tree.
Journal of the Australian Mathematical Society \textbf{17}(2) (1974), 197--213.}

\bibitem{GL}
{{\sc J. Guo and Y. Liu,}
The associated graded algebras of Brauer graph algebras I: finite representation type.
Communications in Algebra \textbf{49}(3) (2021), 1071--1103.}

\bibitem{KZ}
{{\sc S. K{\"o}nig and A. Zimmermann,}
{\it Derived equivalences for group rings}.
Lecture Notes in Mathematics, vol. 1685.
Springer, 1998.}

\bibitem{Kra}
{{\sc H. Krause,}
Maps between tree and band modules.
Journal of Algebra \textbf{137}(1) (1991), 186--194.}

\bibitem{CLLX}
{{\sc N. Li, P. Chen, Y. Liu and B. Xing,}
Classification of Brauer graph algebras under stable equivalence of Morita type.
In preparation.}

\bibitem{LL2}
{{\sc N. Li and Y. Liu,}
The liftability question for stable equivalences between representation-finite self-injective algebras.
Algebras and Representation Theory \textbf{26}(5) (2023), 1519--1547.}

\bibitem{LL}
{{\sc N. Li and Y. Liu,}
A generalization of Dugas' construction on stable auto-equivalences for symmetric algebras.
Journal of Algebra \textbf{678} (2025), 502--542.}

\bibitem{L}
{{\sc Y. Liu,}
Summands of stable equivalences of Morita type.
Communications in Algebra \textbf{36}(10) (2008), 3778--3782.}

\bibitem{OTY}
{{\sc Y. Ohnuki, K. Takeda and K. Yamagata,}
Symmetric Hochschild extension algebras.
Colloquium Mathematicae \textbf{80}(2) (1999), 155--174.}

\bibitem{OZ}
{{\sc S. Opper and A. Zvonareva,} 
Derived equivalence classification of Brauer graph algebras.
Advances in Mathematics \textbf{402} (2022), Article 108341.}

\bibitem{P}
{{\sc Z. Pogorza{\l}y,}
Algebras stably equivalent to selfinjective special biserial algebras.
Communications in Algebra \textbf{22}(4) (1994), 1127--1160.}

\bibitem{Ric}
{{\sc J. Rickard,}
Derived equivalences as derived functors.
Journal of the London Mathematical Society (2) \textbf{43}(1) (1991), 37--48.}

\bibitem{RR}
{{\sc J. Rickard and R. Rouquier,}
Stable categories and reconstruction.
Journal of Algebra \textbf{475} (2017), 287--307.}

\bibitem{Sch99}
{{\sc J. Schr{\"o}er,}
Modules without self-extensions over gentle algebras.
Journal of Algebra \textbf{216}(1) (1999), 178--189.}

\bibitem{Sch}
{{\sc S. Schroll,}
Trivial extensions of gentle algebras and Brauer graph algebras.
Journal of Algebra \textbf{444} (2015), 183--200.}

\bibitem{S}
{{\sc S. Schroll,}
Brauer graph algebras.
In: I. Assem and S. Trepode (eds.),
{\it Homological methods, representation theory, and cluster algebras}.
CRM Short Courses.
Springer, 2018, 177--223.}

\bibitem{ZZ}
{{\sc G. Zhou and A. Zimmermann,}
Classifying tame blocks and related algebras up to stable equivalences of Morita type.
Journal of Pure and Applied Algebra \textbf{215}(12) (2011), 2969--2986.}

\bibitem{Z}
{{\sc A. Zimmermann,}
{\it Representation theory: a homological algebra point of view}.
Algebra and Applications, vol. 19.
Springer, 2014.}

\end{thebibliography}

\end{document}
